\documentclass[letterpaper,twocolumn,openany,nodeprecatedcode,hidelinks]{dndbook}
\usepackage[english]{babel}
\usepackage[utf8]{inputenc}
\usepackage{pdfpages}
\usepackage{graphicx}
\usepackage{wrapfig}
\usepackage{color}
\usepackage{tikz}
\usepackage{lipsum}
\usepackage{hyperref}
\usepackage{titling}
\usepackage{pagecolor}
\usepackage{afterpage}
\usepackage{titletoc}
\usepackage{everyshi}
\usepackage{changepage}
\usepackage{geometry}
\usepackage{xstring}
\usepackage{afterpage}
\usepackage{multicol}

\usepackage{tabularray}
\usepackage{amsmath}
\usepackage{amsfonts}
\usepackage{enumitem}
\usepackage{tipa}
\usepackage{tablefootnote}
\usepackage{linguex}
\usepackage{expex}

\usepackage{makecell}

\usepackage{float}

\renewcommand\glt{} % removes extra vertical space from translation line
\newcommand{\redcolor}[1]{\color{red}#1\normalcolor}

% \catcode`_=\active
% \def_#1{\ifmmode\sb{#1}\else$\sb{#1}$\fi}


\usetikzlibrary{calc}

% \usepackage[unicodeparmaite]{tengwarscript}
% \usepackage{tengwarscript}
% \usepackage[annataritalic]{tengwarscript}
% \pdfmapfile{=tengwarscript.map}


\usepackage{fontspec}
% \newfontface{\tovian}{Ancient-Tovian.ttf}
\newfontface{\tovian}{Tovian.ttf}
\newfontface{\tengwar}{tngani.ttf}
\newcommand*{\red}[1]{\color{red}#1\normalcolor}

\def\vc#1{$\vcenter{\vspace{-30pt}\hbox{\vspace*{-50pt}#1}}$}

\newenvironment{nopagebreak}
  {\par\nobreak\vfil\penalty0\vfilneg
   \vtop\bgroup}
  {\par\xdef\tpd{\the\prevdepth}\egroup
   \prevdepth=\tpd}

\begin{document}
% \setmainfont{Ancient-Tovian.ttf}

% \tengwarannataritalic[1.5]



% IDEAS:
% Speak vs stand means performance versus factual
% to tree: stay in once place long time (good thing)


% Fethrdriel
\title{
    % \Huge\Tformen\Tthuule\TTdotbelow\Troomen\Tanna\Tando\TTdot\Ttelco\TTdot\Tlambe\TTdotbelow\\
    % {\fontsize{60pt}{30pt}\tovian ««ishmahluorniara»»}\\ % ˈiʃmahɫuorɳiara
    {\fontsize{80pt}{30pt}\tovian «««fethrydriel»»»}\\ % ˈapɸeθjdrieɫ
    % {\Huge\tengwar fethrydriel}\\
    \vspace{1in}
    \large \textit{Beginner's Guide to Tovian} \\
    % \Tando\Tsilme\TTacute\Ttinco\TTdoubler\TTthreedots\Tyanta
     % \Tquesse\TTacute\Ttinco\TTdoubler\TTthreedots\\
     % \Tquesse\TTacute\Ttinco\TTdoubler\TTthreedots\\
    % \huge{Tovis}\\
    % \begin{tikzpicture}[remember picture]\\%
        % \node (main) {\includegraphics[width=7in]{img/tovis2.png}};\\%
    % \end{tikzpicture}
    % \vspace{-1in}
    }

% \author{\large Jackson Dean}
\date{}

\newcommand{\continue}[0]{
% \textit{continued on the next page...}
\newpage
}

\newcommand\blfootnote[1]{%
  \begingroup
  \renewcommand\thefootnote{}\footnotetext{#1}%
  %\addtocounter{footnote}{-1}% not needed
  \endgroup
}
\titlespacing*{\subsubsection}
{0pt}{0pt}{0pt}

\newcommand{\ra}[0]{*}
\newcommand{\linkplaceholder}[0]{\subsubsubsection{}}

\newcounter{tablectr}

\newcommand*{\itemrow}[1]{
    \stepcounter{tablectr}%
    \StrSubstitute{#1}{ }{}[\temp] % Replace spaces in the name with nothing
    \StrSubstitute{\temp}{'}{}[\templabel] % Replace apostrophe
    \StrSubstitute{\templabel}{,}{}[\templabel] % Replace apostrophe
    \label{\templabel}
    \arabic{tablectr}\subsubsubsection{}\stepcounter{tablectr}-\arabic{tablectr} & \textbf{#1}
}

\newcommand{\rarity}[3]{%
    \small{\textit{
    \newline #1, #2 %
    \ifbool{#3}{(*)}%
    }}
}

\newcolumntype{B}{X}
\newcolumntype{s}{>{\hsize=0.20\hsize}X}





\newcommand{\multicolinterrupt}[2]{% Stuff to span both rows
\setcounter{tempcolnum}{\col@number}
\end{multicols}
#1%
\begin{multicols}{#2}
}

% \newcommand{\todo}[1]{\color{red}#1\normalcolor}
\newcommand{\todo}[1]{\textit{#1}}

\renewcommand{\chaptermark}[1]{\markboth{#1}{}}

\newcommand{\hanging}[0]{\hspace{0.5in}}


\setcounter{page}{1}
\maketitle



\newcommand{\CLan}[0]{la} % animate class
\newcommand{\CLin}[0]{lo} % inanimate class
\newcommand{\CLab}[0]{le} % abstract class

% cases
\newcommand{\CAnom}[0]{i} % nominative case (default, often omitted)
\newcommand{\CAacc}[0]{yu} % accusative case
\newcommand{\CAgen}[0]{hi} % genitive case
\newcommand{\CAdat}[0]{mi} % dative case
\newcommand{\CAloc}[0]{ti} % locative case (place)
\newcommand{\CAtem}[0]{tu} % temporal case (time)
\newcommand{\CAins}[0]{si} % instrumental case
\newcommand{\CApur}[0]{lhu} % purpose case
\newcommand{\CApri}[0]{ni} % privitive case
\newcommand{\CAcom}[0]{yi} % comlative case
\newcommand{\CAexc}[0]{fi} % excessive case
\newcommand{\CAcomp}[0]{asl} % comparative case
\newcommand{\CAade}[0]{asl} % adessive case
\newcommand{\CAall}[0]{su} % allative case
\newcommand{\CAess}[0]{pu} % allative case


\newcommand{\CAabl}[0]{di}

\newcommand{\casestable}[0]{
\subsection*{Noun Classes}
\begin{itemize}
    \item \textbf{Animate (living beings, includes plants)}
    \begin{itemize}
        \item \textbf{Prefix:} \textit{\CLan}
        \item \textbf{Example:} \textit{\CLan-driel} (person)
    \end{itemize}
    \item \textbf{Inanimate (objects and places)}
    \begin{itemize}
        \item \textbf{Prefix:} \textit{\CLin-}
        \item \textbf{Example:} \textit{\CLin-pelor} (city)
    \end{itemize}
    \item \textbf{Abstract (concepts and ideas)}
    \begin{itemize}
        \item \textbf{Prefix:} \textit{\CLab-}
        \item \textbf{Example:} \textit{\CLab-ara} (wisdom)
    \end{itemize}
\end{itemize}

\subsection*{Noun Cases}
\begin{itemize}
    \item \textbf{Nominative (subject)}
    \begin{itemize}
        \item \textbf{Prefix:} \textit{\CAnom-} (default case, often omitted)
        \item \textbf{Example:} \textit{\CAnom-\CLan-driel} (the person as the subject)
    \end{itemize}
    \item \textbf{Accusative (direct object)}
    \begin{itemize}
        \item \textbf{Prefix:} \textit{\CAacc-}
        \item \textbf{Example:} \textit{\CAacc-\CLan-driel} (the person as the direct object)
        \item \textbf{Sentence}: \textit{\CAacc-\CLin-tel afa muril \CAnom-\CLan-driel} (the person sees the lake)
    \end{itemize}
    \item \textbf{Genitive (possession)}
    \begin{itemize}
        \item \textbf{Prefix:} \textit{\CAgen-}
        \item \textbf{Example:} \textit{\CAgen-\CLan-driel} (of the person)
    \end{itemize}
    \item \textbf{Dative (indirect object)}
    \begin{itemize}
        \item \textbf{Prefix:} \textit{\CAdat-}
        \item \textbf{Example:} \textit{\CAdat-\CLan-driel} (to/for the person)
    \end{itemize}
    \item \textbf{Locative (location)}
    \begin{itemize}
        \item \textbf{Prefix:} \textit{\CAloc-}
        \item \textbf{Example:} \textit{\CAloc-\CLin-pelor} (in the city)
    \end{itemize}
        \item \textbf{Temporal (time)}
    \begin{itemize}
        \item \textbf{Prefix:} \textit{\CAtem-}
        \item \textbf{Example:} \textit{\CAtem-\CLab-loe} (at two o'clock)
    \end{itemize}
    \item \textbf{Instrumental (means or tool)}
    \begin{itemize}
        \item \textbf{Prefix:} \textit{\CAins-}
        \item \textbf{Example:} \textit{\CAins-\CLin-rith} (with the air)
    \end{itemize}
        \item \textbf{Purpose (goal or reason)}
    \begin{itemize}
        \item \textbf{Prefix:} \textit{\CApur-}
        \item \textbf{Example:} \textit{\CApur-\CLab-aray} (for the purpose of wisdom)
    \end{itemize}

    \item \textbf{Ablative (away)}
    \begin{itemize}
        \item \textbf{Prefix:} \textit{\CAabl-}
        \item \textbf{Example:} \textit{\CAabl-\CLab-loe} (before 2)
    \end{itemize}

     \item \textbf{Comlative (together)}
    \begin{itemize}
        \item \textbf{Prefix:} \textit{\CAcom-}
        \item \textbf{Example:} \textit{\CAcom-\CLan-ta} (together with him/her/them)

             \item \textbf{Allative (towards)}
    \begin{itemize}
        \item \textbf{Prefix:} \textit{\CAall-}
        \item \textbf{Example:} \textit{\CAall-\CLan-tal} (towards the east)
    \end{itemize}


     \item \textbf{Essive (time or state)}
    \begin{itemize}
        \item \textbf{Prefix:} \textit{\CAess-}
        \item \textbf{Example:} \textit{\CAess-\CLan-ninadriel} (as a child)
    \end{itemize}


    % \end{itemize}
       % \item \textbf{Excessive (too)}
    % \begin{itemize}
        % \item \textbf{Prefix:} \textit{\CAreq-}
        % \item \textbf{Example:} \textit{\CAreq-\CLab-ara} (requires wisdom)
    % \end{itemize}
           % \item \textbf{Relative (relative clause)}
    % \begin{itemize}
        % \item \textbf{Prefix:} \textit{\CAreq-}
        % \item \textbf{Example:} \textit{\CAreq-\CLab-ara} (requires wisdom)
    % \end{itemize}
\end{itemize}


\subsection*{Combining Classes and Cases}
To use these prefixes effectively, the structure would be Case + Class + Noun.
\begin{itemize}
    \item \textbf{Example:} \textit{\CAacc-\Clan-driel} (the animate person as the direct object)
\end{itemize}

\subsection*{Examples}
\begin{itemize}
    \item \textbf{The person (subject)}: \newline \textit{\CLan-driel} or \textit{\CAnom-\CLan-driel}
    \item \textbf{The person (direct object)}: \newline \textit{\CAacc-\CLan-driel}
    \item \textbf{Of the person (possession)}: \newline \textit{\CAgen-\CLan-driel}
    \item \textbf{To/for the person (indirect object)}:\newline  \textit{\CAdat-\CLan-driel}
    \item \textbf{In the city (location)}:\newline  \textit{\CAloc-\CLin-pelor}
    \item \textbf{With the air (instrumental)}:\newline  \textit{\CAins-\CLin-rith}
        \item \textbf{With the air (comlative)}:\newline  \textit{\CAcom-\CLin-rith}
\end{itemize}

}



\include{alphabet}

\twocolumn

\input{syllables}

Lefethr is an \textit{agglutanative} language, meaning morphemes do not change form


\section*{Plurals}
\textit{-e} indicates plurality. Verbs do not agree with number. Singular nouns therefore never end in \textit{-e}.

\begin{center}
\begin{tabular}{|c c|}
     \hline

    \textit{lothr} &  river \\
     \textit{lothre} & rivers \\

     \hline
     \hline

     \textit{nor} & sea \\

     \textit{nore} & seas \\
     \hline

\end{tabular}
\end{center}

\section*{Verbs}

Verbs are constructed around a consonant root that denotes its aspect. This root is combined with a suffix to indicate tense. The tense suffixes are \textit{-e} for past,\textit{ -a} for present, and \textit{-o} for future. Optional prefixes can be added to denote mood and voice. This final conjugated auxiliary verb is placed in front of the infinitive. The verb that indicates the action is never conjugated, only the auxiliary verb changes form to show person, aspect, tense, mood, and voice. The lexical predicate is an unchanged noun root used verbally. Every auxiliary requires a person prefix in formal writing and complex sentences. The order is person + mood + voice + aspect root + tense.

\begin{table}[H]
    \centering
    \begin{tabular}{|c|c|}
        \hline
        \textbf{Tense} & \textbf{Auxiliary Verb Suffix} \\
        \hline
        Past & -e \\
        Present & -a \\
        Future & -o \\
        \hline
    \end{tabular}
    % \caption{Person Suffixes in Tovian}
    \label{tab:verb-tense}
\end{table}


For example, \textit{fethr} is the noun root for speech, used verbally as ``speak.'' The continuous aspect root \textit{sh-} takes present \textit{-a} and third-person \textit{a-}, forming \textit{asha}. The auxiliary precedes the unchanged predicate: \textit{i-la-yel asha fethr}, ``the person is speaking.''

\begin{table}[h!]
    \centering
    \subsubsection{Aspects with present tense examples}
    \begin{tabular}{c|c|c|c}
    \hline
         \textbf{Aspect}& \makecell{\textbf{Root}\\\textbf{Verb}}& \makecell{\textbf{Aux.}\\\textbf{root}} & \makecell{\textbf{Pres. tense}\\\textbf{Example}} \\
\hline
         Simple&\makecell{\textit{fethr} \\ speak}  & f- & \makecell{\textit{afa muril}\\``sees"} \\ \hline
         \red{Imperfective}&\makecell{\textit{thraf} \\ wonder}  & th- & \makecell{\textit{atha muril}\\ ``is seeing"} \\ \hline
         Perfect&\makecell{\textit{lhara} \\ know}  & lh- & \makecell{\textit{alha muril} \\ ``has seen"} \\ \hline
         Near&\makecell{\textit{tlun} \\ share}& tl- & \makecell{\textit{atla muril} \\ ``is about to see" }\\ \hline
         Immediate&\makecell{\textit{kesh} \\ burn}& k- & \makecell{\textit{aka muril} \\ ``is seeing right now" }\\ \hline
         Habitual&\makecell{\textit{mel} \\ live}  & m- & \makecell{\textit{ama muril} \\``sees regularly" }\\ \hline
         Progressive&\makecell{\textit{sethr} \\ run}& s- & \makecell{\textit{asa muril} \\``is seeing now" }\\ \hline
         Continuous&\makecell{\textit{shotl} \\ flow}& sh- & \makecell{\textit{asha muril} \\ ``is seeing" }\\ \hline
         Iterative&\makecell{\textit{nalh} \\ think}& n- & \makecell{\textit{ana muril} \\ ``sees repeatedly" }\\ \hline
         Inceptive&\makecell{\textit{yalor} \\ begin}& y- & \makecell{\textit{aya muril} \\ ``begins  seeing" }\\ \hline
         Cessative&\makecell{\textit{panor} \\ end}& p- & \makecell{\textit{apa muril} \\ ``stops seeing" }\\ \hline
         Remote&\makecell{\textit{huysh} \\ throw}  & h- & \makecell{\textit{aha muril} \\``sees (distant)" \footnotemark[1]} \\
     \end{tabular}
     \\
     \medskip
     \raggedright
     {\footnotesize \footnotemark[1] \raggedleft In past and future tenses, the remote aspect refers to actions that took place long ago or far in the future, but the meaning in present tense is more nuanced. Elves tend to use the rare present remote aspect as an insult, for example \textit{ta-e aha fethr}: ``they speak but it is irrelevant"}
    \label{tab:aspects}
\end{table}

\medskip


\subsection*{Mood and Voice}

% \subsection*{All verb forms}
\begin{tabular}{rl}
&\textbf{Mood}\\
\textbf{unmarked}& $\rightarrow$ indicative \\
\textbf{wi-}& $\rightarrow$ subjunctive \\
\textbf{se-}& $\rightarrow$ imperative \\
\textbf{ko-} & $\rightarrow$ conditional \\
\textbf{xo-}& $\rightarrow$ counterfactual \\
\textbf{hle-}& $\rightarrow$ optative \\
\textbf{tlo-}& $\rightarrow$ obligative \\
\textbf{ef-}& $\rightarrow$ necessative \\
\textbf{sho-}& $\rightarrow$ potential \\



&\textbf{Voices}\\
\textbf{unmarked}& $\rightarrow$ active \\
\textbf{te-}& $\rightarrow$ reflexive \\
\textbf{pa-}& $\rightarrow$ passive \\
\textbf{mo-}& $\rightarrow$ middle \\
\textbf{ke-}& $\rightarrow$ causative \\
\textbf{ra-}& $\rightarrow$ reciprocal \\

% &\textbf{Other (not combined)}\\
% \textbf{li-}& $\rightarrow$ gerund \\

\end{tabular}

\medskip

Mood and voice can be combined by adding the mood prefix before the voice prefix. For example, adding the obligative mood marker \textit{tlo-} and the reflexive voice marker \textit{te-} to the habitual aspect auxilary root stem \textit{m-} with the present tense conjugation \textit{-a}, with third-person \textit{a-}, the final auxiliary verb is \textit{atlotema}.

\ex
\begingl
\gla i-la-yel atlotema sor //
\glb NOM-ANIM-person 3-OBLIG-REFL-HAB.PRES wash //
\glft "She should habitually wash herself." //
\endgl
\xe



\subsection{Unused Combinations}
Tovian is very flexible when it comes to combining different tenses, aspects, moods, and voices. However, certain verb conjugations result in combinations that are unlikely or never used due to their impractical or nonsensical nature. For example, the Past Imperative is very rare because commands most often refer to actions in the present or future.



\section{Copulas}
Copulas are formed by combining verbs and nouns to describe states and qualities. For example the word "mel" (life/living) can be used as a copula: "i-la-yel afa mel" means "The person is alive," or equally "The person lives" This method applies to various attributes by changing the verb root and noun. For instance, "yel afe yeol" translates to "The person was wise."

 \section{Comparative Constructions}

    Tovian uses comparative constructions to express excessiveness, for example: "bigger than appropriate." The ablative case is often used for this comparison: "the tree is too big" could be expressed as \textit{i-la-gelar \CAabl-le-velar afa mot} ("The tree is big away-from-worthiness" or simply "The tree is bigger than what is worthy/right").

\section{Adjectives}

Tovian does not use a separate class of adjectives. Instead, all descriptive qualities are expressed through nouns. These attributive nouns follow the noun they modify in a fixed noun–noun structure. For example, the phrase \textit{yel yeol} translates literally as “person wisdom,” conveying the meaning “the wise person.” To convey "people wisdom" or "the wisdom of people," one would switch the order to: \textit{yeol yel}. This structure does not need inflection or agreement between adjectives and nouns, for example "the wise people" would simply be \textit{yele yeol}. In cases requiring emphasis or clarification, a copular construction can be used instead: \textit{i-la-yel afa yeol} (“The person is wise”). Many commonly used descriptive terms—such as colors, emotional states, or physical traits—exist as standalone nouns and follow the same syntactic pattern. This noun-only system reinforces the conceptual unity of properties and entities, treating attributes not as modifiers but as qualities with independent substance.


\section{Agreement}

Person is marked by a prefix on every auxiliary in formal writing and complex sentences. Stating a subject noun or pronoun does not remove this requirement. Lexical predicates take no person suffix.

\begin{center}
\begin{tabular}{c|c|l}
\textbf{Person} & \textbf{Prefix} & \textbf{Example} \\\hline
1st & i- & \textit{isa nalh} (I/we think) \\
2nd & o- & \textit{osa nalh} (you think) \\
3rd & a- & \textit{asa nalh} (he/she/it/they think) \\
\end{tabular}
\end{center}

Number is expressed on nouns, not by auxiliary agreement. A plural subject such as ``the truths'' therefore uses the same third-person \textit{a-} as a singular subject.

In short, unambiguous casual speech only, third-person \textit{a-} may be omitted: \textit{fa muril} can stand for \textit{afa muril}. Retain person marking in formal writing and complex sentences, including every auxiliary in a chain. First- and second-person prefixes remain required.

Bare \textit{lha} means ``hardness/strength'' and can modify a preceding noun. The formal perfect auxiliary is \textit{alha}; the passive perfect is \textit{apalha}:
\begin{itemize}
\item \textit{i-le-yeol-e lha apafa nalh}: ``The hard truths are pondered.''
\item \textit{i-le-yeol-e apalha nalh}: ``The truths have been pondered.''
\end{itemize}
The person prefix identifies the auxiliary while the lexical root stays unchanged.

\section{Auxiliary Nesting and Scope}

An auxiliary may take the following auxiliary phrase as its complement. Each earlier auxiliary applies to the entire phrase after it. Read a chain from the inside outward. Brackets here show grouping; they add no spoken or written Tovian morphemes.

Within one shared-subject chain, every auxiliary carries the same person prefix. Each layer may contribute its own mood and aspect. Within a single auxiliary, mood takes scope over aspect. Case and class prefixes remain on noun phrases in their existing order.

\subsection*{Scope contrasts}

\begin{itemize}
\item \textit{apa muril}: stops seeing.
\item \textit{ashopa muril}: can [stop seeing].
\item \textit{apa [ashofa muril]}: stops [being able to see].
\item \textit{ashoya muril}: can [begin seeing].
\item \textit{aya [ashofa muril]}: begins [being able to see].
\item \textit{atlopa for}: should [stop writing].
\item \textit{apa [atlofa for]}: ceases [to be expected to write].
\item \textit{ipa [ishofa muril]}: I/we cease [being able to see].
\end{itemize}

Agreement repeats because the grammatical subject is shared. The moods and aspects need not match.

\subsection*{Can no longer be ignored}

The intended scope is cessation over possibility over a passive predicate. With \texttt{[IGNORE]} as an English placeholder for a lexical root still to be chosen:

\begin{quote}
\texttt{apa [ashopafa [IGNORE]]}
\end{quote}

\begin{itemize}
\item \textit{a-p-a}: third person + cessative + present.
\item \textit{a-sho-pa-f-a}: third person + potential + passive + simple + present.
\item \texttt{[IGNORE]}: the unchanged lexical predicate.
\end{itemize}

The truths remain the grammatical subject throughout. The inner phrase means ``can be ignored''; the outer cessative ends that possibility. ``Stop being ignored'' describes the end of actual ignoring, while ``can stop being ignored'' expresses the possibility of stopping. The chain uses one cessative; the bound \textit{pa-} inside \textit{ashopafa} marks passive voice.

``No longer'' presupposes an earlier possibility without implying that the impossibility lasts forever. This construction adds neither a new particle nor an untensed auxiliary form.

\subsection*{Chains and embedded clauses}

``I cease [being able to see]'' has one shared subject. ``I know [you left yesterday]'' contains a clause with its own subject and time. Matching person is required within one shared-subject chain, not across an entire sentence. The temporal frame \textit{tu-le-ngothr: P; Q} specifies one construction with a finite event description P (see Temporal Clauses). General complement clauses such as ``I know [you left]'' still need their own rule.

\subsection*{Temporal dependence: open question}

The examples here use present tense throughout. Matching person does not require Tovian to copy tense endings across every layer. A general rule for inner temporal reference and mixed-tense chains remains to be defined. These examples establish person agreement and scope without settling temporal dependence.


%  Verb stems are added to nouns and always end in ``lore". Different tenses, aspects, moods, and voices can be applied by adding prefixes. The prefixes in the order they should be applied are as follows:\footnote{\noindent Complex verb forms are saved for formal communication.}

% \subsection{Simple Verb Forms}

% \begin{tabular}{rl}
% &\textbf{Tense}\\
% \textbf{} &$\rightarrow$$ present tense \\
% \textbf{vi-} &$\rightarrow$$ future tense \\
% \textbf{in-} &$\rightarrow$$ past tense \\

% &\textbf{Aspect}\\
% \textbf{de-}& $\rightarrow$ perfect \\
% \textbf{me-} & $\rightarrow$ participle \\
% \textbf{shi-}& $\rightarrow$ continuous \\

% &\textbf{Mood}\\
% \textbf{ko-} & $\rightarrow$ conditional \\
% \textbf{se-}& $\rightarrow$ imperative \\
% \textbf{wi-}& $\rightarrow$ subjunctive \\

% &\textbf{Voices}\\
% \textbf{te-}& $\rightarrow$ reflexive \\
% \textbf{pe-}& $\rightarrow$ passive \\

% &\textbf{Other (not combined)}\\
% \textbf{li-}& $\rightarrow$ gerund \\


% \end{tabular}
% % 55 combinations


% \subsection*{All verb forms}
% \begin{tabular}{rl}
% &\textbf{Tense}\\
% \textbf{} &$\rightarrow$$ present tense \\
% \textbf{vi-} &$\rightarrow$$ future tense \\
% \textbf{in-} &$\rightarrow$$ past tense \\
% &\textbf{Aspect}\\
% \textbf{de-}& $\rightarrow$ perfect \\
% \textbf{me-} & $\rightarrow$ participle \\
% \textbf{re-} & $\rightarrow$ habitual \\
% \textbf{ro-}& $\rightarrow$ iterative \\
% \textbf{ve-} & $\rightarrow$ immediate \\
% \textbf{shi-}& $\rightarrow$ continuous \\
% \textbf{wo-}& $\rightarrow$ remote \\
% \textbf{to-}& $\rightarrow$ near \\
% \textbf{yi-}& $\rightarrow$ inceptive \\
% \textbf{hlo-}& $\rightarrow$ cessative \\

% &\textbf{Mood}\\
% \textbf{ko-} & $\rightarrow$ conditional \\
% \textbf{se-}& $\rightarrow$ imperative \\
% \textbf{wi-}& $\rightarrow$ subjunctive \\
% \textbf{xo-}& $\rightarrow$ counterfactual \\

% &\textbf{Voices}\\
% \textbf{te-}& $\rightarrow$ reflexive \\
% \textbf{pa-}& $\rightarrow$ passive \\
% \textbf{mo-}& $\rightarrow$ middle \\
% \textbf{ke-}& $\rightarrow$ causative \\
% \textbf{ra-}& $\rightarrow$ reciprocal \\

% &\textbf{Other (not combined)}\\
% \textbf{li-}& $\rightarrow$ gerund \\


% \end{tabular}
% % 601 combinations

% \medskip

% \subsection*{Verb Conjugations}
% Using "speak" as an example: \textit{fethr}

% \subsubsection*{Tense}
% \begin{itemize}
%     \item \textbf{Present Tense:} \textit{} \\
%           \hanging Speak: \textit{lore-fethr}

%     \item \textbf{Future Tense:} \textit{vi-} \\
%           \hanging Will speak: \textit{vilore-fethr}

%     \item \textbf{Past Tense:} \textit{in-} \\
%           \hanging Spoke: \textit{inlore-fethr}
% \end{itemize}

% \subsubsection*{Aspect}
% \begin{itemize}
%     \item \textbf{Perfect:} \textit{de-} \\
%           \hanging Has spoken: \textit{delore-fethr}

%     \item \textbf{Continuous:} \textit{shi-} \\
%           \hanging Is speaking: \textit{shilore-fethr}

%     \item \textbf{Habitual:} \textit{re-} \\
%           \hanging Speaks regularly: \textit{relore-fethr}

%     \item \textbf{Immediate:} \textit{ve-} \\
%           \hanging Is about to speak: \textit{velore-fethr}

%     \item \textbf{Participle:} \textit{ma-} \\
%           \hanging Speaking (as a participle): \textit{malore-fethr}
% \end{itemize}

% \subsubsection*{Mood}
% \begin{itemize}
%     \item \textbf{Conditional:} \textit{ko-} \\
%           \hanging Would speak: \textit{kolore-fethr}

%     \item \textbf{Imperative:} \textit{se-} \\
%           \hanging Speak!: \textit{selore-fethr}

%     \item \textbf{Subjunctive:} \textit{wi-} \\
%           \hanging May speak: \textit{wilore-fethr}

%     \item \textbf{Counterfactual:} \textit{xo-} \\
%           \hanging Would have spoken: \textit{xolore-fethr}
% \end{itemize}

% \subsubsection*{Voice}
% \begin{itemize}
%     \item \textbf{Reflexive:} \textit{te-} \\
%           \hanging Speaks to oneself: \textit{telore-fethr}

%     \item \textbf{Passive:} \textit{pa-} \\
%           \hanging Is spoken: \textit{palore-fethr}

%     \item \textbf{Middle:} \textit{mo-} \\
%           \hanging Speaks for oneself: \textit{molore-fethr}

%     \item \textbf{Causative:} \textit{ke-} \\
%           \hanging Makes someone speak: \textit{kelore-fethr}

%     \item \textbf{Reciprocal:} \textit{ra-} \\
%           \hanging They speak to each other: \textit{ralore-fethr}
% \end{itemize}

% \subsubsection*{Other}
% \begin{itemize}
%     \item \textbf{Gerund:} \textit{li-} \\
%           \hanging Speaking (as a noun): \textit{lilore-fethr}
% \end{itemize}

% \subsubsection*{Combined Forms}
% \begin{itemize}
%     \item \textbf{Past Perfect:} \textit{de- + in-} \\
%           \hanging Had spoken: \textit{deinlore-fethr}

%     \item \textbf{Future Continuous:} \textit{shi- + vi-} \\
%           \hanging Will be speaking: \textit{shivilore-fethr}

%     \item \textbf{Conditional Habitual:} \textit{re- + ko-} \\
%           \hanging Would speak regularly: \textit{rekolore-fethr}

%     \item \textbf{Immediate Imperative:} \textit{ve- + se-} \\
%           \hanging Speak immediately!: \textit{veselore-fethr}

%     \item \textbf{Subjunctive Participle:} \textit{wi- + ma-} \\
%           \hanging Speaking in a hypothetical or wished-for sense: \textit{wimalore-fethr}
% \end{itemize}


% \subsection*{To be and to be like}
% All verbs with the exception of two follow the above structure. To express that the subject of the sentence is equivalent to the object, use ``'as" as a postfix on the object. To express similarity, use ``'asl" instead. For example:

% \begin{itemize}
%   \item \textbf{To be:} \\
%         \textit{Tel \CAacc-\CLin-helor'as.} (The lake is glass.) \\
%         - \textit{Tel} (subject: the lake) \\
%         - \textit{\CAacc-\CLin-helor'as} (object: glass with ``as" indicating "is")

%  \item \textbf{To be like:} \\
%         \textit{Tel \CAacc-\CLin-helor'asl.} (The lake is like glass.) \\
%         - \textit{Tel} (subject: the lake) \\
%         - \textit{helor'as} (object: glass with ``asl'" indicating "is like")

% \end{itemize}





\section*{Definiteness, Classes, and Cases}

Speakers usually on context to convey definiteness or indefiniteness.
In situations where there is ambiguity or in very formal communication, there are different prefixes:

\begin{table}[H]
    \centering
    \begin{tabular}{ccc}
         &  Subject& Non-subject\\
         Definite&  \textbf{i}-& \textbf{a}-\\
         Indefinite&  \textbf{e}-& \textbf{o}-\\
    \end{tabular}
\end{table}

{
\hspace{-30pt}
\begin{tabular}{l l}
\hline
    The sun shines in the sky  &   \textit{\textbf{i}-\CLin-efol afa ana \textbf{a}\CAloc-rithil}\\
    A sun shines in the sky   &  \textit{\textbf{e}-\CLin-efol afa ana \textbf{a}\CAloc-rithil}\\

        The sun shines in a sky  &   \textit{\textbf{i}-\CLin-efol afa ana \textbf{o}\CAloc-rithil}\\
    A sun shines in a sky   &  \textit{\textbf{e}-\CLin-efol afa ana \textbf{o}\CAloc-rithil}\\
\hline

\end{tabular}
}



\medskip

Nouns are categorized into three distinct \textbf{classes}, each identified by a unique prefix: \textit{\CLan-} for animate beings, \textit{\CLin-} for inanimate objects and places, and \textit{\CLab-} for abstract concepts and ideas. Nouns can also be inflected for grammatical \textbf{cases} to indicate their role within a sentence. Classes are not always used when the meaning is obvious, for example ``the person sees" would be written formally as \textit{\CLan-driel afa muril}, but in informal communication, the class might be dropped to create \textit{driel afa muril}, since the subject is obviously not a dead person (\textit{\CLin-driel}) or the abstract concept of people (\textit{\CLab-drie}l). When both cases and classes are required, cases are added before classes. For example \textit{\CLin-nume \textbf{\CApur}-\textbf{\CLab}-aray }is ``the paths for wisdom", where \textbf{\CApur} represents ``for", \textbf{\CLab} represents ``abstract concept of", and \textbf{ara} is ``wisdom." In formal or educational contexts, hyphens (romanization: ``-", Tovian: ``\vc{\fontsize{50pt}{-200pt}\tovian\textasciitilde-\textasciitilde}") are used to separate prefixes from their nouns. For example, in this document, a sentence could be written as \textit{\CAnom-\CLan-driel \CAacc-\CLin-tel afa muril} (the person sees the lake), but in practice one might use \textit{driel \CAacc\CLin tel afa muril}, dropping the hyphens, the nomative case, and the animate class.



\input{classes_and_cases}

\subsection*{Word Order}
Sentences historically used a Subject-Object-Verb (SOV) word order. In this structure, the subject of the sentence comes first, followed by the object, and then the verb. For instance, "The person lives in the city" would be structured as "\CLan-driel \CAloc-\CLin-pelor afa mel", where "\CLan-driel" (the person) is the subject, "\CAloc-\CLin-pelor" (the city) is the direct (locative) object, and "afa mel" (lives) is the verb.


For emphasis, modern Tovian allows flexibility in word order. Elements can be fronted to the beginning of the sentence to add emphasis or focus. For example, to emphasize the city being lived in by the person, one might say "\CAloc-\CLin-pelor \CLan-driel afa mel," placing the object "\CAloc-\CLin-pelor" (the city) at the beginning of the sentence. This doesn't introduce ambiguity only because of the locative marker (\CAloc-) on the object of the sentence, which indicates it is never the subject despite coming first in word-order.

When changing word order, it is important to leave adjectives (remember, they are unmarked nouns) directly after the noun they modify. For example, "the person lives in the big city" would be either "\CLan-driel \CAloc-\CLin-pelor mot afa mel" or "\CAloc-\CLin-pelor mot \CLan-driel afa mel." Separating "mot" (big) and "pelor" (city) changes the meaning of the sentence. For example: "\CAloc-\CLin-pelor \CLan-driel mot afa mel" would be "the big person lives in the city."


\subsection*{Questions}
\subsubsection{Yes/No questions}
Add \textit{lhan} before a declarative clause. Person agreement remains on the auxiliary:

\begin{itemize}
\item \textit{lhan yu-la-yel ofa muril?}: ``Do you see the person?''
\item \textit{yu-la-yel ifa muril}: ``I see the person.''
\end{itemize}

A short answer may omit a recoverable subject and object, leaving \textit{ifa muril}, ``I see.'' First- and second-person agreement stays explicit. Only third-person \textit{a-} may be omitted in short, unambiguous casual speech; formal writing and complex sentences retain person prefixes on every auxiliary.

\subsubsection{WH-questions}
A case prefix on \textit{lhan} specifies the question:
\begin{itemize}
\item \textit{ti-lhan ta afe fethr?}: ``Where did he/she speak?''
\item \textit{tu-lhan ta afe fethr?}: ``When did he/she speak?''
\item \textit{si-lhan ta afe fethr?}: ``How did he/she speak?''
\end{itemize}

\newpage

% \section{Conjunctions}

% \subsection{Coordinating}
% \begin{itemize}
%     \item{\makebox[2cm][l]{And}  \_\_\_\_\_\_\_\_}
%     \item{\makebox[2cm][l]{But}  \_\_\_\_\_\_\_\_\_}
%     \item{\makebox[2cm][l]{So} \_\_\_\_\_\_\_\_}
% \end{itemize}

% \subsection{Subordinating}
% \begin{itemize}
%     \item{\makebox[2cm][l]{Because} \_\_\_\_\_\_\_\_}
%     \item{\makebox[2cm][l]{If} \_\_\_\_\_\_\_\_}
%     \item{\makebox[2cm][l]{Although} \_\_\_\_\_\_\_\_}
%     \item{\makebox[2cm][l]{When} \_\_\_\_\_\_\_\_}
%     \item{\makebox[2cm][l]{That} \_\_\_\_\_\_\_\_}
% \end{itemize}

% \subsection{Correlative}
% \begin{itemize}
%     \item{\makebox[2cm][l]{Either ... or} \_\_\_\_\_\_\_\_}
%     \item{\makebox[2cm][l]{Neither ... nor} \_\_\_\_\_\_\_\_}
% \end{itemize}



% \include{vocab}

% \include{poem}

% \include{syntax_test}

% \include{vocab_continued}

% \include{inventory}

\input{history}
\newpage

\section{Phrases}
good morning

wishe-lhano yi-aray




\newpage
\section{Dictionary}
\\twocolumn
\\vspace{15pt}
\\begin{nopagebreak}
\\noindent{\\fontsize{20pt}{10pt}\\textbf{keluloresyaleth} } \\textit{1 of autumn/September 1} (N)\\\\
\\noindent {\\tovian \\fontsize{20pt}{10pt} \\textbf{keluloresyaleth} }\\\\
\\noindent /keluloresy{\textprimstress}ale{\texttheta}/\\\\
\\noindent Syllables: kel.ul.or.esj.al.eθ\\\\
\\noindent lit. branch+grain\\\\
\\noindent \\textit{Starrye}\\\\


\\noindent History:

\\vspace{-0pt}
\\hspace{40pt}
\\begin{tabular}{ccc}
\\textit{6000} & /kelumlloresyale{\dh}a/&$\\rightarrow$ & \\textit{6500} & /kelumloresyale{\dh}a/&$\\rightarrow$ & \\textit{7600} & /keluloresyale{\dh}a/&$\\rightarrow$ & \\textit{9500} & /keluloresyale{\dh}/&$\\rightarrow$ & \\textit{12002} & /keluloresyale{\texttheta}/& \\\\
\\end{tabular}

\\vspace{20pt}\\hline

\\end{nopagebreak}
\\filbreak



\\vspace{15pt}
\\begin{nopagebreak}
\\noindent{\\fontsize{20pt}{10pt}\\textbf{yisklumlor} } \\textit{1 of spring/March 1} (N)\\\\
\\noindent {\\tovian \\fontsize{20pt}{10pt} \\textbf{yisklumlor} }\\\\
\\noindent /yiskl{\textprimstress}umlor/\\\\
\\noindent Syllables: jiskl.uml.or\\\\
\\noindent lit. moon+vine\\\\
\\noindent \\textit{Lunevigne}\\\\


\\noindent History:

\\vspace{-0pt}
\\hspace{40pt}
\\begin{tabular}{ccc}
\\textit{6000} & /aliyskluwmllore/&$\\rightarrow$ & \\textit{6300} & /alyskluwmllore/&$\\rightarrow$ & \\textit{6500} & /alyskluwmlore/&$\\rightarrow$ & \\textit{7500} & /lyskluwmlore/&$\\rightarrow$ & \\textit{7501} & /yskluwmlore/&$\\rightarrow$ & \\textit{7600} & /ysklumlore/&$\\rightarrow$ & \\textit{9500} & /ysklumlor/&$\\rightarrow$ & \\textit{15001} & /yisklumlor/& \\\\
\\end{tabular}

\\vspace{20pt}\\hline

\\end{nopagebreak}
\\filbreak



\\vspace{15pt}
\\begin{nopagebreak}
\\noindent{\\fontsize{20pt}{10pt}\\textbf{muloanam} } \\textit{1 of summer/June 1} (N)\\\\
\\noindent {\\tovian \\fontsize{20pt}{10pt} \\textbf{muloanam} }\\\\
\\noindent /mulo{\textprimstress}anam/\\\\
\\noindent Syllables: mul.o.an.am\\\\
\\noindent lit. morning+bloom\\\\
\\noindent \\textit{Dawnbloom}\\\\


\\noindent History:

\\vspace{-0pt}
\\hspace{40pt}
\\begin{tabular}{ccc}
\\textit{6000} & /ploanadrmpfede/&$\\rightarrow$ & \\textit{7600} & /ploanadmpfede/&$\\rightarrow$ & \\textit{8500} & /ploanadmfede/&$\\rightarrow$ & \\textit{9500} & /ploanadmfed/&$\\rightarrow$ & \\textit{11000} & /mloanadmfed/&$\\rightarrow$ & \\textit{12003} & /mloanadmd/&$\\rightarrow$ & \\textit{13001} & /mloanadm/&$\\rightarrow$ & \\textit{13004} & /mloanam/&$\\rightarrow$ & \\textit{15001} & /muloanam/& \\\\
\\end{tabular}

\\vspace{20pt}\\hline

\\end{nopagebreak}
\\filbreak



\\vspace{15pt}
\\begin{nopagebreak}
\\noindent{\\fontsize{20pt}{10pt}\\textbf{lerithdrem} } \\textit{1 of winter/december 1} (N)\\\\
\\noindent {\\tovian \\fontsize{20pt}{10pt} \\textbf{lerithdrem} }\\\\
\\noindent /ler{\textprimstress}i{\texttheta}drem/\\\\
\\noindent Syllables: ler.iθdr.em\\\\
\\noindent lit. breath+flower\\\\
\\noindent \\textit{Rithwildrem}\\\\


\\noindent History:

\\vspace{-0pt}
\\hspace{40pt}
\\begin{tabular}{ccc}
\\textit{6000} & /kwlyri{\texttheta}dremp/&$\\rightarrow$ & \\textit{7501} & /wlyri{\texttheta}dremp/&$\\rightarrow$ & \\textit{7600} & /lri{\texttheta}dremp/&$\\rightarrow$ & \\textit{8750} & /lri{\texttheta}drem/&$\\rightarrow$ & \\textit{15001} & /leri{\texttheta}drem/& \\\\
\\end{tabular}

\\vspace{20pt}\\hline

\\end{nopagebreak}
\\filbreak



\\vspace{15pt}
\\begin{nopagebreak}
\\noindent{\\fontsize{20pt}{10pt}\\textbf{nanathuv} } \\textit{10 of autumn/September 10} (N)\\\\
\\noindent {\\tovian \\fontsize{20pt}{10pt} \\textbf{nanathuv} }\\\\
\\noindent /nan{\textprimstress}a{\texttheta}uv/\\\\
\\noindent Syllables: nan.aθ.uv\\\\
\\noindent lit. star+sword\\\\
\\noindent \\textit{Stareblad}\\\\


\\noindent History:

\\vspace{-0pt}
\\hspace{40pt}
\\begin{tabular}{ccc}
\\textit{6000} & /nnaana{\dh}n{\textbeltl}uv/&$\\rightarrow$ & \\textit{6400} & /nnaana{\dh}{\textbeltl}uv/&$\\rightarrow$ & \\textit{6500} & /naana{\dh}{\textbeltl}uv/&$\\rightarrow$ & \\textit{12002} & /naana{\texttheta}{\textbeltl}uv/&$\\rightarrow$ & \\textit{12005} & /nana{\texttheta}{\textbeltl}uv/&$\\rightarrow$ & \\textit{13000} & /nana{\texttheta}uv/& \\\\
\\end{tabular}

\\vspace{20pt}\\hline

\\end{nopagebreak}
\\filbreak



\\vspace{15pt}
\\begin{nopagebreak}
\\noindent{\\fontsize{20pt}{10pt}\\textbf{todatalh} } \\textit{10 of spring/ March 10} (N)\\\\
\\noindent {\\tovian \\fontsize{20pt}{10pt} \\textbf{todatalh} }\\\\
\\noindent /{\textsubbridge{t}}od{\textprimstress}a{\textsubbridge{t}}a{\textbeltl}/\\\\
\\noindent Syllables: tod.at.aɬ\\\\
\\noindent lit. eternal+love\\\\
\\noindent \\textit{Eternavial}\\\\


\\noindent History:

\\vspace{-0pt}
\\hspace{40pt}
\\begin{tabular}{ccc}
\\textit{6000} & /{\textsubbridge{t}}oda{\textsubbridge{t}}a{\textbeltl}/& \\\\
\\end{tabular}

\\vspace{20pt}\\hline

\\end{nopagebreak}
\\filbreak



\\vspace{15pt}
\\begin{nopagebreak}
\\noindent{\\fontsize{20pt}{10pt}\\textbf{yethilemnan} } \\textit{10 of summer/June 10} (N)\\\\
\\noindent {\\tovian \\fontsize{20pt}{10pt} \\textbf{yethilemnan} }\\\\
\\noindent /ye{\texttheta}il{\textprimstress}emnan/\\\\
\\noindent Syllables: jeθ.il.emn.an\\\\
\\noindent lit. home+protection\\\\
\\noindent \\textit{Nestkeep}\\\\


\\noindent History:

\\vspace{-0pt}
\\hspace{40pt}
\\begin{tabular}{ccc}
\\textit{6000} & /dryel{\texttheta}ilebnani/&$\\rightarrow$ & \\textit{7501} & /ryel{\texttheta}ilebnani/&$\\rightarrow$ & \\textit{7600} & /yel{\texttheta}ilebnani/&$\\rightarrow$ & \\textit{9500} & /yel{\texttheta}ilebnan/&$\\rightarrow$ & \\textit{11000} & /yel{\texttheta}ilemnan/&$\\rightarrow$ & \\textit{12502} & /ye{\texttheta}ilemnan/& \\\\
\\end{tabular}

\\vspace{20pt}\\hline

\\end{nopagebreak}
\\filbreak



\\vspace{15pt}
\\begin{nopagebreak}
\\noindent{\\fontsize{20pt}{10pt}\\textbf{leriskwil} } \\textit{10 of winter/december 10} (N)\\\\
\\noindent {\\tovian \\fontsize{20pt}{10pt} \\textbf{leriskwil} }\\\\
\\noindent /ler{\textprimstress}iskwil/\\\\
\\noindent Syllables: ler.iskw.il\\\\
\\noindent lit. breath+movement\\\\
\\noindent \\textit{Galespad}\\\\


\\noindent History:

\\vspace{-0pt}
\\hspace{40pt}
\\begin{tabular}{ccc}
\\textit{6000} & /kwlyri{\texttheta}kwil/&$\\rightarrow$ & \\textit{7501} & /wlyri{\texttheta}kwil/&$\\rightarrow$ & \\textit{7600} & /lri{\texttheta}kwil/&$\\rightarrow$ & \\textit{10500} & /lriskwil/&$\\rightarrow$ & \\textit{15001} & /leriskwil/& \\\\
\\end{tabular}

\\vspace{20pt}\\hline

\\end{nopagebreak}
\\filbreak



\\vspace{15pt}
\\begin{nopagebreak}
\\noindent{\\fontsize{20pt}{10pt}\\textbf{???} } \\textit{11 of autumn/September 11} (N)\\\\
\\noindent {\\tovian \\fontsize{20pt}{10pt} \\textbf{???} }\\\\
\\noindent /???/\\\\
\\noindent Syllables: ???\\\\
\\noindent lit. ???\\\\
\\noindent \\textit{Ciantro}\\\\


\\noindent History:

\\vspace{-0pt}
\\hspace{40pt}
\\begin{tabular}{ccc}
\\textit{6000} & /???/& \\\\
\\end{tabular}

\\vspace{20pt}\\hline

\\end{nopagebreak}
\\filbreak



\\vspace{15pt}
\\begin{nopagebreak}
\\noindent{\\fontsize{20pt}{10pt}\\textbf{efith} } \\textit{11 of spring/ March 11} (N)\\\\
\\noindent {\\tovian \\fontsize{20pt}{10pt} \\textbf{efith} }\\\\
\\noindent /{\textprimstress}efi{\texttheta}/\\\\
\\noindent Syllables: ef.iθ\\\\
\\noindent lit. sun+storm\\\\
\\noindent \\textit{Sunshower}\\\\


\\noindent History:

\\vspace{-0pt}
\\hspace{40pt}
\\begin{tabular}{ccc}
\\textit{6000} & /eflkwlyri{\texttheta}l/&$\\rightarrow$ & \\textit{7600} & /eflklri{\texttheta}l/&$\\rightarrow$ & \\textit{8750} & /efllri{\texttheta}l/&$\\rightarrow$ & \\textit{12007} & /eflri{\texttheta}l/&$\\rightarrow$ & \\textit{12500} & /efi{\texttheta}/& \\\\
\\end{tabular}

\\vspace{20pt}\\hline

\\end{nopagebreak}
\\filbreak



\\vspace{15pt}
\\begin{nopagebreak}
\\noindent{\\fontsize{20pt}{10pt}\\textbf{thuvglar} } \\textit{11 of summer/June 11} (N)\\\\
\\noindent {\\tovian \\fontsize{20pt}{10pt} \\textbf{thuvglar} }\\\\
\\noindent /{\texttheta}{\textprimstress}uvglar/\\\\
\\noindent Syllables: θuvgl.ar\\\\
\\noindent lit. sword+woods\\\\
\\noindent \\textit{Vekhangor}\\\\


\\noindent History:

\\vspace{-0pt}
\\hspace{40pt}
\\begin{tabular}{ccc}
\\textit{6000} & /{\dh}n{\textbeltl}uvglar/&$\\rightarrow$ & \\textit{6400} & /{\dh}{\textbeltl}uvglar/&$\\rightarrow$ & \\textit{12002} & /{\texttheta}{\textbeltl}uvglar/&$\\rightarrow$ & \\textit{13000} & /{\texttheta}uvglar/& \\\\
\\end{tabular}

\\vspace{20pt}\\hline

\\end{nopagebreak}
\\filbreak



\\vspace{15pt}
\\begin{nopagebreak}
\\noindent{\\fontsize{20pt}{10pt}\\textbf{fedlad} } \\textit{11 of winter/december 11} (N)\\\\
\\noindent {\\tovian \\fontsize{20pt}{10pt} \\textbf{fedlad} }\\\\
\\noindent /f{\textprimstress}edlad/\\\\
\\noindent Syllables: fedl.ad\\\\
\\noindent lit. showing+jealousy\\\\
\\noindent \\textit{Narcissus}\\\\


\\noindent History:

\\vspace{-0pt}
\\hspace{40pt}
\\begin{tabular}{ccc}
\\textit{6000} & /fedeylrladl/&$\\rightarrow$ & \\textit{6300} & /fedylrladl/&$\\rightarrow$ & \\textit{6503} & /fedyrlladl/&$\\rightarrow$ & \\textit{7600} & /fedrladl/&$\\rightarrow$ & \\textit{12505} & /fedladl/&$\\rightarrow$ & \\textit{13001} & /fedlad/& \\\\
\\end{tabular}

\\vspace{20pt}\\hline

\\end{nopagebreak}
\\filbreak



\\vspace{15pt}
\\begin{nopagebreak}
\\noindent{\\fontsize{20pt}{10pt}\\textbf{thilethan} } \\textit{12 of autumn/September 12} (N)\\\\
\\noindent {\\tovian \\fontsize{20pt}{10pt} \\textbf{thilethan} }\\\\
\\noindent /{\texttheta}il{\textprimstress}e{\texttheta}an/\\\\
\\noindent Syllables: θil.eθ.an\\\\
\\noindent lit. place+tool\\\\
\\noindent \\textit{Earchoke}\\\\


\\noindent History:

\\vspace{-0pt}
\\hspace{40pt}
\\begin{tabular}{ccc}
\\textit{6000} & /{\texttheta}ile{\dh}an/&$\\rightarrow$ & \\textit{12002} & /{\texttheta}ile{\texttheta}an/& \\\\
\\end{tabular}

\\vspace{20pt}\\hline

\\end{nopagebreak}
\\filbreak



\\vspace{15pt}
\\begin{nopagebreak}
\\noindent{\\fontsize{20pt}{10pt}\\textbf{muloelhaswanalor} } \\textit{12 of spring/ March 12} (N)\\\\
\\noindent {\\tovian \\fontsize{20pt}{10pt} \\textbf{muloelhaswanalor} }\\\\
\\noindent /muloe{\textbeltl}aswan{\textprimstress}alor/\\\\
\\noindent Syllables: mul.o.eɬ.asw.an.al.or\\\\
\\noindent lit. rebirth+plant\\\\
\\noindent \\textit{Féenixflour}\\\\


\\noindent History:

\\vspace{-0pt}
\\hspace{40pt}
\\begin{tabular}{ccc}
\\textit{6000} & /ploe{\textbeltl}aswanamllore/&$\\rightarrow$ & \\textit{6500} & /ploe{\textbeltl}aswanamlore/&$\\rightarrow$ & \\textit{7600} & /ploe{\textbeltl}aswanalore/&$\\rightarrow$ & \\textit{9500} & /ploe{\textbeltl}aswanalor/&$\\rightarrow$ & \\textit{11000} & /mloe{\textbeltl}aswanalor/&$\\rightarrow$ & \\textit{15001} & /muloe{\textbeltl}aswanalor/& \\\\
\\end{tabular}

\\vspace{20pt}\\hline

\\end{nopagebreak}
\\filbreak



\\vspace{15pt}
\\begin{nopagebreak}
\\noindent{\\fontsize{20pt}{10pt}\\textbf{nador} } \\textit{12 of summer/June 12} (N)\\\\
\\noindent {\\tovian \\fontsize{20pt}{10pt} \\textbf{nador} }\\\\
\\noindent /n{\textprimstress}ador/\\\\
\\noindent Syllables: nad.or\\\\
\\noindent lit. light+blue\\\\
\\noindent \\textit{Gleamwoad}\\\\


\\noindent History:

\\vspace{-0pt}
\\hspace{40pt}
\\begin{tabular}{ccc}
\\textit{6000} & /anadora/&$\\rightarrow$ & \\textit{7500} & /nadora/&$\\rightarrow$ & \\textit{9500} & /nador/& \\\\
\\end{tabular}

\\vspace{20pt}\\hline

\\end{nopagebreak}
\\filbreak



\\vspace{15pt}
\\begin{nopagebreak}
\\noindent{\\fontsize{20pt}{10pt}\\textbf{rithaglar} } \\textit{12 of winter/december 12} (N)\\\\
\\noindent {\\tovian \\fontsize{20pt}{10pt} \\textbf{rithaglar} }\\\\
\\noindent /ri{\texttheta}{\textprimstress}aglar/\\\\
\\noindent Syllables: riθ.agl.ar\\\\
\\noindent lit. air+tree\\\\
\\noindent \\textit{Windelm}\\\\


\\noindent History:

\\vspace{-0pt}
\\hspace{40pt}
\\begin{tabular}{ccc}
\\textit{6000} & /ri{\texttheta}aglare/&$\\rightarrow$ & \\textit{9500} & /ri{\texttheta}aglar/& \\\\
\\end{tabular}

\\vspace{20pt}\\hline

\\end{nopagebreak}
\\filbreak



\\vspace{15pt}
\\begin{nopagebreak}
\\noindent{\\fontsize{20pt}{10pt}\\textbf{shithan} } \\textit{13 of autumn/September 13} (N)\\\\
\\noindent {\\tovian \\fontsize{20pt}{10pt} \\textbf{shithan} }\\\\
\\noindent /{\textesh}{\textprimstress}i{\texttheta}an/\\\\
\\noindent Syllables: ʃiθ.an\\\\
\\noindent lit. sharpness+tool\\\\
\\noindent \\textit{Spicesta}\\\\


\\noindent History:

\\vspace{-0pt}
\\hspace{40pt}
\\begin{tabular}{ccc}
\\textit{6000} & /{\textyogh}i{\texttheta}{\dh}an/&$\\rightarrow$ & \\textit{12001} & /{\textesh}i{\texttheta}{\dh}an/&$\\rightarrow$ & \\textit{12002} & /{\textesh}i{\texttheta}{\texttheta}an/&$\\rightarrow$ & \\textit{12007} & /{\textesh}i{\texttheta}an/& \\\\
\\end{tabular}

\\vspace{20pt}\\hline

\\end{nopagebreak}
\\filbreak



\\vspace{15pt}
\\begin{nopagebreak}
\\noindent{\\fontsize{20pt}{10pt}\\textbf{tushudol} } \\textit{13 of spring/ March 13} (N)\\\\
\\noindent {\\tovian \\fontsize{20pt}{10pt} \\textbf{tushudol} }\\\\
\\noindent /{\textsubbridge{t}}u{\textesh}{\textprimstress}udol/\\\\
\\noindent Syllables: tuʃ.ud.ol\\\\
\\noindent lit. cold+fruit\\\\
\\noindent \\textit{Víntergourd}\\\\


\\noindent History:

\\vspace{-0pt}
\\hspace{40pt}
\\begin{tabular}{ccc}
\\textit{6000} & /{\textsubbridge{t}}u{\textesh}u{\textyogh}ola/&$\\rightarrow$ & \\textit{9500} & /{\textsubbridge{t}}u{\textesh}u{\textyogh}ol/&$\\rightarrow$ & \\textit{12001} & /{\textsubbridge{t}}u{\textesh}u{\textesh}ol/&$\\rightarrow$ & \\textit{12501} & /{\textsubbridge{t}}u{\textesh}udol/& \\\\
\\end{tabular}

\\vspace{20pt}\\hline

\\end{nopagebreak}
\\filbreak



\\vspace{15pt}
\\begin{nopagebreak}
\\noindent{\\fontsize{20pt}{10pt}\\textbf{muralsyaleth} } \\textit{13 of summer/June 13} (N)\\\\
\\noindent {\\tovian \\fontsize{20pt}{10pt} \\textbf{muralsyaleth} }\\\\
\\noindent /muralsy{\textprimstress}ale{\texttheta}/\\\\
\\noindent Syllables: mur.alsj.al.eθ\\\\
\\noindent lit. darkness+food\\\\
\\noindent \\textit{Nightmush}\\\\


\\noindent History:

\\vspace{-0pt}
\\hspace{40pt}
\\begin{tabular}{ccc}
\\textit{6000} & /bralsyale{\dh}a/&$\\rightarrow$ & \\textit{9500} & /bralsyale{\dh}/&$\\rightarrow$ & \\textit{11000} & /mralsyale{\dh}/&$\\rightarrow$ & \\textit{12002} & /mralsyale{\texttheta}/&$\\rightarrow$ & \\textit{15001} & /muralsyale{\texttheta}/& \\\\
\\end{tabular}

\\vspace{20pt}\\hline

\\end{nopagebreak}
\\filbreak



\\vspace{15pt}
\\begin{nopagebreak}
\\noindent{\\fontsize{20pt}{10pt}\\textbf{shotlan} } \\textit{13 of winter/december 13} (N)\\\\
\\noindent {\\tovian \\fontsize{20pt}{10pt} \\textbf{shotlan} }\\\\
\\noindent /{\textesh}{\textprimstress}o{\texttoptiebar{t\textbeltl}}an/\\\\
\\noindent Syllables: ʃoƛ.an\\\\
\\noindent lit. flowing+tool\\\\
\\noindent \\textit{Mistifen}\\\\


\\noindent History:

\\vspace{-0pt}
\\hspace{40pt}
\\begin{tabular}{ccc}
\\textit{6000} & /{\textesh}o{\texttoptiebar{t\textbeltl}}{\dh}an/&$\\rightarrow$ & \\textit{12002} & /{\textesh}o{\texttoptiebar{t\textbeltl}}{\texttheta}an/&$\\rightarrow$ & \\textit{14000} & /{\textesh}o{\texttoptiebar{t\textbeltl}}an/& \\\\
\\end{tabular}

\\vspace{20pt}\\hline

\\end{nopagebreak}
\\filbreak



\\vspace{15pt}
\\begin{nopagebreak}
\\noindent{\\fontsize{20pt}{10pt}\\textbf{tlurinin} } \\textit{14 of autumn/September 14} (N)\\\\
\\noindent {\\tovian \\fontsize{20pt}{10pt} \\textbf{tlurinin} }\\\\
\\noindent /{\texttoptiebar{t\textbeltl}}ur{\textprimstress}inin/\\\\
\\noindent Syllables: ƛur.in.in\\\\
\\noindent lit. mist+???\\\\
\\noindent \\textit{Mistlaven}\\\\


\\noindent History:

\\vspace{-0pt}
\\hspace{40pt}
\\begin{tabular}{ccc}
\\textit{6000} & /i{\ng}g{\texttoptiebar{t\textbeltl}}unri{\texttheta}nina???/&$\\rightarrow$ & \\textit{7500} & /{\ng}g{\texttoptiebar{t\textbeltl}}unri{\texttheta}nina???/&$\\rightarrow$ & \\textit{7501} & /g{\texttoptiebar{t\textbeltl}}unri{\texttheta}nina???/&$\\rightarrow$ & \\textit{7600} & /g{\texttoptiebar{t\textbeltl}}uri{\texttheta}nina???/&$\\rightarrow$ & \\textit{8500} & /g{\texttoptiebar{t\textbeltl}}uri{\texttheta}nina/&$\\rightarrow$ & \\textit{9500} & /g{\texttoptiebar{t\textbeltl}}uri{\texttheta}nin/&$\\rightarrow$ & \\textit{11501} & /{\texttoptiebar{t\textbeltl}}uri{\texttheta}nin/&$\\rightarrow$ & \\textit{13002} & /{\texttoptiebar{t\textbeltl}}urinin/& \\\\
\\end{tabular}

\\vspace{20pt}\\hline

\\end{nopagebreak}
\\filbreak



\\vspace{15pt}
\\begin{nopagebreak}
\\noindent{\\fontsize{20pt}{10pt}\\textbf{menan} } \\textit{14 of spring/ March 14} (N)\\\\
\\noindent {\\tovian \\fontsize{20pt}{10pt} \\textbf{menan} }\\\\
\\noindent /m{\textprimstress}enan/\\\\
\\noindent Syllables: men.an\\\\
\\noindent lit. silver+???\\\\
\\noindent \\textit{Silvreséda}\\\\


\\noindent History:

\\vspace{-0pt}
\\hspace{40pt}
\\begin{tabular}{ccc}
\\textit{6000} & /{\dh}nhmernana???/&$\\rightarrow$ & \\textit{6400} & /{\dh}hmernana???/&$\\rightarrow$ & \\textit{7501} & /hmernana???/&$\\rightarrow$ & \\textit{7600} & /hmenana???/&$\\rightarrow$ & \\textit{8500} & /hmenana/&$\\rightarrow$ & \\textit{9500} & /hmenan/&$\\rightarrow$ & \\textit{13002} & /menan/& \\\\
\\end{tabular}

\\vspace{20pt}\\hline

\\end{nopagebreak}
\\filbreak



\\vspace{15pt}
\\begin{nopagebreak}
\\noindent{\\fontsize{20pt}{10pt}\\textbf{leriskwil} } \\textit{14 of summer/June 14} (N)\\\\
\\noindent {\\tovian \\fontsize{20pt}{10pt} \\textbf{leriskwil} }\\\\
\\noindent /ler{\textprimstress}iskwil/\\\\
\\noindent Syllables: ler.iskw.il\\\\
\\noindent lit. wind+movement\\\\
\\noindent \\textit{Rithwil}\\\\


\\noindent History:

\\vspace{-0pt}
\\hspace{40pt}
\\begin{tabular}{ccc}
\\textit{6000} & /kwlyri{\texttheta}kwil/&$\\rightarrow$ & \\textit{7501} & /wlyri{\texttheta}kwil/&$\\rightarrow$ & \\textit{7600} & /lri{\texttheta}kwil/&$\\rightarrow$ & \\textit{10500} & /lriskwil/&$\\rightarrow$ & \\textit{15001} & /leriskwil/& \\\\
\\end{tabular}

\\vspace{20pt}\\hline

\\end{nopagebreak}
\\filbreak



\\vspace{15pt}
\\begin{nopagebreak}
\\noindent{\\fontsize{20pt}{10pt}\\textbf{fedetush} } \\textit{14 of winter/december 14} (N)\\\\
\\noindent {\\tovian \\fontsize{20pt}{10pt} \\textbf{fedetush} }\\\\
\\noindent /fed{\textprimstress}e{\textsubbridge{t}}u{\textesh}/\\\\
\\noindent Syllables: fed.et.uʃ\\\\
\\noindent lit. showing+cold\\\\
\\noindent \\textit{Galblom}\\\\


\\noindent History:

\\vspace{-0pt}
\\hspace{40pt}
\\begin{tabular}{ccc}
\\textit{6000} & /fede{\textsubbridge{t}}u{\textesh}u/&$\\rightarrow$ & \\textit{9500} & /fede{\textsubbridge{t}}u{\textesh}/& \\\\
\\end{tabular}

\\vspace{20pt}\\hline

\\end{nopagebreak}
\\filbreak



\\vspace{15pt}
\\begin{nopagebreak}
\\noindent{\\fontsize{20pt}{10pt}\\textbf{???} } \\textit{15 of autumn/September 15} (N)\\\\
\\noindent {\\tovian \\fontsize{20pt}{10pt} \\textbf{???} }\\\\
\\noindent /???/\\\\
\\noindent Syllables: ???\\\\
\\noindent lit. ???\\\\
\\noindent \\textit{Chamois}\\\\


\\noindent History:

\\vspace{-0pt}
\\hspace{40pt}
\\begin{tabular}{ccc}
\\textit{6000} & /???/& \\\\
\\end{tabular}

\\vspace{20pt}\\hline

\\end{nopagebreak}
\\filbreak



\\vspace{15pt}
\\begin{nopagebreak}
\\noindent{\\fontsize{20pt}{10pt}\\textbf{rithilekel} } \\textit{15 of spring/ March 15} (N)\\\\
\\noindent {\\tovian \\fontsize{20pt}{10pt} \\textbf{rithilekel} }\\\\
\\noindent /ri{\texttheta}il{\textprimstress}ekel/\\\\
\\noindent Syllables: riθ.il.ek.el\\\\
\\noindent lit. sky+branch\\\\
\\noindent \\textit{Skyaspe}\\\\


\\noindent History:

\\vspace{-0pt}
\\hspace{40pt}
\\begin{tabular}{ccc}
\\textit{6000} & /ri{\texttheta}ilekelu/&$\\rightarrow$ & \\textit{9500} & /ri{\texttheta}ilekel/& \\\\
\\end{tabular}

\\vspace{20pt}\\hline

\\end{nopagebreak}
\\filbreak



\\vspace{15pt}
\\begin{nopagebreak}
\\noindent{\\fontsize{20pt}{10pt}\\textbf{forengothr} } \\textit{15 of summer/June 15} (N)\\\\
\\noindent {\\tovian \\fontsize{20pt}{10pt} \\textbf{forengothr} }\\\\
\\noindent /for{\textprimstress}e{\ng}o{\texttheta}r/\\\\
\\noindent Syllables: for.eŋ.oθr\\\\
\\noindent lit. writing+time\\\\
\\noindent \\textit{Felor Day}\\\\


\\noindent History:

\\vspace{-0pt}
\\hspace{40pt}
\\begin{tabular}{ccc}
\\textit{6000} & /flore{\ng}go{\texttheta}r/&$\\rightarrow$ & \\textit{8750} & /flore{\ng}o{\texttheta}r/&$\\rightarrow$ & \\textit{12500} & /fore{\ng}o{\texttheta}r/& \\\\
\\end{tabular}

\\vspace{20pt}\\hline

\\end{nopagebreak}
\\filbreak



\\vspace{15pt}
\\begin{nopagebreak}
\\noindent{\\fontsize{20pt}{10pt}\\textbf{lerithen} } \\textit{15 of winter/december 15} (N)\\\\
\\noindent {\\tovian \\fontsize{20pt}{10pt} \\textbf{lerithen} }\\\\
\\noindent /ler{\textprimstress}i{\texttheta}en/\\\\
\\noindent Syllables: ler.iθ.en\\\\
\\noindent lit. breath+dream\\\\
\\noindent \\textit{Windrith}\\\\


\\noindent History:

\\vspace{-0pt}
\\hspace{40pt}
\\begin{tabular}{ccc}
\\textit{6000} & /kwlyri{\texttheta}{\textbeltl}ena/&$\\rightarrow$ & \\textit{7501} & /wlyri{\texttheta}{\textbeltl}ena/&$\\rightarrow$ & \\textit{7600} & /lri{\texttheta}{\textbeltl}ena/&$\\rightarrow$ & \\textit{9500} & /lri{\texttheta}{\textbeltl}en/&$\\rightarrow$ & \\textit{13000} & /lri{\texttheta}en/&$\\rightarrow$ & \\textit{15001} & /leri{\texttheta}en/& \\\\
\\end{tabular}

\\vspace{20pt}\\hline

\\end{nopagebreak}
\\filbreak



\\vspace{15pt}
\\begin{nopagebreak}
\\noindent{\\fontsize{20pt}{10pt}\\textbf{rithol} } \\textit{16 of autumn/September 16} (N)\\\\
\\noindent {\\tovian \\fontsize{20pt}{10pt} \\textbf{rithol} }\\\\
\\noindent /r{\textprimstress}i{\texttheta}ol/\\\\
\\noindent Syllables: riθ.ol\\\\
\\noindent lit. air+sweet\\\\
\\noindent \\textit{Smokesweet}\\\\


\\noindent History:

\\vspace{-0pt}
\\hspace{40pt}
\\begin{tabular}{ccc}
\\textit{6000} & /ri{\texttheta}hola/&$\\rightarrow$ & \\textit{9500} & /ri{\texttheta}hol/&$\\rightarrow$ & \\textit{13000} & /ri{\texttheta}ol/& \\\\
\\end{tabular}

\\vspace{20pt}\\hline

\\end{nopagebreak}
\\filbreak



\\vspace{15pt}
\\begin{nopagebreak}
\\noindent{\\fontsize{20pt}{10pt}\\textbf{murawen} } \\textit{16 of spring/ March 16} (N)\\\\
\\noindent {\\tovian \\fontsize{20pt}{10pt} \\textbf{murawen} }\\\\
\\noindent /mur{\textprimstress}awen/\\\\
\\noindent Syllables: mur.aw.en\\\\
\\noindent lit. darkness+beauty\\\\
\\noindent \\textit{Nightbelle}\\\\


\\noindent History:

\\vspace{-0pt}
\\hspace{40pt}
\\begin{tabular}{ccc}
\\textit{6000} & /bralwen/&$\\rightarrow$ & \\textit{7600} & /brawen/&$\\rightarrow$ & \\textit{11000} & /mrawen/&$\\rightarrow$ & \\textit{15001} & /murawen/& \\\\
\\end{tabular}

\\vspace{20pt}\\hline

\\end{nopagebreak}
\\filbreak



\\vspace{15pt}
\\begin{nopagebreak}
\\noindent{\\fontsize{20pt}{10pt}\\textbf{semlor} } \\textit{16 of summer/June 16} (N)\\\\
\\noindent {\\tovian \\fontsize{20pt}{10pt} \\textbf{semlor} }\\\\
\\noindent /s{\textprimstress}emlor/\\\\
\\noindent Syllables: seml.or\\\\
\\noindent lit. water+plant\\\\
\\noindent \\textit{Dewlettuce}\\\\


\\noindent History:

\\vspace{-0pt}
\\hspace{40pt}
\\begin{tabular}{ccc}
\\textit{6000} & /zelmllore/&$\\rightarrow$ & \\textit{6500} & /zelmlore/&$\\rightarrow$ & \\textit{7600} & /zemlore/&$\\rightarrow$ & \\textit{9500} & /zemlor/&$\\rightarrow$ & \\textit{12000} & /semlor/& \\\\
\\end{tabular}

\\vspace{20pt}\\hline

\\end{nopagebreak}
\\filbreak



\\vspace{15pt}
\\begin{nopagebreak}
\\noindent{\\fontsize{20pt}{10pt}\\textbf{loreayeol} } \\textit{16 of winter/december 16} (N)\\\\
\\noindent {\\tovian \\fontsize{20pt}{10pt} \\textbf{loreayeol} }\\\\
\\noindent /loreay{\textprimstress}eol/\\\\
\\noindent Syllables: lor.e.aj.e.ol\\\\
\\noindent lit. dance+wisdom\\\\
\\noindent \\textit{Dansinath}\\\\


\\noindent History:

\\vspace{-0pt}
\\hspace{40pt}
\\begin{tabular}{ccc}
\\textit{6000} & /kwllorearyeol/&$\\rightarrow$ & \\textit{6500} & /kwlorearyeol/&$\\rightarrow$ & \\textit{7501} & /wlorearyeol/&$\\rightarrow$ & \\textit{7600} & /loreayeol/& \\\\
\\end{tabular}

\\vspace{20pt}\\hline

\\end{nopagebreak}
\\filbreak



\\vspace{15pt}
\\begin{nopagebreak}
\\noindent{\\fontsize{20pt}{10pt}\\textbf{sholalhum} } \\textit{17 of autumn/September 17} (N)\\\\
\\noindent {\\tovian \\fontsize{20pt}{10pt} \\textbf{sholalhum} }\\\\
\\noindent /{\textesh}ol{\textprimstress}a{\textbeltl}um/\\\\
\\noindent Syllables: ʃol.aɬ.um\\\\
\\noindent lit. fruit+goodness\\\\
\\noindent \\textit{Curranthum}\\\\


\\noindent History:

\\vspace{-0pt}
\\hspace{40pt}
\\begin{tabular}{ccc}
\\textit{6000} & /{\textyogh}ola{\textbeltl}um/&$\\rightarrow$ & \\textit{12001} & /{\textesh}ola{\textbeltl}um/& \\\\
\\end{tabular}

\\vspace{20pt}\\hline

\\end{nopagebreak}
\\filbreak



\\vspace{15pt}
\\begin{nopagebreak}
\\noindent{\\fontsize{20pt}{10pt}\\textbf{tlafor} } \\textit{17 of spring/ March 17} (N)\\\\
\\noindent {\\tovian \\fontsize{20pt}{10pt} \\textbf{tlafor} }\\\\
\\noindent /{\texttoptiebar{t\textbeltl}}{\textprimstress}afor/\\\\
\\noindent Syllables: ƛaf.or\\\\
\\noindent lit. peace+plant\\\\
\\noindent \\textit{Tlafen}\\\\


\\noindent History:

\\vspace{-0pt}
\\hspace{40pt}
\\begin{tabular}{ccc}
\\textit{6000} & /{\texttoptiebar{t\textbeltl}}afmllore/&$\\rightarrow$ & \\textit{6500} & /{\texttoptiebar{t\textbeltl}}afmlore/&$\\rightarrow$ & \\textit{7600} & /{\texttoptiebar{t\textbeltl}}aflore/&$\\rightarrow$ & \\textit{9500} & /{\texttoptiebar{t\textbeltl}}aflor/&$\\rightarrow$ & \\textit{12500} & /{\texttoptiebar{t\textbeltl}}afor/& \\\\
\\end{tabular}

\\vspace{20pt}\\hline

\\end{nopagebreak}
\\filbreak



\\vspace{15pt}
\\begin{nopagebreak}
\\noindent{\\fontsize{20pt}{10pt}\\textbf{gaglar} } \\textit{17 of summer/June 17} (N)\\\\
\\noindent {\\tovian \\fontsize{20pt}{10pt} \\textbf{gaglar} }\\\\
\\noindent /g{\textprimstress}aglar/\\\\
\\noindent Syllables: gagl.ar\\\\
\\noindent lit. warmth+tree\\\\
\\noindent \\textit{Frostlarch}\\\\


\\noindent History:

\\vspace{-0pt}
\\hspace{40pt}
\\begin{tabular}{ccc}
\\textit{6000} & /agaaglare/&$\\rightarrow$ & \\textit{7500} & /gaaglare/&$\\rightarrow$ & \\textit{9500} & /gaaglar/&$\\rightarrow$ & \\textit{12005} & /gaglar/& \\\\
\\end{tabular}

\\vspace{20pt}\\hline

\\end{nopagebreak}
\\filbreak



\\vspace{15pt}
\\begin{nopagebreak}
\\noindent{\\fontsize{20pt}{10pt}\\textbf{lanam} } \\textit{17 of winter/december 17} (N)\\\\
\\noindent {\\tovian \\fontsize{20pt}{10pt} \\textbf{lanam} }\\\\
\\noindent /l{\textprimstress}anam/\\\\
\\noindent Syllables: lan.am\\\\
\\noindent lit. twilight+bloom\\\\
\\noindent \\textit{Twilibloom}\\\\


\\noindent History:

\\vspace{-0pt}
\\hspace{40pt}
\\begin{tabular}{ccc}
\\textit{6000} & /brlanadrmpfede/&$\\rightarrow$ & \\textit{7501} & /rlanadrmpfede/&$\\rightarrow$ & \\textit{7600} & /lanadmpfede/&$\\rightarrow$ & \\textit{8500} & /lanadmfede/&$\\rightarrow$ & \\textit{9500} & /lanadmfed/&$\\rightarrow$ & \\textit{12003} & /lanadmd/&$\\rightarrow$ & \\textit{13001} & /lanadm/&$\\rightarrow$ & \\textit{13004} & /lanam/& \\\\
\\end{tabular}

\\vspace{20pt}\\hline

\\end{nopagebreak}
\\filbreak



\\vspace{15pt}
\\begin{nopagebreak}
\\noindent{\\fontsize{20pt}{10pt}\\textbf{sholaglar} } \\textit{18 of autumn/September 18} (N)\\\\
\\noindent {\\tovian \\fontsize{20pt}{10pt} \\textbf{sholaglar} }\\\\
\\noindent /{\textesh}ol{\textprimstress}aglar/\\\\
\\noindent Syllables: ʃol.agl.ar\\\\
\\noindent lit. fruit+tree\\\\
\\noindent \\textit{Peawil}\\\\


\\noindent History:

\\vspace{-0pt}
\\hspace{40pt}
\\begin{tabular}{ccc}
\\textit{6000} & /{\textyogh}olaaglare/&$\\rightarrow$ & \\textit{9500} & /{\textyogh}olaaglar/&$\\rightarrow$ & \\textit{12001} & /{\textesh}olaaglar/&$\\rightarrow$ & \\textit{12005} & /{\textesh}olaglar/& \\\\
\\end{tabular}

\\vspace{20pt}\\hline

\\end{nopagebreak}
\\filbreak



\\vspace{15pt}
\\begin{nopagebreak}
\\noindent{\\fontsize{20pt}{10pt}\\textbf{nanamel} } \\textit{18 of spring/ March 18} (N)\\\\
\\noindent {\\tovian \\fontsize{20pt}{10pt} \\textbf{nanamel} }\\\\
\\noindent /nan{\textprimstress}amel/\\\\
\\noindent Syllables: nan.am.el\\\\
\\noindent lit. star+life\\\\
\\noindent \\textit{Starzin}\\\\


\\noindent History:

\\vspace{-0pt}
\\hspace{40pt}
\\begin{tabular}{ccc}
\\textit{6000} & /nnaanamel/&$\\rightarrow$ & \\textit{6500} & /naanamel/&$\\rightarrow$ & \\textit{12005} & /nanamel/& \\\\
\\end{tabular}

\\vspace{20pt}\\hline

\\end{nopagebreak}
\\filbreak



\\vspace{15pt}
\\begin{nopagebreak}
\\noindent{\\fontsize{20pt}{10pt}\\textbf{murallhenadrem} } \\textit{18 of summer/June 18} (N)\\\\
\\noindent {\\tovian \\fontsize{20pt}{10pt} \\textbf{murallhenadrem} }\\\\
\\noindent /mural{\textbeltl}en{\textprimstress}adrem/\\\\
\\noindent Syllables: mur.alɬ.en.adr.em\\\\
\\noindent lit. darkness+dream+flower\\\\
\\noindent \\textit{Darkbloom}\\\\


\\noindent History:

\\vspace{-0pt}
\\hspace{40pt}
\\begin{tabular}{ccc}
\\textit{6000} & /bral{\textbeltl}enadremp/&$\\rightarrow$ & \\textit{8750} & /bral{\textbeltl}enadrem/&$\\rightarrow$ & \\textit{11000} & /mral{\textbeltl}enadrem/&$\\rightarrow$ & \\textit{15001} & /mural{\textbeltl}enadrem/& \\\\
\\end{tabular}

\\vspace{20pt}\\hline

\\end{nopagebreak}
\\filbreak



\\vspace{15pt}
\\begin{nopagebreak}
\\noindent{\\fontsize{20pt}{10pt}\\textbf{nanalor} } \\textit{18 of winter/december 18} (N)\\\\
\\noindent {\\tovian \\fontsize{20pt}{10pt} \\textbf{nanalor} }\\\\
\\noindent /nan{\textprimstress}alor/\\\\
\\noindent Syllables: nan.al.or\\\\
\\noindent lit. star+plant\\\\
\\noindent \\textit{Starsweed}\\\\


\\noindent History:

\\vspace{-0pt}
\\hspace{40pt}
\\begin{tabular}{ccc}
\\textit{6000} & /nnaanamllore/&$\\rightarrow$ & \\textit{6500} & /naanamlore/&$\\rightarrow$ & \\textit{7600} & /naanalore/&$\\rightarrow$ & \\textit{9500} & /naanalor/&$\\rightarrow$ & \\textit{12005} & /nanalor/& \\\\
\\end{tabular}

\\vspace{20pt}\\hline

\\end{nopagebreak}
\\filbreak



\\vspace{15pt}
\\begin{nopagebreak}
\\noindent{\\fontsize{20pt}{10pt}\\textbf{leshanin} } \\textit{19 of autumn/September 19} (N)\\\\
\\noindent {\\tovian \\fontsize{20pt}{10pt} \\textbf{leshanin} }\\\\
\\noindent /le{\textesh}{\textprimstress}anin/\\\\
\\noindent Syllables: leʃ.an.in\\\\
\\noindent lit. sun+berry\\\\
\\noindent \\textit{Sunberry}\\\\


\\noindent History:

\\vspace{-0pt}
\\hspace{40pt}
\\begin{tabular}{ccc}
\\textit{6000} & /efl{\textyogh}lanina/&$\\rightarrow$ & \\textit{7500} & /fl{\textyogh}lanina/&$\\rightarrow$ & \\textit{7501} & /l{\textyogh}lanina/&$\\rightarrow$ & \\textit{9500} & /l{\textyogh}lanin/&$\\rightarrow$ & \\textit{12001} & /l{\textesh}lanin/&$\\rightarrow$ & \\textit{12500} & /l{\textesh}anin/&$\\rightarrow$ & \\textit{15001} & /le{\textesh}anin/& \\\\
\\end{tabular}

\\vspace{20pt}\\hline

\\end{nopagebreak}
\\filbreak



\\vspace{15pt}
\\begin{nopagebreak}
\\noindent{\\fontsize{20pt}{10pt}\\textbf{efem} } \\textit{19 of spring/ March 19} (N)\\\\
\\noindent {\\tovian \\fontsize{20pt}{10pt} \\textbf{efem} }\\\\
\\noindent /{\textprimstress}efem/\\\\
\\noindent Syllables: ef.em\\\\
\\noindent lit. sun+flower\\\\
\\noindent \\textit{Soldrem}\\\\


\\noindent History:

\\vspace{-0pt}
\\hspace{40pt}
\\begin{tabular}{ccc}
\\textit{6000} & /efldremp/&$\\rightarrow$ & \\textit{8750} & /eflrem/&$\\rightarrow$ & \\textit{12500} & /efem/& \\\\
\\end{tabular}

\\vspace{20pt}\\hline

\\end{nopagebreak}
\\filbreak



\\vspace{15pt}
\\begin{nopagebreak}
\\noindent{\\fontsize{20pt}{10pt}\\textbf{tomomen} } \\textit{19 of summer/June 19} (N)\\\\
\\noindent {\\tovian \\fontsize{20pt}{10pt} \\textbf{tomomen} }\\\\
\\noindent /{\textsubbridge{t}}om{\textprimstress}omen/\\\\
\\noindent Syllables: tom.om.en\\\\
\\noindent lit. root+earth\\\\
\\noindent \\textit{Starroot}\\\\


\\noindent History:

\\vspace{-0pt}
\\hspace{40pt}
\\begin{tabular}{ccc}
\\textit{6000} & /{\textsubbridge{t}}omomern/&$\\rightarrow$ & \\textit{7600} & /{\textsubbridge{t}}omomen/& \\\\
\\end{tabular}

\\vspace{20pt}\\hline

\\end{nopagebreak}
\\filbreak



\\vspace{15pt}
\\begin{nopagebreak}
\\noindent{\\fontsize{20pt}{10pt}\\textbf{ayswen} } \\textit{19 of winter/december 19} (N)\\\\
\\noindent {\\tovian \\fontsize{20pt}{10pt} \\textbf{ayswen} }\\\\
\\noindent /{\textprimstress}ayswen/\\\\
\\noindent Syllables: ajsw.en\\\\
\\noindent lit. moon+beauty\\\\
\\noindent \\textit{Mooncheris}\\\\


\\noindent History:

\\vspace{-0pt}
\\hspace{40pt}
\\begin{tabular}{ccc}
\\textit{6000} & /aliyswen/&$\\rightarrow$ & \\textit{6300} & /alyswen/&$\\rightarrow$ & \\textit{7600} & /ayswen/& \\\\
\\end{tabular}

\\vspace{20pt}\\hline

\\end{nopagebreak}
\\filbreak



\\vspace{15pt}
\\begin{nopagebreak}
\\noindent{\\fontsize{20pt}{10pt}\\textbf{yisoresyaleth} } \\textit{2 of autumn/September 2} (N)\\\\
\\noindent {\\tovian \\fontsize{20pt}{10pt} \\textbf{yisoresyaleth} }\\\\
\\noindent /yisoresy{\textprimstress}ale{\texttheta}/\\\\
\\noindent Syllables: jis.or.esj.al.eθ\\\\
\\noindent lit. moon+grain\\\\
\\noindent \\textit{Moonoat}\\\\


\\noindent History:

\\vspace{-0pt}
\\hspace{40pt}
\\begin{tabular}{ccc}
\\textit{6000} & /aliysmlloresyale{\dh}a/&$\\rightarrow$ & \\textit{6300} & /alysmlloresyale{\dh}a/&$\\rightarrow$ & \\textit{6500} & /alysmloresyale{\dh}a/&$\\rightarrow$ & \\textit{7500} & /lysmloresyale{\dh}a/&$\\rightarrow$ & \\textit{7501} & /ysmloresyale{\dh}a/&$\\rightarrow$ & \\textit{7600} & /ysloresyale{\dh}a/&$\\rightarrow$ & \\textit{9500} & /ysloresyale{\dh}/&$\\rightarrow$ & \\textit{12002} & /ysloresyale{\texttheta}/&$\\rightarrow$ & \\textit{12500} & /ysoresyale{\texttheta}/&$\\rightarrow$ & \\textit{15001} & /yisoresyale{\texttheta}/& \\\\
\\end{tabular}

\\vspace{20pt}\\hline

\\end{nopagebreak}
\\filbreak



\\vspace{15pt}
\\begin{nopagebreak}
\\noindent{\\fontsize{20pt}{10pt}\\textbf{efith} } \\textit{2 of spring/March 2} (N)\\\\
\\noindent {\\tovian \\fontsize{20pt}{10pt} \\textbf{efith} }\\\\
\\noindent /{\textprimstress}efi{\texttheta}/\\\\
\\noindent Syllables: ef.iθ\\\\
\\noindent lit. sun+sharpness\\\\
\\noindent \\textit{Solspica}\\\\


\\noindent History:

\\vspace{-0pt}
\\hspace{40pt}
\\begin{tabular}{ccc}
\\textit{6000} & /efl{\textyogh}i{\texttheta}/&$\\rightarrow$ & \\textit{12001} & /efl{\textesh}i{\texttheta}/&$\\rightarrow$ & \\textit{12500} & /ef{\textesh}i{\texttheta}/&$\\rightarrow$ & \\textit{13000} & /efi{\texttheta}/& \\\\
\\end{tabular}

\\vspace{20pt}\\hline

\\end{nopagebreak}
\\filbreak



\\vspace{15pt}
\\begin{nopagebreak}
\\noindent{\\fontsize{20pt}{10pt}\\textbf{rithileaglar} } \\textit{2 of summer/June 2} (N)\\\\
\\noindent {\\tovian \\fontsize{20pt}{10pt} \\textbf{rithileaglar} }\\\\
\\noindent /ri{\texttheta}ile{\textprimstress}aglar/\\\\
\\noindent Syllables: riθ.il.e.agl.ar\\\\
\\noindent lit. sky+tree\\\\
\\noindent \\textit{Skybark}\\\\


\\noindent History:

\\vspace{-0pt}
\\hspace{40pt}
\\begin{tabular}{ccc}
\\textit{6000} & /ri{\texttheta}ileaglare/&$\\rightarrow$ & \\textit{9500} & /ri{\texttheta}ileaglar/& \\\\
\\end{tabular}

\\vspace{20pt}\\hline

\\end{nopagebreak}
\\filbreak



\\vspace{15pt}
\\begin{nopagebreak}
\\noindent{\\fontsize{20pt}{10pt}\\textbf{dolshanin} } \\textit{2 of winter/december 2} (N)\\\\
\\noindent {\\tovian \\fontsize{20pt}{10pt} \\textbf{dolshanin} }\\\\
\\noindent /dol{\textesh}{\textprimstress}anin/\\\\
\\noindent Syllables: dolʃ.an.in\\\\
\\noindent lit. death+berry\\\\
\\noindent \\textit{Galbery}\\\\


\\noindent History:

\\vspace{-0pt}
\\hspace{40pt}
\\begin{tabular}{ccc}
\\textit{6000} & /dol{\textyogh}lanina/&$\\rightarrow$ & \\textit{9500} & /dol{\textyogh}lanin/&$\\rightarrow$ & \\textit{12001} & /dol{\textesh}lanin/&$\\rightarrow$ & \\textit{12500} & /dol{\textesh}anin/& \\\\
\\end{tabular}

\\vspace{20pt}\\hline

\\end{nopagebreak}
\\filbreak



\\vspace{15pt}
\\begin{nopagebreak}
\\noindent{\\fontsize{20pt}{10pt}\\textbf{tilethil} } \\textit{20 of autumn/September 20} (N)\\\\
\\noindent {\\tovian \\fontsize{20pt}{10pt} \\textbf{tilethil} }\\\\
\\noindent /{\textsubbridge{t}}il{\textprimstress}e{\texttheta}il/\\\\
\\noindent Syllables: til.eθ.il\\\\
\\noindent lit. plain+place\\\\
\\noindent \\textit{Meadowrealm}\\\\


\\noindent History:

\\vspace{-0pt}
\\hspace{40pt}
\\begin{tabular}{ccc}
\\textit{6000} & /{\textsubbridge{t}}n{\texttheta}ile{\texttheta}ile/&$\\rightarrow$ & \\textit{6400} & /{\textsubbridge{t}}{\texttheta}ile{\texttheta}ile/&$\\rightarrow$ & \\textit{9500} & /{\textsubbridge{t}}{\texttheta}ile{\texttheta}il/&$\\rightarrow$ & \\textit{10700} & /{\textsubbridge{t}}{\textsubbridge{t}}ile{\texttheta}il/&$\\rightarrow$ & \\textit{12004} & /{\textsubbridge{t}}ile{\texttheta}il/& \\\\
\\end{tabular}

\\vspace{20pt}\\hline

\\end{nopagebreak}
\\filbreak



\\vspace{15pt}
\\begin{nopagebreak}
\\noindent{\\fontsize{20pt}{10pt}\\textbf{nanashith} } \\textit{20 of spring/March 20} (N)\\\\
\\noindent {\\tovian \\fontsize{20pt}{10pt} \\textbf{nanashith} }\\\\
\\noindent /nan{\textprimstress}a{\textesh}i{\texttheta}/\\\\
\\noindent Syllables: nan.aʃ.iθ\\\\
\\noindent lit. star+sharpness\\\\
\\noindent \\textit{Starpress}\\\\


\\noindent History:

\\vspace{-0pt}
\\hspace{40pt}
\\begin{tabular}{ccc}
\\textit{6000} & /nnaana{\textyogh}i{\texttheta}/&$\\rightarrow$ & \\textit{6500} & /naana{\textyogh}i{\texttheta}/&$\\rightarrow$ & \\textit{12001} & /naana{\textesh}i{\texttheta}/&$\\rightarrow$ & \\textit{12005} & /nana{\textesh}i{\texttheta}/& \\\\
\\end{tabular}

\\vspace{20pt}\\hline

\\end{nopagebreak}
\\filbreak



\\vspace{15pt}
\\begin{nopagebreak}
\\noindent{\\fontsize{20pt}{10pt}\\textbf{holadyethil} } \\textit{20 of summer/June 20} (N)\\\\
\\noindent {\\tovian \\fontsize{20pt}{10pt} \\textbf{holadyethil} }\\\\
\\noindent /ħolady{\textprimstress}e{\texttheta}il/\\\\
\\noindent Syllables: ħol.adj.eθ.il\\\\
\\noindent lit. sweet+home\\\\
\\noindent \\textit{Honeyhall}\\\\


\\noindent History:

\\vspace{-0pt}
\\hspace{40pt}
\\begin{tabular}{ccc}
\\textit{6000} & /holadryel{\texttheta}ile/&$\\rightarrow$ & \\textit{7600} & /holadyel{\texttheta}ile/&$\\rightarrow$ & \\textit{9500} & /holadyel{\texttheta}il/&$\\rightarrow$ & \\textit{12502} & /holadye{\texttheta}il/&$\\rightarrow$ & \\textit{14006} & /ħoladye{\texttheta}il/& \\\\
\\end{tabular}

\\vspace{20pt}\\hline

\\end{nopagebreak}
\\filbreak



\\vspace{15pt}
\\begin{nopagebreak}
\\noindent{\\fontsize{20pt}{10pt}\\textbf{lerithun} } \\textit{20 of winter/december 20} (N)\\\\
\\noindent {\\tovian \\fontsize{20pt}{10pt} \\textbf{lerithun} }\\\\
\\noindent /ler{\textprimstress}i{\texttheta}un/\\\\
\\noindent Syllables: ler.iθ.un\\\\
\\noindent lit. breath+joining\\\\
\\noindent \\textit{Windcord}\\\\


\\noindent History:

\\vspace{-0pt}
\\hspace{40pt}
\\begin{tabular}{ccc}
\\textit{6000} & /kwlyri{\texttheta}{\texttoptiebar{t\textbeltl}}una/&$\\rightarrow$ & \\textit{7501} & /wlyri{\texttheta}{\texttoptiebar{t\textbeltl}}una/&$\\rightarrow$ & \\textit{7600} & /lri{\texttheta}{\texttoptiebar{t\textbeltl}}una/&$\\rightarrow$ & \\textit{9500} & /lri{\texttheta}{\texttoptiebar{t\textbeltl}}un/&$\\rightarrow$ & \\textit{14001} & /lri{\texttheta}un/&$\\rightarrow$ & \\textit{15001} & /leri{\texttheta}un/& \\\\
\\end{tabular}

\\vspace{20pt}\\hline

\\end{nopagebreak}
\\filbreak



\\vspace{15pt}
\\begin{nopagebreak}
\\noindent{\\fontsize{20pt}{10pt}\\textbf{shush} } \\textit{20 of winter/January 20} (N)\\\\
\\noindent {\\tovian \\fontsize{20pt}{10pt} \\textbf{shush} }\\\\
\\noindent /{\textesh}{\textprimstress}u{\textesh}/\\\\
\\noindent Syllables: ʃuʃ\\\\
\\noindent lit. ice+???\\\\
\\noindent \\textit{Icetole}\\\\


\\noindent History:

\\vspace{-0pt}
\\hspace{40pt}
\\begin{tabular}{ccc}
\\textit{6000} & /{\textesh}u{\textesh}a???/&$\\rightarrow$ & \\textit{8500} & /{\textesh}u{\textesh}a/&$\\rightarrow$ & \\textit{9500} & /{\textesh}u{\textesh}/& \\\\
\\end{tabular}

\\vspace{20pt}\\hline

\\end{nopagebreak}
\\filbreak



\\vspace{15pt}
\\begin{nopagebreak}
\\noindent{\\fontsize{20pt}{10pt}\\textbf{riskel} } \\textit{21 of autumn/September 21} (N)\\\\
\\noindent {\\tovian \\fontsize{20pt}{10pt} \\textbf{riskel} }\\\\
\\noindent /r{\textprimstress}iskel/\\\\
\\noindent Syllables: risk.el\\\\
\\noindent lit. air+branch\\\\
\\noindent \\textit{Breezemint}\\\\


\\noindent History:

\\vspace{-0pt}
\\hspace{40pt}
\\begin{tabular}{ccc}
\\textit{6000} & /ri{\texttheta}kelu/&$\\rightarrow$ & \\textit{9500} & /ri{\texttheta}kel/&$\\rightarrow$ & \\textit{10500} & /riskel/& \\\\
\\end{tabular}

\\vspace{20pt}\\hline

\\end{nopagebreak}
\\filbreak



\\vspace{15pt}
\\begin{nopagebreak}
\\noindent{\\fontsize{20pt}{10pt}\\textbf{deremlor} } \\textit{21 of spring/March 21} (N)\\\\
\\noindent {\\tovian \\fontsize{20pt}{10pt} \\textbf{deremlor} }\\\\
\\noindent /der{\textprimstress}emlor/\\\\
\\noindent Syllables: der.eml.or\\\\
\\noindent lit. flower+plant\\\\
\\noindent \\textit{Dremhemp}\\\\


\\noindent History:

\\vspace{-0pt}
\\hspace{40pt}
\\begin{tabular}{ccc}
\\textit{6000} & /drempmllore/&$\\rightarrow$ & \\textit{6500} & /drempmlore/&$\\rightarrow$ & \\textit{7600} & /dremplore/&$\\rightarrow$ & \\textit{8500} & /dremlore/&$\\rightarrow$ & \\textit{9500} & /dremlor/&$\\rightarrow$ & \\textit{15001} & /deremlor/& \\\\
\\end{tabular}

\\vspace{20pt}\\hline

\\end{nopagebreak}
\\filbreak



\\vspace{15pt}
\\begin{nopagebreak}
\\noindent{\\fontsize{20pt}{10pt}\\textbf{meloremnan} } \\textit{21 of summer/June 21} (N)\\\\
\\noindent {\\tovian \\fontsize{20pt}{10pt} \\textbf{meloremnan} }\\\\
\\noindent /melor{\textprimstress}emnan/\\\\
\\noindent Syllables: mel.or.emn.an\\\\
\\noindent lit. plant+protection\\\\
\\noindent \\textit{Seedguard}\\\\


\\noindent History:

\\vspace{-0pt}
\\hspace{40pt}
\\begin{tabular}{ccc}
\\textit{6000} & /mllorebnani/&$\\rightarrow$ & \\textit{6500} & /mlorebnani/&$\\rightarrow$ & \\textit{7502} & /melorebnani/&$\\rightarrow$ & \\textit{9500} & /melorebnan/&$\\rightarrow$ & \\textit{11000} & /meloremnan/& \\\\
\\end{tabular}

\\vspace{20pt}\\hline

\\end{nopagebreak}
\\filbreak



\\vspace{15pt}
\\begin{nopagebreak}
\\noindent{\\fontsize{20pt}{10pt}\\textbf{lhilaglar} } \\textit{21 of winter/december 21} (N)\\\\
\\noindent {\\tovian \\fontsize{20pt}{10pt} \\textbf{lhilaglar} }\\\\
\\noindent /{\textbeltl}il{\textprimstress}aglar/\\\\
\\noindent Syllables: ɬil.agl.ar\\\\
\\noindent lit. night+tree\\\\
\\noindent \\textit{Nightroot}\\\\


\\noindent History:

\\vspace{-0pt}
\\hspace{40pt}
\\begin{tabular}{ccc}
\\textit{6000} & /a{\textbeltl}inlaaglare/&$\\rightarrow$ & \\textit{7500} & /{\textbeltl}inlaaglare/&$\\rightarrow$ & \\textit{7600} & /{\textbeltl}ilaaglare/&$\\rightarrow$ & \\textit{9500} & /{\textbeltl}ilaaglar/&$\\rightarrow$ & \\textit{12005} & /{\textbeltl}ilaglar/& \\\\
\\end{tabular}

\\vspace{20pt}\\hline

\\end{nopagebreak}
\\filbreak



\\vspace{15pt}
\\begin{nopagebreak}
\\noindent{\\fontsize{20pt}{10pt}\\textbf{muraltlun} } \\textit{22 of autumn/September 22} (N)\\\\
\\noindent {\\tovian \\fontsize{20pt}{10pt} \\textbf{muraltlun} }\\\\
\\noindent /mur{\textprimstress}al{\texttoptiebar{t\textbeltl}}un/\\\\
\\noindent Syllables: mur.alƛ.un\\\\
\\noindent lit. darkness+joining\\\\
\\noindent \\textit{Nightspice}\\\\


\\noindent History:

\\vspace{-0pt}
\\hspace{40pt}
\\begin{tabular}{ccc}
\\textit{6000} & /bral{\texttoptiebar{t\textbeltl}}una/&$\\rightarrow$ & \\textit{9500} & /bral{\texttoptiebar{t\textbeltl}}un/&$\\rightarrow$ & \\textit{11000} & /mral{\texttoptiebar{t\textbeltl}}un/&$\\rightarrow$ & \\textit{15001} & /mural{\texttoptiebar{t\textbeltl}}un/& \\\\
\\end{tabular}

\\vspace{20pt}\\hline

\\end{nopagebreak}
\\filbreak



\\vspace{15pt}
\\begin{nopagebreak}
\\noindent{\\fontsize{20pt}{10pt}\\textbf{rithiledrem} } \\textit{22 of spring/March 22} (N)\\\\
\\noindent {\\tovian \\fontsize{20pt}{10pt} \\textbf{rithiledrem} }\\\\
\\noindent /ri{\texttheta}il{\textprimstress}edrem/\\\\
\\noindent Syllables: riθ.il.edr.em\\\\
\\noindent lit. sky+flower\\\\
\\noindent \\textit{Skypech}\\\\


\\noindent History:

\\vspace{-0pt}
\\hspace{40pt}
\\begin{tabular}{ccc}
\\textit{6000} & /ri{\texttheta}iledremp/&$\\rightarrow$ & \\textit{8750} & /ri{\texttheta}iledrem/& \\\\
\\end{tabular}

\\vspace{20pt}\\hline

\\end{nopagebreak}
\\filbreak



\\vspace{15pt}
\\begin{nopagebreak}
\\noindent{\\fontsize{20pt}{10pt}\\textbf{shotlathang} } \\textit{22 of summer/June 22} (N)\\\\
\\noindent {\\tovian \\fontsize{20pt}{10pt} \\textbf{shotlathang} }\\\\
\\noindent /{\textesh}o{\texttoptiebar{t\textbeltl}}{\textprimstress}a{\texttheta}a{\ng}/\\\\
\\noindent Syllables: ʃoƛ.aθ.aŋ\\\\
\\noindent lit. flowing+leaf\\\\
\\noindent \\textit{Mistleaf}\\\\


\\noindent History:

\\vspace{-0pt}
\\hspace{40pt}
\\begin{tabular}{ccc}
\\textit{6000} & /{\textesh}o{\texttoptiebar{t\textbeltl}}a{\dh}a{\ng}ge/&$\\rightarrow$ & \\textit{8750} & /{\textesh}o{\texttoptiebar{t\textbeltl}}a{\dh}a{\ng}e/&$\\rightarrow$ & \\textit{9500} & /{\textesh}o{\texttoptiebar{t\textbeltl}}a{\dh}a{\ng}/&$\\rightarrow$ & \\textit{12002} & /{\textesh}o{\texttoptiebar{t\textbeltl}}a{\texttheta}a{\ng}/& \\\\
\\end{tabular}

\\vspace{20pt}\\hline

\\end{nopagebreak}
\\filbreak



\\vspace{15pt}
\\begin{nopagebreak}
\\noindent{\\fontsize{20pt}{10pt}\\textbf{lan} } \\textit{22 of winter/december 22} (N)\\\\
\\noindent {\\tovian \\fontsize{20pt}{10pt} \\textbf{lan} }\\\\
\\noindent /l{\textprimstress}an/\\\\
\\noindent Syllables: lan\\\\
\\noindent lit. twilight+???\\\\
\\noindent \\textit{Duspars}\\\\


\\noindent History:

\\vspace{-0pt}
\\hspace{40pt}
\\begin{tabular}{ccc}
\\textit{6000} & /brlana???/&$\\rightarrow$ & \\textit{7501} & /rlana???/&$\\rightarrow$ & \\textit{7600} & /lana???/&$\\rightarrow$ & \\textit{8500} & /lana/&$\\rightarrow$ & \\textit{9500} & /lan/& \\\\
\\end{tabular}

\\vspace{20pt}\\hline

\\end{nopagebreak}
\\filbreak



\\vspace{15pt}
\\begin{nopagebreak}
\\noindent{\\fontsize{20pt}{10pt}\\textbf{kelushol} } \\textit{23 of autumn/September 23} (N)\\\\
\\noindent {\\tovian \\fontsize{20pt}{10pt} \\textbf{kelushol} }\\\\
\\noindent /kel{\textprimstress}u{\textesh}ol/\\\\
\\noindent Syllables: kel.uʃ.ol\\\\
\\noindent lit. branch+fruit\\\\
\\noindent \\textit{Starbean}\\\\


\\noindent History:

\\vspace{-0pt}
\\hspace{40pt}
\\begin{tabular}{ccc}
\\textit{6000} & /kelu{\textyogh}ola/&$\\rightarrow$ & \\textit{9500} & /kelu{\textyogh}ol/&$\\rightarrow$ & \\textit{12001} & /kelu{\textesh}ol/& \\\\
\\end{tabular}

\\vspace{20pt}\\hline

\\end{nopagebreak}
\\filbreak



\\vspace{15pt}
\\begin{nopagebreak}
\\noindent{\\fontsize{20pt}{10pt}\\textbf{shasith} } \\textit{23 of spring/March 23} (N)\\\\
\\noindent {\\tovian \\fontsize{20pt}{10pt} \\textbf{shasith} }\\\\
\\noindent /{\textesh}{\textprimstress}asi{\texttheta}/\\\\
\\noindent Syllables: ʃas.iθ\\\\
\\noindent lit. ice+ship\\\\
\\noindent \\textit{Frostnave}\\\\


\\noindent History:

\\vspace{-0pt}
\\hspace{40pt}
\\begin{tabular}{ccc}
\\textit{6000} & /{\textesh}u{\textesh}azli{\texttheta}a/&$\\rightarrow$ & \\textit{9200} & /{\textesh}{\textschwa}{\textesh}azli{\texttheta}a/&$\\rightarrow$ & \\textit{9500} & /{\textesh}{\textschwa}{\textesh}azli{\texttheta}/&$\\rightarrow$ & \\textit{11990} & /{\textesh}{\textesh}azli{\texttheta}/&$\\rightarrow$ & \\textit{12000} & /{\textesh}{\textesh}asli{\texttheta}/&$\\rightarrow$ & \\textit{12007} & /{\textesh}asli{\texttheta}/&$\\rightarrow$ & \\textit{12500} & /{\textesh}asi{\texttheta}/& \\\\
\\end{tabular}

\\vspace{20pt}\\hline

\\end{nopagebreak}
\\filbreak



\\vspace{15pt}
\\begin{nopagebreak}
\\noindent{\\fontsize{20pt}{10pt}\\textbf{melaglar} } \\textit{23 of summer/June 23} (N)\\\\
\\noindent {\\tovian \\fontsize{20pt}{10pt} \\textbf{melaglar} }\\\\
\\noindent /mel{\textprimstress}aglar/\\\\
\\noindent Syllables: mel.agl.ar\\\\
\\noindent lit. life+tree\\\\
\\noindent \\textit{Heartree}\\\\


\\noindent History:

\\vspace{-0pt}
\\hspace{40pt}
\\begin{tabular}{ccc}
\\textit{6000} & /melaglare/&$\\rightarrow$ & \\textit{9500} & /melaglar/& \\\\
\\end{tabular}

\\vspace{20pt}\\hline

\\end{nopagebreak}
\\filbreak



\\vspace{15pt}
\\begin{nopagebreak}
\\noindent{\\fontsize{20pt}{10pt}\\textbf{lanasyaleth} } \\textit{23 of winter/december 23} (N)\\\\
\\noindent {\\tovian \\fontsize{20pt}{10pt} \\textbf{lanasyaleth} }\\\\
\\noindent /lanasy{\textprimstress}ale{\texttheta}/\\\\
\\noindent Syllables: lan.asj.al.eθ\\\\
\\noindent lit. twilight+food\\\\
\\noindent \\textit{Twilispoon}\\\\


\\noindent History:

\\vspace{-0pt}
\\hspace{40pt}
\\begin{tabular}{ccc}
\\textit{6000} & /brlanasyale{\dh}a/&$\\rightarrow$ & \\textit{7501} & /rlanasyale{\dh}a/&$\\rightarrow$ & \\textit{7600} & /lanasyale{\dh}a/&$\\rightarrow$ & \\textit{9500} & /lanasyale{\dh}/&$\\rightarrow$ & \\textit{12002} & /lanasyale{\texttheta}/& \\\\
\\end{tabular}

\\vspace{20pt}\\hline

\\end{nopagebreak}
\\filbreak



\\vspace{15pt}
\\begin{nopagebreak}
\\noindent{\\fontsize{20pt}{10pt}\\textbf{muloanamen} } \\textit{24 of autumn/September 24} (N)\\\\
\\noindent {\\tovian \\fontsize{20pt}{10pt} \\textbf{muloanamen} }\\\\
\\noindent /muloan{\textprimstress}amen/\\\\
\\noindent Syllables: mul.o.an.am.en\\\\
\\noindent lit. morning+earth\\\\
\\noindent \\textit{Dawnborage}\\\\


\\noindent History:

\\vspace{-0pt}
\\hspace{40pt}
\\begin{tabular}{ccc}
\\textit{6000} & /ploanamern/&$\\rightarrow$ & \\textit{7600} & /ploanamen/&$\\rightarrow$ & \\textit{11000} & /mloanamen/&$\\rightarrow$ & \\textit{15001} & /muloanamen/& \\\\
\\end{tabular}

\\vspace{20pt}\\hline

\\end{nopagebreak}
\\filbreak



\\vspace{15pt}
\\begin{nopagebreak}
\\noindent{\\fontsize{20pt}{10pt}\\textbf{nanadrem} } \\textit{24 of spring/March 24} (N)\\\\
\\noindent {\\tovian \\fontsize{20pt}{10pt} \\textbf{nanadrem} }\\\\
\\noindent /nan{\textprimstress}adrem/\\\\
\\noindent Syllables: nan.adr.em\\\\
\\noindent lit. star+flower\\\\
\\noindent \\textit{Starlilis}\\\\


\\noindent History:

\\vspace{-0pt}
\\hspace{40pt}
\\begin{tabular}{ccc}
\\textit{6000} & /nnaanadremp/&$\\rightarrow$ & \\textit{6500} & /naanadremp/&$\\rightarrow$ & \\textit{8750} & /naanadrem/&$\\rightarrow$ & \\textit{12005} & /nanadrem/& \\\\
\\end{tabular}

\\vspace{20pt}\\hline

\\end{nopagebreak}
\\filbreak



\\vspace{15pt}
\\begin{nopagebreak}
\\noindent{\\fontsize{20pt}{10pt}\\textbf{tleshor} } \\textit{24 of summer/June 24} (N)\\\\
\\noindent {\\tovian \\fontsize{20pt}{10pt} \\textbf{tleshor} }\\\\
\\noindent /{\texttoptiebar{t\textbeltl}}{\textprimstress}e{\textesh}or/\\\\
\\noindent Syllables: ƛeʃ.or\\\\
\\noindent lit. speed+plant\\\\
\\noindent \\textit{Swiftgreen}\\\\


\\noindent History:

\\vspace{-0pt}
\\hspace{40pt}
\\begin{tabular}{ccc}
\\textit{6000} & /{\texttoptiebar{t\textbeltl}}e{\textyogh}mllore/&$\\rightarrow$ & \\textit{6500} & /{\texttoptiebar{t\textbeltl}}e{\textyogh}mlore/&$\\rightarrow$ & \\textit{7600} & /{\texttoptiebar{t\textbeltl}}e{\textyogh}lore/&$\\rightarrow$ & \\textit{9500} & /{\texttoptiebar{t\textbeltl}}e{\textyogh}lor/&$\\rightarrow$ & \\textit{12001} & /{\texttoptiebar{t\textbeltl}}e{\textesh}lor/&$\\rightarrow$ & \\textit{12500} & /{\texttoptiebar{t\textbeltl}}e{\textesh}or/& \\\\
\\end{tabular}

\\vspace{20pt}\\hline

\\end{nopagebreak}
\\filbreak



\\vspace{15pt}
\\begin{nopagebreak}
\\noindent{\\fontsize{20pt}{10pt}\\textbf{lhonrem} } \\textit{24 of winter/december 24} (N)\\\\
\\noindent {\\tovian \\fontsize{20pt}{10pt} \\textbf{lhonrem} }\\\\
\\noindent /{\textbeltl}{\textprimstress}onrem/\\\\
\\noindent Syllables: ɬonr.em\\\\
\\noindent lit. day+flower\\\\
\\noindent \\textit{Daydaisy}\\\\


\\noindent History:

\\vspace{-0pt}
\\hspace{40pt}
\\begin{tabular}{ccc}
\\textit{6000} & /{\textbeltl}onmdremp/&$\\rightarrow$ & \\textit{6000} & /{\textbeltl}onndremp/&$\\rightarrow$ & \\textit{6500} & /{\textbeltl}ondremp/&$\\rightarrow$ & \\textit{8500} & /{\textbeltl}onremp/&$\\rightarrow$ & \\textit{8750} & /{\textbeltl}onrem/& \\\\
\\end{tabular}

\\vspace{20pt}\\hline

\\end{nopagebreak}
\\filbreak



\\vspace{15pt}
\\begin{nopagebreak}
\\noindent{\\fontsize{20pt}{10pt}\\textbf{???} } \\textit{25 of autumn/September 25} (N)\\\\
\\noindent {\\tovian \\fontsize{20pt}{10pt} \\textbf{???} }\\\\
\\noindent /???/\\\\
\\noindent Syllables: ???\\\\
\\noindent lit. ???\\\\
\\noindent \\textit{Skuinea}\\\\


\\noindent History:

\\vspace{-0pt}
\\hspace{40pt}
\\begin{tabular}{ccc}
\\textit{6000} & /???/& \\\\
\\end{tabular}

\\vspace{20pt}\\hline

\\end{nopagebreak}
\\filbreak



\\vspace{15pt}
\\begin{nopagebreak}
\\noindent{\\fontsize{20pt}{10pt}\\textbf{lerithan} } \\textit{25 of spring/March 25} (N)\\\\
\\noindent {\\tovian \\fontsize{20pt}{10pt} \\textbf{lerithan} }\\\\
\\noindent /ler{\textprimstress}i{\texttheta}an/\\\\
\\noindent Syllables: ler.iθ.an\\\\
\\noindent lit. storm+animal\\\\
\\noindent \\textit{Stormbull}\\\\


\\noindent History:

\\vspace{-0pt}
\\hspace{40pt}
\\begin{tabular}{ccc}
\\textit{6000} & /kwlyri{\texttheta}l{\textsubbridge{t}}lamn/&$\\rightarrow$ & \\textit{6401} & /kwlyri{\texttheta}l{\texttoptiebar{t\textbeltl}}amn/&$\\rightarrow$ & \\textit{7501} & /wlyri{\texttheta}l{\texttoptiebar{t\textbeltl}}amn/&$\\rightarrow$ & \\textit{7600} & /lri{\texttheta}l{\texttoptiebar{t\textbeltl}}an/&$\\rightarrow$ & \\textit{12500} & /lri{\texttheta}{\texttoptiebar{t\textbeltl}}an/&$\\rightarrow$ & \\textit{14001} & /lri{\texttheta}an/&$\\rightarrow$ & \\textit{15001} & /leri{\texttheta}an/& \\\\
\\end{tabular}

\\vspace{20pt}\\hline

\\end{nopagebreak}
\\filbreak



\\vspace{15pt}
\\begin{nopagebreak}
\\noindent{\\fontsize{20pt}{10pt}\\textbf{rithileay} } \\textit{25 of summer/June 25} (N)\\\\
\\noindent {\\tovian \\fontsize{20pt}{10pt} \\textbf{rithileay} }\\\\
\\noindent /ri{\texttheta}il{\textprimstress}eay/\\\\
\\noindent Syllables: riθ.il.e.aj\\\\
\\noindent lit. sky+feather\\\\
\\noindent \\textit{Skyfeather}\\\\


\\noindent History:

\\vspace{-0pt}
\\hspace{40pt}
\\begin{tabular}{ccc}
\\textit{6000} & /ri{\texttheta}ilearya/&$\\rightarrow$ & \\textit{7600} & /ri{\texttheta}ileaya/&$\\rightarrow$ & \\textit{9500} & /ri{\texttheta}ileay/& \\\\
\\end{tabular}

\\vspace{20pt}\\hline

\\end{nopagebreak}
\\filbreak



\\vspace{15pt}
\\begin{nopagebreak}
\\noindent{\\fontsize{20pt}{10pt}\\textbf{nan} } \\textit{25 of winter/december 25} (N)\\\\
\\noindent {\\tovian \\fontsize{20pt}{10pt} \\textbf{nan} }\\\\
\\noindent /n{\textprimstress}an/\\\\
\\noindent Syllables: nan\\\\
\\noindent lit. star+???\\\\
\\noindent \\textit{Starfis}\\\\


\\noindent History:

\\vspace{-0pt}
\\hspace{40pt}
\\begin{tabular}{ccc}
\\textit{6000} & /nnaana???/&$\\rightarrow$ & \\textit{6500} & /naana???/&$\\rightarrow$ & \\textit{8500} & /naana/&$\\rightarrow$ & \\textit{9500} & /naan/&$\\rightarrow$ & \\textit{12005} & /nan/& \\\\
\\end{tabular}

\\vspace{20pt}\\hline

\\end{nopagebreak}
\\filbreak



\\vspace{15pt}
\\begin{nopagebreak}
\\noindent{\\fontsize{20pt}{10pt}\\textbf{melorelor} } \\textit{26 of autumn/September 26} (N)\\\\
\\noindent {\\tovian \\fontsize{20pt}{10pt} \\textbf{melorelor} }\\\\
\\noindent /melor{\textprimstress}elor/\\\\
\\noindent Syllables: mel.or.el.or\\\\
\\noindent lit. plant+plant\\\\
\\noindent \\textit{Hayherb}\\\\


\\noindent History:

\\vspace{-0pt}
\\hspace{40pt}
\\begin{tabular}{ccc}
\\textit{6000} & /mlloremllore/&$\\rightarrow$ & \\textit{6500} & /mloremlore/&$\\rightarrow$ & \\textit{7502} & /meloremlore/&$\\rightarrow$ & \\textit{7600} & /melorelore/&$\\rightarrow$ & \\textit{9500} & /melorelor/& \\\\
\\end{tabular}

\\vspace{20pt}\\hline

\\end{nopagebreak}
\\filbreak



\\vspace{15pt}
\\begin{nopagebreak}
\\noindent{\\fontsize{20pt}{10pt}\\textbf{muramlorelhen} } \\textit{26 of spring/March 26} (N)\\\\
\\noindent {\\tovian \\fontsize{20pt}{10pt} \\textbf{muramlorelhen} }\\\\
\\noindent /muramlor{\textprimstress}e{\textbeltl}en/\\\\
\\noindent Syllables: mur.aml.or.eɬ.en\\\\
\\noindent lit. darkness+plant+dream\\\\
\\noindent \\textit{Nightshade}\\\\


\\noindent History:

\\vspace{-0pt}
\\hspace{40pt}
\\begin{tabular}{ccc}
\\textit{6000} & /bralmllore{\textbeltl}ena/&$\\rightarrow$ & \\textit{6500} & /bralmlore{\textbeltl}ena/&$\\rightarrow$ & \\textit{7600} & /bramlore{\textbeltl}ena/&$\\rightarrow$ & \\textit{9500} & /bramlore{\textbeltl}en/&$\\rightarrow$ & \\textit{11000} & /mramlore{\textbeltl}en/&$\\rightarrow$ & \\textit{15001} & /muramlore{\textbeltl}en/& \\\\
\\end{tabular}

\\vspace{20pt}\\hline

\\end{nopagebreak}
\\filbreak



\\vspace{15pt}
\\begin{nopagebreak}
\\noindent{\\fontsize{20pt}{10pt}\\textbf{dolithor} } \\textit{26 of summer/June 26} (N)\\\\
\\noindent {\\tovian \\fontsize{20pt}{10pt} \\textbf{dolithor} }\\\\
\\noindent /dol{\textprimstress}i{\texttheta}or/\\\\
\\noindent Syllables: dol.iθ.or\\\\
\\noindent lit. death+smell\\\\
\\noindent \\textit{Duskscent}\\\\


\\noindent History:

\\vspace{-0pt}
\\hspace{40pt}
\\begin{tabular}{ccc}
\\textit{6000} & /dolri{\texttheta}lore/&$\\rightarrow$ & \\textit{6503} & /dorli{\texttheta}lore/&$\\rightarrow$ & \\textit{7600} & /doli{\texttheta}lore/&$\\rightarrow$ & \\textit{9500} & /doli{\texttheta}lor/&$\\rightarrow$ & \\textit{12500} & /doli{\texttheta}or/& \\\\
\\end{tabular}

\\vspace{20pt}\\hline

\\end{nopagebreak}
\\filbreak



\\vspace{15pt}
\\begin{nopagebreak}
\\noindent{\\fontsize{20pt}{10pt}\\textbf{shotlodrem} } \\textit{26 of winter/december 26} (N)\\\\
\\noindent {\\tovian \\fontsize{20pt}{10pt} \\textbf{shotlodrem} }\\\\
\\noindent /{\textesh}o{\texttoptiebar{t\textbeltl}}{\textprimstress}odrem/\\\\
\\noindent Syllables: ʃoƛ.odr.em\\\\
\\noindent lit. flowing+flower\\\\
\\noindent \\textit{Mistdande}\\\\


\\noindent History:

\\vspace{-0pt}
\\hspace{40pt}
\\begin{tabular}{ccc}
\\textit{6000} & /{\textesh}o{\texttoptiebar{t\textbeltl}}dremp/&$\\rightarrow$ & \\textit{8750} & /{\textesh}o{\texttoptiebar{t\textbeltl}}drem/&$\\rightarrow$ & \\textit{9300} & /{\textesh}o{\texttoptiebar{t\textbeltl}}odrem/& \\\\
\\end{tabular}

\\vspace{20pt}\\hline

\\end{nopagebreak}
\\filbreak



\\vspace{15pt}
\\begin{nopagebreak}
\\noindent{\\fontsize{20pt}{10pt}\\textbf{thullhen} } \\textit{27 of autumn/September 27} (N)\\\\
\\noindent {\\tovian \\fontsize{20pt}{10pt} \\textbf{thullhen} }\\\\
\\noindent /{\texttheta}{\textprimstress}ul{\textbeltl}en/\\\\
\\noindent Syllables: θulɬ.en\\\\
\\noindent lit. top+dream\\\\
\\noindent \\textit{Sarlic}\\\\


\\noindent History:

\\vspace{-0pt}
\\hspace{40pt}
\\begin{tabular}{ccc}
\\textit{6000} & /{\dh}ul{\textbeltl}ena/&$\\rightarrow$ & \\textit{9500} & /{\dh}ul{\textbeltl}en/&$\\rightarrow$ & \\textit{12002} & /{\texttheta}ul{\textbeltl}en/& \\\\
\\end{tabular}

\\vspace{20pt}\\hline

\\end{nopagebreak}
\\filbreak



\\vspace{15pt}
\\begin{nopagebreak}
\\noindent{\\fontsize{20pt}{10pt}\\textbf{gadrem} } \\textit{27 of spring/March 27} (N)\\\\
\\noindent {\\tovian \\fontsize{20pt}{10pt} \\textbf{gadrem} }\\\\
\\noindent /g{\textprimstress}adrem/\\\\
\\noindent Syllables: gadr.em\\\\
\\noindent lit. hotness+flower\\\\
\\noindent \\textit{Ascbloom}\\\\


\\noindent History:

\\vspace{-0pt}
\\hspace{40pt}
\\begin{tabular}{ccc}
\\textit{6000} & /agadremp/&$\\rightarrow$ & \\textit{7500} & /gadremp/&$\\rightarrow$ & \\textit{8750} & /gadrem/& \\\\
\\end{tabular}

\\vspace{20pt}\\hline

\\end{nopagebreak}
\\filbreak



\\vspace{15pt}
\\begin{nopagebreak}
\\noindent{\\fontsize{20pt}{10pt}\\textbf{lerithdrem} } \\textit{27 of summer/June 27} (N)\\\\
\\noindent {\\tovian \\fontsize{20pt}{10pt} \\textbf{lerithdrem} }\\\\
\\noindent /ler{\textprimstress}i{\texttheta}drem/\\\\
\\noindent Syllables: ler.iθdr.em\\\\
\\noindent lit. wind+flower\\\\
\\noindent \\textit{Windflower}\\\\


\\noindent History:

\\vspace{-0pt}
\\hspace{40pt}
\\begin{tabular}{ccc}
\\textit{6000} & /kwlyri{\texttheta}dremp/&$\\rightarrow$ & \\textit{7501} & /wlyri{\texttheta}dremp/&$\\rightarrow$ & \\textit{7600} & /lri{\texttheta}dremp/&$\\rightarrow$ & \\textit{8750} & /lri{\texttheta}drem/&$\\rightarrow$ & \\textit{15001} & /leri{\texttheta}drem/& \\\\
\\end{tabular}

\\vspace{20pt}\\hline

\\end{nopagebreak}
\\filbreak



\\vspace{15pt}
\\begin{nopagebreak}
\\noindent{\\fontsize{20pt}{10pt}\\textbf{doraklrith} } \\textit{27 of winter/december 27} (N)\\\\
\\noindent {\\tovian \\fontsize{20pt}{10pt} \\textbf{doraklrith} }\\\\
\\noindent /dor{\textprimstress}aklri{\texttheta}/\\\\
\\noindent Syllables: dor.aklr.iθ\\\\
\\noindent lit. blue+breath\\\\
\\noindent \\textit{Windwoad}\\\\


\\noindent History:

\\vspace{-0pt}
\\hspace{40pt}
\\begin{tabular}{ccc}
\\textit{6000} & /dorakwlyri{\texttheta}/&$\\rightarrow$ & \\textit{7600} & /doraklri{\texttheta}/& \\\\
\\end{tabular}

\\vspace{20pt}\\hline

\\end{nopagebreak}
\\filbreak



\\vspace{15pt}
\\begin{nopagebreak}
\\noindent{\\fontsize{20pt}{10pt}\\textbf{???} } \\textit{28 of autumn/September 28} (N)\\\\
\\noindent {\\tovian \\fontsize{20pt}{10pt} \\textbf{???} }\\\\
\\noindent /???/\\\\
\\noindent Syllables: ???\\\\
\\noindent lit. ???\\\\
\\noindent \\textit{Whivetch}\\\\


\\noindent History:

\\vspace{-0pt}
\\hspace{40pt}
\\begin{tabular}{ccc}
\\textit{6000} & /???/& \\\\
\\end{tabular}

\\vspace{20pt}\\hline

\\end{nopagebreak}
\\filbreak



\\vspace{15pt}
\\begin{nopagebreak}
\\noindent{\\fontsize{20pt}{10pt}\\textbf{leshol} } \\textit{28 of spring/March 28} (N)\\\\
\\noindent {\\tovian \\fontsize{20pt}{10pt} \\textbf{leshol} }\\\\
\\noindent /l{\textprimstress}e{\textesh}ol/\\\\
\\noindent Syllables: leʃ.ol\\\\
\\noindent lit. sun+fruit\\\\
\\noindent \\textit{Sunfruit}\\\\


\\noindent History:

\\vspace{-0pt}
\\hspace{40pt}
\\begin{tabular}{ccc}
\\textit{6000} & /efl{\textyogh}ola/&$\\rightarrow$ & \\textit{7500} & /fl{\textyogh}ola/&$\\rightarrow$ & \\textit{7501} & /l{\textyogh}ola/&$\\rightarrow$ & \\textit{9500} & /l{\textyogh}ol/&$\\rightarrow$ & \\textit{12001} & /l{\textesh}ol/&$\\rightarrow$ & \\textit{15001} & /le{\textesh}ol/& \\\\
\\end{tabular}

\\vspace{20pt}\\hline

\\end{nopagebreak}
\\filbreak



\\vspace{15pt}
\\begin{nopagebreak}
\\noindent{\\fontsize{20pt}{10pt}\\textbf{ithrem} } \\textit{28 of summer/June 28} (N)\\\\
\\noindent {\\tovian \\fontsize{20pt}{10pt} \\textbf{ithrem} }\\\\
\\noindent /{\textprimstress}i{\texttheta}rem/\\\\
\\noindent Syllables: iθr.em\\\\
\\noindent lit. thought+flower\\\\
\\noindent \\textit{Thoughtbloom}\\\\


\\noindent History:

\\vspace{-0pt}
\\hspace{40pt}
\\begin{tabular}{ccc}
\\textit{6000} & /i{\texttheta}rdremp/&$\\rightarrow$ & \\textit{8750} & /i{\texttheta}rrem/&$\\rightarrow$ & \\textit{12007} & /i{\texttheta}rem/& \\\\
\\end{tabular}

\\vspace{20pt}\\hline

\\end{nopagebreak}
\\filbreak



\\vspace{15pt}
\\begin{nopagebreak}
\\noindent{\\fontsize{20pt}{10pt}\\textbf{kelumloremlan} } \\textit{28 of winter/december 28} (N)\\\\
\\noindent {\\tovian \\fontsize{20pt}{10pt} \\textbf{kelumloremlan} }\\\\
\\noindent /kelumlor{\textprimstress}emlan/\\\\
\\noindent Syllables: kel.uml.or.eml.an\\\\
\\noindent lit. vine+twilight\\\\
\\noindent \\textit{Twilifern}\\\\


\\noindent History:

\\vspace{-0pt}
\\hspace{40pt}
\\begin{tabular}{ccc}
\\textit{6000} & /kluwmllorebrlana/&$\\rightarrow$ & \\textit{6500} & /kluwmlorebrlana/&$\\rightarrow$ & \\textit{7600} & /klumloreblana/&$\\rightarrow$ & \\textit{9500} & /klumloreblan/&$\\rightarrow$ & \\textit{11000} & /klumloremlan/&$\\rightarrow$ & \\textit{15001} & /kelumloremlan/& \\\\
\\end{tabular}

\\vspace{20pt}\\hline

\\end{nopagebreak}
\\filbreak



\\vspace{15pt}
\\begin{nopagebreak}
\\noindent{\\fontsize{20pt}{10pt}\\textbf{vareloresyaleth} } \\textit{29 of autumn/September 29} (N)\\\\
\\noindent {\\tovian \\fontsize{20pt}{10pt} \\textbf{vareloresyaleth} }\\\\
\\noindent /vareloresy{\textprimstress}ale{\texttheta}/\\\\
\\noindent Syllables: var.el.or.esj.al.eθ\\\\
\\noindent lit. value+grain\\\\
\\noindent \\textit{Goldwheat}\\\\


\\noindent History:

\\vspace{-0pt}
\\hspace{40pt}
\\begin{tabular}{ccc}
\\textit{6000} & /vlaremlloresyale{\dh}a/&$\\rightarrow$ & \\textit{6500} & /vlaremloresyale{\dh}a/&$\\rightarrow$ & \\textit{7600} & /vlareloresyale{\dh}a/&$\\rightarrow$ & \\textit{9500} & /vlareloresyale{\dh}/&$\\rightarrow$ & \\textit{12002} & /vlareloresyale{\texttheta}/&$\\rightarrow$ & \\textit{12500} & /vareloresyale{\texttheta}/& \\\\
\\end{tabular}

\\vspace{20pt}\\hline

\\end{nopagebreak}
\\filbreak



\\vspace{15pt}
\\begin{nopagebreak}
\\noindent{\\fontsize{20pt}{10pt}\\textbf{nanaloresyaleth} } \\textit{29 of spring/March 29} (N)\\\\
\\noindent {\\tovian \\fontsize{20pt}{10pt} \\textbf{nanaloresyaleth} }\\\\
\\noindent /nanaloresy{\textprimstress}ale{\texttheta}/\\\\
\\noindent Syllables: nan.al.or.esj.al.eθ\\\\
\\noindent lit. star+grain\\\\
\\noindent \\textit{Starbarley}\\\\


\\noindent History:

\\vspace{-0pt}
\\hspace{40pt}
\\begin{tabular}{ccc}
\\textit{6000} & /nnaanamlloresyale{\dh}a/&$\\rightarrow$ & \\textit{6500} & /naanamloresyale{\dh}a/&$\\rightarrow$ & \\textit{7600} & /naanaloresyale{\dh}a/&$\\rightarrow$ & \\textit{9500} & /naanaloresyale{\dh}/&$\\rightarrow$ & \\textit{12002} & /naanaloresyale{\texttheta}/&$\\rightarrow$ & \\textit{12005} & /nanaloresyale{\texttheta}/& \\\\
\\end{tabular}

\\vspace{20pt}\\hline

\\end{nopagebreak}
\\filbreak



\\vspace{15pt}
\\begin{nopagebreak}
\\noindent{\\fontsize{20pt}{10pt}\\textbf{muralshol} } \\textit{29 of summer/June 29} (N)\\\\
\\noindent {\\tovian \\fontsize{20pt}{10pt} \\textbf{muralshol} }\\\\
\\noindent /mur{\textprimstress}al{\textesh}ol/\\\\
\\noindent Syllables: mur.alʃ.ol\\\\
\\noindent lit. darkness+fruit\\\\
\\noindent \\textit{Nightberry}\\\\


\\noindent History:

\\vspace{-0pt}
\\hspace{40pt}
\\begin{tabular}{ccc}
\\textit{6000} & /bral{\textyogh}ola/&$\\rightarrow$ & \\textit{9500} & /bral{\textyogh}ol/&$\\rightarrow$ & \\textit{11000} & /mral{\textyogh}ol/&$\\rightarrow$ & \\textit{12001} & /mral{\textesh}ol/&$\\rightarrow$ & \\textit{15001} & /mural{\textesh}ol/& \\\\
\\end{tabular}

\\vspace{20pt}\\hline

\\end{nopagebreak}
\\filbreak



\\vspace{15pt}
\\begin{nopagebreak}
\\noindent{\\fontsize{20pt}{10pt}\\textbf{ayskesh} } \\textit{29 of winter/december 29} (N)\\\\
\\noindent {\\tovian \\fontsize{20pt}{10pt} \\textbf{ayskesh} }\\\\
\\noindent /{\textprimstress}ayske{\textesh}/\\\\
\\noindent Syllables: ajsk.eʃ\\\\
\\noindent lit. moon+fire\\\\
\\noindent \\textit{Moonash}\\\\


\\noindent History:

\\vspace{-0pt}
\\hspace{40pt}
\\begin{tabular}{ccc}
\\textit{6000} & /aliyske{\textesh}/&$\\rightarrow$ & \\textit{6300} & /alyske{\textesh}/&$\\rightarrow$ & \\textit{7600} & /ayske{\textesh}/& \\\\
\\end{tabular}

\\vspace{20pt}\\hline

\\end{nopagebreak}
\\filbreak



\\vspace{15pt}
\\begin{nopagebreak}
\\noindent{\\fontsize{20pt}{10pt}\\textbf{muralaysm} } \\textit{3 of autumn/September 3} (N)\\\\
\\noindent {\\tovian \\fontsize{20pt}{10pt} \\textbf{muralaysm} }\\\\
\\noindent /mur{\textprimstress}alaysm/\\\\
\\noindent Syllables: mur.al.ajsm\\\\
\\noindent lit. darkness+moon+bloom\\\\
\\noindent \\textit{Nightbulb}\\\\


\\noindent History:

\\vspace{-0pt}
\\hspace{40pt}
\\begin{tabular}{ccc}
\\textit{6000} & /bralaliysdrmpfede/&$\\rightarrow$ & \\textit{6300} & /bralalysdrmpfede/&$\\rightarrow$ & \\textit{7600} & /bralaysdmpfede/&$\\rightarrow$ & \\textit{8500} & /bralaysdmfede/&$\\rightarrow$ & \\textit{9500} & /bralaysdmfed/&$\\rightarrow$ & \\textit{11000} & /mralaysdmfed/&$\\rightarrow$ & \\textit{12003} & /mralaysdmd/&$\\rightarrow$ & \\textit{13001} & /mralaysdm/&$\\rightarrow$ & \\textit{13004} & /mralaysm/&$\\rightarrow$ & \\textit{15001} & /muralaysm/& \\\\
\\end{tabular}

\\vspace{20pt}\\hline

\\end{nopagebreak}
\\filbreak



\\vspace{15pt}
\\begin{nopagebreak}
\\noindent{\\fontsize{20pt}{10pt}\\textbf{thanshol} } \\textit{3 of spring/March 3} (N)\\\\
\\noindent {\\tovian \\fontsize{20pt}{10pt} \\textbf{thanshol} }\\\\
\\noindent /{\texttheta}{\textprimstress}an{\textesh}ol/\\\\
\\noindent Syllables: θanʃ.ol\\\\
\\noindent lit. tool+fruit\\\\
\\noindent \\textit{Stellnut}\\\\


\\noindent History:

\\vspace{-0pt}
\\hspace{40pt}
\\begin{tabular}{ccc}
\\textit{6000} & /{\dh}an{\textyogh}ola/&$\\rightarrow$ & \\textit{9500} & /{\dh}an{\textyogh}ol/&$\\rightarrow$ & \\textit{12001} & /{\dh}an{\textesh}ol/&$\\rightarrow$ & \\textit{12002} & /{\texttheta}an{\textesh}ol/& \\\\
\\end{tabular}

\\vspace{20pt}\\hline

\\end{nopagebreak}
\\filbreak



\\vspace{15pt}
\\begin{nopagebreak}
\\noindent{\\fontsize{20pt}{10pt}\\textbf{aysan} } \\textit{3 of summer/June 3} (N)\\\\
\\noindent {\\tovian \\fontsize{20pt}{10pt} \\textbf{aysan} }\\\\
\\noindent /{\textprimstress}aysan/\\\\
\\noindent Syllables: ajs.an\\\\
\\noindent lit. moon+tool\\\\
\\noindent \\textit{Moonspear}\\\\


\\noindent History:

\\vspace{-0pt}
\\hspace{40pt}
\\begin{tabular}{ccc}
\\textit{6000} & /aliys{\dh}an/&$\\rightarrow$ & \\textit{6300} & /alys{\dh}an/&$\\rightarrow$ & \\textit{7600} & /ays{\dh}an/&$\\rightarrow$ & \\textit{12002} & /ays{\texttheta}an/&$\\rightarrow$ & \\textit{13000} & /aysan/& \\\\
\\end{tabular}

\\vspace{20pt}\\hline

\\end{nopagebreak}
\\filbreak



\\vspace{15pt}
\\begin{nopagebreak}
\\noindent{\\fontsize{20pt}{10pt}\\textbf{???} } \\textit{3 of winter/december 3} (N)\\\\
\\noindent {\\tovian \\fontsize{20pt}{10pt} \\textbf{???} }\\\\
\\noindent /???/\\\\
\\noindent Syllables: ???\\\\
\\noindent lit. ???\\\\
\\noindent \\textit{Brezgilly}\\\\


\\noindent History:

\\vspace{-0pt}
\\hspace{40pt}
\\begin{tabular}{ccc}
\\textit{6000} & /???/& \\\\
\\end{tabular}

\\vspace{20pt}\\hline

\\end{nopagebreak}
\\filbreak



\\vspace{15pt}
\\begin{nopagebreak}
\\noindent{\\fontsize{20pt}{10pt}\\textbf{rithtesh} } \\textit{30 of autumn/September 30} (N)\\\\
\\noindent {\\tovian \\fontsize{20pt}{10pt} \\textbf{rithtesh} }\\\\
\\noindent /r{\textprimstress}i{\texttheta}{\textsubbridge{t}}e{\textesh}/\\\\
\\noindent Syllables: riθt.eʃ\\\\
\\noindent lit. air+sound\\\\
\\noindent \\textit{Windflute}\\\\


\\noindent History:

\\vspace{-0pt}
\\hspace{40pt}
\\begin{tabular}{ccc}
\\textit{6000} & /ri{\texttheta}{\textsubbridge{t}}e{\textyogh}/&$\\rightarrow$ & \\textit{12001} & /ri{\texttheta}{\textsubbridge{t}}e{\textesh}/& \\\\
\\end{tabular}

\\vspace{20pt}\\hline

\\end{nopagebreak}
\\filbreak



\\vspace{15pt}
\\begin{nopagebreak}
\\noindent{\\fontsize{20pt}{10pt}\\textbf{shimilhom} } \\textit{30 of spring/March 30} (N)\\\\
\\noindent {\\tovian \\fontsize{20pt}{10pt} \\textbf{shimilhom} }\\\\
\\noindent /{\textesh}im{\textprimstress}i{\textbeltl}om/\\\\
\\noindent Syllables: ʃim.iɬ.om\\\\
\\noindent lit. mystery+container\\\\
\\noindent \\textit{Mystbarrel}\\\\


\\noindent History:

\\vspace{-0pt}
\\hspace{40pt}
\\begin{tabular}{ccc}
\\textit{6000} & /{\textesh}imi{\textbeltl}oma/&$\\rightarrow$ & \\textit{9500} & /{\textesh}imi{\textbeltl}om/& \\\\
\\end{tabular}

\\vspace{20pt}\\hline

\\end{nopagebreak}
\\filbreak



\\vspace{15pt}
\\begin{nopagebreak}
\\noindent{\\fontsize{20pt}{10pt}\\textbf{melor} } \\textit{30 of summer/June 30} (N)\\\\
\\noindent {\\tovian \\fontsize{20pt}{10pt} \\textbf{melor} }\\\\
\\noindent /m{\textprimstress}elor/\\\\
\\noindent Syllables: mel.or\\\\
\\noindent lit. plant+???\\\\
\\noindent \\textit{Sproutbles}\\\\


\\noindent History:

\\vspace{-0pt}
\\hspace{40pt}
\\begin{tabular}{ccc}
\\textit{6000} & /mllore???/&$\\rightarrow$ & \\textit{6500} & /mlore???/&$\\rightarrow$ & \\textit{7502} & /melore???/&$\\rightarrow$ & \\textit{8500} & /melore/&$\\rightarrow$ & \\textit{9500} & /melor/& \\\\
\\end{tabular}

\\vspace{20pt}\\hline

\\end{nopagebreak}
\\filbreak



\\vspace{15pt}
\\begin{nopagebreak}
\\noindent{\\fontsize{20pt}{10pt}\\textbf{lerith} } \\textit{30 of winter/december 30} (N)\\\\
\\noindent {\\tovian \\fontsize{20pt}{10pt} \\textbf{lerith} }\\\\
\\noindent /l{\textprimstress}eri{\texttheta}/\\\\
\\noindent Syllables: ler.iθ\\\\
\\noindent lit. breath+???\\\\
\\noindent \\textit{Galedibble}\\\\


\\noindent History:

\\vspace{-0pt}
\\hspace{40pt}
\\begin{tabular}{ccc}
\\textit{6000} & /kwlyri{\texttheta}???/&$\\rightarrow$ & \\textit{7501} & /wlyri{\texttheta}???/&$\\rightarrow$ & \\textit{7600} & /lri{\texttheta}???/&$\\rightarrow$ & \\textit{8500} & /lri{\texttheta}/&$\\rightarrow$ & \\textit{15001} & /leri{\texttheta}/& \\\\
\\end{tabular}

\\vspace{20pt}\\hline

\\end{nopagebreak}
\\filbreak



\\vspace{15pt}
\\begin{nopagebreak}
\\noindent{\\fontsize{20pt}{10pt}\\textbf{leshol} } \\textit{31 of autumn/October 1} (N)\\\\
\\noindent {\\tovian \\fontsize{20pt}{10pt} \\textbf{leshol} }\\\\
\\noindent /l{\textprimstress}e{\textesh}ol/\\\\
\\noindent Syllables: leʃ.ol\\\\
\\noindent lit. harvest+fruit\\\\
\\noindent \\textit{Starplum}\\\\


\\noindent History:

\\vspace{-0pt}
\\hspace{40pt}
\\begin{tabular}{ccc}
\\textit{6000} & /mllre{\textyogh}{\textyogh}ola/&$\\rightarrow$ & \\textit{6500} & /mlre{\textyogh}ola/&$\\rightarrow$ & \\textit{6503} & /mrle{\textyogh}ola/&$\\rightarrow$ & \\textit{7501} & /rle{\textyogh}ola/&$\\rightarrow$ & \\textit{7600} & /le{\textyogh}ola/&$\\rightarrow$ & \\textit{9500} & /le{\textyogh}ol/&$\\rightarrow$ & \\textit{12001} & /le{\textesh}ol/& \\\\
\\end{tabular}

\\vspace{20pt}\\hline

\\end{nopagebreak}
\\filbreak



\\vspace{15pt}
\\begin{nopagebreak}
\\noindent{\\fontsize{20pt}{10pt}\\textbf{rithileshol} } \\textit{31 of spring/April 1} (N)\\\\
\\noindent {\\tovian \\fontsize{20pt}{10pt} \\textbf{rithileshol} }\\\\
\\noindent /ri{\texttheta}il{\textprimstress}e{\textesh}ol/\\\\
\\noindent Syllables: riθ.il.eʃ.ol\\\\
\\noindent lit. sky+fruit\\\\
\\noindent \\textit{Skyfruit}\\\\


\\noindent History:

\\vspace{-0pt}
\\hspace{40pt}
\\begin{tabular}{ccc}
\\textit{6000} & /ri{\texttheta}ile{\textyogh}ola/&$\\rightarrow$ & \\textit{9500} & /ri{\texttheta}ile{\textyogh}ol/&$\\rightarrow$ & \\textit{12001} & /ri{\texttheta}ile{\textesh}ol/& \\\\
\\end{tabular}

\\vspace{20pt}\\hline

\\end{nopagebreak}
\\filbreak



\\vspace{15pt}
\\begin{nopagebreak}
\\noindent{\\fontsize{20pt}{10pt}\\textbf{lhenananaloresyaleth} } \\textit{31 of summer/July 1} (N)\\\\
\\noindent {\\tovian \\fontsize{20pt}{10pt} \\textbf{lhenananaloresyaleth} }\\\\
\\noindent /{\textbeltl}enananaloresy{\textprimstress}ale{\texttheta}/\\\\
\\noindent Syllables: ɬen.an.an.al.or.esj.al.eθ\\\\
\\noindent lit. dream+star+grain\\\\
\\noindent \\textit{Startriticale}\\\\


\\noindent History:

\\vspace{-0pt}
\\hspace{40pt}
\\begin{tabular}{ccc}
\\textit{6000} & /{\textbeltl}enannaanamlloresyale{\dh}a/&$\\rightarrow$ & \\textit{6500} & /{\textbeltl}enanaanamloresyale{\dh}a/&$\\rightarrow$ & \\textit{7600} & /{\textbeltl}enanaanaloresyale{\dh}a/&$\\rightarrow$ & \\textit{9500} & /{\textbeltl}enanaanaloresyale{\dh}/&$\\rightarrow$ & \\textit{12002} & /{\textbeltl}enanaanaloresyale{\texttheta}/&$\\rightarrow$ & \\textit{12005} & /{\textbeltl}enananaloresyale{\texttheta}/& \\\\
\\end{tabular}

\\vspace{20pt}\\hline

\\end{nopagebreak}
\\filbreak



\\vspace{15pt}
\\begin{nopagebreak}
\\noindent{\\fontsize{20pt}{10pt}\\textbf{???} } \\textit{31 of winter/January 31} (N)\\\\
\\noindent {\\tovian \\fontsize{20pt}{10pt} \\textbf{???} }\\\\
\\noindent /???/\\\\
\\noindent Syllables: ???\\\\
\\noindent lit. ???\\\\
\\noindent \\textit{Frohstlef}\\\\


\\noindent History:

\\vspace{-0pt}
\\hspace{40pt}
\\begin{tabular}{ccc}
\\textit{6000} & /???/& \\\\
\\end{tabular}

\\vspace{20pt}\\hline

\\end{nopagebreak}
\\filbreak



\\vspace{15pt}
\\begin{nopagebreak}
\\noindent{\\fontsize{20pt}{10pt}\\textbf{kelumloreloresyaleth} } \\textit{32 of autumn/October 2} (N)\\\\
\\noindent {\\tovian \\fontsize{20pt}{10pt} \\textbf{kelumloreloresyaleth} }\\\\
\\noindent /kelumloreloresy{\textprimstress}ale{\texttheta}/\\\\
\\noindent Syllables: kel.uml.or.el.or.esj.al.eθ\\\\
\\noindent lit. vine+grain\\\\
\\noindent \\textit{Duskmillet}\\\\


\\noindent History:

\\vspace{-0pt}
\\hspace{40pt}
\\begin{tabular}{ccc}
\\textit{6000} & /kluwmlloremlloresyale{\dh}a/&$\\rightarrow$ & \\textit{6500} & /kluwmloremloresyale{\dh}a/&$\\rightarrow$ & \\textit{7600} & /klumloreloresyale{\dh}a/&$\\rightarrow$ & \\textit{9500} & /klumloreloresyale{\dh}/&$\\rightarrow$ & \\textit{12002} & /klumloreloresyale{\texttheta}/&$\\rightarrow$ & \\textit{15001} & /kelumloreloresyale{\texttheta}/& \\\\
\\end{tabular}

\\vspace{20pt}\\hline

\\end{nopagebreak}
\\filbreak



\\vspace{15pt}
\\begin{nopagebreak}
\\noindent{\\fontsize{20pt}{10pt}\\textbf{aysith} } \\textit{32 of spring/April 2} (N)\\\\
\\noindent {\\tovian \\fontsize{20pt}{10pt} \\textbf{aysith} }\\\\
\\noindent /{\textprimstress}aysi{\texttheta}/\\\\
\\noindent Syllables: ajs.iθ\\\\
\\noindent lit. moon+sharpness\\\\
\\noindent \\textit{Mooncrísp}\\\\


\\noindent History:

\\vspace{-0pt}
\\hspace{40pt}
\\begin{tabular}{ccc}
\\textit{6000} & /aliys{\textyogh}i{\texttheta}/&$\\rightarrow$ & \\textit{6300} & /alys{\textyogh}i{\texttheta}/&$\\rightarrow$ & \\textit{7600} & /ays{\textyogh}i{\texttheta}/&$\\rightarrow$ & \\textit{12001} & /ays{\textesh}i{\texttheta}/&$\\rightarrow$ & \\textit{13000} & /aysi{\texttheta}/& \\\\
\\end{tabular}

\\vspace{20pt}\\hline

\\end{nopagebreak}
\\filbreak



\\vspace{15pt}
\\begin{nopagebreak}
\\noindent{\\fontsize{20pt}{10pt}\\textbf{yeomlor} } \\textit{32 of summer/July 2} (N)\\\\
\\noindent {\\tovian \\fontsize{20pt}{10pt} \\textbf{yeomlor} }\\\\
\\noindent /ye{\textprimstress}omlor/\\\\
\\noindent Syllables: je.oml.or\\\\
\\noindent lit. wisdom+plant\\\\
\\noindent \\textit{Moonmullen}\\\\


\\noindent History:

\\vspace{-0pt}
\\hspace{40pt}
\\begin{tabular}{ccc}
\\textit{6000} & /aryeolmllore/&$\\rightarrow$ & \\textit{6500} & /aryeolmlore/&$\\rightarrow$ & \\textit{7500} & /ryeolmlore/&$\\rightarrow$ & \\textit{7600} & /yeomlore/&$\\rightarrow$ & \\textit{9500} & /yeomlor/& \\\\
\\end{tabular}

\\vspace{20pt}\\hline

\\end{nopagebreak}
\\filbreak



\\vspace{15pt}
\\begin{nopagebreak}
\\noindent{\\fontsize{20pt}{10pt}\\textbf{tushuthok} } \\textit{32 of winter/January 32} (N)\\\\
\\noindent {\\tovian \\fontsize{20pt}{10pt} \\textbf{tushuthok} }\\\\
\\noindent /{\textsubbridge{t}}u{\textesh}{\textprimstress}u{\texttheta}ok/\\\\
\\noindent Syllables: tuʃ.uθ.ok\\\\
\\noindent lit. cold+ground\\\\
\\noindent \\textit{Wintersnap}\\\\


\\noindent History:

\\vspace{-0pt}
\\hspace{40pt}
\\begin{tabular}{ccc}
\\textit{6000} & /{\textsubbridge{t}}u{\textesh}u{\texttheta}ok{\textesh}/&$\\rightarrow$ & \\textit{13001} & /{\textsubbridge{t}}u{\textesh}u{\texttheta}ok/& \\\\
\\end{tabular}

\\vspace{20pt}\\hline

\\end{nopagebreak}
\\filbreak



\\vspace{15pt}
\\begin{nopagebreak}
\\noindent{\\fontsize{20pt}{10pt}\\textbf{???} } \\textit{33 of autumn/October 3} (N)\\\\
\\noindent {\\tovian \\fontsize{20pt}{10pt} \\textbf{???} }\\\\
\\noindent /???/\\\\
\\noindent Syllables: ???\\\\
\\noindent lit. ???\\\\
\\noindent \\textit{Munpuf}\\\\


\\noindent History:

\\vspace{-0pt}
\\hspace{40pt}
\\begin{tabular}{ccc}
\\textit{6000} & /???/& \\\\
\\end{tabular}

\\vspace{20pt}\\hline

\\end{nopagebreak}
\\filbreak



\\vspace{15pt}
\\begin{nopagebreak}
\\noindent{\\fontsize{20pt}{10pt}\\textbf{lerithwen} } \\textit{33 of spring/April 3} (N)\\\\
\\noindent {\\tovian \\fontsize{20pt}{10pt} \\textbf{lerithwen} }\\\\
\\noindent /ler{\textprimstress}i{\texttheta}wen/\\\\
\\noindent Syllables: ler.iθw.en\\\\
\\noindent lit. storm+beauty\\\\
\\noindent \\textit{Thunderpearl}\\\\


\\noindent History:

\\vspace{-0pt}
\\hspace{40pt}
\\begin{tabular}{ccc}
\\textit{6000} & /kwlyri{\texttheta}lwen/&$\\rightarrow$ & \\textit{7501} & /wlyri{\texttheta}lwen/&$\\rightarrow$ & \\textit{7600} & /lri{\texttheta}wen/&$\\rightarrow$ & \\textit{15001} & /leri{\texttheta}wen/& \\\\
\\end{tabular}

\\vspace{20pt}\\hline

\\end{nopagebreak}
\\filbreak



\\vspace{15pt}
\\begin{nopagebreak}
\\noindent{\\fontsize{20pt}{10pt}\\textbf{lelhenashol} } \\textit{33 of summer/July 3} (N)\\\\
\\noindent {\\tovian \\fontsize{20pt}{10pt} \\textbf{lelhenashol} }\\\\
\\noindent /le{\textbeltl}en{\textprimstress}a{\textesh}ol/\\\\
\\noindent Syllables: leɬ.en.aʃ.ol\\\\
\\noindent lit. sun+dream+fruit\\\\
\\noindent \\textit{Sanfruta}\\\\


\\noindent History:

\\vspace{-0pt}
\\hspace{40pt}
\\begin{tabular}{ccc}
\\textit{6000} & /efl{\textbeltl}ena{\textyogh}ola/&$\\rightarrow$ & \\textit{7500} & /fl{\textbeltl}ena{\textyogh}ola/&$\\rightarrow$ & \\textit{7501} & /l{\textbeltl}ena{\textyogh}ola/&$\\rightarrow$ & \\textit{9500} & /l{\textbeltl}ena{\textyogh}ol/&$\\rightarrow$ & \\textit{12001} & /l{\textbeltl}ena{\textesh}ol/&$\\rightarrow$ & \\textit{15001} & /le{\textbeltl}ena{\textesh}ol/& \\\\
\\end{tabular}

\\vspace{20pt}\\hline

\\end{nopagebreak}
\\filbreak



\\vspace{15pt}
\\begin{nopagebreak}
\\noindent{\\fontsize{20pt}{10pt}\\textbf{muralrem} } \\textit{33 of winter/January 33} (N)\\\\
\\noindent {\\tovian \\fontsize{20pt}{10pt} \\textbf{muralrem} }\\\\
\\noindent /mur{\textprimstress}alrem/\\\\
\\noindent Syllables: mur.alr.em\\\\
\\noindent lit. darkness+flower\\\\
\\noindent \\textit{Níghtbloom}\\\\


\\noindent History:

\\vspace{-0pt}
\\hspace{40pt}
\\begin{tabular}{ccc}
\\textit{6000} & /braldremp/&$\\rightarrow$ & \\textit{8750} & /bralrem/&$\\rightarrow$ & \\textit{11000} & /mralrem/&$\\rightarrow$ & \\textit{15001} & /muralrem/& \\\\
\\end{tabular}

\\vspace{20pt}\\hline

\\end{nopagebreak}
\\filbreak



\\vspace{15pt}
\\begin{nopagebreak}
\\noindent{\\fontsize{20pt}{10pt}\\textbf{lhilaloresyaleth} } \\textit{34 of autumn/October 4} (N)\\\\
\\noindent {\\tovian \\fontsize{20pt}{10pt} \\textbf{lhilaloresyaleth} }\\\\
\\noindent /{\textbeltl}ilaloresy{\textprimstress}ale{\texttheta}/\\\\
\\noindent Syllables: ɬil.al.or.esj.al.eθ\\\\
\\noindent lit. night+grain\\\\
\\noindent \\textit{Twilbar}\\\\


\\noindent History:

\\vspace{-0pt}
\\hspace{40pt}
\\begin{tabular}{ccc}
\\textit{6000} & /a{\textbeltl}inlamlloresyale{\dh}a/&$\\rightarrow$ & \\textit{6500} & /a{\textbeltl}inlamloresyale{\dh}a/&$\\rightarrow$ & \\textit{7500} & /{\textbeltl}inlamloresyale{\dh}a/&$\\rightarrow$ & \\textit{7600} & /{\textbeltl}ilaloresyale{\dh}a/&$\\rightarrow$ & \\textit{9500} & /{\textbeltl}ilaloresyale{\dh}/&$\\rightarrow$ & \\textit{12002} & /{\textbeltl}ilaloresyale{\texttheta}/& \\\\
\\end{tabular}

\\vspace{20pt}\\hline

\\end{nopagebreak}
\\filbreak



\\vspace{15pt}
\\begin{nopagebreak}
\\noindent{\\fontsize{20pt}{10pt}\\textbf{delelaglar} } \\textit{34 of spring/April 4} (N)\\\\
\\noindent {\\tovian \\fontsize{20pt}{10pt} \\textbf{delelaglar} }\\\\
\\noindent /delel{\textprimstress}aglar/\\\\
\\noindent Syllables: del.el.agl.ar\\\\
\\noindent lit. body+tree\\\\
\\noindent \\textit{Heartroot}\\\\


\\noindent History:

\\vspace{-0pt}
\\hspace{40pt}
\\begin{tabular}{ccc}
\\textit{6000} & /adllenlaglare/&$\\rightarrow$ & \\textit{6500} & /adlenlaglare/&$\\rightarrow$ & \\textit{7500} & /dlenlaglare/&$\\rightarrow$ & \\textit{7600} & /dlelaglare/&$\\rightarrow$ & \\textit{9500} & /dlelaglar/&$\\rightarrow$ & \\textit{15001} & /delelaglar/& \\\\
\\end{tabular}

\\vspace{20pt}\\hline

\\end{nopagebreak}
\\filbreak



\\vspace{15pt}
\\begin{nopagebreak}
\\noindent{\\fontsize{20pt}{10pt}\\textbf{melorethanglar} } \\textit{34 of summer/July 4} (N)\\\\
\\noindent {\\tovian \\fontsize{20pt}{10pt} \\textbf{melorethanglar} }\\\\
\\noindent /melore{\texttheta}{\textprimstress}a{\ng}lar/\\\\
\\noindent Syllables: mel.or.eθ.aŋl.ar\\\\
\\noindent lit. plant+tool+woods\\\\
\\noindent \\textit{Melorethangor}\\\\


\\noindent History:

\\vspace{-0pt}
\\hspace{40pt}
\\begin{tabular}{ccc}
\\textit{6000} & /mllore{\dh}anglar/&$\\rightarrow$ & \\textit{6000} & /mllore{\dh}a{\ng}glar/&$\\rightarrow$ & \\textit{6500} & /mlore{\dh}a{\ng}glar/&$\\rightarrow$ & \\textit{7502} & /melore{\dh}a{\ng}glar/&$\\rightarrow$ & \\textit{8500} & /melore{\dh}a{\ng}lar/&$\\rightarrow$ & \\textit{12002} & /melore{\texttheta}a{\ng}lar/& \\\\
\\end{tabular}

\\vspace{20pt}\\hline

\\end{nopagebreak}
\\filbreak



\\vspace{15pt}
\\begin{nopagebreak}
\\noindent{\\fontsize{20pt}{10pt}\\textbf{shashanin} } \\textit{34 of winter/January 34} (N)\\\\
\\noindent {\\tovian \\fontsize{20pt}{10pt} \\textbf{shashanin} }\\\\
\\noindent /{\textesh}a{\textesh}{\textprimstress}anin/\\\\
\\noindent Syllables: ʃaʃ.an.in\\\\
\\noindent lit. ice+berry\\\\
\\noindent \\textit{Icepear}\\\\


\\noindent History:

\\vspace{-0pt}
\\hspace{40pt}
\\begin{tabular}{ccc}
\\textit{6000} & /{\textesh}u{\textesh}a{\textyogh}lanina/&$\\rightarrow$ & \\textit{9200} & /{\textesh}{\textschwa}{\textesh}a{\textyogh}lanina/&$\\rightarrow$ & \\textit{9500} & /{\textesh}{\textschwa}{\textesh}a{\textyogh}lanin/&$\\rightarrow$ & \\textit{11990} & /{\textesh}{\textesh}a{\textyogh}lanin/&$\\rightarrow$ & \\textit{12001} & /{\textesh}{\textesh}a{\textesh}lanin/&$\\rightarrow$ & \\textit{12004} & /{\textesh}a{\textesh}lanin/&$\\rightarrow$ & \\textit{12500} & /{\textesh}a{\textesh}anin/& \\\\
\\end{tabular}

\\vspace{20pt}\\hline

\\end{nopagebreak}
\\filbreak



\\vspace{15pt}
\\begin{nopagebreak}
\\noindent{\\fontsize{20pt}{10pt}\\textbf{shotlan} } \\textit{35 of autumn/October 5} (N)\\\\
\\noindent {\\tovian \\fontsize{20pt}{10pt} \\textbf{shotlan} }\\\\
\\noindent /{\textesh}{\textprimstress}o{\texttoptiebar{t\textbeltl}}an/\\\\
\\noindent Syllables: ʃoƛ.an\\\\
\\noindent lit. river+light\\\\
\\noindent \\textit{Rivshine}\\\\


\\noindent History:

\\vspace{-0pt}
\\hspace{40pt}
\\begin{tabular}{ccc}
\\textit{6000} & /{\textesh}o{\texttoptiebar{t\textbeltl}}zlana/&$\\rightarrow$ & \\textit{9500} & /{\textesh}o{\texttoptiebar{t\textbeltl}}zlan/&$\\rightarrow$ & \\textit{12000} & /{\textesh}o{\texttoptiebar{t\textbeltl}}slan/&$\\rightarrow$ & \\textit{12500} & /{\textesh}o{\texttoptiebar{t\textbeltl}}san/&$\\rightarrow$ & \\textit{14000} & /{\textesh}o{\texttoptiebar{t\textbeltl}}an/& \\\\
\\end{tabular}

\\vspace{20pt}\\hline

\\end{nopagebreak}
\\filbreak



\\vspace{15pt}
\\begin{nopagebreak}
\\noindent{\\fontsize{20pt}{10pt}\\textbf{lerithay} } \\textit{35 of spring/April 5} (N)\\\\
\\noindent {\\tovian \\fontsize{20pt}{10pt} \\textbf{lerithay} }\\\\
\\noindent /ler{\textprimstress}i{\texttheta}ay/\\\\
\\noindent Syllables: ler.iθ.aj\\\\
\\noindent lit. storm+feather\\\\
\\noindent \\textit{Stormféther}\\\\


\\noindent History:

\\vspace{-0pt}
\\hspace{40pt}
\\begin{tabular}{ccc}
\\textit{6000} & /kwlyri{\texttheta}larya/&$\\rightarrow$ & \\textit{7501} & /wlyri{\texttheta}larya/&$\\rightarrow$ & \\textit{7600} & /lri{\texttheta}laya/&$\\rightarrow$ & \\textit{9500} & /lri{\texttheta}lay/&$\\rightarrow$ & \\textit{12500} & /lri{\texttheta}ay/&$\\rightarrow$ & \\textit{15001} & /leri{\texttheta}ay/& \\\\
\\end{tabular}

\\vspace{20pt}\\hline

\\end{nopagebreak}
\\filbreak



\\vspace{15pt}
\\begin{nopagebreak}
\\noindent{\\fontsize{20pt}{10pt}\\textbf{lhomtlan} } \\textit{35 of summer/July 5} (N)\\\\
\\noindent {\\tovian \\fontsize{20pt}{10pt} \\textbf{lhomtlan} }\\\\
\\noindent /{\textbeltl}{\textprimstress}om{\texttoptiebar{t\textbeltl}}an/\\\\
\\noindent Syllables: ɬomƛ.an\\\\
\\noindent lit. day+animal\\\\
\\noindent \\textit{Dayram}\\\\


\\noindent History:

\\vspace{-0pt}
\\hspace{40pt}
\\begin{tabular}{ccc}
\\textit{6000} & /{\textbeltl}onm{\textsubbridge{t}}lamn/&$\\rightarrow$ & \\textit{6401} & /{\textbeltl}onm{\texttoptiebar{t\textbeltl}}amn/&$\\rightarrow$ & \\textit{7600} & /{\textbeltl}om{\texttoptiebar{t\textbeltl}}an/& \\\\
\\end{tabular}

\\vspace{20pt}\\hline

\\end{nopagebreak}
\\filbreak



\\vspace{15pt}
\\begin{nopagebreak}
\\noindent{\\fontsize{20pt}{10pt}\\textbf{selshith} } \\textit{35 of winter/January 35} (N)\\\\
\\noindent {\\tovian \\fontsize{20pt}{10pt} \\textbf{selshith} }\\\\
\\noindent /s{\textprimstress}el{\textesh}i{\texttheta}/\\\\
\\noindent Syllables: selʃ.iθ\\\\
\\noindent lit. water+sharpness\\\\
\\noindent \\textit{Snowbristle}\\\\


\\noindent History:

\\vspace{-0pt}
\\hspace{40pt}
\\begin{tabular}{ccc}
\\textit{6000} & /zel{\textyogh}i{\texttheta}/&$\\rightarrow$ & \\textit{12000} & /sel{\textyogh}i{\texttheta}/&$\\rightarrow$ & \\textit{12001} & /sel{\textesh}i{\texttheta}/& \\\\
\\end{tabular}

\\vspace{20pt}\\hline

\\end{nopagebreak}
\\filbreak



\\vspace{15pt}
\\begin{nopagebreak}
\\noindent{\\fontsize{20pt}{10pt}\\textbf{yisnal} } \\textit{36 of autumn/October 6} (N)\\\\
\\noindent {\\tovian \\fontsize{20pt}{10pt} \\textbf{yisnal} }\\\\
\\noindent /y{\textprimstress}isnal/\\\\
\\noindent Syllables: jisn.al\\\\
\\noindent lit. moon+sharing\\\\
\\noindent \\textit{Tlunalis}\\\\


\\noindent History:

\\vspace{-0pt}
\\hspace{40pt}
\\begin{tabular}{ccc}
\\textit{6000} & /aliys{\texttoptiebar{t\textbeltl}}naels/&$\\rightarrow$ & \\textit{6300} & /alys{\texttoptiebar{t\textbeltl}}naels/&$\\rightarrow$ & \\textit{7500} & /lys{\texttoptiebar{t\textbeltl}}naels/&$\\rightarrow$ & \\textit{7501} & /ys{\texttoptiebar{t\textbeltl}}naels/&$\\rightarrow$ & \\textit{10000} & /ys{\texttoptiebar{t\textbeltl}}nals/&$\\rightarrow$ & \\textit{13001} & /ys{\texttoptiebar{t\textbeltl}}nal/&$\\rightarrow$ & \\textit{14001} & /ysnal/&$\\rightarrow$ & \\textit{15001} & /yisnal/& \\\\
\\end{tabular}

\\vspace{20pt}\\hline

\\end{nopagebreak}
\\filbreak



\\vspace{15pt}
\\begin{nopagebreak}
\\noindent{\\fontsize{20pt}{10pt}\\textbf{teshglar} } \\textit{36 of spring/April 6} (N)\\\\
\\noindent {\\tovian \\fontsize{20pt}{10pt} \\textbf{teshglar} }\\\\
\\noindent /{\textsubbridge{t}}{\textprimstress}e{\textesh}glar/\\\\
\\noindent Syllables: teʃgl.ar\\\\
\\noindent lit. sound+woods\\\\
\\noindent \\textit{Teshangor}\\\\


\\noindent History:

\\vspace{-0pt}
\\hspace{40pt}
\\begin{tabular}{ccc}
\\textit{6000} & /{\textsubbridge{t}}e{\textyogh}glar/&$\\rightarrow$ & \\textit{12001} & /{\textsubbridge{t}}e{\textesh}glar/& \\\\
\\end{tabular}

\\vspace{20pt}\\hline

\\end{nopagebreak}
\\filbreak



\\vspace{15pt}
\\begin{nopagebreak}
\\noindent{\\fontsize{20pt}{10pt}\\textbf{muloanalor} } \\textit{36 of summer/July 6} (N)\\\\
\\noindent {\\tovian \\fontsize{20pt}{10pt} \\textbf{muloanalor} }\\\\
\\noindent /muloan{\textprimstress}alor/\\\\
\\noindent Syllables: mul.o.an.al.or\\\\
\\noindent lit. dawn+plant\\\\
\\noindent \\textit{Dawnfern}\\\\


\\noindent History:

\\vspace{-0pt}
\\hspace{40pt}
\\begin{tabular}{ccc}
\\textit{6000} & /ploanamllore/&$\\rightarrow$ & \\textit{6500} & /ploanamlore/&$\\rightarrow$ & \\textit{7600} & /ploanalore/&$\\rightarrow$ & \\textit{9500} & /ploanalor/&$\\rightarrow$ & \\textit{11000} & /mloanalor/&$\\rightarrow$ & \\textit{15001} & /muloanalor/& \\\\
\\end{tabular}

\\vspace{20pt}\\hline

\\end{nopagebreak}
\\filbreak



\\vspace{15pt}
\\begin{nopagebreak}
\\noindent{\\fontsize{20pt}{10pt}\\textbf{kelumloreshud} } \\textit{36 of winter/January 36} (N)\\\\
\\noindent {\\tovian \\fontsize{20pt}{10pt} \\textbf{kelumloreshud} }\\\\
\\noindent /kelumlor{\textprimstress}e{\textesh}ud/\\\\
\\noindent Syllables: kel.uml.or.eʃ.ud\\\\
\\noindent lit. vine+ice\\\\
\\noindent \\textit{Frostfern}\\\\


\\noindent History:

\\vspace{-0pt}
\\hspace{40pt}
\\begin{tabular}{ccc}
\\textit{6000} & /kluwmllore{\textesh}u{\textesh}a/&$\\rightarrow$ & \\textit{6500} & /kluwmlore{\textesh}u{\textesh}a/&$\\rightarrow$ & \\textit{7600} & /klumlore{\textesh}u{\textesh}a/&$\\rightarrow$ & \\textit{9500} & /klumlore{\textesh}u{\textesh}/&$\\rightarrow$ & \\textit{12501} & /klumlore{\textesh}ud/&$\\rightarrow$ & \\textit{15001} & /kelumlore{\textesh}ud/& \\\\
\\end{tabular}

\\vspace{20pt}\\hline

\\end{nopagebreak}
\\filbreak



\\vspace{15pt}
\\begin{nopagebreak}
\\noindent{\\fontsize{20pt}{10pt}\\textbf{hol} } \\textit{37 of autumn/October 7} (N)\\\\
\\noindent {\\tovian \\fontsize{20pt}{10pt} \\textbf{hol} }\\\\
\\noindent /ħ{\textprimstress}ol/\\\\
\\noindent Syllables: ħol\\\\
\\noindent lit. sweet+???\\\\
\\noindent \\textit{Sweetspore}\\\\


\\noindent History:

\\vspace{-0pt}
\\hspace{40pt}
\\begin{tabular}{ccc}
\\textit{6000} & /hola???/&$\\rightarrow$ & \\textit{8500} & /hola/&$\\rightarrow$ & \\textit{9500} & /hol/&$\\rightarrow$ & \\textit{14006} & /ħol/& \\\\
\\end{tabular}

\\vspace{20pt}\\hline

\\end{nopagebreak}
\\filbreak



\\vspace{15pt}
\\begin{nopagebreak}
\\noindent{\\fontsize{20pt}{10pt}\\textbf{deremshol} } \\textit{37 of spring/April 7} (N)\\\\
\\noindent {\\tovian \\fontsize{20pt}{10pt} \\textbf{deremshol} }\\\\
\\noindent /der{\textprimstress}em{\textesh}ol/\\\\
\\noindent Syllables: der.emʃ.ol\\\\
\\noindent lit. flower+fruit\\\\
\\noindent \\textit{Drémfig}\\\\


\\noindent History:

\\vspace{-0pt}
\\hspace{40pt}
\\begin{tabular}{ccc}
\\textit{6000} & /dremp{\textyogh}ola/&$\\rightarrow$ & \\textit{8500} & /drem{\textyogh}ola/&$\\rightarrow$ & \\textit{9500} & /drem{\textyogh}ol/&$\\rightarrow$ & \\textit{12001} & /drem{\textesh}ol/&$\\rightarrow$ & \\textit{15001} & /derem{\textesh}ol/& \\\\
\\end{tabular}

\\vspace{20pt}\\hline

\\end{nopagebreak}
\\filbreak



\\vspace{15pt}
\\begin{nopagebreak}
\\noindent{\\fontsize{20pt}{10pt}\\textbf{lhenamral} } \\textit{37 of summer/July 7} (N)\\\\
\\noindent {\\tovian \\fontsize{20pt}{10pt} \\textbf{lhenamral} }\\\\
\\noindent /{\textbeltl}en{\textprimstress}amral/\\\\
\\noindent Syllables: ɬen.amr.al\\\\
\\noindent lit. dream+darkness\\\\
\\noindent \\textit{Nightwood}\\\\


\\noindent History:

\\vspace{-0pt}
\\hspace{40pt}
\\begin{tabular}{ccc}
\\textit{6000} & /{\textbeltl}enabral/&$\\rightarrow$ & \\textit{11000} & /{\textbeltl}enamral/& \\\\
\\end{tabular}

\\vspace{20pt}\\hline

\\end{nopagebreak}
\\filbreak



\\vspace{15pt}
\\begin{nopagebreak}
\\noindent{\\fontsize{20pt}{10pt}\\textbf{aysel} } \\textit{37 of winter/January 37} (N)\\\\
\\noindent {\\tovian \\fontsize{20pt}{10pt} \\textbf{aysel} }\\\\
\\noindent /{\textprimstress}aysel/\\\\
\\noindent Syllables: ajs.el\\\\
\\noindent lit. moon+center\\\\
\\noindent \\textit{Mooncrést}\\\\


\\noindent History:

\\vspace{-0pt}
\\hspace{40pt}
\\begin{tabular}{ccc}
\\textit{6000} & /aliys{\textbeltl}el/&$\\rightarrow$ & \\textit{6300} & /alys{\textbeltl}el/&$\\rightarrow$ & \\textit{7600} & /ays{\textbeltl}el/&$\\rightarrow$ & \\textit{13000} & /aysel/& \\\\
\\end{tabular}

\\vspace{20pt}\\hline

\\end{nopagebreak}
\\filbreak



\\vspace{15pt}
\\begin{nopagebreak}
\\noindent{\\fontsize{20pt}{10pt}\\textbf{lhenadrem} } \\textit{38 of autumn/October 8} (N)\\\\
\\noindent {\\tovian \\fontsize{20pt}{10pt} \\textbf{lhenadrem} }\\\\
\\noindent /{\textbeltl}en{\textprimstress}adrem/\\\\
\\noindent Syllables: ɬen.adr.em\\\\
\\noindent lit. dream+flower\\\\
\\noindent \\textit{Dremblossom}\\\\


\\noindent History:

\\vspace{-0pt}
\\hspace{40pt}
\\begin{tabular}{ccc}
\\textit{6000} & /{\textbeltl}enadremp/&$\\rightarrow$ & \\textit{8750} & /{\textbeltl}enadrem/& \\\\
\\end{tabular}

\\vspace{20pt}\\hline

\\end{nopagebreak}
\\filbreak



\\vspace{15pt}
\\begin{nopagebreak}
\\noindent{\\fontsize{20pt}{10pt}\\textbf{muraladlel} } \\textit{38 of spring/April 8} (N)\\\\
\\noindent {\\tovian \\fontsize{20pt}{10pt} \\textbf{muraladlel} }\\\\
\\noindent /mural{\textprimstress}adlel/\\\\
\\noindent Syllables: mur.al.adl.el\\\\
\\noindent lit. darkness+body\\\\
\\noindent \\textit{Níghtstalk}\\\\


\\noindent History:

\\vspace{-0pt}
\\hspace{40pt}
\\begin{tabular}{ccc}
\\textit{6000} & /braladllenl/&$\\rightarrow$ & \\textit{6500} & /braladlenl/&$\\rightarrow$ & \\textit{7600} & /braladlel/&$\\rightarrow$ & \\textit{11000} & /mraladlel/&$\\rightarrow$ & \\textit{15001} & /muraladlel/& \\\\
\\end{tabular}

\\vspace{20pt}\\hline

\\end{nopagebreak}
\\filbreak



\\vspace{15pt}
\\begin{nopagebreak}
\\noindent{\\fontsize{20pt}{10pt}\\textbf{lelhenalor} } \\textit{38 of summer/July 8} (N)\\\\
\\noindent {\\tovian \\fontsize{20pt}{10pt} \\textbf{lelhenalor} }\\\\
\\noindent /le{\textbeltl}en{\textprimstress}alor/\\\\
\\noindent Syllables: leɬ.en.al.or\\\\
\\noindent lit. sun+dream+plant\\\\
\\noindent \\textit{Sunthistle}\\\\


\\noindent History:

\\vspace{-0pt}
\\hspace{40pt}
\\begin{tabular}{ccc}
\\textit{6000} & /efl{\textbeltl}enamllore/&$\\rightarrow$ & \\textit{6500} & /efl{\textbeltl}enamlore/&$\\rightarrow$ & \\textit{7500} & /fl{\textbeltl}enamlore/&$\\rightarrow$ & \\textit{7501} & /l{\textbeltl}enamlore/&$\\rightarrow$ & \\textit{7600} & /l{\textbeltl}enalore/&$\\rightarrow$ & \\textit{9500} & /l{\textbeltl}enalor/&$\\rightarrow$ & \\textit{15001} & /le{\textbeltl}enalor/& \\\\
\\end{tabular}

\\vspace{20pt}\\hline

\\end{nopagebreak}
\\filbreak



\\vspace{15pt}
\\begin{nopagebreak}
\\noindent{\\fontsize{20pt}{10pt}\\textbf{fol} } \\textit{38 of winter/January 38} (N)\\\\
\\noindent {\\tovian \\fontsize{20pt}{10pt} \\textbf{fol} }\\\\
\\noindent /f{\textprimstress}ol/\\\\
\\noindent Syllables: fol\\\\
\\noindent lit. sun+sweet\\\\
\\noindent \\textit{Sunnectar}\\\\


\\noindent History:

\\vspace{-0pt}
\\hspace{40pt}
\\begin{tabular}{ccc}
\\textit{6000} & /eflhola/&$\\rightarrow$ & \\textit{6402} & /ef{\textbeltl}ola/&$\\rightarrow$ & \\textit{7500} & /f{\textbeltl}ola/&$\\rightarrow$ & \\textit{9500} & /f{\textbeltl}ol/&$\\rightarrow$ & \\textit{13000} & /fol/& \\\\
\\end{tabular}

\\vspace{20pt}\\hline

\\end{nopagebreak}
\\filbreak



\\vspace{15pt}
\\begin{nopagebreak}
\\noindent{\\fontsize{20pt}{10pt}\\textbf{muralaglar} } \\textit{39 of autumn/October 9} (N)\\\\
\\noindent {\\tovian \\fontsize{20pt}{10pt} \\textbf{muralaglar} }\\\\
\\noindent /mural{\textprimstress}aglar/\\\\
\\noindent Syllables: mur.al.agl.ar\\\\
\\noindent lit. darkness+tree\\\\
\\noindent \\textit{Shadroot}\\\\


\\noindent History:

\\vspace{-0pt}
\\hspace{40pt}
\\begin{tabular}{ccc}
\\textit{6000} & /bralaglare/&$\\rightarrow$ & \\textit{9500} & /bralaglar/&$\\rightarrow$ & \\textit{11000} & /mralaglar/&$\\rightarrow$ & \\textit{15001} & /muralaglar/& \\\\
\\end{tabular}

\\vspace{20pt}\\hline

\\end{nopagebreak}
\\filbreak



\\vspace{15pt}
\\begin{nopagebreak}
\\noindent{\\fontsize{20pt}{10pt}\\textbf{yeolaglar} } \\textit{39 of spring/April 9} (N)\\\\
\\noindent {\\tovian \\fontsize{20pt}{10pt} \\textbf{yeolaglar} }\\\\
\\noindent /yeol{\textprimstress}aglar/\\\\
\\noindent Syllables: je.ol.agl.ar\\\\
\\noindent lit. truth+tree\\\\
\\noindent \\textit{Wísertre}\\\\


\\noindent History:

\\vspace{-0pt}
\\hspace{40pt}
\\begin{tabular}{ccc}
\\textit{6000} & /aryeolaglare/&$\\rightarrow$ & \\textit{7500} & /ryeolaglare/&$\\rightarrow$ & \\textit{7600} & /yeolaglare/&$\\rightarrow$ & \\textit{9500} & /yeolaglar/& \\\\
\\end{tabular}

\\vspace{20pt}\\hline

\\end{nopagebreak}
\\filbreak



\\vspace{15pt}
\\begin{nopagebreak}
\\noindent{\\fontsize{20pt}{10pt}\\textbf{shaninalhil} } \\textit{39 of summer/July 9} (N)\\\\
\\noindent {\\tovian \\fontsize{20pt}{10pt} \\textbf{shaninalhil} }\\\\
\\noindent /{\textesh}anin{\textprimstress}a{\textbeltl}il/\\\\
\\noindent Syllables: ʃan.in.aɬ.il\\\\
\\noindent lit. berry+night\\\\
\\noindent \\textit{Twilightberry}\\\\


\\noindent History:

\\vspace{-0pt}
\\hspace{40pt}
\\begin{tabular}{ccc}
\\textit{6000} & /{\textyogh}laninaa{\textbeltl}inla/&$\\rightarrow$ & \\textit{7600} & /{\textyogh}laninaa{\textbeltl}ila/&$\\rightarrow$ & \\textit{9500} & /{\textyogh}laninaa{\textbeltl}il/&$\\rightarrow$ & \\textit{12001} & /{\textesh}laninaa{\textbeltl}il/&$\\rightarrow$ & \\textit{12005} & /{\textesh}lanina{\textbeltl}il/&$\\rightarrow$ & \\textit{12500} & /{\textesh}anina{\textbeltl}il/& \\\\
\\end{tabular}

\\vspace{20pt}\\hline

\\end{nopagebreak}
\\filbreak



\\vspace{15pt}
\\begin{nopagebreak}
\\noindent{\\fontsize{20pt}{10pt}\\textbf{lhenashanin} } \\textit{39 of winter/January 39} (N)\\\\
\\noindent {\\tovian \\fontsize{20pt}{10pt} \\textbf{lhenashanin} }\\\\
\\noindent /{\textbeltl}ena{\textesh}{\textprimstress}anin/\\\\
\\noindent Syllables: ɬen.aʃ.an.in\\\\
\\noindent lit. spirit+berry\\\\
\\noindent \\textit{Spirbery}\\\\


\\noindent History:

\\vspace{-0pt}
\\hspace{40pt}
\\begin{tabular}{ccc}
\\textit{6000} & /{\textbeltl}ena{\textyogh}lanina/&$\\rightarrow$ & \\textit{9500} & /{\textbeltl}ena{\textyogh}lanin/&$\\rightarrow$ & \\textit{12001} & /{\textbeltl}ena{\textesh}lanin/&$\\rightarrow$ & \\textit{12500} & /{\textbeltl}ena{\textesh}anin/& \\\\
\\end{tabular}

\\vspace{20pt}\\hline

\\end{nopagebreak}
\\filbreak



\\vspace{15pt}
\\begin{nopagebreak}
\\noindent{\\fontsize{20pt}{10pt}\\textbf{???} } \\textit{4 of autumn/September 4} (N)\\\\
\\noindent {\\tovian \\fontsize{20pt}{10pt} \\textbf{???} }\\\\
\\noindent /???/\\\\
\\noindent Syllables: ???\\\\
\\noindent lit. ???\\\\
\\noindent \\textit{Daeronica}\\\\


\\noindent History:

\\vspace{-0pt}
\\hspace{40pt}
\\begin{tabular}{ccc}
\\textit{6000} & /???/& \\\\
\\end{tabular}

\\vspace{20pt}\\hline

\\end{nopagebreak}
\\filbreak



\\vspace{15pt}
\\begin{nopagebreak}
\\noindent{\\fontsize{20pt}{10pt}\\textbf{deremrem} } \\textit{4 of spring/March 4} (N)\\\\
\\noindent {\\tovian \\fontsize{20pt}{10pt} \\textbf{deremrem} }\\\\
\\noindent /der{\textprimstress}emrem/\\\\
\\noindent Syllables: der.emr.em\\\\
\\noindent lit. flower+flower\\\\
\\noindent \\textit{Drembloom}\\\\


\\noindent History:

\\vspace{-0pt}
\\hspace{40pt}
\\begin{tabular}{ccc}
\\textit{6000} & /drempdremp/&$\\rightarrow$ & \\textit{8500} & /dremremp/&$\\rightarrow$ & \\textit{8750} & /dremrem/&$\\rightarrow$ & \\textit{15001} & /deremrem/& \\\\
\\end{tabular}

\\vspace{20pt}\\hline

\\end{nopagebreak}
\\filbreak



\\vspace{15pt}
\\begin{nopagebreak}
\\noindent{\\fontsize{20pt}{10pt}\\textbf{lel} } \\textit{4 of summer/June 4} (N)\\\\
\\noindent {\\tovian \\fontsize{20pt}{10pt} \\textbf{lel} }\\\\
\\noindent /l{\textprimstress}el/\\\\
\\noindent Syllables: lel\\\\
\\noindent lit. sun+branch\\\\
\\noindent \\textit{Sunpetal}\\\\


\\noindent History:

\\vspace{-0pt}
\\hspace{40pt}
\\begin{tabular}{ccc}
\\textit{6000} & /eflkelu/&$\\rightarrow$ & \\textit{7500} & /flkelu/&$\\rightarrow$ & \\textit{7501} & /lkelu/&$\\rightarrow$ & \\textit{8750} & /lelu/&$\\rightarrow$ & \\textit{9500} & /lel/& \\\\
\\end{tabular}

\\vspace{20pt}\\hline

\\end{nopagebreak}
\\filbreak



\\vspace{15pt}
\\begin{nopagebreak}
\\noindent{\\fontsize{20pt}{10pt}\\textbf{???} } \\textit{4 of winter/december 4} (N)\\\\
\\noindent {\\tovian \\fontsize{20pt}{10pt} \\textbf{???} }\\\\
\\noindent /???/\\\\
\\noindent Syllables: ???\\\\
\\noindent lit. ???\\\\
\\noindent \\textit{Stormprive}\\\\


\\noindent History:

\\vspace{-0pt}
\\hspace{40pt}
\\begin{tabular}{ccc}
\\textit{6000} & /???/& \\\\
\\end{tabular}

\\vspace{20pt}\\hline

\\end{nopagebreak}
\\filbreak



\\vspace{15pt}
\\begin{nopagebreak}
\\noindent{\\fontsize{20pt}{10pt}\\textbf{nanathan} } \\textit{40 of autumn/October 10} (N)\\\\
\\noindent {\\tovian \\fontsize{20pt}{10pt} \\textbf{nanathan} }\\\\
\\noindent /nan{\textprimstress}a{\texttheta}an/\\\\
\\noindent Syllables: nan.aθ.an\\\\
\\noindent lit. star+tool\\\\
\\noindent \\textit{Starladder}\\\\


\\noindent History:

\\vspace{-0pt}
\\hspace{40pt}
\\begin{tabular}{ccc}
\\textit{6000} & /nnaana{\dh}an/&$\\rightarrow$ & \\textit{6500} & /naana{\dh}an/&$\\rightarrow$ & \\textit{12002} & /naana{\texttheta}an/&$\\rightarrow$ & \\textit{12005} & /nana{\texttheta}an/& \\\\
\\end{tabular}

\\vspace{20pt}\\hline

\\end{nopagebreak}
\\filbreak



\\vspace{15pt}
\\begin{nopagebreak}
\\noindent{\\fontsize{20pt}{10pt}\\textbf{rithilenum} } \\textit{40 of spring} (N)\\\\
\\noindent {\\tovian \\fontsize{20pt}{10pt} \\textbf{rithilenum} }\\\\
\\noindent /ri{\texttheta}il{\textprimstress}enum/\\\\
\\noindent Syllables: riθ.il.en.um\\\\
\\noindent lit. sky+path\\\\
\\noindent \\textit{Cloudfurrow}\\\\


\\noindent History:

\\vspace{-0pt}
\\hspace{40pt}
\\begin{tabular}{ccc}
\\textit{6000} & /ri{\texttheta}ilenum/& \\\\
\\end{tabular}

\\vspace{20pt}\\hline

\\end{nopagebreak}
\\filbreak



\\vspace{15pt}
\\begin{nopagebreak}
\\noindent{\\fontsize{20pt}{10pt}\\textbf{selnan} } \\textit{40 of summer/July 10} (N)\\\\
\\noindent {\\tovian \\fontsize{20pt}{10pt} \\textbf{selnan} }\\\\
\\noindent /s{\textprimstress}elnan/\\\\
\\noindent Syllables: seln.an\\\\
\\noindent lit. water+protection\\\\
\\noindent \\textit{Rainkeeper}\\\\


\\noindent History:

\\vspace{-0pt}
\\hspace{40pt}
\\begin{tabular}{ccc}
\\textit{6000} & /zelbnani/&$\\rightarrow$ & \\textit{8750} & /zelnani/&$\\rightarrow$ & \\textit{9500} & /zelnan/&$\\rightarrow$ & \\textit{12000} & /selnan/& \\\\
\\end{tabular}

\\vspace{20pt}\\hline

\\end{nopagebreak}
\\filbreak



\\vspace{15pt}
\\begin{nopagebreak}
\\noindent{\\fontsize{20pt}{10pt}\\textbf{shaingtlurith} } \\textit{40 of winter/January 40} (N)\\\\
\\noindent {\\tovian \\fontsize{20pt}{10pt} \\textbf{shaingtlurith} }\\\\
\\noindent /{\textesh}ai{\ng}{\texttoptiebar{t\textbeltl}}{\textprimstress}uri{\texttheta}/\\\\
\\noindent Syllables: ʃa.iŋƛ.ur.iθ\\\\
\\noindent lit. ice+fog\\\\
\\noindent \\textit{Frostfoge}\\\\


\\noindent History:

\\vspace{-0pt}
\\hspace{40pt}
\\begin{tabular}{ccc}
\\textit{6000} & /{\textesh}u{\textesh}ai{\ng}g{\texttoptiebar{t\textbeltl}}unri{\texttheta}/&$\\rightarrow$ & \\textit{7600} & /{\textesh}u{\textesh}ai{\ng}g{\texttoptiebar{t\textbeltl}}uri{\texttheta}/&$\\rightarrow$ & \\textit{8500} & /{\textesh}u{\textesh}ai{\ng}{\texttoptiebar{t\textbeltl}}uri{\texttheta}/&$\\rightarrow$ & \\textit{9200} & /{\textesh}{\textschwa}{\textesh}ai{\ng}{\texttoptiebar{t\textbeltl}}uri{\texttheta}/&$\\rightarrow$ & \\textit{11990} & /{\textesh}{\textesh}ai{\ng}{\texttoptiebar{t\textbeltl}}uri{\texttheta}/&$\\rightarrow$ & \\textit{12004} & /{\textesh}ai{\ng}{\texttoptiebar{t\textbeltl}}uri{\texttheta}/& \\\\
\\end{tabular}

\\vspace{20pt}\\hline

\\end{nopagebreak}
\\filbreak



\\vspace{15pt}
\\begin{nopagebreak}
\\noindent{\\fontsize{20pt}{10pt}\\textbf{yisanin} } \\textit{41 of autumn/October 11} (N)\\\\
\\noindent {\\tovian \\fontsize{20pt}{10pt} \\textbf{yisanin} }\\\\
\\noindent /yis{\textprimstress}anin/\\\\
\\noindent Syllables: jis.an.in\\\\
\\noindent lit. moon+berry\\\\
\\noindent \\textit{Moonmelon}\\\\


\\noindent History:

\\vspace{-0pt}
\\hspace{40pt}
\\begin{tabular}{ccc}
\\textit{6000} & /aliys{\textyogh}lanina/&$\\rightarrow$ & \\textit{6300} & /alys{\textyogh}lanina/&$\\rightarrow$ & \\textit{7500} & /lys{\textyogh}lanina/&$\\rightarrow$ & \\textit{7501} & /ys{\textyogh}lanina/&$\\rightarrow$ & \\textit{9500} & /ys{\textyogh}lanin/&$\\rightarrow$ & \\textit{12001} & /ys{\textesh}lanin/&$\\rightarrow$ & \\textit{12500} & /ys{\textesh}anin/&$\\rightarrow$ & \\textit{13000} & /ysanin/&$\\rightarrow$ & \\textit{15001} & /yisanin/& \\\\
\\end{tabular}

\\vspace{20pt}\\hline

\\end{nopagebreak}
\\filbreak



\\vspace{15pt}
\\begin{nopagebreak}
\\noindent{\\fontsize{20pt}{10pt}\\textbf{lanatan} } \\textit{41 of spring/April 11} (N)\\\\
\\noindent {\\tovian \\fontsize{20pt}{10pt} \\textbf{lanatan} }\\\\
\\noindent /lan{\textprimstress}a{\textsubbridge{t}}an/\\\\
\\noindent Syllables: lan.at.an\\\\
\\noindent lit. twilight+face\\\\
\\noindent \\textit{Duskwhísker}\\\\


\\noindent History:

\\vspace{-0pt}
\\hspace{40pt}
\\begin{tabular}{ccc}
\\textit{6000} & /brlana{\textsubbridge{t}}arn/&$\\rightarrow$ & \\textit{7501} & /rlana{\textsubbridge{t}}arn/&$\\rightarrow$ & \\textit{7600} & /lana{\textsubbridge{t}}an/& \\\\
\\end{tabular}

\\vspace{20pt}\\hline

\\end{nopagebreak}
\\filbreak



\\vspace{15pt}
\\begin{nopagebreak}
\\noindent{\\fontsize{20pt}{10pt}\\textbf{melsyaleth} } \\textit{41 of summer/July 11} (N)\\\\
\\noindent {\\tovian \\fontsize{20pt}{10pt} \\textbf{melsyaleth} }\\\\
\\noindent /melsy{\textprimstress}ale{\texttheta}/\\\\
\\noindent Syllables: melsj.al.eθ\\\\
\\noindent lit. life+food\\\\
\\noindent \\textit{Heartbread}\\\\


\\noindent History:

\\vspace{-0pt}
\\hspace{40pt}
\\begin{tabular}{ccc}
\\textit{6000} & /melsyale{\dh}a/&$\\rightarrow$ & \\textit{9500} & /melsyale{\dh}/&$\\rightarrow$ & \\textit{12002} & /melsyale{\texttheta}/& \\\\
\\end{tabular}

\\vspace{20pt}\\hline

\\end{nopagebreak}
\\filbreak



\\vspace{15pt}
\\begin{nopagebreak}
\\noindent{\\fontsize{20pt}{10pt}\\textbf{nanaingtlun} } \\textit{41 of winter/January 41} (N)\\\\
\\noindent {\\tovian \\fontsize{20pt}{10pt} \\textbf{nanaingtlun} }\\\\
\\noindent /nana{\textprimstress}i{\ng}{\texttoptiebar{t\textbeltl}}un/\\\\
\\noindent Syllables: nan.a.iŋƛ.un\\\\
\\noindent lit. star+tar\\\\
\\noindent \\textit{Starwax}\\\\


\\noindent History:

\\vspace{-0pt}
\\hspace{40pt}
\\begin{tabular}{ccc}
\\textit{6000} & /nnaanai{\ng}g{\texttoptiebar{t\textbeltl}}una/&$\\rightarrow$ & \\textit{6500} & /naanai{\ng}g{\texttoptiebar{t\textbeltl}}una/&$\\rightarrow$ & \\textit{8500} & /naanai{\ng}{\texttoptiebar{t\textbeltl}}una/&$\\rightarrow$ & \\textit{9500} & /naanai{\ng}{\texttoptiebar{t\textbeltl}}un/&$\\rightarrow$ & \\textit{12005} & /nanai{\ng}{\texttoptiebar{t\textbeltl}}un/& \\\\
\\end{tabular}

\\vspace{20pt}\\hline

\\end{nopagebreak}
\\filbreak



\\vspace{15pt}
\\begin{nopagebreak}
\\noindent{\\fontsize{20pt}{10pt}\\textbf{kelumral} } \\textit{42 of autumn/October 12} (N)\\\\
\\noindent {\\tovian \\fontsize{20pt}{10pt} \\textbf{kelumral} }\\\\
\\noindent /kel{\textprimstress}umral/\\\\
\\noindent Syllables: kel.umr.al\\\\
\\noindent lit. branch+darkness\\\\
\\noindent \\textit{Starfennel}\\\\


\\noindent History:

\\vspace{-0pt}
\\hspace{40pt}
\\begin{tabular}{ccc}
\\textit{6000} & /kelubral/&$\\rightarrow$ & \\textit{11000} & /kelumral/& \\\\
\\end{tabular}

\\vspace{20pt}\\hline

\\end{nopagebreak}
\\filbreak



\\vspace{15pt}
\\begin{nopagebreak}
\\noindent{\\fontsize{20pt}{10pt}\\textbf{yisotel} } \\textit{42 of spring/April 12} (N)\\\\
\\noindent {\\tovian \\fontsize{20pt}{10pt} \\textbf{yisotel} }\\\\
\\noindent /yis{\textprimstress}o{\textsubbridge{t}}el/\\\\
\\noindent Syllables: jis.ot.el\\\\
\\noindent lit. moon+lake\\\\
\\noindent \\textit{Moonpool}\\\\


\\noindent History:

\\vspace{-0pt}
\\hspace{40pt}
\\begin{tabular}{ccc}
\\textit{6000} & /aliyso{\textsubbridge{t}}ele/&$\\rightarrow$ & \\textit{6300} & /alyso{\textsubbridge{t}}ele/&$\\rightarrow$ & \\textit{7500} & /lyso{\textsubbridge{t}}ele/&$\\rightarrow$ & \\textit{7501} & /yso{\textsubbridge{t}}ele/&$\\rightarrow$ & \\textit{9500} & /yso{\textsubbridge{t}}el/&$\\rightarrow$ & \\textit{15001} & /yiso{\textsubbridge{t}}el/& \\\\
\\end{tabular}

\\vspace{20pt}\\hline

\\end{nopagebreak}
\\filbreak



\\vspace{15pt}
\\begin{nopagebreak}
\\noindent{\\fontsize{20pt}{10pt}\\textbf{nomen} } \\textit{42 of summer/July 12} (N)\\\\
\\noindent {\\tovian \\fontsize{20pt}{10pt} \\textbf{nomen} }\\\\
\\noindent /n{\textprimstress}omen/\\\\
\\noindent Syllables: nom.en\\\\
\\noindent lit. sea+stone\\\\
\\noindent \\textit{Seaspire}\\\\


\\noindent History:

\\vspace{-0pt}
\\hspace{40pt}
\\begin{tabular}{ccc}
\\textit{6000} & /normern/&$\\rightarrow$ & \\textit{7600} & /nomen/& \\\\
\\end{tabular}

\\vspace{20pt}\\hline

\\end{nopagebreak}
\\filbreak



\\vspace{15pt}
\\begin{nopagebreak}
\\noindent{\\fontsize{20pt}{10pt}\\textbf{???} } \\textit{42 of winter/January 42} (N)\\\\
\\noindent {\\tovian \\fontsize{20pt}{10pt} \\textbf{???} }\\\\
\\noindent /???/\\\\
\\noindent Syllables: ???\\\\
\\noindent lit. ???\\\\
\\noindent \\textit{Namiloray}\\\\


\\noindent History:

\\vspace{-0pt}
\\hspace{40pt}
\\begin{tabular}{ccc}
\\textit{6000} & /???/& \\\\
\\end{tabular}

\\vspace{20pt}\\hline

\\end{nopagebreak}
\\filbreak



\\vspace{15pt}
\\begin{nopagebreak}
\\noindent{\\fontsize{20pt}{10pt}\\textbf{shithol} } \\textit{43 of autumn/October 13} (N)\\\\
\\noindent {\\tovian \\fontsize{20pt}{10pt} \\textbf{shithol} }\\\\
\\noindent /{\textesh}{\textprimstress}i{\texttheta}ol/\\\\
\\noindent Syllables: ʃiθ.ol\\\\
\\noindent lit. sharpness+sweet\\\\
\\noindent \\textit{Thornsweet}\\\\


\\noindent History:

\\vspace{-0pt}
\\hspace{40pt}
\\begin{tabular}{ccc}
\\textit{6000} & /{\textyogh}i{\texttheta}hola/&$\\rightarrow$ & \\textit{9500} & /{\textyogh}i{\texttheta}hol/&$\\rightarrow$ & \\textit{12001} & /{\textesh}i{\texttheta}hol/&$\\rightarrow$ & \\textit{13000} & /{\textesh}i{\texttheta}ol/& \\\\
\\end{tabular}

\\vspace{20pt}\\hline

\\end{nopagebreak}
\\filbreak



\\vspace{15pt}
\\begin{nopagebreak}
\\noindent{\\fontsize{20pt}{10pt}\\textbf{murilaglar} } \\textit{43 of spring/April 13} (N)\\\\
\\noindent {\\tovian \\fontsize{20pt}{10pt} \\textbf{murilaglar} }\\\\
\\noindent /muril{\textprimstress}aglar/\\\\
\\noindent Syllables: mur.il.agl.ar\\\\
\\noindent lit. sight+tree\\\\
\\noindent \\textit{Mrilaple}\\\\


\\noindent History:

\\vspace{-0pt}
\\hspace{40pt}
\\begin{tabular}{ccc}
\\textit{6000} & /emrilaaglare/&$\\rightarrow$ & \\textit{7500} & /mrilaaglare/&$\\rightarrow$ & \\textit{9500} & /mrilaaglar/&$\\rightarrow$ & \\textit{12005} & /mrilaglar/&$\\rightarrow$ & \\textit{15001} & /murilaglar/& \\\\
\\end{tabular}

\\vspace{20pt}\\hline

\\end{nopagebreak}
\\filbreak



\\vspace{15pt}
\\begin{nopagebreak}
\\noindent{\\fontsize{20pt}{10pt}\\textbf{muloanafor} } \\textit{43 of summer/July 13} (N)\\\\
\\noindent {\\tovian \\fontsize{20pt}{10pt} \\textbf{muloanafor} }\\\\
\\noindent /muloan{\textprimstress}afor/\\\\
\\noindent Syllables: mul.o.an.af.or\\\\
\\noindent lit. morning+writing\\\\
\\noindent \\textit{Dawnfelor}\\\\


\\noindent History:

\\vspace{-0pt}
\\hspace{40pt}
\\begin{tabular}{ccc}
\\textit{6000} & /ploanaflore/&$\\rightarrow$ & \\textit{9500} & /ploanaflor/&$\\rightarrow$ & \\textit{11000} & /mloanaflor/&$\\rightarrow$ & \\textit{12500} & /mloanafor/&$\\rightarrow$ & \\textit{15001} & /muloanafor/& \\\\
\\end{tabular}

\\vspace{20pt}\\hline

\\end{nopagebreak}
\\filbreak



\\vspace{15pt}
\\begin{nopagebreak}
\\noindent{\\fontsize{20pt}{10pt}\\textbf{nanador} } \\textit{43 of winter/January 43} (N)\\\\
\\noindent {\\tovian \\fontsize{20pt}{10pt} \\textbf{nanador} }\\\\
\\noindent /nan{\textprimstress}ador/\\\\
\\noindent Syllables: nan.ad.or\\\\
\\noindent lit. star+blue\\\\
\\noindent \\textit{Starwoad}\\\\


\\noindent History:

\\vspace{-0pt}
\\hspace{40pt}
\\begin{tabular}{ccc}
\\textit{6000} & /nnaanadora/&$\\rightarrow$ & \\textit{6500} & /naanadora/&$\\rightarrow$ & \\textit{9500} & /naanador/&$\\rightarrow$ & \\textit{12005} & /nanador/& \\\\
\\end{tabular}

\\vspace{20pt}\\hline

\\end{nopagebreak}
\\filbreak



\\vspace{15pt}
\\begin{nopagebreak}
\\noindent{\\fontsize{20pt}{10pt}\\textbf{munorhanashol} } \\textit{44 of autumn/October 14} (N)\\\\
\\noindent {\\tovian \\fontsize{20pt}{10pt} \\textbf{munorhanashol} }\\\\
\\noindent /munorħan{\textprimstress}a{\textesh}ol/\\\\
\\noindent Syllables: mun.orħ.an.aʃ.ol\\\\
\\noindent lit. dusk+fruit\\\\
\\noindent \\textit{Dusknut}\\\\


\\noindent History:

\\vspace{-0pt}
\\hspace{40pt}
\\begin{tabular}{ccc}
\\textit{6000} & /pnorhana{\textyogh}ola/&$\\rightarrow$ & \\textit{9500} & /pnorhana{\textyogh}ol/&$\\rightarrow$ & \\textit{11000} & /mnorhana{\textyogh}ol/&$\\rightarrow$ & \\textit{11996} & /munorhana{\textyogh}ol/&$\\rightarrow$ & \\textit{12001} & /munorhana{\textesh}ol/&$\\rightarrow$ & \\textit{14006} & /munorħana{\textesh}ol/& \\\\
\\end{tabular}

\\vspace{20pt}\\hline

\\end{nopagebreak}
\\filbreak



\\vspace{15pt}
\\begin{nopagebreak}
\\noindent{\\fontsize{20pt}{10pt}\\textbf{rithilefethr} } \\textit{44 of spring/April 14} (N)\\\\
\\noindent {\\tovian \\fontsize{20pt}{10pt} \\textbf{rithilefethr} }\\\\
\\noindent /ri{\texttheta}il{\textprimstress}efe{\texttheta}r/\\\\
\\noindent Syllables: riθ.il.ef.eθr\\\\
\\noindent lit. sky+speech\\\\
\\noindent \\textit{Skyfethr}\\\\


\\noindent History:

\\vspace{-0pt}
\\hspace{40pt}
\\begin{tabular}{ccc}
\\textit{6000} & /ri{\texttheta}ilefe{\dh}ra/&$\\rightarrow$ & \\textit{9500} & /ri{\texttheta}ilefe{\dh}r/&$\\rightarrow$ & \\textit{12002} & /ri{\texttheta}ilefe{\texttheta}r/& \\\\
\\end{tabular}

\\vspace{20pt}\\hline

\\end{nopagebreak}
\\filbreak



\\vspace{15pt}
\\begin{nopagebreak}
\\noindent{\\fontsize{20pt}{10pt}\\textbf{thangekel} } \\textit{44 of summer/July 14} (N)\\\\
\\noindent {\\tovian \\fontsize{20pt}{10pt} \\textbf{thangekel} }\\\\
\\noindent /{\texttheta}a{\ng}{\textprimstress}ekel/\\\\
\\noindent Syllables: θaŋ.ek.el\\\\
\\noindent lit. leaf+branch\\\\
\\noindent \\textit{Starherb}\\\\


\\noindent History:

\\vspace{-0pt}
\\hspace{40pt}
\\begin{tabular}{ccc}
\\textit{6000} & /a{\dh}a{\ng}gekelu/&$\\rightarrow$ & \\textit{7500} & /{\dh}a{\ng}gekelu/&$\\rightarrow$ & \\textit{8750} & /{\dh}a{\ng}ekelu/&$\\rightarrow$ & \\textit{9500} & /{\dh}a{\ng}ekel/&$\\rightarrow$ & \\textit{12002} & /{\texttheta}a{\ng}ekel/& \\\\
\\end{tabular}

\\vspace{20pt}\\hline

\\end{nopagebreak}
\\filbreak



\\vspace{15pt}
\\begin{nopagebreak}
\\noindent{\\fontsize{20pt}{10pt}\\textbf{shakel} } \\textit{44 of winter/January 44} (N)\\\\
\\noindent {\\tovian \\fontsize{20pt}{10pt} \\textbf{shakel} }\\\\
\\noindent /{\textesh}{\textprimstress}akel/\\\\
\\noindent Syllables: ʃak.el\\\\
\\noindent lit. ice+branch\\\\
\\noindent \\textit{Frostpine}\\\\


\\noindent History:

\\vspace{-0pt}
\\hspace{40pt}
\\begin{tabular}{ccc}
\\textit{6000} & /{\textesh}u{\textesh}akelu/&$\\rightarrow$ & \\textit{9200} & /{\textesh}{\textschwa}{\textesh}akelu/&$\\rightarrow$ & \\textit{9500} & /{\textesh}{\textschwa}{\textesh}akel/&$\\rightarrow$ & \\textit{11990} & /{\textesh}{\textesh}akel/&$\\rightarrow$ & \\textit{12007} & /{\textesh}akel/& \\\\
\\end{tabular}

\\vspace{20pt}\\hline

\\end{nopagebreak}
\\filbreak



\\vspace{15pt}
\\begin{nopagebreak}
\\noindent{\\fontsize{20pt}{10pt}\\textbf{???} } \\textit{45 of autumn/October 15} (N)\\\\
\\noindent {\\tovian \\fontsize{20pt}{10pt} \\textbf{???} }\\\\
\\noindent /???/\\\\
\\noindent Syllables: ???\\\\
\\noindent lit. ???\\\\
\\noindent \\textit{Stamglith}\\\\


\\noindent History:

\\vspace{-0pt}
\\hspace{40pt}
\\begin{tabular}{ccc}
\\textit{6000} & /???/& \\\\
\\end{tabular}

\\vspace{20pt}\\hline

\\end{nopagebreak}
\\filbreak



\\vspace{15pt}
\\begin{nopagebreak}
\\noindent{\\fontsize{20pt}{10pt}\\textbf{lerithitan} } \\textit{45 of spring/April 15} (N)\\\\
\\noindent {\\tovian \\fontsize{20pt}{10pt} \\textbf{lerithitan} }\\\\
\\noindent /leri{\texttheta}{\textprimstress}i{\textsubbridge{t}}an/\\\\
\\noindent Syllables: ler.iθ.it.an\\\\
\\noindent lit. storm+bird\\\\
\\noindent \\textit{Stormbird}\\\\


\\noindent History:

\\vspace{-0pt}
\\hspace{40pt}
\\begin{tabular}{ccc}
\\textit{6000} & /kwlyri{\texttheta}lri{\texttheta}lamn/&$\\rightarrow$ & \\textit{6503} & /kwlyri{\texttheta}rli{\texttheta}lamn/&$\\rightarrow$ & \\textit{7501} & /wlyri{\texttheta}rli{\texttheta}lamn/&$\\rightarrow$ & \\textit{7600} & /lri{\texttheta}li{\texttheta}lan/&$\\rightarrow$ & \\textit{12500} & /lri{\texttheta}i{\texttheta}an/&$\\rightarrow$ & \\textit{12501} & /lri{\texttheta}i{\textsubbridge{t}}an/&$\\rightarrow$ & \\textit{15001} & /leri{\texttheta}i{\textsubbridge{t}}an/& \\\\
\\end{tabular}

\\vspace{20pt}\\hline

\\end{nopagebreak}
\\filbreak



\\vspace{15pt}
\\begin{nopagebreak}
\\noindent{\\fontsize{20pt}{10pt}\\textbf{lhenaysan} } \\textit{45 of summer/July 15} (N)\\\\
\\noindent {\\tovian \\fontsize{20pt}{10pt} \\textbf{lhenaysan} }\\\\
\\noindent /{\textbeltl}en{\textprimstress}aysan/\\\\
\\noindent Syllables: ɬen.ajs.an\\\\
\\noindent lit. dream+moon+animal\\\\
\\noindent \\textit{Moonewe}\\\\


\\noindent History:

\\vspace{-0pt}
\\hspace{40pt}
\\begin{tabular}{ccc}
\\textit{6000} & /{\textbeltl}enaaliys{\textsubbridge{t}}lamn/&$\\rightarrow$ & \\textit{6300} & /{\textbeltl}enaalys{\textsubbridge{t}}lamn/&$\\rightarrow$ & \\textit{6401} & /{\textbeltl}enaalys{\texttoptiebar{t\textbeltl}}amn/&$\\rightarrow$ & \\textit{7600} & /{\textbeltl}enaays{\texttoptiebar{t\textbeltl}}an/&$\\rightarrow$ & \\textit{12005} & /{\textbeltl}enays{\texttoptiebar{t\textbeltl}}an/&$\\rightarrow$ & \\textit{14001} & /{\textbeltl}enaysan/& \\\\
\\end{tabular}

\\vspace{20pt}\\hline

\\end{nopagebreak}
\\filbreak



\\vspace{15pt}
\\begin{nopagebreak}
\\noindent{\\fontsize{20pt}{10pt}\\textbf{rishwen} } \\textit{45 of winter/January 45} (N)\\\\
\\noindent {\\tovian \\fontsize{20pt}{10pt} \\textbf{rishwen} }\\\\
\\noindent /r{\textprimstress}i{\textesh}wen/\\\\
\\noindent Syllables: riʃw.en\\\\
\\noindent lit. kindness+beauty\\\\
\\noindent \\textit{Rashelwen}\\\\


\\noindent History:

\\vspace{-0pt}
\\hspace{40pt}
\\begin{tabular}{ccc}
\\textit{6000} & /ri{\textesh}lwen/&$\\rightarrow$ & \\textit{7600} & /ri{\textesh}wen/& \\\\
\\end{tabular}

\\vspace{20pt}\\hline

\\end{nopagebreak}
\\filbreak



\\vspace{15pt}
\\begin{nopagebreak}
\\noindent{\\fontsize{20pt}{10pt}\\textbf{for} } \\textit{46 of autumn/October 16} (N)\\\\
\\noindent {\\tovian \\fontsize{20pt}{10pt} \\textbf{for} }\\\\
\\noindent /f{\textprimstress}or/\\\\
\\noindent Syllables: for\\\\
\\noindent lit. sun+action\\\\
\\noindent \\textit{Sunzest}\\\\


\\noindent History:

\\vspace{-0pt}
\\hspace{40pt}
\\begin{tabular}{ccc}
\\textit{6000} & /efllore/&$\\rightarrow$ & \\textit{6500} & /eflore/&$\\rightarrow$ & \\textit{7500} & /flore/&$\\rightarrow$ & \\textit{9500} & /flor/&$\\rightarrow$ & \\textit{12500} & /for/& \\\\
\\end{tabular}

\\vspace{20pt}\\hline

\\end{nopagebreak}
\\filbreak



\\vspace{15pt}
\\begin{nopagebreak}
\\noindent{\\fontsize{20pt}{10pt}\\textbf{nanatom} } \\textit{46 of spring/April 16} (N)\\\\
\\noindent {\\tovian \\fontsize{20pt}{10pt} \\textbf{nanatom} }\\\\
\\noindent /nan{\textprimstress}a{\textsubbridge{t}}om/\\\\
\\noindent Syllables: nan.at.om\\\\
\\noindent lit. star+root\\\\
\\noindent \\textit{Starroot}\\\\


\\noindent History:

\\vspace{-0pt}
\\hspace{40pt}
\\begin{tabular}{ccc}
\\textit{6000} & /nnaana{\textsubbridge{t}}omo/&$\\rightarrow$ & \\textit{6500} & /naana{\textsubbridge{t}}omo/&$\\rightarrow$ & \\textit{9500} & /naana{\textsubbridge{t}}om/&$\\rightarrow$ & \\textit{12005} & /nana{\textsubbridge{t}}om/& \\\\
\\end{tabular}

\\vspace{20pt}\\hline

\\end{nopagebreak}
\\filbreak



\\vspace{15pt}
\\begin{nopagebreak}
\\noindent{\\fontsize{20pt}{10pt}\\textbf{lhenalor} } \\textit{46 of summer/July 16} (N)\\\\
\\noindent {\\tovian \\fontsize{20pt}{10pt} \\textbf{lhenalor} }\\\\
\\noindent /{\textbeltl}en{\textprimstress}alor/\\\\
\\noindent Syllables: ɬen.al.or\\\\
\\noindent lit. dream+plant\\\\
\\noindent \\textit{Dreammallow}\\\\


\\noindent History:

\\vspace{-0pt}
\\hspace{40pt}
\\begin{tabular}{ccc}
\\textit{6000} & /{\textbeltl}enamllore/&$\\rightarrow$ & \\textit{6500} & /{\textbeltl}enamlore/&$\\rightarrow$ & \\textit{7600} & /{\textbeltl}enalore/&$\\rightarrow$ & \\textit{9500} & /{\textbeltl}enalor/& \\\\
\\end{tabular}

\\vspace{20pt}\\hline

\\end{nopagebreak}
\\filbreak



\\vspace{15pt}
\\begin{nopagebreak}
\\noindent{\\fontsize{20pt}{10pt}\\textbf{thatlun} } \\textit{46 of winter/January 46} (N)\\\\
\\noindent {\\tovian \\fontsize{20pt}{10pt} \\textbf{thatlun} }\\\\
\\noindent /{\texttheta}{\textprimstress}a{\texttoptiebar{t\textbeltl}}un/\\\\
\\noindent Syllables: θaƛ.un\\\\
\\noindent lit. tool+joining\\\\
\\noindent \\textit{Thornweave}\\\\


\\noindent History:

\\vspace{-0pt}
\\hspace{40pt}
\\begin{tabular}{ccc}
\\textit{6000} & /{\dh}an{\texttoptiebar{t\textbeltl}}una/&$\\rightarrow$ & \\textit{6400} & /{\dh}a{\texttoptiebar{t\textbeltl}}una/&$\\rightarrow$ & \\textit{9500} & /{\dh}a{\texttoptiebar{t\textbeltl}}un/&$\\rightarrow$ & \\textit{12002} & /{\texttheta}a{\texttoptiebar{t\textbeltl}}un/& \\\\
\\end{tabular}

\\vspace{20pt}\\hline

\\end{nopagebreak}
\\filbreak



\\vspace{15pt}
\\begin{nopagebreak}
\\noindent{\\fontsize{20pt}{10pt}\\textbf{lenan} } \\textit{47 of autumn/October 17} (N)\\\\
\\noindent {\\tovian \\fontsize{20pt}{10pt} \\textbf{lenan} }\\\\
\\noindent /l{\textprimstress}enan/\\\\
\\noindent Syllables: len.an\\\\
\\noindent lit. hoe+covering\\\\
\\noindent \\textit{Melorethangor}\\\\


\\noindent History:

\\vspace{-0pt}
\\hspace{40pt}
\\begin{tabular}{ccc}
\\textit{6000} & /mllre{\dh}nbnani/&$\\rightarrow$ & \\textit{6500} & /mlre{\dh}nbnani/&$\\rightarrow$ & \\textit{6503} & /mrle{\dh}nbnani/&$\\rightarrow$ & \\textit{7501} & /rle{\dh}nbnani/&$\\rightarrow$ & \\textit{7600} & /le{\dh}nbnani/&$\\rightarrow$ & \\textit{8500} & /le{\dh}nnani/&$\\rightarrow$ & \\textit{9500} & /le{\dh}nnan/&$\\rightarrow$ & \\textit{12002} & /le{\texttheta}nnan/&$\\rightarrow$ & \\textit{12007} & /le{\texttheta}nan/&$\\rightarrow$ & \\textit{13002} & /lenan/& \\\\
\\end{tabular}

\\vspace{20pt}\\hline

\\end{nopagebreak}
\\filbreak



\\vspace{15pt}
\\begin{nopagebreak}
\\noindent{\\fontsize{20pt}{10pt}\\textbf{shasel} } \\textit{47 of spring/April 17} (N)\\\\
\\noindent {\\tovian \\fontsize{20pt}{10pt} \\textbf{shasel} }\\\\
\\noindent /{\textesh}{\textprimstress}asel/\\\\
\\noindent Syllables: ʃas.el\\\\
\\noindent lit. ice+water\\\\
\\noindent \\textit{Mooncress}\\\\


\\noindent History:

\\vspace{-0pt}
\\hspace{40pt}
\\begin{tabular}{ccc}
\\textit{6000} & /{\textesh}u{\textesh}azel/&$\\rightarrow$ & \\textit{9200} & /{\textesh}{\textschwa}{\textesh}azel/&$\\rightarrow$ & \\textit{11990} & /{\textesh}{\textesh}azel/&$\\rightarrow$ & \\textit{12000} & /{\textesh}{\textesh}asel/&$\\rightarrow$ & \\textit{12007} & /{\textesh}asel/& \\\\
\\end{tabular}

\\vspace{20pt}\\hline

\\end{nopagebreak}
\\filbreak



\\vspace{15pt}
\\begin{nopagebreak}
\\noindent{\\fontsize{20pt}{10pt}\\textbf{shaninakel} } \\textit{47 of summer/July 17} (N)\\\\
\\noindent {\\tovian \\fontsize{20pt}{10pt} \\textbf{shaninakel} }\\\\
\\noindent /{\textesh}anin{\textprimstress}akel/\\\\
\\noindent Syllables: ʃan.in.ak.el\\\\
\\noindent lit. berry+branch\\\\
\\noindent \\textit{Starflax}\\\\


\\noindent History:

\\vspace{-0pt}
\\hspace{40pt}
\\begin{tabular}{ccc}
\\textit{6000} & /{\textyogh}laninakelu/&$\\rightarrow$ & \\textit{9500} & /{\textyogh}laninakel/&$\\rightarrow$ & \\textit{12001} & /{\textesh}laninakel/&$\\rightarrow$ & \\textit{12500} & /{\textesh}aninakel/& \\\\
\\end{tabular}

\\vspace{20pt}\\hline

\\end{nopagebreak}
\\filbreak



\\vspace{15pt}
\\begin{nopagebreak}
\\noindent{\\fontsize{20pt}{10pt}\\textbf{???} } \\textit{47 of winter/January 47} (N)\\\\
\\noindent {\\tovian \\fontsize{20pt}{10pt} \\textbf{???} }\\\\
\\noindent /???/\\\\
\\noindent Syllables: ???\\\\
\\noindent lit. ???\\\\
\\noindent \\textit{Nicypres}\\\\


\\noindent History:

\\vspace{-0pt}
\\hspace{40pt}
\\begin{tabular}{ccc}
\\textit{6000} & /???/& \\\\
\\end{tabular}

\\vspace{20pt}\\hline

\\end{nopagebreak}
\\filbreak



\\vspace{15pt}
\\begin{nopagebreak}
\\noindent{\\fontsize{20pt}{10pt}\\textbf{muloanashith} } \\textit{48 of autumn/October 18} (N)\\\\
\\noindent {\\tovian \\fontsize{20pt}{10pt} \\textbf{muloanashith} }\\\\
\\noindent /muloan{\textprimstress}a{\textesh}i{\texttheta}/\\\\
\\noindent Syllables: mul.o.an.aʃ.iθ\\\\
\\noindent lit. morning+sharpness\\\\
\\noindent \\textit{Dawnthor}\\\\


\\noindent History:

\\vspace{-0pt}
\\hspace{40pt}
\\begin{tabular}{ccc}
\\textit{6000} & /ploana{\textyogh}i{\texttheta}/&$\\rightarrow$ & \\textit{11000} & /mloana{\textyogh}i{\texttheta}/&$\\rightarrow$ & \\textit{12001} & /mloana{\textesh}i{\texttheta}/&$\\rightarrow$ & \\textit{15001} & /muloana{\textesh}i{\texttheta}/& \\\\
\\end{tabular}

\\vspace{20pt}\\hline

\\end{nopagebreak}
\\filbreak



\\vspace{15pt}
\\begin{nopagebreak}
\\noindent{\\fontsize{20pt}{10pt}\\textbf{shagang} } \\textit{48 of spring/April 18} (N)\\\\
\\noindent {\\tovian \\fontsize{20pt}{10pt} \\textbf{shagang} }\\\\
\\noindent /{\textesh}{\textprimstress}aga{\ng}/\\\\
\\noindent Syllables: ʃag.aŋ\\\\
\\noindent lit. ice+tooth\\\\
\\noindent \\textit{Frostdent}\\\\


\\noindent History:

\\vspace{-0pt}
\\hspace{40pt}
\\begin{tabular}{ccc}
\\textit{6000} & /{\textesh}u{\textesh}aga{\ng}ga/&$\\rightarrow$ & \\textit{8750} & /{\textesh}u{\textesh}aga{\ng}a/&$\\rightarrow$ & \\textit{9200} & /{\textesh}{\textschwa}{\textesh}aga{\ng}a/&$\\rightarrow$ & \\textit{9500} & /{\textesh}{\textschwa}{\textesh}aga{\ng}/&$\\rightarrow$ & \\textit{11990} & /{\textesh}{\textesh}aga{\ng}/&$\\rightarrow$ & \\textit{12007} & /{\textesh}aga{\ng}/& \\\\
\\end{tabular}

\\vspace{20pt}\\hline

\\end{nopagebreak}
\\filbreak



\\vspace{15pt}
\\begin{nopagebreak}
\\noindent{\\fontsize{20pt}{10pt}\\textbf{holashol} } \\textit{48 of summer/July 18} (N)\\\\
\\noindent {\\tovian \\fontsize{20pt}{10pt} \\textbf{holashol} }\\\\
\\noindent /ħol{\textprimstress}a{\textesh}ol/\\\\
\\noindent Syllables: ħol.aʃ.ol\\\\
\\noindent lit. sweet+fruit\\\\
\\noindent \\textit{Sweetnut}\\\\


\\noindent History:

\\vspace{-0pt}
\\hspace{40pt}
\\begin{tabular}{ccc}
\\textit{6000} & /hola{\textyogh}ola/&$\\rightarrow$ & \\textit{9500} & /hola{\textyogh}ol/&$\\rightarrow$ & \\textit{12001} & /hola{\textesh}ol/&$\\rightarrow$ & \\textit{14006} & /ħola{\textesh}ol/& \\\\
\\end{tabular}

\\vspace{20pt}\\hline

\\end{nopagebreak}
\\filbreak



\\vspace{15pt}
\\begin{nopagebreak}
\\noindent{\\fontsize{20pt}{10pt}\\textbf{kelumloretush} } \\textit{48 of winter/January 48} (N)\\\\
\\noindent {\\tovian \\fontsize{20pt}{10pt} \\textbf{kelumloretush} }\\\\
\\noindent /kelumlor{\textprimstress}e{\textsubbridge{t}}u{\textesh}/\\\\
\\noindent Syllables: kel.uml.or.et.uʃ\\\\
\\noindent lit. vine+cold\\\\
\\noindent \\textit{Winterivy}\\\\


\\noindent History:

\\vspace{-0pt}
\\hspace{40pt}
\\begin{tabular}{ccc}
\\textit{6000} & /kluwmllore{\textsubbridge{t}}u{\textesh}u/&$\\rightarrow$ & \\textit{6500} & /kluwmlore{\textsubbridge{t}}u{\textesh}u/&$\\rightarrow$ & \\textit{7600} & /klumlore{\textsubbridge{t}}u{\textesh}u/&$\\rightarrow$ & \\textit{9500} & /klumlore{\textsubbridge{t}}u{\textesh}/&$\\rightarrow$ & \\textit{15001} & /kelumlore{\textsubbridge{t}}u{\textesh}/& \\\\
\\end{tabular}

\\vspace{20pt}\\hline

\\end{nopagebreak}
\\filbreak



\\vspace{15pt}
\\begin{nopagebreak}
\\noindent{\\fontsize{20pt}{10pt}\\textbf{varethan} } \\textit{49 of autumn/October 19} (N)\\\\
\\noindent {\\tovian \\fontsize{20pt}{10pt} \\textbf{varethan} }\\\\
\\noindent /var{\textprimstress}e{\texttheta}an/\\\\
\\noindent Syllables: var.eθ.an\\\\
\\noindent lit. value+tool\\\\
\\noindent \\textit{Golace}\\\\


\\noindent History:

\\vspace{-0pt}
\\hspace{40pt}
\\begin{tabular}{ccc}
\\textit{6000} & /vlare{\dh}an/&$\\rightarrow$ & \\textit{12002} & /vlare{\texttheta}an/&$\\rightarrow$ & \\textit{12500} & /vare{\texttheta}an/& \\\\
\\end{tabular}

\\vspace{20pt}\\hline

\\end{nopagebreak}
\\filbreak



\\vspace{15pt}
\\begin{nopagebreak}
\\noindent{\\fontsize{20pt}{10pt}\\textbf{nanashol} } \\textit{49 of spring/April 19} (N)\\\\
\\noindent {\\tovian \\fontsize{20pt}{10pt} \\textbf{nanashol} }\\\\
\\noindent /nan{\textprimstress}a{\textesh}ol/\\\\
\\noindent Syllables: nan.aʃ.ol\\\\
\\noindent lit. star+fruit\\\\
\\noindent \\textit{Starfruit}\\\\


\\noindent History:

\\vspace{-0pt}
\\hspace{40pt}
\\begin{tabular}{ccc}
\\textit{6000} & /nnaana{\textyogh}ola/&$\\rightarrow$ & \\textit{6500} & /naana{\textyogh}ola/&$\\rightarrow$ & \\textit{9500} & /naana{\textyogh}ol/&$\\rightarrow$ & \\textit{12001} & /naana{\textesh}ol/&$\\rightarrow$ & \\textit{12005} & /nana{\textesh}ol/& \\\\
\\end{tabular}

\\vspace{20pt}\\hline

\\end{nopagebreak}
\\filbreak



\\vspace{15pt}
\\begin{nopagebreak}
\\noindent{\\fontsize{20pt}{10pt}\\textbf{dolaglar} } \\textit{49 of summer/July 19} (N)\\\\
\\noindent {\\tovian \\fontsize{20pt}{10pt} \\textbf{dolaglar} }\\\\
\\noindent /dol{\textprimstress}aglar/\\\\
\\noindent Syllables: dol.agl.ar\\\\
\\noindent lit. death+tree\\\\
\\noindent \\textit{Twilightroot}\\\\


\\noindent History:

\\vspace{-0pt}
\\hspace{40pt}
\\begin{tabular}{ccc}
\\textit{6000} & /dolaglare/&$\\rightarrow$ & \\textit{9500} & /dolaglar/& \\\\
\\end{tabular}

\\vspace{20pt}\\hline

\\end{nopagebreak}
\\filbreak



\\vspace{15pt}
\\begin{nopagebreak}
\\noindent{\\fontsize{20pt}{10pt}\\textbf{aymen} } \\textit{49 of winter/January 49} (N)\\\\
\\noindent {\\tovian \\fontsize{20pt}{10pt} \\textbf{aymen} }\\\\
\\noindent /{\textprimstress}aymen/\\\\
\\noindent Syllables: ajm.en\\\\
\\noindent lit. moon+earth\\\\
\\noindent \\textit{Moonjuniper}\\\\


\\noindent History:

\\vspace{-0pt}
\\hspace{40pt}
\\begin{tabular}{ccc}
\\textit{6000} & /aliysmern/&$\\rightarrow$ & \\textit{6300} & /alysmern/&$\\rightarrow$ & \\textit{7600} & /aysmen/&$\\rightarrow$ & \\textit{13002} & /aymen/& \\\\
\\end{tabular}

\\vspace{20pt}\\hline

\\end{nopagebreak}
\\filbreak



\\vspace{15pt}
\\begin{nopagebreak}
\\noindent{\\fontsize{20pt}{10pt}\\textbf{nanamen} } \\textit{5 of autumn/September 5} (N)\\\\
\\noindent {\\tovian \\fontsize{20pt}{10pt} \\textbf{nanamen} }\\\\
\\noindent /nan{\textprimstress}amen/\\\\
\\noindent Syllables: nan.am.en\\\\
\\noindent lit. star+earth\\\\
\\noindent \\textit{Starmul}\\\\


\\noindent History:

\\vspace{-0pt}
\\hspace{40pt}
\\begin{tabular}{ccc}
\\textit{6000} & /nnaanamern/&$\\rightarrow$ & \\textit{6500} & /naanamern/&$\\rightarrow$ & \\textit{7600} & /naanamen/&$\\rightarrow$ & \\textit{12005} & /nanamen/& \\\\
\\end{tabular}

\\vspace{20pt}\\hline

\\end{nopagebreak}
\\filbreak



\\vspace{15pt}
\\begin{nopagebreak}
\\noindent{\\fontsize{20pt}{10pt}\\textbf{lerithan} } \\textit{5 of spring/March 5} (N)\\\\
\\noindent {\\tovian \\fontsize{20pt}{10pt} \\textbf{lerithan} }\\\\
\\noindent /ler{\textprimstress}i{\texttheta}an/\\\\
\\noindent Syllables: ler.iθ.an\\\\
\\noindent lit. wind+animal\\\\
\\noindent \\textit{Vindsteed}\\\\


\\noindent History:

\\vspace{-0pt}
\\hspace{40pt}
\\begin{tabular}{ccc}
\\textit{6000} & /kwlyri{\texttheta}{\textsubbridge{t}}lamn/&$\\rightarrow$ & \\textit{6401} & /kwlyri{\texttheta}{\texttoptiebar{t\textbeltl}}amn/&$\\rightarrow$ & \\textit{7501} & /wlyri{\texttheta}{\texttoptiebar{t\textbeltl}}amn/&$\\rightarrow$ & \\textit{7600} & /lri{\texttheta}{\texttoptiebar{t\textbeltl}}an/&$\\rightarrow$ & \\textit{14001} & /lri{\texttheta}an/&$\\rightarrow$ & \\textit{15001} & /leri{\texttheta}an/& \\\\
\\end{tabular}

\\vspace{20pt}\\hline

\\end{nopagebreak}
\\filbreak



\\vspace{15pt}
\\begin{nopagebreak}
\\noindent{\\fontsize{20pt}{10pt}\\textbf{muloanatlan} } \\textit{5 of summer/June 5} (N)\\\\
\\noindent {\\tovian \\fontsize{20pt}{10pt} \\textbf{muloanatlan} }\\\\
\\noindent /muloan{\textprimstress}a{\texttoptiebar{t\textbeltl}}an/\\\\
\\noindent Syllables: mul.o.an.aƛ.an\\\\
\\noindent lit. morning+animal\\\\
\\noindent \\textit{Dawncluck}\\\\


\\noindent History:

\\vspace{-0pt}
\\hspace{40pt}
\\begin{tabular}{ccc}
\\textit{6000} & /ploana{\textsubbridge{t}}lamn/&$\\rightarrow$ & \\textit{6401} & /ploana{\texttoptiebar{t\textbeltl}}amn/&$\\rightarrow$ & \\textit{7600} & /ploana{\texttoptiebar{t\textbeltl}}an/&$\\rightarrow$ & \\textit{11000} & /mloana{\texttoptiebar{t\textbeltl}}an/&$\\rightarrow$ & \\textit{15001} & /muloana{\texttoptiebar{t\textbeltl}}an/& \\\\
\\end{tabular}

\\vspace{20pt}\\hline

\\end{nopagebreak}
\\filbreak



\\vspace{15pt}
\\begin{nopagebreak}
\\noindent{\\fontsize{20pt}{10pt}\\textbf{lerithan} } \\textit{5 of winter/december 5} (N)\\\\
\\noindent {\\tovian \\fontsize{20pt}{10pt} \\textbf{lerithan} }\\\\
\\noindent /ler{\textprimstress}i{\texttheta}an/\\\\
\\noindent Syllables: ler.iθ.an\\\\
\\noindent lit. breath+tool\\\\
\\noindent \\textit{Windhorn}\\\\


\\noindent History:

\\vspace{-0pt}
\\hspace{40pt}
\\begin{tabular}{ccc}
\\textit{6000} & /kwlyri{\texttheta}{\dh}an/&$\\rightarrow$ & \\textit{7501} & /wlyri{\texttheta}{\dh}an/&$\\rightarrow$ & \\textit{7600} & /lri{\texttheta}{\dh}an/&$\\rightarrow$ & \\textit{12002} & /lri{\texttheta}{\texttheta}an/&$\\rightarrow$ & \\textit{12007} & /lri{\texttheta}an/&$\\rightarrow$ & \\textit{15001} & /leri{\texttheta}an/& \\\\
\\end{tabular}

\\vspace{20pt}\\hline

\\end{nopagebreak}
\\filbreak



\\vspace{15pt}
\\begin{nopagebreak}
\\noindent{\\fontsize{20pt}{10pt}\\textbf{harathan} } \\textit{50 of autumn/October 20} (N)\\\\
\\noindent {\\tovian \\fontsize{20pt}{10pt} \\textbf{harathan} }\\\\
\\noindent /ħar{\textprimstress}a{\texttheta}an/\\\\
\\noindent Syllables: ħar.aθ.an\\\\
\\noindent lit. pain+tool\\\\
\\noindent \\textit{Harsack}\\\\


\\noindent History:

\\vspace{-0pt}
\\hspace{40pt}
\\begin{tabular}{ccc}
\\textit{6000} & /hara{\dh}an/&$\\rightarrow$ & \\textit{12002} & /hara{\texttheta}an/&$\\rightarrow$ & \\textit{14006} & /ħara{\texttheta}an/& \\\\
\\end{tabular}

\\vspace{20pt}\\hline

\\end{nopagebreak}
\\filbreak



\\vspace{15pt}
\\begin{nopagebreak}
\\noindent{\\fontsize{20pt}{10pt}\\textbf{nanamnan} } \\textit{50 of spring/April 20} (N)\\\\
\\noindent {\\tovian \\fontsize{20pt}{10pt} \\textbf{nanamnan} }\\\\
\\noindent /nan{\textprimstress}amnan/\\\\
\\noindent Syllables: nan.amn.an\\\\
\\noindent lit. star+covering\\\\
\\noindent \\textit{Starharrow}\\\\


\\noindent History:

\\vspace{-0pt}
\\hspace{40pt}
\\begin{tabular}{ccc}
\\textit{6000} & /nnaanabnani/&$\\rightarrow$ & \\textit{6500} & /naanabnani/&$\\rightarrow$ & \\textit{9500} & /naanabnan/&$\\rightarrow$ & \\textit{11000} & /naanamnan/&$\\rightarrow$ & \\textit{12005} & /nanamnan/& \\\\
\\end{tabular}

\\vspace{20pt}\\hline

\\end{nopagebreak}
\\filbreak



\\vspace{15pt}
\\begin{nopagebreak}
\\noindent{\\fontsize{20pt}{10pt}\\textbf{selyethil} } \\textit{50 of summer/July 20} (N)\\\\
\\noindent {\\tovian \\fontsize{20pt}{10pt} \\textbf{selyethil} }\\\\
\\noindent /sely{\textprimstress}e{\texttheta}il/\\\\
\\noindent Syllables: selj.eθ.il\\\\
\\noindent lit. water+home\\\\
\\noindent \\textit{Wetgate}\\\\


\\noindent History:

\\vspace{-0pt}
\\hspace{40pt}
\\begin{tabular}{ccc}
\\textit{6000} & /zeldryel{\texttheta}ile/&$\\rightarrow$ & \\textit{7600} & /zeldyel{\texttheta}ile/&$\\rightarrow$ & \\textit{8750} & /zelyel{\texttheta}ile/&$\\rightarrow$ & \\textit{9500} & /zelyel{\texttheta}il/&$\\rightarrow$ & \\textit{12000} & /selyel{\texttheta}il/&$\\rightarrow$ & \\textit{12502} & /selye{\texttheta}il/& \\\\
\\end{tabular}

\\vspace{20pt}\\hline

\\end{nopagebreak}
\\filbreak



\\vspace{15pt}
\\begin{nopagebreak}
\\noindent{\\fontsize{20pt}{10pt}\\textbf{lethiledrem} } \\textit{51 of autumn/October 21} (N)\\\\
\\noindent {\\tovian \\fontsize{20pt}{10pt} \\textbf{lethiledrem} }\\\\
\\noindent /le{\texttheta}il{\textprimstress}edrem/\\\\
\\noindent Syllables: leθ.il.edr.em\\\\
\\noindent lit. forest+flower\\\\
\\noindent \\textit{Wildrose}\\\\


\\noindent History:

\\vspace{-0pt}
\\hspace{40pt}
\\begin{tabular}{ccc}
\\textit{6000} & /aglre{\texttheta}iledremp/&$\\rightarrow$ & \\textit{6503} & /agrle{\texttheta}iledremp/&$\\rightarrow$ & \\textit{7500} & /grle{\texttheta}iledremp/&$\\rightarrow$ & \\textit{7501} & /rle{\texttheta}iledremp/&$\\rightarrow$ & \\textit{7600} & /le{\texttheta}iledremp/&$\\rightarrow$ & \\textit{8750} & /le{\texttheta}iledrem/& \\\\
\\end{tabular}

\\vspace{20pt}\\hline

\\end{nopagebreak}
\\filbreak



\\vspace{15pt}
\\begin{nopagebreak}
\\noindent{\\fontsize{20pt}{10pt}\\textbf{shimisel} } \\textit{51 of spring/April 21} (N)\\\\
\\noindent {\\tovian \\fontsize{20pt}{10pt} \\textbf{shimisel} }\\\\
\\noindent /{\textesh}im{\textprimstress}isel/\\\\
\\noindent Syllables: ʃim.is.el\\\\
\\noindent lit. mystery+water\\\\
\\noindent \\textit{Mystbach}\\\\


\\noindent History:

\\vspace{-0pt}
\\hspace{40pt}
\\begin{tabular}{ccc}
\\textit{6000} & /{\textesh}imizel/&$\\rightarrow$ & \\textit{12000} & /{\textesh}imisel/& \\\\
\\end{tabular}

\\vspace{20pt}\\hline

\\end{nopagebreak}
\\filbreak



\\vspace{15pt}
\\begin{nopagebreak}
\\noindent{\\fontsize{20pt}{10pt}\\textbf{felakorim} } \\textit{51 of summer/July 21} (N)\\\\
\\noindent {\\tovian \\fontsize{20pt}{10pt} \\textbf{felakorim} }\\\\
\\noindent /felak{\textprimstress}orim/\\\\
\\noindent Syllables: fel.ak.or.im\\\\
\\noindent lit. sun+cycle\\\\
\\noindent \\textit{Sunwheel}\\\\


\\noindent History:

\\vspace{-0pt}
\\hspace{40pt}
\\begin{tabular}{ccc}
\\textit{6000} & /eflelakorim/&$\\rightarrow$ & \\textit{7500} & /flelakorim/&$\\rightarrow$ & \\textit{12500} & /felakorim/& \\\\
\\end{tabular}

\\vspace{20pt}\\hline

\\end{nopagebreak}
\\filbreak



\\vspace{15pt}
\\begin{nopagebreak}
\\noindent{\\fontsize{20pt}{10pt}\\textbf{nanaglar} } \\textit{51 of winter/January 51} (N)\\\\
\\noindent {\\tovian \\fontsize{20pt}{10pt} \\textbf{nanaglar} }\\\\
\\noindent /nan{\textprimstress}aglar/\\\\
\\noindent Syllables: nan.agl.ar\\\\
\\noindent lit. star+tree\\\\
\\noindent \\textit{Starmaple}\\\\


\\noindent History:

\\vspace{-0pt}
\\hspace{40pt}
\\begin{tabular}{ccc}
\\textit{6000} & /nnaanaaglare/&$\\rightarrow$ & \\textit{6500} & /naanaaglare/&$\\rightarrow$ & \\textit{9500} & /naanaaglar/&$\\rightarrow$ & \\textit{12005} & /nanaglar/& \\\\
\\end{tabular}

\\vspace{20pt}\\hline

\\end{nopagebreak}
\\filbreak



\\vspace{15pt}
\\begin{nopagebreak}
\\noindent{\\fontsize{20pt}{10pt}\\textbf{lhenaysol} } \\textit{52 of autumn/October 22} (N)\\\\
\\noindent {\\tovian \\fontsize{20pt}{10pt} \\textbf{lhenaysol} }\\\\
\\noindent /{\textbeltl}en{\textprimstress}aysol/\\\\
\\noindent Syllables: ɬen.ajs.ol\\\\
\\noindent lit. dream+moon+fruit\\\\
\\noindent \\textit{Monut}\\\\


\\noindent History:

\\vspace{-0pt}
\\hspace{40pt}
\\begin{tabular}{ccc}
\\textit{6000} & /{\textbeltl}enaaliys{\textyogh}ola/&$\\rightarrow$ & \\textit{6300} & /{\textbeltl}enaalys{\textyogh}ola/&$\\rightarrow$ & \\textit{7600} & /{\textbeltl}enaays{\textyogh}ola/&$\\rightarrow$ & \\textit{9500} & /{\textbeltl}enaays{\textyogh}ol/&$\\rightarrow$ & \\textit{12001} & /{\textbeltl}enaays{\textesh}ol/&$\\rightarrow$ & \\textit{12005} & /{\textbeltl}enays{\textesh}ol/&$\\rightarrow$ & \\textit{13000} & /{\textbeltl}enaysol/& \\\\
\\end{tabular}

\\vspace{20pt}\\hline

\\end{nopagebreak}
\\filbreak



\\vspace{15pt}
\\begin{nopagebreak}
\\noindent{\\fontsize{20pt}{10pt}\\textbf{muralshanin} } \\textit{52 of spring/April 22} (N)\\\\
\\noindent {\\tovian \\fontsize{20pt}{10pt} \\textbf{muralshanin} }\\\\
\\noindent /mural{\textesh}{\textprimstress}anin/\\\\
\\noindent Syllables: mur.alʃ.an.in\\\\
\\noindent lit. darkness+berry\\\\
\\noindent \\textit{Nightberry}\\\\


\\noindent History:

\\vspace{-0pt}
\\hspace{40pt}
\\begin{tabular}{ccc}
\\textit{6000} & /bral{\textyogh}lanina/&$\\rightarrow$ & \\textit{9500} & /bral{\textyogh}lanin/&$\\rightarrow$ & \\textit{11000} & /mral{\textyogh}lanin/&$\\rightarrow$ & \\textit{12001} & /mral{\textesh}lanin/&$\\rightarrow$ & \\textit{12500} & /mral{\textesh}anin/&$\\rightarrow$ & \\textit{15001} & /mural{\textesh}anin/& \\\\
\\end{tabular}

\\vspace{20pt}\\hline

\\end{nopagebreak}
\\filbreak



\\vspace{15pt}
\\begin{nopagebreak}
\\noindent{\\fontsize{20pt}{10pt}\\textbf{muralshol} } \\textit{52 of summer/July 22} (N)\\\\
\\noindent {\\tovian \\fontsize{20pt}{10pt} \\textbf{muralshol} }\\\\
\\noindent /mur{\textprimstress}al{\textesh}ol/\\\\
\\noindent Syllables: mur.alʃ.ol\\\\
\\noindent lit. dark+fruit\\\\
\\noindent \\textit{Duskberry}\\\\


\\noindent History:

\\vspace{-0pt}
\\hspace{40pt}
\\begin{tabular}{ccc}
\\textit{6000} & /bral{\textyogh}ola/&$\\rightarrow$ & \\textit{9500} & /bral{\textyogh}ol/&$\\rightarrow$ & \\textit{11000} & /mral{\textyogh}ol/&$\\rightarrow$ & \\textit{12001} & /mral{\textesh}ol/&$\\rightarrow$ & \\textit{15001} & /mural{\textesh}ol/& \\\\
\\end{tabular}

\\vspace{20pt}\\hline

\\end{nopagebreak}
\\filbreak



\\vspace{15pt}
\\begin{nopagebreak}
\\noindent{\\fontsize{20pt}{10pt}\\textbf{lanalor} } \\textit{52 of winter/January 52} (N)\\\\
\\noindent {\\tovian \\fontsize{20pt}{10pt} \\textbf{lanalor} }\\\\
\\noindent /lan{\textprimstress}alor/\\\\
\\noindent Syllables: lan.al.or\\\\
\\noindent lit. twilight+plant\\\\
\\noindent \\textit{Duskheather}\\\\


\\noindent History:

\\vspace{-0pt}
\\hspace{40pt}
\\begin{tabular}{ccc}
\\textit{6000} & /brlanamllore/&$\\rightarrow$ & \\textit{6500} & /brlanamlore/&$\\rightarrow$ & \\textit{7501} & /rlanamlore/&$\\rightarrow$ & \\textit{7600} & /lanalore/&$\\rightarrow$ & \\textit{9500} & /lanalor/& \\\\
\\end{tabular}

\\vspace{20pt}\\hline

\\end{nopagebreak}
\\filbreak



\\vspace{15pt}
\\begin{nopagebreak}
\\noindent{\\fontsize{20pt}{10pt}\\textbf{thamlor} } \\textit{53 of autumn/October 23} (N)\\\\
\\noindent {\\tovian \\fontsize{20pt}{10pt} \\textbf{thamlor} }\\\\
\\noindent /{\texttheta}{\textprimstress}amlor/\\\\
\\noindent Syllables: θaml.or\\\\
\\noindent lit. tool+plant\\\\
\\noindent \\textit{Brewvine}\\\\


\\noindent History:

\\vspace{-0pt}
\\hspace{40pt}
\\begin{tabular}{ccc}
\\textit{6000} & /{\dh}anmllore/&$\\rightarrow$ & \\textit{6500} & /{\dh}anmlore/&$\\rightarrow$ & \\textit{7600} & /{\dh}amlore/&$\\rightarrow$ & \\textit{9500} & /{\dh}amlor/&$\\rightarrow$ & \\textit{12002} & /{\texttheta}amlor/& \\\\
\\end{tabular}

\\vspace{20pt}\\hline

\\end{nopagebreak}
\\filbreak



\\vspace{15pt}
\\begin{nopagebreak}
\\noindent{\\fontsize{20pt}{10pt}\\textbf{femlor} } \\textit{53 of spring/April 23} (N)\\\\
\\noindent {\\tovian \\fontsize{20pt}{10pt} \\textbf{femlor} }\\\\
\\noindent /f{\textprimstress}emlor/\\\\
\\noindent Syllables: feml.or\\\\
\\noindent lit. speaking+plant\\\\
\\noindent \\textit{Duskmadder}\\\\


\\noindent History:

\\vspace{-0pt}
\\hspace{40pt}
\\begin{tabular}{ccc}
\\textit{6000} & /fe{\texttheta}rmllore/&$\\rightarrow$ & \\textit{6500} & /fe{\texttheta}rmlore/&$\\rightarrow$ & \\textit{7600} & /fe{\texttheta}mlore/&$\\rightarrow$ & \\textit{9500} & /fe{\texttheta}mlor/&$\\rightarrow$ & \\textit{13002} & /femlor/& \\\\
\\end{tabular}

\\vspace{20pt}\\hline

\\end{nopagebreak}
\\filbreak



\\vspace{15pt}
\\begin{nopagebreak}
\\noindent{\\fontsize{20pt}{10pt}\\textbf{yisayeol} } \\textit{53 of summer/July 23} (N)\\\\
\\noindent {\\tovian \\fontsize{20pt}{10pt} \\textbf{yisayeol} }\\\\
\\noindent /yisay{\textprimstress}eol/\\\\
\\noindent Syllables: jis.aj.e.ol\\\\
\\noindent lit. moon+wisdom\\\\
\\noindent \\textit{Moonlent}\\\\


\\noindent History:

\\vspace{-0pt}
\\hspace{40pt}
\\begin{tabular}{ccc}
\\textit{6000} & /aliysaryeol/&$\\rightarrow$ & \\textit{6300} & /alysaryeol/&$\\rightarrow$ & \\textit{7500} & /lysaryeol/&$\\rightarrow$ & \\textit{7501} & /ysaryeol/&$\\rightarrow$ & \\textit{7600} & /ysayeol/&$\\rightarrow$ & \\textit{15001} & /yisayeol/& \\\\
\\end{tabular}

\\vspace{20pt}\\hline

\\end{nopagebreak}
\\filbreak



\\vspace{15pt}
\\begin{nopagebreak}
\\noindent{\\fontsize{20pt}{10pt}\\textbf{shalor} } \\textit{53 of winter/January 53} (N)\\\\
\\noindent {\\tovian \\fontsize{20pt}{10pt} \\textbf{shalor} }\\\\
\\noindent /{\textesh}{\textprimstress}alor/\\\\
\\noindent Syllables: ʃal.or\\\\
\\noindent lit. ice+plant\\\\
\\noindent \\textit{Frostreed}\\\\


\\noindent History:

\\vspace{-0pt}
\\hspace{40pt}
\\begin{tabular}{ccc}
\\textit{6000} & /{\textesh}u{\textesh}amllore/&$\\rightarrow$ & \\textit{6500} & /{\textesh}u{\textesh}amlore/&$\\rightarrow$ & \\textit{7600} & /{\textesh}u{\textesh}alore/&$\\rightarrow$ & \\textit{9200} & /{\textesh}{\textschwa}{\textesh}alore/&$\\rightarrow$ & \\textit{9500} & /{\textesh}{\textschwa}{\textesh}alor/&$\\rightarrow$ & \\textit{11990} & /{\textesh}{\textesh}alor/&$\\rightarrow$ & \\textit{12007} & /{\textesh}alor/& \\\\
\\end{tabular}

\\vspace{20pt}\\hline

\\end{nopagebreak}
\\filbreak



\\vspace{15pt}
\\begin{nopagebreak}
\\noindent{\\fontsize{20pt}{10pt}\\textbf{meloresyalethavar} } \\textit{54 of autumn/October 24} (N)\\\\
\\noindent {\\tovian \\fontsize{20pt}{10pt} \\textbf{meloresyalethavar} }\\\\
\\noindent /meloresyale{\texttheta}{\textprimstress}avar/\\\\
\\noindent Syllables: mel.or.esj.al.eθ.av.ar\\\\
\\noindent lit. grain+value\\\\
\\noindent \\textit{Goldgrain}\\\\


\\noindent History:

\\vspace{-0pt}
\\hspace{40pt}
\\begin{tabular}{ccc}
\\textit{6000} & /mlloresyale{\dh}avlare/&$\\rightarrow$ & \\textit{6500} & /mloresyale{\dh}avlare/&$\\rightarrow$ & \\textit{7502} & /meloresyale{\dh}avlare/&$\\rightarrow$ & \\textit{9500} & /meloresyale{\dh}avlar/&$\\rightarrow$ & \\textit{12002} & /meloresyale{\texttheta}avlar/&$\\rightarrow$ & \\textit{12500} & /meloresyale{\texttheta}avar/& \\\\
\\end{tabular}

\\vspace{20pt}\\hline

\\end{nopagebreak}
\\filbreak



\\vspace{15pt}
\\begin{nopagebreak}
\\noindent{\\fontsize{20pt}{10pt}\\textbf{lunum} } \\textit{54 of spring/April 24} (N)\\\\
\\noindent {\\tovian \\fontsize{20pt}{10pt} \\textbf{lunum} }\\\\
\\noindent /l{\textprimstress}unum/\\\\
\\noindent Syllables: lun.um\\\\
\\noindent lit. sun+TEMP+present\\\\
\\noindent \\textit{Sunorb}\\\\


\\noindent History:

\\vspace{-0pt}
\\hspace{40pt}
\\begin{tabular}{ccc}
\\textit{6000} & /efl{\textsubbridge{t}}u-num{\textbeltl}le/&$\\rightarrow$ & \\textit{7500} & /fl{\textsubbridge{t}}u-num{\textbeltl}le/&$\\rightarrow$ & \\textit{7501} & /l{\textsubbridge{t}}u-num{\textbeltl}le/&$\\rightarrow$ & \\textit{8500} & /l{\textsubbridge{t}}unum{\textbeltl}le/&$\\rightarrow$ & \\textit{8750} & /lunum{\textbeltl}le/&$\\rightarrow$ & \\textit{9500} & /lunum{\textbeltl}l/&$\\rightarrow$ & \\textit{12500} & /lunum{\textbeltl}/&$\\rightarrow$ & \\textit{13001} & /lunum/& \\\\
\\end{tabular}

\\vspace{20pt}\\hline

\\end{nopagebreak}
\\filbreak



\\vspace{15pt}
\\begin{nopagebreak}
\\noindent{\\fontsize{20pt}{10pt}\\textbf{kelulor} } \\textit{54 of summer/July 24} (N)\\\\
\\noindent {\\tovian \\fontsize{20pt}{10pt} \\textbf{kelulor} }\\\\
\\noindent /kel{\textprimstress}ulor/\\\\
\\noindent Syllables: kel.ul.or\\\\
\\noindent lit. branch+plant\\\\
\\noindent \\textit{Staranise}\\\\


\\noindent History:

\\vspace{-0pt}
\\hspace{40pt}
\\begin{tabular}{ccc}
\\textit{6000} & /kelumllore/&$\\rightarrow$ & \\textit{6500} & /kelumlore/&$\\rightarrow$ & \\textit{7600} & /kelulore/&$\\rightarrow$ & \\textit{9500} & /kelulor/& \\\\
\\end{tabular}

\\vspace{20pt}\\hline

\\end{nopagebreak}
\\filbreak



\\vspace{15pt}
\\begin{nopagebreak}
\\noindent{\\fontsize{20pt}{10pt}\\textbf{kelusyaleth} } \\textit{54 of winter/January 54} (N)\\\\
\\noindent {\\tovian \\fontsize{20pt}{10pt} \\textbf{kelusyaleth} }\\\\
\\noindent /kelusy{\textprimstress}ale{\texttheta}/\\\\
\\noindent Syllables: kel.usj.al.eθ\\\\
\\noindent lit. branch+food\\\\
\\noindent \\textit{Startherb}\\\\


\\noindent History:

\\vspace{-0pt}
\\hspace{40pt}
\\begin{tabular}{ccc}
\\textit{6000} & /kelusyale{\dh}a/&$\\rightarrow$ & \\textit{9500} & /kelusyale{\dh}/&$\\rightarrow$ & \\textit{12002} & /kelusyale{\texttheta}/& \\\\
\\end{tabular}

\\vspace{20pt}\\hline

\\end{nopagebreak}
\\filbreak



\\vspace{15pt}
\\begin{nopagebreak}
\\noindent{\\fontsize{20pt}{10pt}\\textbf{shotlesh} } \\textit{55 of autumn/October 25} (N)\\\\
\\noindent {\\tovian \\fontsize{20pt}{10pt} \\textbf{shotlesh} }\\\\
\\noindent /{\textesh}{\textprimstress}o{\texttoptiebar{t\textbeltl}}e{\textesh}/\\\\
\\noindent Syllables: ʃoƛ.eʃ\\\\
\\noindent lit. river+sound\\\\
\\noindent \\textit{Rivercaw}\\\\


\\noindent History:

\\vspace{-0pt}
\\hspace{40pt}
\\begin{tabular}{ccc}
\\textit{6000} & /{\textesh}o{\texttoptiebar{t\textbeltl}}zl{\textsubbridge{t}}e{\textyogh}/&$\\rightarrow$ & \\textit{8750} & /{\textesh}o{\texttoptiebar{t\textbeltl}}zle{\textyogh}/&$\\rightarrow$ & \\textit{12000} & /{\textesh}o{\texttoptiebar{t\textbeltl}}sle{\textyogh}/&$\\rightarrow$ & \\textit{12001} & /{\textesh}o{\texttoptiebar{t\textbeltl}}sle{\textesh}/&$\\rightarrow$ & \\textit{12500} & /{\textesh}o{\texttoptiebar{t\textbeltl}}se{\textesh}/&$\\rightarrow$ & \\textit{14000} & /{\textesh}o{\texttoptiebar{t\textbeltl}}e{\textesh}/& \\\\
\\end{tabular}

\\vspace{20pt}\\hline

\\end{nopagebreak}
\\filbreak



\\vspace{15pt}
\\begin{nopagebreak}
\\noindent{\\fontsize{20pt}{10pt}\\textbf{rithiletlan} } \\textit{55 of spring/April 25} (N)\\\\
\\noindent {\\tovian \\fontsize{20pt}{10pt} \\textbf{rithiletlan} }\\\\
\\noindent /ri{\texttheta}il{\textprimstress}e{\texttoptiebar{t\textbeltl}}an/\\\\
\\noindent Syllables: riθ.il.eƛ.an\\\\
\\noindent lit. sky+animal\\\\
\\noindent \\textit{Skyfowl}\\\\


\\noindent History:

\\vspace{-0pt}
\\hspace{40pt}
\\begin{tabular}{ccc}
\\textit{6000} & /ri{\texttheta}ile{\textsubbridge{t}}lamn/&$\\rightarrow$ & \\textit{6401} & /ri{\texttheta}ile{\texttoptiebar{t\textbeltl}}amn/&$\\rightarrow$ & \\textit{7600} & /ri{\texttheta}ile{\texttoptiebar{t\textbeltl}}an/& \\\\
\\end{tabular}

\\vspace{20pt}\\hline

\\end{nopagebreak}
\\filbreak



\\vspace{15pt}
\\begin{nopagebreak}
\\noindent{\\fontsize{20pt}{10pt}\\textbf{sethan} } \\textit{55 of summer/July 25} (N)\\\\
\\noindent {\\tovian \\fontsize{20pt}{10pt} \\textbf{sethan} }\\\\
\\noindent /s{\textprimstress}e{\texttheta}an/\\\\
\\noindent Syllables: seθ.an\\\\
\\noindent lit. water+tool\\\\
\\noindent \\textit{Waterwhisker}\\\\


\\noindent History:

\\vspace{-0pt}
\\hspace{40pt}
\\begin{tabular}{ccc}
\\textit{6000} & /zel{\dh}an/&$\\rightarrow$ & \\textit{12000} & /sel{\dh}an/&$\\rightarrow$ & \\textit{12002} & /sel{\texttheta}an/&$\\rightarrow$ & \\textit{12502} & /se{\texttheta}an/& \\\\
\\end{tabular}

\\vspace{20pt}\\hline

\\end{nopagebreak}
\\filbreak



\\vspace{15pt}
\\begin{nopagebreak}
\\noindent{\\fontsize{20pt}{10pt}\\textbf{muralesh} } \\textit{55 of winter/January 55} (N)\\\\
\\noindent {\\tovian \\fontsize{20pt}{10pt} \\textbf{muralesh} }\\\\
\\noindent /mur{\textprimstress}ale{\textesh}/\\\\
\\noindent Syllables: mur.al.eʃ\\\\
\\noindent lit. darkness+sound\\\\
\\noindent \\textit{Nightchirp}\\\\


\\noindent History:

\\vspace{-0pt}
\\hspace{40pt}
\\begin{tabular}{ccc}
\\textit{6000} & /bral{\textsubbridge{t}}e{\textyogh}/&$\\rightarrow$ & \\textit{8750} & /brale{\textyogh}/&$\\rightarrow$ & \\textit{11000} & /mrale{\textyogh}/&$\\rightarrow$ & \\textit{12001} & /mrale{\textesh}/&$\\rightarrow$ & \\textit{15001} & /murale{\textesh}/& \\\\
\\end{tabular}

\\vspace{20pt}\\hline

\\end{nopagebreak}
\\filbreak



\\vspace{15pt}
\\begin{nopagebreak}
\\noindent{\\fontsize{20pt}{10pt}\\textbf{haraon} } \\textit{56 of autumn/October 26} (N)\\\\
\\noindent {\\tovian \\fontsize{20pt}{10pt} \\textbf{haraon} }\\\\
\\noindent /ħar{\textprimstress}aon/\\\\
\\noindent Syllables: ħar.a.on\\\\
\\noindent lit. pain+mouth\\\\
\\noindent \\textit{Bitterpearl}\\\\


\\noindent History:

\\vspace{-0pt}
\\hspace{40pt}
\\begin{tabular}{ccc}
\\textit{6000} & /haraeoln/&$\\rightarrow$ & \\textit{7600} & /haraeon/&$\\rightarrow$ & \\textit{10000} & /haraon/&$\\rightarrow$ & \\textit{14006} & /ħaraon/& \\\\
\\end{tabular}

\\vspace{20pt}\\hline

\\end{nopagebreak}
\\filbreak



\\vspace{15pt}
\\begin{nopagebreak}
\\noindent{\\fontsize{20pt}{10pt}\\textbf{muralsyalethashol} } \\textit{56 of spring/April 26} (N)\\\\
\\noindent {\\tovian \\fontsize{20pt}{10pt} \\textbf{muralsyalethashol} }\\\\
\\noindent /muralsyale{\texttheta}{\textprimstress}a{\textesh}ol/\\\\
\\noindent Syllables: mur.alsj.al.eθ.aʃ.ol\\\\
\\noindent lit. darkness+food+fruit\\\\
\\noindent \\textit{Nightnut}\\\\


\\noindent History:

\\vspace{-0pt}
\\hspace{40pt}
\\begin{tabular}{ccc}
\\textit{6000} & /bralsyale{\dh}a{\textyogh}ola/&$\\rightarrow$ & \\textit{9500} & /bralsyale{\dh}a{\textyogh}ol/&$\\rightarrow$ & \\textit{11000} & /mralsyale{\dh}a{\textyogh}ol/&$\\rightarrow$ & \\textit{12001} & /mralsyale{\dh}a{\textesh}ol/&$\\rightarrow$ & \\textit{12002} & /mralsyale{\texttheta}a{\textesh}ol/&$\\rightarrow$ & \\textit{15001} & /muralsyale{\texttheta}a{\textesh}ol/& \\\\
\\end{tabular}

\\vspace{20pt}\\hline

\\end{nopagebreak}
\\filbreak



\\vspace{15pt}
\\begin{nopagebreak}
\\noindent{\\fontsize{20pt}{10pt}\\textbf{muramlor} } \\textit{56 of summer/July 26} (N)\\\\
\\noindent {\\tovian \\fontsize{20pt}{10pt} \\textbf{muramlor} }\\\\
\\noindent /mur{\textprimstress}amlor/\\\\
\\noindent Syllables: mur.aml.or\\\\
\\noindent lit. darkness+plant\\\\
\\noindent \\textit{Nightmyrtle}\\\\


\\noindent History:

\\vspace{-0pt}
\\hspace{40pt}
\\begin{tabular}{ccc}
\\textit{6000} & /bralmllore/&$\\rightarrow$ & \\textit{6500} & /bralmlore/&$\\rightarrow$ & \\textit{7600} & /bramlore/&$\\rightarrow$ & \\textit{9500} & /bramlor/&$\\rightarrow$ & \\textit{11000} & /mramlor/&$\\rightarrow$ & \\textit{15001} & /muramlor/& \\\\
\\end{tabular}

\\vspace{20pt}\\hline

\\end{nopagebreak}
\\filbreak



\\vspace{15pt}
\\begin{nopagebreak}
\\noindent{\\fontsize{20pt}{10pt}\\textbf{nanad} } \\textit{56 of winter/January 56} (N)\\\\
\\noindent {\\tovian \\fontsize{20pt}{10pt} \\textbf{nanad} }\\\\
\\noindent /n{\textprimstress}anad/\\\\
\\noindent Syllables: nan.ad\\\\
\\noindent lit. star+showing\\\\
\\noindent \\textit{Stardush}\\\\


\\noindent History:

\\vspace{-0pt}
\\hspace{40pt}
\\begin{tabular}{ccc}
\\textit{6000} & /nnaanafede/&$\\rightarrow$ & \\textit{6500} & /naanafede/&$\\rightarrow$ & \\textit{9500} & /naanafed/&$\\rightarrow$ & \\textit{12003} & /naanad/&$\\rightarrow$ & \\textit{12005} & /nanad/& \\\\
\\end{tabular}

\\vspace{20pt}\\hline

\\end{nopagebreak}
\\filbreak



\\vspace{15pt}
\\begin{nopagebreak}
\\noindent{\\fontsize{20pt}{10pt}\\textbf{lethilekel} } \\textit{57 of autumn/October 27} (N)\\\\
\\noindent {\\tovian \\fontsize{20pt}{10pt} \\textbf{lethilekel} }\\\\
\\noindent /le{\texttheta}il{\textprimstress}ekel/\\\\
\\noindent Syllables: leθ.il.ek.el\\\\
\\noindent lit. forest+branch\\\\
\\noindent \\textit{Starblum}\\\\


\\noindent History:

\\vspace{-0pt}
\\hspace{40pt}
\\begin{tabular}{ccc}
\\textit{6000} & /aglre{\texttheta}ilekelu/&$\\rightarrow$ & \\textit{6503} & /agrle{\texttheta}ilekelu/&$\\rightarrow$ & \\textit{7500} & /grle{\texttheta}ilekelu/&$\\rightarrow$ & \\textit{7501} & /rle{\texttheta}ilekelu/&$\\rightarrow$ & \\textit{7600} & /le{\texttheta}ilekelu/&$\\rightarrow$ & \\textit{9500} & /le{\texttheta}ilekel/& \\\\
\\end{tabular}

\\vspace{20pt}\\hline

\\end{nopagebreak}
\\filbreak



\\vspace{15pt}
\\begin{nopagebreak}
\\noindent{\\fontsize{20pt}{10pt}\\textbf{yelathan} } \\textit{57 of spring/April 27} (N)\\\\
\\noindent {\\tovian \\fontsize{20pt}{10pt} \\textbf{yelathan} }\\\\
\\noindent /yel{\textprimstress}a{\texttheta}an/\\\\
\\noindent Syllables: jel.aθ.an\\\\
\\noindent lit. person+tool\\\\
\\noindent \\textit{Drielthorn}\\\\


\\noindent History:

\\vspace{-0pt}
\\hspace{40pt}
\\begin{tabular}{ccc}
\\textit{6000} & /dryela{\dh}an/&$\\rightarrow$ & \\textit{7501} & /ryela{\dh}an/&$\\rightarrow$ & \\textit{7600} & /yela{\dh}an/&$\\rightarrow$ & \\textit{12002} & /yela{\texttheta}an/& \\\\
\\end{tabular}

\\vspace{20pt}\\hline

\\end{nopagebreak}
\\filbreak



\\vspace{15pt}
\\begin{nopagebreak}
\\noindent{\\fontsize{20pt}{10pt}\\textbf{varedrem} } \\textit{57 of summer/July 27} (N)\\\\
\\noindent {\\tovian \\fontsize{20pt}{10pt} \\textbf{varedrem} }\\\\
\\noindent /var{\textprimstress}edrem/\\\\
\\noindent Syllables: var.edr.em\\\\
\\noindent lit. value+flower\\\\
\\noindent \\textit{Goldbloom}\\\\


\\noindent History:

\\vspace{-0pt}
\\hspace{40pt}
\\begin{tabular}{ccc}
\\textit{6000} & /vlaredremp/&$\\rightarrow$ & \\textit{8750} & /vlaredrem/&$\\rightarrow$ & \\textit{12500} & /varedrem/& \\\\
\\end{tabular}

\\vspace{20pt}\\hline

\\end{nopagebreak}
\\filbreak



\\vspace{15pt}
\\begin{nopagebreak}
\\noindent{\\fontsize{20pt}{10pt}\\textbf{yisaglar} } \\textit{57 of winter/January 57} (N)\\\\
\\noindent {\\tovian \\fontsize{20pt}{10pt} \\textbf{yisaglar} }\\\\
\\noindent /yis{\textprimstress}aglar/\\\\
\\noindent Syllables: jis.agl.ar\\\\
\\noindent lit. moon+tree\\\\
\\noindent \\textit{Moonbark}\\\\


\\noindent History:

\\vspace{-0pt}
\\hspace{40pt}
\\begin{tabular}{ccc}
\\textit{6000} & /aliysaglare/&$\\rightarrow$ & \\textit{6300} & /alysaglare/&$\\rightarrow$ & \\textit{7500} & /lysaglare/&$\\rightarrow$ & \\textit{7501} & /ysaglare/&$\\rightarrow$ & \\textit{9500} & /ysaglar/&$\\rightarrow$ & \\textit{15001} & /yisaglar/& \\\\
\\end{tabular}

\\vspace{20pt}\\hline

\\end{nopagebreak}
\\filbreak



\\vspace{15pt}
\\begin{nopagebreak}
\\noindent{\\fontsize{20pt}{10pt}\\textbf{mulor} } \\textit{58 of autumn/October 28} (N)\\\\
\\noindent {\\tovian \\fontsize{20pt}{10pt} \\textbf{mulor} }\\\\
\\noindent /m{\textprimstress}ulor/\\\\
\\noindent Syllables: mul.or\\\\
\\noindent lit. sun+plant\\\\
\\noindent \\textit{Sunstalk}\\\\


\\noindent History:

\\vspace{-0pt}
\\hspace{40pt}
\\begin{tabular}{ccc}
\\textit{6000} & /eflmllore/&$\\rightarrow$ & \\textit{6500} & /eflmlore/&$\\rightarrow$ & \\textit{7500} & /flmlore/&$\\rightarrow$ & \\textit{7501} & /lmlore/&$\\rightarrow$ & \\textit{7600} & /mlore/&$\\rightarrow$ & \\textit{9500} & /mlor/&$\\rightarrow$ & \\textit{15001} & /mulor/& \\\\
\\end{tabular}

\\vspace{20pt}\\hline

\\end{nopagebreak}
\\filbreak



\\vspace{15pt}
\\begin{nopagebreak}
\\noindent{\\fontsize{20pt}{10pt}\\textbf{yisol} } \\textit{58 of spring/April 28} (N)\\\\
\\noindent {\\tovian \\fontsize{20pt}{10pt} \\textbf{yisol} }\\\\
\\noindent /y{\textprimstress}isol/\\\\
\\noindent Syllables: jis.ol\\\\
\\noindent lit. moon+fruit\\\\
\\noindent \\textit{Moonquince}\\\\


\\noindent History:

\\vspace{-0pt}
\\hspace{40pt}
\\begin{tabular}{ccc}
\\textit{6000} & /aliys{\textyogh}ola/&$\\rightarrow$ & \\textit{6300} & /alys{\textyogh}ola/&$\\rightarrow$ & \\textit{7500} & /lys{\textyogh}ola/&$\\rightarrow$ & \\textit{7501} & /ys{\textyogh}ola/&$\\rightarrow$ & \\textit{9500} & /ys{\textyogh}ol/&$\\rightarrow$ & \\textit{12001} & /ys{\textesh}ol/&$\\rightarrow$ & \\textit{13000} & /ysol/&$\\rightarrow$ & \\textit{15001} & /yisol/& \\\\
\\end{tabular}

\\vspace{20pt}\\hline

\\end{nopagebreak}
\\filbreak



\\vspace{15pt}
\\begin{nopagebreak}
\\noindent{\\fontsize{20pt}{10pt}\\textbf{yisdmfed} } \\textit{58 of summer/July 28} (N)\\\\
\\noindent {\\tovian \\fontsize{20pt}{10pt} \\textbf{yisdmfed} }\\\\
\\noindent /y{\textprimstress}isdmfed/\\\\
\\noindent Syllables: jisdmf.ed\\\\
\\noindent lit. moon+bloom\\\\
\\noindent \\textit{Moonbloom}\\\\


\\noindent History:

\\vspace{-0pt}
\\hspace{40pt}
\\begin{tabular}{ccc}
\\textit{6000} & /aliysdrmpfede/&$\\rightarrow$ & \\textit{6300} & /alysdrmpfede/&$\\rightarrow$ & \\textit{7500} & /lysdrmpfede/&$\\rightarrow$ & \\textit{7501} & /ysdrmpfede/&$\\rightarrow$ & \\textit{7600} & /ysdmpfede/&$\\rightarrow$ & \\textit{8500} & /ysdmfede/&$\\rightarrow$ & \\textit{9500} & /ysdmfed/&$\\rightarrow$ & \\textit{15001} & /yisdmfed/& \\\\
\\end{tabular}

\\vspace{20pt}\\hline

\\end{nopagebreak}
\\filbreak



\\vspace{15pt}
\\begin{nopagebreak}
\\noindent{\\fontsize{20pt}{10pt}\\textbf{mr} } \\textit{58 of winter/January 58} (N)\\\\
\\noindent {\\tovian \\fontsize{20pt}{10pt} \\textbf{mr} }\\\\
\\noindent /mr/\\\\
\\noindent Syllables: mr\\\\
\\noindent lit. earth+showing\\\\
\\noindent \\textit{Earthem}\\\\


\\noindent History:

\\vspace{-0pt}
\\hspace{40pt}
\\begin{tabular}{ccc}
\\textit{6000} & /mernfede/&$\\rightarrow$ & \\textit{6400} & /merfede/&$\\rightarrow$ & \\textit{9500} & /merfed/&$\\rightarrow$ & \\textit{12003} & /merd/&$\\rightarrow$ & \\textit{13001} & /mer/&$\\rightarrow$ & \\textit{16002} & /mr/& \\\\
\\end{tabular}

\\vspace{20pt}\\hline

\\end{nopagebreak}
\\filbreak



\\vspace{15pt}
\\begin{nopagebreak}
\\noindent{\\fontsize{20pt}{10pt}\\textbf{gelareshol} } \\textit{59 of autumn/October 29} (N)\\\\
\\noindent {\\tovian \\fontsize{20pt}{10pt} \\textbf{gelareshol} }\\\\
\\noindent /gelar{\textprimstress}e{\textesh}ol/\\\\
\\noindent Syllables: gel.ar.eʃ.ol\\\\
\\noindent lit. tree+fruit\\\\
\\noindent \\textit{Hovenut}\\\\


\\noindent History:

\\vspace{-0pt}
\\hspace{40pt}
\\begin{tabular}{ccc}
\\textit{6000} & /aglare{\textyogh}ola/&$\\rightarrow$ & \\textit{7500} & /glare{\textyogh}ola/&$\\rightarrow$ & \\textit{9500} & /glare{\textyogh}ol/&$\\rightarrow$ & \\textit{12001} & /glare{\textesh}ol/&$\\rightarrow$ & \\textit{15001} & /gelare{\textesh}ol/& \\\\
\\end{tabular}

\\vspace{20pt}\\hline

\\end{nopagebreak}
\\filbreak



\\vspace{15pt}
\\begin{nopagebreak}
\\noindent{\\fontsize{20pt}{10pt}\\textbf{keluaglar} } \\textit{59 of spring/April 29} (N)\\\\
\\noindent {\\tovian \\fontsize{20pt}{10pt} \\textbf{keluaglar} }\\\\
\\noindent /kelu{\textprimstress}aglar/\\\\
\\noindent Syllables: kel.u.agl.ar\\\\
\\noindent lit. branch+tree\\\\
\\noindent \\textit{Starwood}\\\\


\\noindent History:

\\vspace{-0pt}
\\hspace{40pt}
\\begin{tabular}{ccc}
\\textit{6000} & /keluaglare/&$\\rightarrow$ & \\textit{9500} & /keluaglar/& \\\\
\\end{tabular}

\\vspace{20pt}\\hline

\\end{nopagebreak}
\\filbreak



\\vspace{15pt}
\\begin{nopagebreak}
\\noindent{\\fontsize{20pt}{10pt}\\textbf{rithiledyel} } \\textit{59 of summer/July 29} (N)\\\\
\\noindent {\\tovian \\fontsize{20pt}{10pt} \\textbf{rithiledyel} }\\\\
\\noindent /ri{\texttheta}il{\textprimstress}edyel/\\\\
\\noindent Syllables: riθ.il.edj.el\\\\
\\noindent lit. sky+person\\\\
\\noindent \\textit{Cloudmun}\\\\


\\noindent History:

\\vspace{-0pt}
\\hspace{40pt}
\\begin{tabular}{ccc}
\\textit{6000} & /ri{\texttheta}iledryela/&$\\rightarrow$ & \\textit{7600} & /ri{\texttheta}iledyela/&$\\rightarrow$ & \\textit{9500} & /ri{\texttheta}iledyel/& \\\\
\\end{tabular}

\\vspace{20pt}\\hline

\\end{nopagebreak}
\\filbreak



\\vspace{15pt}
\\begin{nopagebreak}
\\noindent{\\fontsize{20pt}{10pt}\\textbf{varefethr} } \\textit{59 of winter/January 59} (N)\\\\
\\noindent {\\tovian \\fontsize{20pt}{10pt} \\textbf{varefethr} }\\\\
\\noindent /var{\textprimstress}efe{\texttheta}r/\\\\
\\noindent Syllables: var.ef.eθr\\\\
\\noindent lit. value+speaking\\\\
\\noindent \\textit{Velarum}\\\\


\\noindent History:

\\vspace{-0pt}
\\hspace{40pt}
\\begin{tabular}{ccc}
\\textit{6000} & /vlarefe{\texttheta}r/&$\\rightarrow$ & \\textit{12500} & /varefe{\texttheta}r/& \\\\
\\end{tabular}

\\vspace{20pt}\\hline

\\end{nopagebreak}
\\filbreak



\\vspace{15pt}
\\begin{nopagebreak}
\\noindent{\\fontsize{20pt}{10pt}\\textbf{mar} } \\textit{6 of autumn/September 6} (N)\\\\
\\noindent {\\tovian \\fontsize{20pt}{10pt} \\textbf{mar} }\\\\
\\noindent /m{\textprimstress}ar/\\\\
\\noindent Syllables: mar\\\\
\\noindent lit. sun+mother\\\\
\\noindent \\textit{Sunmary}\\\\


\\noindent History:

\\vspace{-0pt}
\\hspace{40pt}
\\begin{tabular}{ccc}
\\textit{6000} & /eflmari/&$\\rightarrow$ & \\textit{7500} & /flmari/&$\\rightarrow$ & \\textit{7501} & /lmari/&$\\rightarrow$ & \\textit{7600} & /mari/&$\\rightarrow$ & \\textit{9500} & /mar/& \\\\
\\end{tabular}

\\vspace{20pt}\\hline

\\end{nopagebreak}
\\filbreak



\\vspace{15pt}
\\begin{nopagebreak}
\\noindent{\\fontsize{20pt}{10pt}\\textbf{shimilor} } \\textit{6 of spring/March 6} (N)\\\\
\\noindent {\\tovian \\fontsize{20pt}{10pt} \\textbf{shimilor} }\\\\
\\noindent /{\textesh}im{\textprimstress}ilor/\\\\
\\noindent Syllables: ʃim.il.or\\\\
\\noindent lit. mystery+plant\\\\
\\noindent \\textit{Mystsap}\\\\


\\noindent History:

\\vspace{-0pt}
\\hspace{40pt}
\\begin{tabular}{ccc}
\\textit{6000} & /{\textesh}imimllore/&$\\rightarrow$ & \\textit{6500} & /{\textesh}imimlore/&$\\rightarrow$ & \\textit{7600} & /{\textesh}imilore/&$\\rightarrow$ & \\textit{9500} & /{\textesh}imilor/& \\\\
\\end{tabular}

\\vspace{20pt}\\hline

\\end{nopagebreak}
\\filbreak



\\vspace{15pt}
\\begin{nopagebreak}
\\noindent{\\fontsize{20pt}{10pt}\\textbf{meloreathang} } \\textit{6 of summer/June 6} (N)\\\\
\\noindent {\\tovian \\fontsize{20pt}{10pt} \\textbf{meloreathang} }\\\\
\\noindent /melore{\textprimstress}a{\texttheta}a{\ng}/\\\\
\\noindent Syllables: mel.or.e.aθ.aŋ\\\\
\\noindent lit. plant+leaf\\\\
\\noindent \\textit{Melorleaf}\\\\


\\noindent History:

\\vspace{-0pt}
\\hspace{40pt}
\\begin{tabular}{ccc}
\\textit{6000} & /mllorea{\dh}a{\ng}ge/&$\\rightarrow$ & \\textit{6500} & /mlorea{\dh}a{\ng}ge/&$\\rightarrow$ & \\textit{7502} & /melorea{\dh}a{\ng}ge/&$\\rightarrow$ & \\textit{8750} & /melorea{\dh}a{\ng}e/&$\\rightarrow$ & \\textit{9500} & /melorea{\dh}a{\ng}/&$\\rightarrow$ & \\textit{12002} & /melorea{\texttheta}a{\ng}/& \\\\
\\end{tabular}

\\vspace{20pt}\\hline

\\end{nopagebreak}
\\filbreak



\\vspace{15pt}
\\begin{nopagebreak}
\\noindent{\\fontsize{20pt}{10pt}\\textbf{shotltom} } \\textit{6 of winter/december 6} (N)\\\\
\\noindent {\\tovian \\fontsize{20pt}{10pt} \\textbf{shotltom} }\\\\
\\noindent /{\textesh}{\textprimstress}o{\texttoptiebar{t\textbeltl}}{\textsubbridge{t}}om/\\\\
\\noindent Syllables: ʃoƛt.om\\\\
\\noindent lit. flowing+root\\\\
\\noindent \\textit{Mistroot}\\\\


\\noindent History:

\\vspace{-0pt}
\\hspace{40pt}
\\begin{tabular}{ccc}
\\textit{6000} & /{\textesh}o{\texttoptiebar{t\textbeltl}}{\textsubbridge{t}}omo/&$\\rightarrow$ & \\textit{9500} & /{\textesh}o{\texttoptiebar{t\textbeltl}}{\textsubbridge{t}}om/& \\\\
\\end{tabular}

\\vspace{20pt}\\hline

\\end{nopagebreak}
\\filbreak



\\vspace{15pt}
\\begin{nopagebreak}
\\noindent{\\fontsize{20pt}{10pt}\\textbf{leshor} } \\textit{60 of autumn/October 30} (N)\\\\
\\noindent {\\tovian \\fontsize{20pt}{10pt} \\textbf{leshor} }\\\\
\\noindent /l{\textprimstress}e{\textesh}or/\\\\
\\noindent Syllables: leʃ.or\\\\
\\noindent lit. harvest+action\\\\
\\noindent \\textit{Harvesh}\\\\


\\noindent History:

\\vspace{-0pt}
\\hspace{40pt}
\\begin{tabular}{ccc}
\\textit{6000} & /mllre{\textyogh}lore/&$\\rightarrow$ & \\textit{6500} & /mlre{\textyogh}lore/&$\\rightarrow$ & \\textit{6503} & /mrle{\textyogh}lore/&$\\rightarrow$ & \\textit{7501} & /rle{\textyogh}lore/&$\\rightarrow$ & \\textit{7600} & /le{\textyogh}lore/&$\\rightarrow$ & \\textit{9500} & /le{\textyogh}lor/&$\\rightarrow$ & \\textit{12001} & /le{\textesh}lor/&$\\rightarrow$ & \\textit{12500} & /le{\textesh}or/& \\\\
\\end{tabular}

\\vspace{20pt}\\hline

\\end{nopagebreak}
\\filbreak



\\vspace{15pt}
\\begin{nopagebreak}
\\noindent{\\fontsize{20pt}{10pt}\\textbf{rithilekwil} } \\textit{60 of spring/April 30} (N)\\\\
\\noindent {\\tovian \\fontsize{20pt}{10pt} \\textbf{rithilekwil} }\\\\
\\noindent /ri{\texttheta}il{\textprimstress}ekwil/\\\\
\\noindent Syllables: riθ.il.ekw.il\\\\
\\noindent lit. sky+movement\\\\
\\noindent \\textit{Cloudroll}\\\\


\\noindent History:

\\vspace{-0pt}
\\hspace{40pt}
\\begin{tabular}{ccc}
\\textit{6000} & /ri{\texttheta}ilekwil/& \\\\
\\end{tabular}

\\vspace{20pt}\\hline

\\end{nopagebreak}
\\filbreak



\\vspace{15pt}
\\begin{nopagebreak}
\\noindent{\\fontsize{20pt}{10pt}\\textbf{lhenumlhen} } \\textit{60 of summer/July 30} (N)\\\\
\\noindent {\\tovian \\fontsize{20pt}{10pt} \\textbf{lhenumlhen} }\\\\
\\noindent /{\textbeltl}en{\textprimstress}um{\textbeltl}en/\\\\
\\noindent Syllables: ɬen.umɬ.en\\\\
\\noindent lit. center+path+dream\\\\
\\noindent \\textit{Lhelnumly}\\\\


\\noindent History:

\\vspace{-0pt}
\\hspace{40pt}
\\begin{tabular}{ccc}
\\textit{6000} & /{\textbeltl}elnum{\textbeltl}ena/&$\\rightarrow$ & \\textit{7600} & /{\textbeltl}enum{\textbeltl}ena/&$\\rightarrow$ & \\textit{9500} & /{\textbeltl}enum{\textbeltl}en/& \\\\
\\end{tabular}

\\vspace{20pt}\\hline

\\end{nopagebreak}
\\filbreak



\\vspace{15pt}
\\begin{nopagebreak}
\\noindent{\\fontsize{20pt}{10pt}\\textbf{nanadyethil} } \\textit{60 of winter/January 60} (N)\\\\
\\noindent {\\tovian \\fontsize{20pt}{10pt} \\textbf{nanadyethil} }\\\\
\\noindent /nanady{\textprimstress}e{\texttheta}il/\\\\
\\noindent Syllables: nan.adj.eθ.il\\\\
\\noindent lit. star+home\\\\
\\noindent \\textit{Starhovel}\\\\


\\noindent History:

\\vspace{-0pt}
\\hspace{40pt}
\\begin{tabular}{ccc}
\\textit{6000} & /nnaanadryel{\texttheta}ile/&$\\rightarrow$ & \\textit{6500} & /naanadryel{\texttheta}ile/&$\\rightarrow$ & \\textit{7600} & /naanadyel{\texttheta}ile/&$\\rightarrow$ & \\textit{9500} & /naanadyel{\texttheta}il/&$\\rightarrow$ & \\textit{12005} & /nanadyel{\texttheta}il/&$\\rightarrow$ & \\textit{12502} & /nanadye{\texttheta}il/& \\\\
\\end{tabular}

\\vspace{20pt}\\hline

\\end{nopagebreak}
\\filbreak



\\vspace{15pt}
\\begin{nopagebreak}
\\noindent{\\fontsize{20pt}{10pt}\\textbf{nanasyaleth} } \\textit{61 of autumn/November 1} (N)\\\\
\\noindent {\\tovian \\fontsize{20pt}{10pt} \\textbf{nanasyaleth} }\\\\
\\noindent /nanasy{\textprimstress}ale{\texttheta}/\\\\
\\noindent Syllables: nan.asj.al.eθ\\\\
\\noindent lit. star+food\\\\
\\noindent \\textit{Starfodder}\\\\


\\noindent History:

\\vspace{-0pt}
\\hspace{40pt}
\\begin{tabular}{ccc}
\\textit{6000} & /nnaanasyale{\dh}a/&$\\rightarrow$ & \\textit{6500} & /naanasyale{\dh}a/&$\\rightarrow$ & \\textit{9500} & /naanasyale{\dh}/&$\\rightarrow$ & \\textit{12002} & /naanasyale{\texttheta}/&$\\rightarrow$ & \\textit{12005} & /nanasyale{\texttheta}/& \\\\
\\end{tabular}

\\vspace{20pt}\\hline

\\end{nopagebreak}
\\filbreak



\\vspace{15pt}
\\begin{nopagebreak}
\\noindent{\\fontsize{20pt}{10pt}\\textbf{selaglar} } \\textit{61 of spring/May 1} (N)\\\\
\\noindent {\\tovian \\fontsize{20pt}{10pt} \\textbf{selaglar} }\\\\
\\noindent /sel{\textprimstress}aglar/\\\\
\\noindent Syllables: sel.agl.ar\\\\
\\noindent lit. water+tree\\\\
\\noindent \\textit{Rainlaurel}\\\\


\\noindent History:

\\vspace{-0pt}
\\hspace{40pt}
\\begin{tabular}{ccc}
\\textit{6000} & /zelaglare/&$\\rightarrow$ & \\textit{9500} & /zelaglar/&$\\rightarrow$ & \\textit{12000} & /selaglar/& \\\\
\\end{tabular}

\\vspace{20pt}\\hline

\\end{nopagebreak}
\\filbreak



\\vspace{15pt}
\\begin{nopagebreak}
\\noindent{\\fontsize{20pt}{10pt}\\textbf{muloanadrem} } \\textit{61 of summer/August 1} (N)\\\\
\\noindent {\\tovian \\fontsize{20pt}{10pt} \\textbf{muloanadrem} }\\\\
\\noindent /muloan{\textprimstress}adrem/\\\\
\\noindent Syllables: mul.o.an.adr.em\\\\
\\noindent lit. morning+flower\\\\
\\noindent \\textit{Dawnrose}\\\\


\\noindent History:

\\vspace{-0pt}
\\hspace{40pt}
\\begin{tabular}{ccc}
\\textit{6000} & /ploanadremp/&$\\rightarrow$ & \\textit{8750} & /ploanadrem/&$\\rightarrow$ & \\textit{11000} & /mloanadrem/&$\\rightarrow$ & \\textit{15001} & /muloanadrem/& \\\\
\\end{tabular}

\\vspace{20pt}\\hline

\\end{nopagebreak}
\\filbreak



\\vspace{15pt}
\\begin{nopagebreak}
\\noindent{\\fontsize{20pt}{10pt}\\textbf{tushulor} } \\textit{61 of winter/February 1} (N)\\\\
\\noindent {\\tovian \\fontsize{20pt}{10pt} \\textbf{tushulor} }\\\\
\\noindent /{\textsubbridge{t}}u{\textesh}{\textprimstress}ulor/\\\\
\\noindent Syllables: tuʃ.ul.or\\\\
\\noindent lit. cold+plant\\\\
\\noindent \\textit{Froshpea}\\\\


\\noindent History:

\\vspace{-0pt}
\\hspace{40pt}
\\begin{tabular}{ccc}
\\textit{6000} & /{\textsubbridge{t}}u{\textesh}umllore/&$\\rightarrow$ & \\textit{6500} & /{\textsubbridge{t}}u{\textesh}umlore/&$\\rightarrow$ & \\textit{7600} & /{\textsubbridge{t}}u{\textesh}ulore/&$\\rightarrow$ & \\textit{9500} & /{\textsubbridge{t}}u{\textesh}ulor/& \\\\
\\end{tabular}

\\vspace{20pt}\\hline

\\end{nopagebreak}
\\filbreak



\\vspace{15pt}
\\begin{nopagebreak}
\\noindent{\\fontsize{20pt}{10pt}\\textbf{lhoman} } \\textit{62 of autumn/November 2} (N)\\\\
\\noindent {\\tovian \\fontsize{20pt}{10pt} \\textbf{lhoman} }\\\\
\\noindent /{\textbeltl}{\textprimstress}oman/\\\\
\\noindent Syllables: ɬom.an\\\\
\\noindent lit. day+light\\\\
\\noindent \\textit{Dayglem}\\\\


\\noindent History:

\\vspace{-0pt}
\\hspace{40pt}
\\begin{tabular}{ccc}
\\textit{6000} & /{\textbeltl}onmana/&$\\rightarrow$ & \\textit{7600} & /{\textbeltl}omana/&$\\rightarrow$ & \\textit{9500} & /{\textbeltl}oman/& \\\\
\\end{tabular}

\\vspace{20pt}\\hline

\\end{nopagebreak}
\\filbreak



\\vspace{15pt}
\\begin{nopagebreak}
\\noindent{\\fontsize{20pt}{10pt}\\textbf{meloresel} } \\textit{62 of spring/May 2} (N)\\\\
\\noindent {\\tovian \\fontsize{20pt}{10pt} \\textbf{meloresel} }\\\\
\\noindent /melor{\textprimstress}esel/\\\\
\\noindent Syllables: mel.or.es.el\\\\
\\noindent lit. plant+water\\\\
\\noindent \\textit{Dewmoss}\\\\


\\noindent History:

\\vspace{-0pt}
\\hspace{40pt}
\\begin{tabular}{ccc}
\\textit{6000} & /mllorezel/&$\\rightarrow$ & \\textit{6500} & /mlorezel/&$\\rightarrow$ & \\textit{7502} & /melorezel/&$\\rightarrow$ & \\textit{12000} & /meloresel/& \\\\
\\end{tabular}

\\vspace{20pt}\\hline

\\end{nopagebreak}
\\filbreak



\\vspace{15pt}
\\begin{nopagebreak}
\\noindent{\\fontsize{20pt}{10pt}\\textbf{teshaglar} } \\textit{62 of summer/August 2} (N)\\\\
\\noindent {\\tovian \\fontsize{20pt}{10pt} \\textbf{teshaglar} }\\\\
\\noindent /{\textsubbridge{t}}e{\textesh}{\textprimstress}aglar/\\\\
\\noindent Syllables: teʃ.agl.ar\\\\
\\noindent lit. sound+tree\\\\
\\noindent \\textit{Thunderoak}\\\\


\\noindent History:

\\vspace{-0pt}
\\hspace{40pt}
\\begin{tabular}{ccc}
\\textit{6000} & /{\textsubbridge{t}}e{\textyogh}aglare/&$\\rightarrow$ & \\textit{9500} & /{\textsubbridge{t}}e{\textyogh}aglar/&$\\rightarrow$ & \\textit{12001} & /{\textsubbridge{t}}e{\textesh}aglar/& \\\\
\\end{tabular}

\\vspace{20pt}\\hline

\\end{nopagebreak}
\\filbreak



\\vspace{15pt}
\\begin{nopagebreak}
\\noindent{\\fontsize{20pt}{10pt}\\textbf{thokmramen} } \\textit{62 of winter/February 2} (N)\\\\
\\noindent {\\tovian \\fontsize{20pt}{10pt} \\textbf{thokmramen} }\\\\
\\noindent /{\texttheta}okmr{\textprimstress}amen/\\\\
\\noindent Syllables: θokmr.am.en\\\\
\\noindent lit. bottom+darkness+stone\\\\
\\noindent \\textit{Nightcoal}\\\\


\\noindent History:

\\vspace{-0pt}
\\hspace{40pt}
\\begin{tabular}{ccc}
\\textit{6000} & /{\texttheta}okbralmern/&$\\rightarrow$ & \\textit{7600} & /{\texttheta}okbramen/&$\\rightarrow$ & \\textit{11000} & /{\texttheta}okmramen/& \\\\
\\end{tabular}

\\vspace{20pt}\\hline

\\end{nopagebreak}
\\filbreak



\\vspace{15pt}
\\begin{nopagebreak}
\\noindent{\\fontsize{20pt}{10pt}\\textbf{haralor} } \\textit{63 of autumn/November 3} (N)\\\\
\\noindent {\\tovian \\fontsize{20pt}{10pt} \\textbf{haralor} }\\\\
\\noindent /ħar{\textprimstress}alor/\\\\
\\noindent Syllables: ħar.al.or\\\\
\\noindent lit. pain+plant\\\\
\\noindent \\textit{Harclover}\\\\


\\noindent History:

\\vspace{-0pt}
\\hspace{40pt}
\\begin{tabular}{ccc}
\\textit{6000} & /haramllore/&$\\rightarrow$ & \\textit{6500} & /haramlore/&$\\rightarrow$ & \\textit{7600} & /haralore/&$\\rightarrow$ & \\textit{9500} & /haralor/&$\\rightarrow$ & \\textit{14006} & /ħaralor/& \\\\
\\end{tabular}

\\vspace{20pt}\\hline

\\end{nopagebreak}
\\filbreak



\\vspace{15pt}
\\begin{nopagebreak}
\\noindent{\\fontsize{20pt}{10pt}\\textbf{shotlen} } \\textit{63 of spring/May 3} (N)\\\\
\\noindent {\\tovian \\fontsize{20pt}{10pt} \\textbf{shotlen} }\\\\
\\noindent /{\textesh}{\textprimstress}o{\texttoptiebar{t\textbeltl}}en/\\\\
\\noindent Syllables: ʃoƛ.en\\\\
\\noindent lit. flowing+dream\\\\
\\noindent \\textit{Mistthorn}\\\\


\\noindent History:

\\vspace{-0pt}
\\hspace{40pt}
\\begin{tabular}{ccc}
\\textit{6000} & /{\textesh}o{\texttoptiebar{t\textbeltl}}{\textbeltl}ena/&$\\rightarrow$ & \\textit{9500} & /{\textesh}o{\texttoptiebar{t\textbeltl}}{\textbeltl}en/&$\\rightarrow$ & \\textit{14000} & /{\textesh}o{\texttoptiebar{t\textbeltl}}en/& \\\\
\\end{tabular}

\\vspace{20pt}\\hline

\\end{nopagebreak}
\\filbreak



\\vspace{15pt}
\\begin{nopagebreak}
\\noindent{\\fontsize{20pt}{10pt}\\textbf{munanilor} } \\textit{63 of summer/August 3} (N)\\\\
\\noindent {\\tovian \\fontsize{20pt}{10pt} \\textbf{munanilor} }\\\\
\\noindent /munan{\textprimstress}ilor/\\\\
\\noindent Syllables: mun.an.il.or\\\\
\\noindent lit. covering+plant\\\\
\\noindent \\textit{Veilfern}\\\\


\\noindent History:

\\vspace{-0pt}
\\hspace{40pt}
\\begin{tabular}{ccc}
\\textit{6000} & /bnanimllore/&$\\rightarrow$ & \\textit{6500} & /bnanimlore/&$\\rightarrow$ & \\textit{7600} & /bnanilore/&$\\rightarrow$ & \\textit{9500} & /bnanilor/&$\\rightarrow$ & \\textit{11000} & /mnanilor/&$\\rightarrow$ & \\textit{11996} & /munanilor/& \\\\
\\end{tabular}

\\vspace{20pt}\\hline

\\end{nopagebreak}
\\filbreak



\\vspace{15pt}
\\begin{nopagebreak}
\\noindent{\\fontsize{20pt}{10pt}\\textbf{muralingtlun} } \\textit{63 of winter/February 3} (N)\\\\
\\noindent {\\tovian \\fontsize{20pt}{10pt} \\textbf{muralingtlun} }\\\\
\\noindent /mural{\textprimstress}i{\ng}{\texttoptiebar{t\textbeltl}}un/\\\\
\\noindent Syllables: mur.al.iŋƛ.un\\\\
\\noindent lit. darkness+tar\\\\
\\noindent \\textit{Shadowtar}\\\\


\\noindent History:

\\vspace{-0pt}
\\hspace{40pt}
\\begin{tabular}{ccc}
\\textit{6000} & /brali{\ng}g{\texttoptiebar{t\textbeltl}}una/&$\\rightarrow$ & \\textit{8500} & /brali{\ng}{\texttoptiebar{t\textbeltl}}una/&$\\rightarrow$ & \\textit{9500} & /brali{\ng}{\texttoptiebar{t\textbeltl}}un/&$\\rightarrow$ & \\textit{11000} & /mrali{\ng}{\texttoptiebar{t\textbeltl}}un/&$\\rightarrow$ & \\textit{15001} & /murali{\ng}{\texttoptiebar{t\textbeltl}}un/& \\\\
\\end{tabular}

\\vspace{20pt}\\hline

\\end{nopagebreak}
\\filbreak



\\vspace{15pt}
\\begin{nopagebreak}
\\noindent{\\fontsize{20pt}{10pt}\\textbf{ay} } \\textit{64 of autumn/November 4} (N)\\\\
\\noindent {\\tovian \\fontsize{20pt}{10pt} \\textbf{ay} }\\\\
\\noindent /{\textprimstress}ay/\\\\
\\noindent Syllables: aj\\\\
\\noindent lit. moon+???\\\\
\\noindent \\textit{Moonerb}\\\\


\\noindent History:

\\vspace{-0pt}
\\hspace{40pt}
\\begin{tabular}{ccc}
\\textit{6000} & /aliys???/&$\\rightarrow$ & \\textit{6300} & /alys???/&$\\rightarrow$ & \\textit{7600} & /ays???/&$\\rightarrow$ & \\textit{8500} & /ays/&$\\rightarrow$ & \\textit{13001} & /ay/& \\\\
\\end{tabular}

\\vspace{20pt}\\hline

\\end{nopagebreak}
\\filbreak



\\vspace{15pt}
\\begin{nopagebreak}
\\noindent{\\fontsize{20pt}{10pt}\\textbf{shashithan} } \\textit{64 of spring/May 4} (N)\\\\
\\noindent {\\tovian \\fontsize{20pt}{10pt} \\textbf{shashithan} }\\\\
\\noindent /{\textesh}a{\textesh}{\textprimstress}i{\texttheta}an/\\\\
\\noindent Syllables: ʃaʃ.iθ.an\\\\
\\noindent lit. ice+sharpness+tool\\\\
\\noindent \\textit{Snowpierce}\\\\


\\noindent History:

\\vspace{-0pt}
\\hspace{40pt}
\\begin{tabular}{ccc}
\\textit{6000} & /{\textesh}u{\textesh}a{\textyogh}i{\texttheta}{\dh}an/&$\\rightarrow$ & \\textit{9200} & /{\textesh}{\textschwa}{\textesh}a{\textyogh}i{\texttheta}{\dh}an/&$\\rightarrow$ & \\textit{11990} & /{\textesh}{\textesh}a{\textyogh}i{\texttheta}{\dh}an/&$\\rightarrow$ & \\textit{12001} & /{\textesh}{\textesh}a{\textesh}i{\texttheta}{\dh}an/&$\\rightarrow$ & \\textit{12002} & /{\textesh}{\textesh}a{\textesh}i{\texttheta}{\texttheta}an/&$\\rightarrow$ & \\textit{12004} & /{\textesh}a{\textesh}i{\texttheta}{\texttheta}an/&$\\rightarrow$ & \\textit{12007} & /{\textesh}a{\textesh}i{\texttheta}an/& \\\\
\\end{tabular}

\\vspace{20pt}\\hline

\\end{nopagebreak}
\\filbreak



\\vspace{15pt}
\\begin{nopagebreak}
\\noindent{\\fontsize{20pt}{10pt}\\textbf{kelushith} } \\textit{64 of summer/August 4} (N)\\\\
\\noindent {\\tovian \\fontsize{20pt}{10pt} \\textbf{kelushith} }\\\\
\\noindent /kel{\textprimstress}u{\textesh}i{\texttheta}/\\\\
\\noindent Syllables: kel.uʃ.iθ\\\\
\\noindent lit. branch+sharpness\\\\
\\noindent \\textit{Starthorn}\\\\


\\noindent History:

\\vspace{-0pt}
\\hspace{40pt}
\\begin{tabular}{ccc}
\\textit{6000} & /kelu{\textyogh}i{\texttheta}/&$\\rightarrow$ & \\textit{12001} & /kelu{\textesh}i{\texttheta}/& \\\\
\\end{tabular}

\\vspace{20pt}\\hline

\\end{nopagebreak}
\\filbreak



\\vspace{15pt}
\\begin{nopagebreak}
\\noindent{\\fontsize{20pt}{10pt}\\textbf{nanamen} } \\textit{64 of winter/February 4} (N)\\\\
\\noindent {\\tovian \\fontsize{20pt}{10pt} \\textbf{nanamen} }\\\\
\\noindent /nan{\textprimstress}amen/\\\\
\\noindent Syllables: nan.am.en\\\\
\\noindent lit. star+stone\\\\
\\noindent \\textit{Starbrimstone}\\\\


\\noindent History:

\\vspace{-0pt}
\\hspace{40pt}
\\begin{tabular}{ccc}
\\textit{6000} & /nnaanamern/&$\\rightarrow$ & \\textit{6500} & /naanamern/&$\\rightarrow$ & \\textit{7600} & /naanamen/&$\\rightarrow$ & \\textit{12005} & /nanamen/& \\\\
\\end{tabular}

\\vspace{20pt}\\hline

\\end{nopagebreak}
\\filbreak



\\vspace{15pt}
\\begin{nopagebreak}
\\noindent{\\fontsize{20pt}{10pt}\\textbf{teleay} } \\textit{65 of autumn/November 5} (N)\\\\
\\noindent {\\tovian \\fontsize{20pt}{10pt} \\textbf{teleay} }\\\\
\\noindent /{\textsubbridge{t}}el{\textprimstress}eay/\\\\
\\noindent Syllables: tel.e.aj\\\\
\\noindent lit. lake+feather\\\\
\\noindent \\textit{Lakequill}\\\\


\\noindent History:

\\vspace{-0pt}
\\hspace{40pt}
\\begin{tabular}{ccc}
\\textit{6000} & /o{\textsubbridge{t}}elearya/&$\\rightarrow$ & \\textit{7500} & /{\textsubbridge{t}}elearya/&$\\rightarrow$ & \\textit{7600} & /{\textsubbridge{t}}eleaya/&$\\rightarrow$ & \\textit{9500} & /{\textsubbridge{t}}eleay/& \\\\
\\end{tabular}

\\vspace{20pt}\\hline

\\end{nopagebreak}
\\filbreak



\\vspace{15pt}
\\begin{nopagebreak}
\\noindent{\\fontsize{20pt}{10pt}\\textbf{doltlan} } \\textit{65 of spring/May 5} (N)\\\\
\\noindent {\\tovian \\fontsize{20pt}{10pt} \\textbf{doltlan} }\\\\
\\noindent /d{\textprimstress}ol{\texttoptiebar{t\textbeltl}}an/\\\\
\\noindent Syllables: dolƛ.an\\\\
\\noindent lit. death+animal\\\\
\\noindent \\textit{Stambul}\\\\


\\noindent History:

\\vspace{-0pt}
\\hspace{40pt}
\\begin{tabular}{ccc}
\\textit{6000} & /dol{\textsubbridge{t}}lamn/&$\\rightarrow$ & \\textit{6401} & /dol{\texttoptiebar{t\textbeltl}}amn/&$\\rightarrow$ & \\textit{7600} & /dol{\texttoptiebar{t\textbeltl}}an/& \\\\
\\end{tabular}

\\vspace{20pt}\\hline

\\end{nopagebreak}
\\filbreak



\\vspace{15pt}
\\begin{nopagebreak}
\\noindent{\\fontsize{20pt}{10pt}\\textbf{mural} } \\textit{65 of summer/August 5} (N)\\\\
\\noindent {\\tovian \\fontsize{20pt}{10pt} \\textbf{mural} }\\\\
\\noindent /m{\textprimstress}ural/\\\\
\\noindent Syllables: mur.al\\\\
\\noindent lit. darkness+???\\\\
\\noindent \\textit{Nightahn}\\\\


\\noindent History:

\\vspace{-0pt}
\\hspace{40pt}
\\begin{tabular}{ccc}
\\textit{6000} & /bral???/&$\\rightarrow$ & \\textit{8500} & /bral/&$\\rightarrow$ & \\textit{11000} & /mral/&$\\rightarrow$ & \\textit{15001} & /mural/& \\\\
\\end{tabular}

\\vspace{20pt}\\hline

\\end{nopagebreak}
\\filbreak



\\vspace{15pt}
\\begin{nopagebreak}
\\noindent{\\fontsize{20pt}{10pt}\\textbf{yisanon} } \\textit{65 of winter/February 5} (N)\\\\
\\noindent {\\tovian \\fontsize{20pt}{10pt} \\textbf{yisanon} }\\\\
\\noindent /yis{\textprimstress}anon/\\\\
\\noindent Syllables: jis.an.on\\\\
\\noindent lit. moon+dog\\\\
\\noindent \\textit{Moonhound}\\\\


\\noindent History:

\\vspace{-0pt}
\\hspace{40pt}
\\begin{tabular}{ccc}
\\textit{6000} & /aliysandorn/&$\\rightarrow$ & \\textit{6300} & /alysandorn/&$\\rightarrow$ & \\textit{7500} & /lysandorn/&$\\rightarrow$ & \\textit{7501} & /ysandorn/&$\\rightarrow$ & \\textit{7600} & /ysandon/&$\\rightarrow$ & \\textit{8750} & /ysanon/&$\\rightarrow$ & \\textit{15001} & /yisanon/& \\\\
\\end{tabular}

\\vspace{20pt}\\hline

\\end{nopagebreak}
\\filbreak



\\vspace{15pt}
\\begin{nopagebreak}
\\noindent{\\fontsize{20pt}{10pt}\\textbf{holalor} } \\textit{66 of autumn/November 6} (N)\\\\
\\noindent {\\tovian \\fontsize{20pt}{10pt} \\textbf{holalor} }\\\\
\\noindent /ħol{\textprimstress}alor/\\\\
\\noindent Syllables: ħol.al.or\\\\
\\noindent lit. sweet+plant\\\\
\\noindent \\textit{Sweetbalm}\\\\


\\noindent History:

\\vspace{-0pt}
\\hspace{40pt}
\\begin{tabular}{ccc}
\\textit{6000} & /holamllore/&$\\rightarrow$ & \\textit{6500} & /holamlore/&$\\rightarrow$ & \\textit{7600} & /holalore/&$\\rightarrow$ & \\textit{9500} & /holalor/&$\\rightarrow$ & \\textit{14006} & /ħolalor/& \\\\
\\end{tabular}

\\vspace{20pt}\\hline

\\end{nopagebreak}
\\filbreak



\\vspace{15pt}
\\begin{nopagebreak}
\\noindent{\\fontsize{20pt}{10pt}\\textbf{rithor} } \\textit{66 of spring/May 6} (N)\\\\
\\noindent {\\tovian \\fontsize{20pt}{10pt} \\textbf{rithor} }\\\\
\\noindent /r{\textprimstress}i{\texttheta}or/\\\\
\\noindent Syllables: riθ.or\\\\
\\noindent lit. air+plant\\\\
\\noindent \\textit{Rinthyme}\\\\


\\noindent History:

\\vspace{-0pt}
\\hspace{40pt}
\\begin{tabular}{ccc}
\\textit{6000} & /ri{\texttheta}mllore/&$\\rightarrow$ & \\textit{6500} & /ri{\texttheta}mlore/&$\\rightarrow$ & \\textit{7600} & /ri{\texttheta}lore/&$\\rightarrow$ & \\textit{9500} & /ri{\texttheta}lor/&$\\rightarrow$ & \\textit{12500} & /ri{\texttheta}or/& \\\\
\\end{tabular}

\\vspace{20pt}\\hline

\\end{nopagebreak}
\\filbreak



\\vspace{15pt}
\\begin{nopagebreak}
\\noindent{\\fontsize{20pt}{10pt}\\textbf{lhenathan} } \\textit{66 of summer/August 6} (N)\\\\
\\noindent {\\tovian \\fontsize{20pt}{10pt} \\textbf{lhenathan} }\\\\
\\noindent /{\textbeltl}en{\textprimstress}a{\texttheta}an/\\\\
\\noindent Syllables: ɬen.aθ.an\\\\
\\noindent lit. dream+tool\\\\
\\noindent \\textit{Drembell}\\\\


\\noindent History:

\\vspace{-0pt}
\\hspace{40pt}
\\begin{tabular}{ccc}
\\textit{6000} & /{\textbeltl}ena{\dh}an/&$\\rightarrow$ & \\textit{12002} & /{\textbeltl}ena{\texttheta}an/& \\\\
\\end{tabular}

\\vspace{20pt}\\hline

\\end{nopagebreak}
\\filbreak



\\vspace{15pt}
\\begin{nopagebreak}
\\noindent{\\fontsize{20pt}{10pt}\\textbf{lhenashotl} } \\textit{66 of winter/February 6} (N)\\\\
\\noindent {\\tovian \\fontsize{20pt}{10pt} \\textbf{lhenashotl} }\\\\
\\noindent /{\textbeltl}en{\textprimstress}a{\textesh}o{\texttoptiebar{t\textbeltl}}/\\\\
\\noindent Syllables: ɬen.aʃ.oƛ\\\\
\\noindent lit. dream+flow\\\\
\\noindent \\textit{Dreamlava}\\\\


\\noindent History:

\\vspace{-0pt}
\\hspace{40pt}
\\begin{tabular}{ccc}
\\textit{6000} & /{\textbeltl}ena{\textesh}o{\texttoptiebar{t\textbeltl}}/& \\\\
\\end{tabular}

\\vspace{20pt}\\hline

\\end{nopagebreak}
\\filbreak



\\vspace{15pt}
\\begin{nopagebreak}
\\noindent{\\fontsize{20pt}{10pt}\\textbf{lhuloresyaleth} } \\textit{67 of autumn/November 7} (N)\\\\
\\noindent {\\tovian \\fontsize{20pt}{10pt} \\textbf{lhuloresyaleth} }\\\\
\\noindent /{\textbeltl}uloresy{\textprimstress}ale{\texttheta}/\\\\
\\noindent Syllables: ɬul.or.esj.al.eθ\\\\
\\noindent lit. goodness+grain\\\\
\\noindent \\textit{Sugargrain}\\\\


\\noindent History:

\\vspace{-0pt}
\\hspace{40pt}
\\begin{tabular}{ccc}
\\textit{6000} & /{\textbeltl}ummlloresyale{\dh}a/&$\\rightarrow$ & \\textit{6500} & /{\textbeltl}umloresyale{\dh}a/&$\\rightarrow$ & \\textit{7600} & /{\textbeltl}uloresyale{\dh}a/&$\\rightarrow$ & \\textit{9500} & /{\textbeltl}uloresyale{\dh}/&$\\rightarrow$ & \\textit{12002} & /{\textbeltl}uloresyale{\texttheta}/& \\\\
\\end{tabular}

\\vspace{20pt}\\hline

\\end{nopagebreak}
\\filbreak



\\vspace{15pt}
\\begin{nopagebreak}
\\noindent{\\fontsize{20pt}{10pt}\\textbf{munaniay} } \\textit{67 of spring/May 7} (N)\\\\
\\noindent {\\tovian \\fontsize{20pt}{10pt} \\textbf{munaniay} }\\\\
\\noindent /munan{\textprimstress}iay/\\\\
\\noindent Syllables: mun.an.i.aj\\\\
\\noindent lit. protection+moon\\\\
\\noindent \\textit{Moonsponge}\\\\


\\noindent History:

\\vspace{-0pt}
\\hspace{40pt}
\\begin{tabular}{ccc}
\\textit{6000} & /bnanialiys/&$\\rightarrow$ & \\textit{6300} & /bnanialys/&$\\rightarrow$ & \\textit{7600} & /bnaniays/&$\\rightarrow$ & \\textit{11000} & /mnaniays/&$\\rightarrow$ & \\textit{11996} & /munaniays/&$\\rightarrow$ & \\textit{13001} & /munaniay/& \\\\
\\end{tabular}

\\vspace{20pt}\\hline

\\end{nopagebreak}
\\filbreak



\\vspace{15pt}
\\begin{nopagebreak}
\\noindent{\\fontsize{20pt}{10pt}\\textbf{aysdrem} } \\textit{67 of summer/August 7} (N)\\\\
\\noindent {\\tovian \\fontsize{20pt}{10pt} \\textbf{aysdrem} }\\\\
\\noindent /{\textprimstress}aysdrem/\\\\
\\noindent Syllables: ajsdr.em\\\\
\\noindent lit. moon+flower\\\\
\\noindent \\textit{Moonlily}\\\\


\\noindent History:

\\vspace{-0pt}
\\hspace{40pt}
\\begin{tabular}{ccc}
\\textit{6000} & /aliysdremp/&$\\rightarrow$ & \\textit{6300} & /alysdremp/&$\\rightarrow$ & \\textit{7600} & /aysdremp/&$\\rightarrow$ & \\textit{8750} & /aysdrem/& \\\\
\\end{tabular}

\\vspace{20pt}\\hline

\\end{nopagebreak}
\\filbreak



\\vspace{15pt}
\\begin{nopagebreak}
\\noindent{\\fontsize{20pt}{10pt}\\textbf{menlor} } \\textit{67 of winter/February 7} (N)\\\\
\\noindent {\\tovian \\fontsize{20pt}{10pt} \\textbf{menlor} }\\\\
\\noindent /m{\textprimstress}enlor/\\\\
\\noindent Syllables: menl.or\\\\
\\noindent lit. earth+plant\\\\
\\noindent \\textit{Earthmoss}\\\\


\\noindent History:

\\vspace{-0pt}
\\hspace{40pt}
\\begin{tabular}{ccc}
\\textit{6000} & /mernmllore/&$\\rightarrow$ & \\textit{6500} & /mernmlore/&$\\rightarrow$ & \\textit{7600} & /menlore/&$\\rightarrow$ & \\textit{9500} & /menlor/& \\\\
\\end{tabular}

\\vspace{20pt}\\hline

\\end{nopagebreak}
\\filbreak



\\vspace{15pt}
\\begin{nopagebreak}
\\noindent{\\fontsize{20pt}{10pt}\\textbf{dolag} } \\textit{68 of autumn/November 8} (N)\\\\
\\noindent {\\tovian \\fontsize{20pt}{10pt} \\textbf{dolag} }\\\\
\\noindent /d{\textprimstress}olag/\\\\
\\noindent Syllables: dol.ag\\\\
\\noindent lit. death+hotness\\\\
\\noindent \\textit{Twilily}\\\\


\\noindent History:

\\vspace{-0pt}
\\hspace{40pt}
\\begin{tabular}{ccc}
\\textit{6000} & /dolaga/&$\\rightarrow$ & \\textit{9500} & /dolag/& \\\\
\\end{tabular}

\\vspace{20pt}\\hline

\\end{nopagebreak}
\\filbreak



\\vspace{15pt}
\\begin{nopagebreak}
\\noindent{\\fontsize{20pt}{10pt}\\textbf{munorhanalor} } \\textit{68 of spring/May 8} (N)\\\\
\\noindent {\\tovian \\fontsize{20pt}{10pt} \\textbf{munorhanalor} }\\\\
\\noindent /munorħan{\textprimstress}alor/\\\\
\\noindent Syllables: mun.orħ.an.al.or\\\\
\\noindent lit. dusk+plant\\\\
\\noindent \\textit{Duskdaphne}\\\\


\\noindent History:

\\vspace{-0pt}
\\hspace{40pt}
\\begin{tabular}{ccc}
\\textit{6000} & /pnorhanamllore/&$\\rightarrow$ & \\textit{6500} & /pnorhanamlore/&$\\rightarrow$ & \\textit{7600} & /pnorhanalore/&$\\rightarrow$ & \\textit{9500} & /pnorhanalor/&$\\rightarrow$ & \\textit{11000} & /mnorhanalor/&$\\rightarrow$ & \\textit{11996} & /munorhanalor/&$\\rightarrow$ & \\textit{14006} & /munorħanalor/& \\\\
\\end{tabular}

\\vspace{20pt}\\hline

\\end{nopagebreak}
\\filbreak



\\vspace{15pt}
\\begin{nopagebreak}
\\noindent{\\fontsize{20pt}{10pt}\\textbf{menrunan} } \\textit{68 of summer/August 8} (N)\\\\
\\noindent {\\tovian \\fontsize{20pt}{10pt} \\textbf{menrunan} }\\\\
\\noindent /menr{\textprimstress}unan/\\\\
\\noindent Syllables: menr.un.an\\\\
\\noindent lit. earth+hat\\\\
\\noindent \\textit{Earthcap}\\\\


\\noindent History:

\\vspace{-0pt}
\\hspace{40pt}
\\begin{tabular}{ccc}
\\textit{6000} & /merndrunbnani/&$\\rightarrow$ & \\textit{7600} & /mendrunbnani/&$\\rightarrow$ & \\textit{8500} & /menrunnani/&$\\rightarrow$ & \\textit{9500} & /menrunnan/&$\\rightarrow$ & \\textit{12007} & /menrunan/& \\\\
\\end{tabular}

\\vspace{20pt}\\hline

\\end{nopagebreak}
\\filbreak



\\vspace{15pt}
\\begin{nopagebreak}
\\noindent{\\fontsize{20pt}{10pt}\\textbf{mrsyaleth} } \\textit{68 of winter/February 8} (N)\\\\
\\noindent {\\tovian \\fontsize{20pt}{10pt} \\textbf{mrsyaleth} }\\\\
\\noindent /mrsy{\textprimstress}ale{\texttheta}/\\\\
\\noindent Syllables: mrsjal.eθ\\\\
\\noindent lit. earth+food\\\\
\\noindent \\textit{Starfern}\\\\


\\noindent History:

\\vspace{-0pt}
\\hspace{40pt}
\\begin{tabular}{ccc}
\\textit{6000} & /mernsyale{\dh}a/&$\\rightarrow$ & \\textit{6400} & /mersyale{\dh}a/&$\\rightarrow$ & \\textit{9500} & /mersyale{\dh}/&$\\rightarrow$ & \\textit{12002} & /mersyale{\texttheta}/&$\\rightarrow$ & \\textit{16002} & /mrsyale{\texttheta}/& \\\\
\\end{tabular}

\\vspace{20pt}\\hline

\\end{nopagebreak}
\\filbreak



\\vspace{15pt}
\\begin{nopagebreak}
\\noindent{\\fontsize{20pt}{10pt}\\textbf{rithtom} } \\textit{69 of autumn/November 9} (N)\\\\
\\noindent {\\tovian \\fontsize{20pt}{10pt} \\textbf{rithtom} }\\\\
\\noindent /r{\textprimstress}i{\texttheta}{\textsubbridge{t}}om/\\\\
\\noindent Syllables: riθt.om\\\\
\\noindent lit. air+root\\\\
\\noindent \\textit{Windthyme}\\\\


\\noindent History:

\\vspace{-0pt}
\\hspace{40pt}
\\begin{tabular}{ccc}
\\textit{6000} & /ri{\texttheta}{\textsubbridge{t}}omo/&$\\rightarrow$ & \\textit{9500} & /ri{\texttheta}{\textsubbridge{t}}om/& \\\\
\\end{tabular}

\\vspace{20pt}\\hline

\\end{nopagebreak}
\\filbreak



\\vspace{15pt}
\\begin{nopagebreak}
\\noindent{\\fontsize{20pt}{10pt}\\textbf{teleaglar} } \\textit{69 of spring/May 9} (N)\\\\
\\noindent {\\tovian \\fontsize{20pt}{10pt} \\textbf{teleaglar} }\\\\
\\noindent /{\textsubbridge{t}}ele{\textprimstress}aglar/\\\\
\\noindent Syllables: tel.e.agl.ar\\\\
\\noindent lit. lake+tree\\\\
\\noindent \\textit{Telpoplar}\\\\


\\noindent History:

\\vspace{-0pt}
\\hspace{40pt}
\\begin{tabular}{ccc}
\\textit{6000} & /o{\textsubbridge{t}}eleaglare/&$\\rightarrow$ & \\textit{7500} & /{\textsubbridge{t}}eleaglare/&$\\rightarrow$ & \\textit{9500} & /{\textsubbridge{t}}eleaglar/& \\\\
\\end{tabular}

\\vspace{20pt}\\hline

\\end{nopagebreak}
\\filbreak



\\vspace{15pt}
\\begin{nopagebreak}
\\noindent{\\fontsize{20pt}{10pt}\\textbf{selrem} } \\textit{69 of summer/August 9} (N)\\\\
\\noindent {\\tovian \\fontsize{20pt}{10pt} \\textbf{selrem} }\\\\
\\noindent /s{\textprimstress}elrem/\\\\
\\noindent Syllables: selr.em\\\\
\\noindent lit. water+flower\\\\
\\noindent \\textit{Rainblossom}\\\\


\\noindent History:

\\vspace{-0pt}
\\hspace{40pt}
\\begin{tabular}{ccc}
\\textit{6000} & /zeldremp/&$\\rightarrow$ & \\textit{8750} & /zelrem/&$\\rightarrow$ & \\textit{12000} & /selrem/& \\\\
\\end{tabular}

\\vspace{20pt}\\hline

\\end{nopagebreak}
\\filbreak



\\vspace{15pt}
\\begin{nopagebreak}
\\noindent{\\fontsize{20pt}{10pt}\\textbf{shalhen} } \\textit{69 of winter/February 9} (N)\\\\
\\noindent {\\tovian \\fontsize{20pt}{10pt} \\textbf{shalhen} }\\\\
\\noindent /{\textesh}{\textprimstress}a{\textbeltl}en/\\\\
\\noindent Syllables: ʃaɬ.en\\\\
\\noindent lit. ice+dream\\\\
\\noindent \\textit{Frostniter}\\\\


\\noindent History:

\\vspace{-0pt}
\\hspace{40pt}
\\begin{tabular}{ccc}
\\textit{6000} & /{\textesh}u{\textesh}a{\textbeltl}ena/&$\\rightarrow$ & \\textit{9200} & /{\textesh}{\textschwa}{\textesh}a{\textbeltl}ena/&$\\rightarrow$ & \\textit{9500} & /{\textesh}{\textschwa}{\textesh}a{\textbeltl}en/&$\\rightarrow$ & \\textit{11990} & /{\textesh}{\textesh}a{\textbeltl}en/&$\\rightarrow$ & \\textit{12007} & /{\textesh}a{\textbeltl}en/& \\\\
\\end{tabular}

\\vspace{20pt}\\hline

\\end{nopagebreak}
\\filbreak



\\vspace{15pt}
\\begin{nopagebreak}
\\noindent{\\fontsize{20pt}{10pt}\\textbf{ay} } \\textit{7 of autumn/September 7} (N)\\\\
\\noindent {\\tovian \\fontsize{20pt}{10pt} \\textbf{ay} }\\\\
\\noindent /{\textprimstress}ay/\\\\
\\noindent Syllables: aj\\\\
\\noindent lit. moon+???\\\\
\\noindent \\textit{Mooncoke}\\\\


\\noindent History:

\\vspace{-0pt}
\\hspace{40pt}
\\begin{tabular}{ccc}
\\textit{6000} & /aliys???/&$\\rightarrow$ & \\textit{6300} & /alys???/&$\\rightarrow$ & \\textit{7600} & /ays???/&$\\rightarrow$ & \\textit{8500} & /ays/&$\\rightarrow$ & \\textit{13001} & /ay/& \\\\
\\end{tabular}

\\vspace{20pt}\\hline

\\end{nopagebreak}
\\filbreak



\\vspace{15pt}
\\begin{nopagebreak}
\\noindent{\\fontsize{20pt}{10pt}\\textbf{murilamen} } \\textit{7 of spring/March 7} (N)\\\\
\\noindent {\\tovian \\fontsize{20pt}{10pt} \\textbf{murilamen} }\\\\
\\noindent /muril{\textprimstress}amen/\\\\
\\noindent Syllables: mur.il.am.en\\\\
\\noindent lit. sight+stone\\\\
\\noindent \\textit{Mrilgem}\\\\


\\noindent History:

\\vspace{-0pt}
\\hspace{40pt}
\\begin{tabular}{ccc}
\\textit{6000} & /emrilamern/&$\\rightarrow$ & \\textit{7500} & /mrilamern/&$\\rightarrow$ & \\textit{7600} & /mrilamen/&$\\rightarrow$ & \\textit{15001} & /murilamen/& \\\\
\\end{tabular}

\\vspace{20pt}\\hline

\\end{nopagebreak}
\\filbreak



\\vspace{15pt}
\\begin{nopagebreak}
\\noindent{\\fontsize{20pt}{10pt}\\textbf{gelarenan} } \\textit{7 of summer/June 7} (N)\\\\
\\noindent {\\tovian \\fontsize{20pt}{10pt} \\textbf{gelarenan} }\\\\
\\noindent /gelar{\textprimstress}enan/\\\\
\\noindent Syllables: gel.ar.en.an\\\\
\\noindent lit. tree+star\\\\
\\noindent \\textit{Starbirch}\\\\


\\noindent History:

\\vspace{-0pt}
\\hspace{40pt}
\\begin{tabular}{ccc}
\\textit{6000} & /aglarennaana/&$\\rightarrow$ & \\textit{6500} & /aglarenaana/&$\\rightarrow$ & \\textit{7500} & /glarenaana/&$\\rightarrow$ & \\textit{9500} & /glarenaan/&$\\rightarrow$ & \\textit{12005} & /glarenan/&$\\rightarrow$ & \\textit{15001} & /gelarenan/& \\\\
\\end{tabular}

\\vspace{20pt}\\hline

\\end{nopagebreak}
\\filbreak



\\vspace{15pt}
\\begin{nopagebreak}
\\noindent{\\fontsize{20pt}{10pt}\\textbf{muloanathan} } \\textit{7 of winter/december 7} (N)\\\\
\\noindent {\\tovian \\fontsize{20pt}{10pt} \\textbf{muloanathan} }\\\\
\\noindent /muloan{\textprimstress}a{\texttheta}an/\\\\
\\noindent Syllables: mul.o.an.aθ.an\\\\
\\noindent lit. morning+tool\\\\
\\noindent \\textit{Dawnthorn}\\\\


\\noindent History:

\\vspace{-0pt}
\\hspace{40pt}
\\begin{tabular}{ccc}
\\textit{6000} & /ploana{\dh}an/&$\\rightarrow$ & \\textit{11000} & /mloana{\dh}an/&$\\rightarrow$ & \\textit{12002} & /mloana{\texttheta}an/&$\\rightarrow$ & \\textit{15001} & /muloana{\texttheta}an/& \\\\
\\end{tabular}

\\vspace{20pt}\\hline

\\end{nopagebreak}
\\filbreak



\\vspace{15pt}
\\begin{nopagebreak}
\\noindent{\\fontsize{20pt}{10pt}\\textbf{murallhenathan} } \\textit{70 of autumn/November 10} (N)\\\\
\\noindent {\\tovian \\fontsize{20pt}{10pt} \\textbf{murallhenathan} }\\\\
\\noindent /mural{\textbeltl}en{\textprimstress}a{\texttheta}an/\\\\
\\noindent Syllables: mur.alɬ.en.aθ.an\\\\
\\noindent lit. darkness+dream+tool\\\\
\\noindent \\textit{Shedscythe}\\\\


\\noindent History:

\\vspace{-0pt}
\\hspace{40pt}
\\begin{tabular}{ccc}
\\textit{6000} & /bral{\textbeltl}ena{\dh}an/&$\\rightarrow$ & \\textit{11000} & /mral{\textbeltl}ena{\dh}an/&$\\rightarrow$ & \\textit{12002} & /mral{\textbeltl}ena{\texttheta}an/&$\\rightarrow$ & \\textit{15001} & /mural{\textbeltl}ena{\texttheta}an/& \\\\
\\end{tabular}

\\vspace{20pt}\\hline

\\end{nopagebreak}
\\filbreak



\\vspace{15pt}
\\begin{nopagebreak}
\\noindent{\\fontsize{20pt}{10pt}\\textbf{shotlol} } \\textit{70 of spring/May 10} (N)\\\\
\\noindent {\\tovian \\fontsize{20pt}{10pt} \\textbf{shotlol} }\\\\
\\noindent /{\textesh}{\textprimstress}o{\texttoptiebar{t\textbeltl}}ol/\\\\
\\noindent Syllables: ʃoƛ.ol\\\\
\\noindent lit. flowing+fruit\\\\
\\noindent \\textit{Mistquince}\\\\


\\noindent History:

\\vspace{-0pt}
\\hspace{40pt}
\\begin{tabular}{ccc}
\\textit{6000} & /{\textesh}o{\texttoptiebar{t\textbeltl}}{\textyogh}ola/&$\\rightarrow$ & \\textit{9500} & /{\textesh}o{\texttoptiebar{t\textbeltl}}{\textyogh}ol/&$\\rightarrow$ & \\textit{12001} & /{\textesh}o{\texttoptiebar{t\textbeltl}}{\textesh}ol/&$\\rightarrow$ & \\textit{14000} & /{\textesh}o{\texttoptiebar{t\textbeltl}}ol/& \\\\
\\end{tabular}

\\vspace{20pt}\\hline

\\end{nopagebreak}
\\filbreak



\\vspace{15pt}
\\begin{nopagebreak}
\\noindent{\\fontsize{20pt}{10pt}\\textbf{nanalor} } \\textit{70 of summer/August 10} (N)\\\\
\\noindent {\\tovian \\fontsize{20pt}{10pt} \\textbf{nanalor} }\\\\
\\noindent /nan{\textprimstress}alor/\\\\
\\noindent Syllables: nan.al.or\\\\
\\noindent lit. star+action\\\\
\\noindent \\textit{Starharvest}\\\\


\\noindent History:

\\vspace{-0pt}
\\hspace{40pt}
\\begin{tabular}{ccc}
\\textit{6000} & /nnaanalore/&$\\rightarrow$ & \\textit{6500} & /naanalore/&$\\rightarrow$ & \\textit{9500} & /naanalor/&$\\rightarrow$ & \\textit{12005} & /nanalor/& \\\\
\\end{tabular}

\\vspace{20pt}\\hline

\\end{nopagebreak}
\\filbreak



\\vspace{15pt}
\\begin{nopagebreak}
\\noindent{\\fontsize{20pt}{10pt}\\textbf{tushuthan} } \\textit{70 of winter/February 10} (N)\\\\
\\noindent {\\tovian \\fontsize{20pt}{10pt} \\textbf{tushuthan} }\\\\
\\noindent /{\textsubbridge{t}}u{\textesh}{\textprimstress}u{\texttheta}an/\\\\
\\noindent Syllables: tuʃ.uθ.an\\\\
\\noindent lit. cold+tool\\\\
\\noindent \\textit{Winterflail}\\\\


\\noindent History:

\\vspace{-0pt}
\\hspace{40pt}
\\begin{tabular}{ccc}
\\textit{6000} & /{\textsubbridge{t}}u{\textesh}u{\dh}an/&$\\rightarrow$ & \\textit{12002} & /{\textsubbridge{t}}u{\textesh}u{\texttheta}an/& \\\\
\\end{tabular}

\\vspace{20pt}\\hline

\\end{nopagebreak}
\\filbreak



\\vspace{15pt}
\\begin{nopagebreak}
\\noindent{\\fontsize{20pt}{10pt}\\textbf{shithanin} } \\textit{71 of autumn/November 11} (N)\\\\
\\noindent {\\tovian \\fontsize{20pt}{10pt} \\textbf{shithanin} }\\\\
\\noindent /{\textesh}i{\texttheta}{\textprimstress}anin/\\\\
\\noindent Syllables: ʃiθ.an.in\\\\
\\noindent lit. sharpness+berry\\\\
\\noindent \\textit{Stefbery}\\\\


\\noindent History:

\\vspace{-0pt}
\\hspace{40pt}
\\begin{tabular}{ccc}
\\textit{6000} & /{\textyogh}i{\texttheta}{\textyogh}lanina/&$\\rightarrow$ & \\textit{9500} & /{\textyogh}i{\texttheta}{\textyogh}lanin/&$\\rightarrow$ & \\textit{12001} & /{\textesh}i{\texttheta}{\textesh}lanin/&$\\rightarrow$ & \\textit{12500} & /{\textesh}i{\texttheta}{\textesh}anin/&$\\rightarrow$ & \\textit{13000} & /{\textesh}i{\texttheta}anin/& \\\\
\\end{tabular}

\\vspace{20pt}\\hline

\\end{nopagebreak}
\\filbreak



\\vspace{15pt}
\\begin{nopagebreak}
\\noindent{\\fontsize{20pt}{10pt}\\textbf{muramlorean} } \\textit{71 of spring/May 11} (N)\\\\
\\noindent {\\tovian \\fontsize{20pt}{10pt} \\textbf{muramlorean} }\\\\
\\noindent /muramlor{\textprimstress}ean/\\\\
\\noindent Syllables: mur.aml.or.e.an\\\\
\\noindent lit. darkness+plant+light\\\\
\\noindent \\textit{Nightshade}\\\\


\\noindent History:

\\vspace{-0pt}
\\hspace{40pt}
\\begin{tabular}{ccc}
\\textit{6000} & /bralmlloreana/&$\\rightarrow$ & \\textit{6500} & /bralmloreana/&$\\rightarrow$ & \\textit{7600} & /bramloreana/&$\\rightarrow$ & \\textit{9500} & /bramlorean/&$\\rightarrow$ & \\textit{11000} & /mramlorean/&$\\rightarrow$ & \\textit{15001} & /muramlorean/& \\\\
\\end{tabular}

\\vspace{20pt}\\hline

\\end{nopagebreak}
\\filbreak



\\vspace{15pt}
\\begin{nopagebreak}
\\noindent{\\fontsize{20pt}{10pt}\\textbf{lanathan} } \\textit{71 of summer/August 11} (N)\\\\
\\noindent {\\tovian \\fontsize{20pt}{10pt} \\textbf{lanathan} }\\\\
\\noindent /lan{\textprimstress}a{\texttheta}an/\\\\
\\noindent Syllables: lan.aθ.an\\\\
\\noindent lit. twilight+tool\\\\
\\noindent \\textit{Dusktang}\\\\


\\noindent History:

\\vspace{-0pt}
\\hspace{40pt}
\\begin{tabular}{ccc}
\\textit{6000} & /brlana{\dh}an/&$\\rightarrow$ & \\textit{7501} & /rlana{\dh}an/&$\\rightarrow$ & \\textit{7600} & /lana{\dh}an/&$\\rightarrow$ & \\textit{12002} & /lana{\texttheta}an/& \\\\
\\end{tabular}

\\vspace{20pt}\\hline

\\end{nopagebreak}
\\filbreak



\\vspace{15pt}
\\begin{nopagebreak}
\\noindent{\\fontsize{20pt}{10pt}\\textbf{thanramen} } \\textit{71 of winter/February 11} (N)\\\\
\\noindent {\\tovian \\fontsize{20pt}{10pt} \\textbf{thanramen} }\\\\
\\noindent /{\texttheta}anr{\textprimstress}amen/\\\\
\\noindent Syllables: θanr.am.en\\\\
\\noindent lit. tool+darkness+stone\\\\
\\noindent \\textit{Nightstone}\\\\


\\noindent History:

\\vspace{-0pt}
\\hspace{40pt}
\\begin{tabular}{ccc}
\\textit{6000} & /{\dh}anbralmern/&$\\rightarrow$ & \\textit{7600} & /{\dh}anbramen/&$\\rightarrow$ & \\textit{8500} & /{\dh}anramen/&$\\rightarrow$ & \\textit{12002} & /{\texttheta}anramen/& \\\\
\\end{tabular}

\\vspace{20pt}\\hline

\\end{nopagebreak}
\\filbreak



\\vspace{15pt}
\\begin{nopagebreak}
\\noindent{\\fontsize{20pt}{10pt}\\textbf{lethileaglar} } \\textit{72 of autumn/November 12} (N)\\\\
\\noindent {\\tovian \\fontsize{20pt}{10pt} \\textbf{lethileaglar} }\\\\
\\noindent /le{\texttheta}ile{\textprimstress}aglar/\\\\
\\noindent Syllables: leθ.il.e.agl.ar\\\\
\\noindent lit. forest+tree\\\\
\\noindent \\textit{Wildroot}\\\\


\\noindent History:

\\vspace{-0pt}
\\hspace{40pt}
\\begin{tabular}{ccc}
\\textit{6000} & /aglre{\texttheta}ileaglare/&$\\rightarrow$ & \\textit{6503} & /agrle{\texttheta}ileaglare/&$\\rightarrow$ & \\textit{7500} & /grle{\texttheta}ileaglare/&$\\rightarrow$ & \\textit{7501} & /rle{\texttheta}ileaglare/&$\\rightarrow$ & \\textit{7600} & /le{\texttheta}ileaglare/&$\\rightarrow$ & \\textit{9500} & /le{\texttheta}ileaglar/& \\\\
\\end{tabular}

\\vspace{20pt}\\hline

\\end{nopagebreak}
\\filbreak



\\vspace{15pt}
\\begin{nopagebreak}
\\noindent{\\fontsize{20pt}{10pt}\\textbf{deleleon} } \\textit{72 of spring/May 12} (N)\\\\
\\noindent {\\tovian \\fontsize{20pt}{10pt} \\textbf{deleleon} }\\\\
\\noindent /delel{\textprimstress}eon/\\\\
\\noindent Syllables: del.el.e.on\\\\
\\noindent lit. trunk+mouth\\\\
\\noindent \\textit{Starstalk}\\\\


\\noindent History:

\\vspace{-0pt}
\\hspace{40pt}
\\begin{tabular}{ccc}
\\textit{6000} & /adllenleoln/&$\\rightarrow$ & \\textit{6500} & /adlenleoln/&$\\rightarrow$ & \\textit{7500} & /dlenleoln/&$\\rightarrow$ & \\textit{7600} & /dleleon/&$\\rightarrow$ & \\textit{15001} & /deleleon/& \\\\
\\end{tabular}

\\vspace{20pt}\\hline

\\end{nopagebreak}
\\filbreak



\\vspace{15pt}
\\begin{nopagebreak}
\\noindent{\\fontsize{20pt}{10pt}\\textbf{muloanatesh} } \\textit{72 of summer/August 12} (N)\\\\
\\noindent {\\tovian \\fontsize{20pt}{10pt} \\textbf{muloanatesh} }\\\\
\\noindent /muloan{\textprimstress}a{\textsubbridge{t}}e{\textesh}/\\\\
\\noindent Syllables: mul.o.an.at.eʃ\\\\
\\noindent lit. morning+sound\\\\
\\noindent \\textit{Dawnvox}\\\\


\\noindent History:

\\vspace{-0pt}
\\hspace{40pt}
\\begin{tabular}{ccc}
\\textit{6000} & /ploana{\textsubbridge{t}}e{\textyogh}/&$\\rightarrow$ & \\textit{11000} & /mloana{\textsubbridge{t}}e{\textyogh}/&$\\rightarrow$ & \\textit{12001} & /mloana{\textsubbridge{t}}e{\textesh}/&$\\rightarrow$ & \\textit{15001} & /muloana{\textsubbridge{t}}e{\textesh}/& \\\\
\\end{tabular}

\\vspace{20pt}\\hline

\\end{nopagebreak}
\\filbreak



\\vspace{15pt}
\\begin{nopagebreak}
\\noindent{\\fontsize{20pt}{10pt}\\textbf{numvar} } \\textit{72 of winter/February 12} (N)\\\\
\\noindent {\\tovian \\fontsize{20pt}{10pt} \\textbf{numvar} }\\\\
\\noindent /n{\textprimstress}umvar/\\\\
\\noindent Syllables: numv.ar\\\\
\\noindent lit. path+value\\\\
\\noindent \\textit{Numkar}\\\\


\\noindent History:

\\vspace{-0pt}
\\hspace{40pt}
\\begin{tabular}{ccc}
\\textit{6000} & /numvlare/&$\\rightarrow$ & \\textit{9500} & /numvlar/&$\\rightarrow$ & \\textit{12500} & /numvar/& \\\\
\\end{tabular}

\\vspace{20pt}\\hline

\\end{nopagebreak}
\\filbreak



\\vspace{15pt}
\\begin{nopagebreak}
\\noindent{\\fontsize{20pt}{10pt}\\textbf{kanaithr} } \\textit{73 of autumn/November 13} (N)\\\\
\\noindent {\\tovian \\fontsize{20pt}{10pt} \\textbf{kanaithr} }\\\\
\\noindent /kan{\textprimstress}ai{\texttheta}r/\\\\
\\noindent Syllables: kan.a.iθr\\\\
\\noindent lit. change+thought\\\\
\\noindent \\textit{Kanalh}\\\\


\\noindent History:

\\vspace{-0pt}
\\hspace{40pt}
\\begin{tabular}{ccc}
\\textit{6000} & /kanai{\texttheta}r/& \\\\
\\end{tabular}

\\vspace{20pt}\\hline

\\end{nopagebreak}
\\filbreak



\\vspace{15pt}
\\begin{nopagebreak}
\\noindent{\\fontsize{20pt}{10pt}\\textbf{???} } \\textit{73 of spring/May 13} (N)\\\\
\\noindent {\\tovian \\fontsize{20pt}{10pt} \\textbf{???} }\\\\
\\noindent /???/\\\\
\\noindent Syllables: ???\\\\
\\noindent lit. ???\\\\
\\noindent \\textit{Bayra}\\\\


\\noindent History:

\\vspace{-0pt}
\\hspace{40pt}
\\begin{tabular}{ccc}
\\textit{6000} & /???/& \\\\
\\end{tabular}

\\vspace{20pt}\\hline

\\end{nopagebreak}
\\filbreak



\\vspace{15pt}
\\begin{nopagebreak}
\\noindent{\\fontsize{20pt}{10pt}\\textbf{efan} } \\textit{73 of summer/August 13} (N)\\\\
\\noindent {\\tovian \\fontsize{20pt}{10pt} \\textbf{efan} }\\\\
\\noindent /{\textprimstress}efan/\\\\
\\noindent Syllables: ef.an\\\\
\\noindent lit. sun+tool\\\\
\\noindent \\textit{Sunstaff}\\\\


\\noindent History:

\\vspace{-0pt}
\\hspace{40pt}
\\begin{tabular}{ccc}
\\textit{6000} & /efl{\dh}an/&$\\rightarrow$ & \\textit{12002} & /efl{\texttheta}an/&$\\rightarrow$ & \\textit{12500} & /ef{\texttheta}an/&$\\rightarrow$ & \\textit{13000} & /efan/& \\\\
\\end{tabular}

\\vspace{20pt}\\hline

\\end{nopagebreak}
\\filbreak



\\vspace{15pt}
\\begin{nopagebreak}
\\noindent{\\fontsize{20pt}{10pt}\\textbf{muramen} } \\textit{73 of winter/February 13} (N)\\\\
\\noindent {\\tovian \\fontsize{20pt}{10pt} \\textbf{muramen} }\\\\
\\noindent /mur{\textprimstress}amen/\\\\
\\noindent Syllables: mur.am.en\\\\
\\noindent lit. darkness+stone\\\\
\\noindent \\textit{Shedowslate}\\\\


\\noindent History:

\\vspace{-0pt}
\\hspace{40pt}
\\begin{tabular}{ccc}
\\textit{6000} & /bralmern/&$\\rightarrow$ & \\textit{7600} & /bramen/&$\\rightarrow$ & \\textit{11000} & /mramen/&$\\rightarrow$ & \\textit{15001} & /muramen/& \\\\
\\end{tabular}

\\vspace{20pt}\\hline

\\end{nopagebreak}
\\filbreak



\\vspace{15pt}
\\begin{nopagebreak}
\\noindent{\\fontsize{20pt}{10pt}\\textbf{sillhum} } \\textit{74 of autumn/November 14} (N)\\\\
\\noindent {\\tovian \\fontsize{20pt}{10pt} \\textbf{sillhum} }\\\\
\\noindent /s{\textprimstress}il{\textbeltl}um/\\\\
\\noindent Syllables: silɬ.um\\\\
\\noindent lit. existence+goodness\\\\
\\noindent \\textit{Silshede}\\\\


\\noindent History:

\\vspace{-0pt}
\\hspace{40pt}
\\begin{tabular}{ccc}
\\textit{6000} & /sil{\textbeltl}um/& \\\\
\\end{tabular}

\\vspace{20pt}\\hline

\\end{nopagebreak}
\\filbreak



\\vspace{15pt}
\\begin{nopagebreak}
\\noindent{\\fontsize{20pt}{10pt}\\textbf{shotl} } \\textit{74 of spring/May 14} (N)\\\\
\\noindent {\\tovian \\fontsize{20pt}{10pt} \\textbf{shotl} }\\\\
\\noindent /{\textesh}{\textprimstress}o{\texttoptiebar{t\textbeltl}}/\\\\
\\noindent Syllables: ʃoƛ\\\\
\\noindent lit. flowing+???\\\\
\\noindent \\textit{Misazel}\\\\


\\noindent History:

\\vspace{-0pt}
\\hspace{40pt}
\\begin{tabular}{ccc}
\\textit{6000} & /{\textesh}o{\texttoptiebar{t\textbeltl}}???/&$\\rightarrow$ & \\textit{8500} & /{\textesh}o{\texttoptiebar{t\textbeltl}}/& \\\\
\\end{tabular}

\\vspace{20pt}\\hline

\\end{nopagebreak}
\\filbreak



\\vspace{15pt}
\\begin{nopagebreak}
\\noindent{\\fontsize{20pt}{10pt}\\textbf{gelaretesh} } \\textit{74 of summer/August 14} (N)\\\\
\\noindent {\\tovian \\fontsize{20pt}{10pt} \\textbf{gelaretesh} }\\\\
\\noindent /gelar{\textprimstress}e{\textsubbridge{t}}e{\textesh}/\\\\
\\noindent Syllables: gel.ar.et.eʃ\\\\
\\noindent lit. tree+sound\\\\
\\noindent \\textit{Palmwhisper}\\\\


\\noindent History:

\\vspace{-0pt}
\\hspace{40pt}
\\begin{tabular}{ccc}
\\textit{6000} & /aglare{\textsubbridge{t}}e{\textyogh}/&$\\rightarrow$ & \\textit{7500} & /glare{\textsubbridge{t}}e{\textyogh}/&$\\rightarrow$ & \\textit{12001} & /glare{\textsubbridge{t}}e{\textesh}/&$\\rightarrow$ & \\textit{15001} & /gelare{\textsubbridge{t}}e{\textesh}/& \\\\
\\end{tabular}

\\vspace{20pt}\\hline

\\end{nopagebreak}
\\filbreak



\\vspace{15pt}
\\begin{nopagebreak}
\\noindent{\\fontsize{20pt}{10pt}\\textbf{mrfor} } \\textit{74 of winter/February 14} (N)\\\\
\\noindent {\\tovian \\fontsize{20pt}{10pt} \\textbf{mrfor} }\\\\
\\noindent /mrf{\textprimstress}or/\\\\
\\noindent Syllables: mrfor\\\\
\\noindent lit. earth+writing\\\\
\\noindent \\textit{Peatstone}\\\\


\\noindent History:

\\vspace{-0pt}
\\hspace{40pt}
\\begin{tabular}{ccc}
\\textit{6000} & /mernflore/&$\\rightarrow$ & \\textit{6400} & /merflore/&$\\rightarrow$ & \\textit{9500} & /merflor/&$\\rightarrow$ & \\textit{12500} & /merfor/&$\\rightarrow$ & \\textit{16002} & /mrfor/& \\\\
\\end{tabular}

\\vspace{20pt}\\hline

\\end{nopagebreak}
\\filbreak



\\vspace{15pt}
\\begin{nopagebreak}
\\noindent{\\fontsize{20pt}{10pt}\\textbf{muloanarithan} } \\textit{75 of autumn/November 15} (N)\\\\
\\noindent {\\tovian \\fontsize{20pt}{10pt} \\textbf{muloanarithan} }\\\\
\\noindent /muloanar{\textprimstress}i{\texttheta}an/\\\\
\\noindent Syllables: mul.o.an.ar.iθ.an\\\\
\\noindent lit. morning+bird\\\\
\\noindent \\textit{Dawnquail}\\\\


\\noindent History:

\\vspace{-0pt}
\\hspace{40pt}
\\begin{tabular}{ccc}
\\textit{6000} & /ploanari{\texttheta}lamn/&$\\rightarrow$ & \\textit{7600} & /ploanari{\texttheta}lan/&$\\rightarrow$ & \\textit{11000} & /mloanari{\texttheta}lan/&$\\rightarrow$ & \\textit{12500} & /mloanari{\texttheta}an/&$\\rightarrow$ & \\textit{15001} & /muloanari{\texttheta}an/& \\\\
\\end{tabular}

\\vspace{20pt}\\hline

\\end{nopagebreak}
\\filbreak



\\vspace{15pt}
\\begin{nopagebreak}
\\noindent{\\fontsize{20pt}{10pt}\\textbf{aysan} } \\textit{75 of spring/May 15} (N)\\\\
\\noindent {\\tovian \\fontsize{20pt}{10pt} \\textbf{aysan} }\\\\
\\noindent /{\textprimstress}aysan/\\\\
\\noindent Syllables: ajs.an\\\\
\\noindent lit. moon+animal\\\\
\\noindent \\textit{Mooncalf}\\\\


\\noindent History:

\\vspace{-0pt}
\\hspace{40pt}
\\begin{tabular}{ccc}
\\textit{6000} & /aliys{\textsubbridge{t}}lamn/&$\\rightarrow$ & \\textit{6300} & /alys{\textsubbridge{t}}lamn/&$\\rightarrow$ & \\textit{6401} & /alys{\texttoptiebar{t\textbeltl}}amn/&$\\rightarrow$ & \\textit{7600} & /ays{\texttoptiebar{t\textbeltl}}an/&$\\rightarrow$ & \\textit{14001} & /aysan/& \\\\
\\end{tabular}

\\vspace{20pt}\\hline

\\end{nopagebreak}
\\filbreak



\\vspace{15pt}
\\begin{nopagebreak}
\\noindent{\\fontsize{20pt}{10pt}\\textbf{holalor} } \\textit{75 of summer/August 15} (N)\\\\
\\noindent {\\tovian \\fontsize{20pt}{10pt} \\textbf{holalor} }\\\\
\\noindent /ħol{\textprimstress}alor/\\\\
\\noindent Syllables: ħol.al.or\\\\
\\noindent lit. sweet+action\\\\
\\noindent \\textit{Silkspinner}\\\\


\\noindent History:

\\vspace{-0pt}
\\hspace{40pt}
\\begin{tabular}{ccc}
\\textit{6000} & /holalore/&$\\rightarrow$ & \\textit{9500} & /holalor/&$\\rightarrow$ & \\textit{14006} & /ħolalor/& \\\\
\\end{tabular}

\\vspace{20pt}\\hline

\\end{nopagebreak}
\\filbreak



\\vspace{15pt}
\\begin{nopagebreak}
\\noindent{\\fontsize{20pt}{10pt}\\textbf{shashol} } \\textit{75 of winter/February 15} (N)\\\\
\\noindent {\\tovian \\fontsize{20pt}{10pt} \\textbf{shashol} }\\\\
\\noindent /{\textesh}{\textprimstress}a{\textesh}ol/\\\\
\\noindent Syllables: ʃaʃ.ol\\\\
\\noindent lit. ice+fruit\\\\
\\noindent \\textit{Frostbun}\\\\


\\noindent History:

\\vspace{-0pt}
\\hspace{40pt}
\\begin{tabular}{ccc}
\\textit{6000} & /{\textesh}u{\textesh}a{\textyogh}ola/&$\\rightarrow$ & \\textit{9200} & /{\textesh}{\textschwa}{\textesh}a{\textyogh}ola/&$\\rightarrow$ & \\textit{9500} & /{\textesh}{\textschwa}{\textesh}a{\textyogh}ol/&$\\rightarrow$ & \\textit{11990} & /{\textesh}{\textesh}a{\textyogh}ol/&$\\rightarrow$ & \\textit{12001} & /{\textesh}{\textesh}a{\textesh}ol/&$\\rightarrow$ & \\textit{12007} & /{\textesh}a{\textesh}ol/& \\\\
\\end{tabular}

\\vspace{20pt}\\hline

\\end{nopagebreak}
\\filbreak



\\vspace{15pt}
\\begin{nopagebreak}
\\noindent{\\fontsize{20pt}{10pt}\\textbf{fedemralrem} } \\textit{76 of autumn/November 16} (N)\\\\
\\noindent {\\tovian \\fontsize{20pt}{10pt} \\textbf{fedemralrem} }\\\\
\\noindent /fedemr{\textprimstress}alrem/\\\\
\\noindent Syllables: fed.emr.alr.em\\\\
\\noindent lit. showing+darkness+flower\\\\
\\noindent \\textit{Duflor}\\\\


\\noindent History:

\\vspace{-0pt}
\\hspace{40pt}
\\begin{tabular}{ccc}
\\textit{6000} & /fedebraldremp/&$\\rightarrow$ & \\textit{8750} & /fedebralrem/&$\\rightarrow$ & \\textit{11000} & /fedemralrem/& \\\\
\\end{tabular}

\\vspace{20pt}\\hline

\\end{nopagebreak}
\\filbreak



\\vspace{15pt}
\\begin{nopagebreak}
\\noindent{\\fontsize{20pt}{10pt}\\textbf{rithilesyaleth} } \\textit{76 of spring/May 16} (N)\\\\
\\noindent {\\tovian \\fontsize{20pt}{10pt} \\textbf{rithilesyaleth} }\\\\
\\noindent /ri{\texttheta}ilesy{\textprimstress}ale{\texttheta}/\\\\
\\noindent Syllables: riθ.il.esj.al.eθ\\\\
\\noindent lit. sky+food\\\\
\\noindent \\textit{Sky's Feed}\\\\


\\noindent History:

\\vspace{-0pt}
\\hspace{40pt}
\\begin{tabular}{ccc}
\\textit{6000} & /ri{\texttheta}ilesyale{\dh}a/&$\\rightarrow$ & \\textit{9500} & /ri{\texttheta}ilesyale{\dh}/&$\\rightarrow$ & \\textit{12002} & /ri{\texttheta}ilesyale{\texttheta}/& \\\\
\\end{tabular}

\\vspace{20pt}\\hline

\\end{nopagebreak}
\\filbreak



\\vspace{15pt}
\\begin{nopagebreak}
\\noindent{\\fontsize{20pt}{10pt}\\textbf{melrem} } \\textit{76 of summer/August 16} (N)\\\\
\\noindent {\\tovian \\fontsize{20pt}{10pt} \\textbf{melrem} }\\\\
\\noindent /m{\textprimstress}elrem/\\\\
\\noindent Syllables: melr.em\\\\
\\noindent lit. life+flower\\\\
\\noindent \\textit{Healbloom}\\\\


\\noindent History:

\\vspace{-0pt}
\\hspace{40pt}
\\begin{tabular}{ccc}
\\textit{6000} & /meldremp/&$\\rightarrow$ & \\textit{8750} & /melrem/& \\\\
\\end{tabular}

\\vspace{20pt}\\hline

\\end{nopagebreak}
\\filbreak



\\vspace{15pt}
\\begin{nopagebreak}
\\noindent{\\fontsize{20pt}{10pt}\\textbf{namen} } \\textit{76 of winter/February 16} (N)\\\\
\\noindent {\\tovian \\fontsize{20pt}{10pt} \\textbf{namen} }\\\\
\\noindent /n{\textprimstress}amen/\\\\
\\noindent Syllables: nam.en\\\\
\\noindent lit. light+stone\\\\
\\noindent \\textit{Sparkstone}\\\\


\\noindent History:

\\vspace{-0pt}
\\hspace{40pt}
\\begin{tabular}{ccc}
\\textit{6000} & /anamern/&$\\rightarrow$ & \\textit{7500} & /namern/&$\\rightarrow$ & \\textit{7600} & /namen/& \\\\
\\end{tabular}

\\vspace{20pt}\\hline

\\end{nopagebreak}
\\filbreak



\\vspace{15pt}
\\begin{nopagebreak}
\\noindent{\\fontsize{20pt}{10pt}\\textbf{lhiladrem} } \\textit{77 of autumn/November 17} (N)\\\\
\\noindent {\\tovian \\fontsize{20pt}{10pt} \\textbf{lhiladrem} }\\\\
\\noindent /{\textbeltl}il{\textprimstress}adrem/\\\\
\\noindent Syllables: ɬil.adr.em\\\\
\\noindent lit. night+flower\\\\
\\noindent \\textit{Niteblum}\\\\


\\noindent History:

\\vspace{-0pt}
\\hspace{40pt}
\\begin{tabular}{ccc}
\\textit{6000} & /a{\textbeltl}inladremp/&$\\rightarrow$ & \\textit{7500} & /{\textbeltl}inladremp/&$\\rightarrow$ & \\textit{7600} & /{\textbeltl}iladremp/&$\\rightarrow$ & \\textit{8750} & /{\textbeltl}iladrem/& \\\\
\\end{tabular}

\\vspace{20pt}\\hline

\\end{nopagebreak}
\\filbreak



\\vspace{15pt}
\\begin{nopagebreak}
\\noindent{\\fontsize{20pt}{10pt}\\textbf{kelueon} } \\textit{77 of spring/May 17} (N)\\\\
\\noindent {\\tovian \\fontsize{20pt}{10pt} \\textbf{kelueon} }\\\\
\\noindent /kelu{\textprimstress}eon/\\\\
\\noindent Syllables: kel.u.e.on\\\\
\\noindent lit. branch+mouth\\\\
\\noindent \\textit{Starlichen}\\\\


\\noindent History:

\\vspace{-0pt}
\\hspace{40pt}
\\begin{tabular}{ccc}
\\textit{6000} & /kelueoln/&$\\rightarrow$ & \\textit{7600} & /kelueon/& \\\\
\\end{tabular}

\\vspace{20pt}\\hline

\\end{nopagebreak}
\\filbreak



\\vspace{15pt}
\\begin{nopagebreak}
\\noindent{\\fontsize{20pt}{10pt}\\textbf{keluthan} } \\textit{77 of summer/August 17} (N)\\\\
\\noindent {\\tovian \\fontsize{20pt}{10pt} \\textbf{keluthan} }\\\\
\\noindent /kel{\textprimstress}u{\texttheta}an/\\\\
\\noindent Syllables: kel.uθ.an\\\\
\\noindent lit. branch+tool\\\\
\\noindent \\textit{Starpim}\\\\


\\noindent History:

\\vspace{-0pt}
\\hspace{40pt}
\\begin{tabular}{ccc}
\\textit{6000} & /kelu{\dh}an/&$\\rightarrow$ & \\textit{12002} & /kelu{\texttheta}an/& \\\\
\\end{tabular}

\\vspace{20pt}\\hline

\\end{nopagebreak}
\\filbreak



\\vspace{15pt}
\\begin{nopagebreak}
\\noindent{\\fontsize{20pt}{10pt}\\textbf{lhenamen} } \\textit{77 of winter/February 17} (N)\\\\
\\noindent {\\tovian \\fontsize{20pt}{10pt} \\textbf{lhenamen} }\\\\
\\noindent /{\textbeltl}en{\textprimstress}amen/\\\\
\\noindent Syllables: ɬen.am.en\\\\
\\noindent lit. dream+earth\\\\
\\noindent \\textit{Dreamclay}\\\\


\\noindent History:

\\vspace{-0pt}
\\hspace{40pt}
\\begin{tabular}{ccc}
\\textit{6000} & /{\textbeltl}enamern/&$\\rightarrow$ & \\textit{7600} & /{\textbeltl}enamen/& \\\\
\\end{tabular}

\\vspace{20pt}\\hline

\\end{nopagebreak}
\\filbreak



\\vspace{15pt}
\\begin{nopagebreak}
\\noindent{\\fontsize{20pt}{10pt}\\textbf{lhenakel} } \\textit{78 of autumn/November 18} (N)\\\\
\\noindent {\\tovian \\fontsize{20pt}{10pt} \\textbf{lhenakel} }\\\\
\\noindent /{\textbeltl}en{\textprimstress}akel/\\\\
\\noindent Syllables: ɬen.ak.el\\\\
\\noindent lit. dream+branch\\\\
\\noindent \\textit{Dremsweed}\\\\


\\noindent History:

\\vspace{-0pt}
\\hspace{40pt}
\\begin{tabular}{ccc}
\\textit{6000} & /{\textbeltl}enakelu/&$\\rightarrow$ & \\textit{9500} & /{\textbeltl}enakel/& \\\\
\\end{tabular}

\\vspace{20pt}\\hline

\\end{nopagebreak}
\\filbreak



\\vspace{15pt}
\\begin{nopagebreak}
\\noindent{\\fontsize{20pt}{10pt}\\textbf{muralel} } \\textit{78 of spring/May 18} (N)\\\\
\\noindent {\\tovian \\fontsize{20pt}{10pt} \\textbf{muralel} }\\\\
\\noindent /mur{\textprimstress}alel/\\\\
\\noindent Syllables: mur.al.el\\\\
\\noindent lit. darkness+branch\\\\
\\noindent \\textit{Nightyew}\\\\


\\noindent History:

\\vspace{-0pt}
\\hspace{40pt}
\\begin{tabular}{ccc}
\\textit{6000} & /bralkelu/&$\\rightarrow$ & \\textit{8750} & /bralelu/&$\\rightarrow$ & \\textit{9500} & /bralel/&$\\rightarrow$ & \\textit{11000} & /mralel/&$\\rightarrow$ & \\textit{15001} & /muralel/& \\\\
\\end{tabular}

\\vspace{20pt}\\hline

\\end{nopagebreak}
\\filbreak



\\vspace{15pt}
\\begin{nopagebreak}
\\noindent{\\fontsize{20pt}{10pt}\\textbf{varetlun} } \\textit{78 of summer/August 18} (N)\\\\
\\noindent {\\tovian \\fontsize{20pt}{10pt} \\textbf{varetlun} }\\\\
\\noindent /var{\textprimstress}e{\texttoptiebar{t\textbeltl}}un/\\\\
\\noindent Syllables: var.eƛ.un\\\\
\\noindent lit. value+joining\\\\
\\noindent \\textit{Goldweave}\\\\


\\noindent History:

\\vspace{-0pt}
\\hspace{40pt}
\\begin{tabular}{ccc}
\\textit{6000} & /vlare{\texttoptiebar{t\textbeltl}}una/&$\\rightarrow$ & \\textit{9500} & /vlare{\texttoptiebar{t\textbeltl}}un/&$\\rightarrow$ & \\textit{12500} & /vare{\texttoptiebar{t\textbeltl}}un/& \\\\
\\end{tabular}

\\vspace{20pt}\\hline

\\end{nopagebreak}
\\filbreak



\\vspace{15pt}
\\begin{nopagebreak}
\\noindent{\\fontsize{20pt}{10pt}\\textbf{yisdeshol} } \\textit{78 of winter/February 18} (N)\\\\
\\noindent {\\tovian \\fontsize{20pt}{10pt} \\textbf{yisdeshol} }\\\\
\\noindent /yisd{\textprimstress}e{\textesh}ol/\\\\
\\noindent Syllables: jisd.eʃ.ol\\\\
\\noindent lit. moon+showing+fruit\\\\
\\noindent \\textit{Moonlime}\\\\


\\noindent History:

\\vspace{-0pt}
\\hspace{40pt}
\\begin{tabular}{ccc}
\\textit{6000} & /aliysfede{\textyogh}ola/&$\\rightarrow$ & \\textit{6300} & /alysfede{\textyogh}ola/&$\\rightarrow$ & \\textit{7500} & /lysfede{\textyogh}ola/&$\\rightarrow$ & \\textit{7501} & /ysfede{\textyogh}ola/&$\\rightarrow$ & \\textit{9500} & /ysfede{\textyogh}ol/&$\\rightarrow$ & \\textit{12001} & /ysfede{\textesh}ol/&$\\rightarrow$ & \\textit{12003} & /ysde{\textesh}ol/&$\\rightarrow$ & \\textit{15001} & /yisde{\textesh}ol/& \\\\
\\end{tabular}

\\vspace{20pt}\\hline

\\end{nopagebreak}
\\filbreak



\\vspace{15pt}
\\begin{nopagebreak}
\\noindent{\\fontsize{20pt}{10pt}\\textbf{nanathang} } \\textit{79 of autumn/November 19} (N)\\\\
\\noindent {\\tovian \\fontsize{20pt}{10pt} \\textbf{nanathang} }\\\\
\\noindent /nan{\textprimstress}a{\texttheta}a{\ng}/\\\\
\\noindent Syllables: nan.aθ.aŋ\\\\
\\noindent lit. star+leaf\\\\
\\noindent \\textit{Starlef}\\\\


\\noindent History:

\\vspace{-0pt}
\\hspace{40pt}
\\begin{tabular}{ccc}
\\textit{6000} & /nnaanaa{\dh}a{\ng}ge/&$\\rightarrow$ & \\textit{6500} & /naanaa{\dh}a{\ng}ge/&$\\rightarrow$ & \\textit{8750} & /naanaa{\dh}a{\ng}e/&$\\rightarrow$ & \\textit{9500} & /naanaa{\dh}a{\ng}/&$\\rightarrow$ & \\textit{12002} & /naanaa{\texttheta}a{\ng}/&$\\rightarrow$ & \\textit{12005} & /nana{\texttheta}a{\ng}/& \\\\
\\end{tabular}

\\vspace{20pt}\\hline

\\end{nopagebreak}
\\filbreak



\\vspace{15pt}
\\begin{nopagebreak}
\\noindent{\\fontsize{20pt}{10pt}\\textbf{yisor} } \\textit{79 of spring/May 19} (N)\\\\
\\noindent {\\tovian \\fontsize{20pt}{10pt} \\textbf{yisor} }\\\\
\\noindent /y{\textprimstress}isor/\\\\
\\noindent Syllables: jis.or\\\\
\\noindent lit. moon+plant\\\\
\\noindent \\textit{Moongrass}\\\\


\\noindent History:

\\vspace{-0pt}
\\hspace{40pt}
\\begin{tabular}{ccc}
\\textit{6000} & /aliysmllore/&$\\rightarrow$ & \\textit{6300} & /alysmllore/&$\\rightarrow$ & \\textit{6500} & /alysmlore/&$\\rightarrow$ & \\textit{7500} & /lysmlore/&$\\rightarrow$ & \\textit{7501} & /ysmlore/&$\\rightarrow$ & \\textit{7600} & /yslore/&$\\rightarrow$ & \\textit{9500} & /yslor/&$\\rightarrow$ & \\textit{12500} & /ysor/&$\\rightarrow$ & \\textit{15001} & /yisor/& \\\\
\\end{tabular}

\\vspace{20pt}\\hline

\\end{nopagebreak}
\\filbreak



\\vspace{15pt}
\\begin{nopagebreak}
\\noindent{\\fontsize{20pt}{10pt}\\textbf{muralaysor} } \\textit{79 of summer/August 19} (N)\\\\
\\noindent {\\tovian \\fontsize{20pt}{10pt} \\textbf{muralaysor} }\\\\
\\noindent /mural{\textprimstress}aysor/\\\\
\\noindent Syllables: mur.al.ajs.or\\\\
\\noindent lit. darkness+moon+plant\\\\
\\noindent \\textit{Shadegreen}\\\\


\\noindent History:

\\vspace{-0pt}
\\hspace{40pt}
\\begin{tabular}{ccc}
\\textit{6000} & /bralaliysmllore/&$\\rightarrow$ & \\textit{6300} & /bralalysmllore/&$\\rightarrow$ & \\textit{6500} & /bralalysmlore/&$\\rightarrow$ & \\textit{7600} & /bralayslore/&$\\rightarrow$ & \\textit{9500} & /bralayslor/&$\\rightarrow$ & \\textit{11000} & /mralayslor/&$\\rightarrow$ & \\textit{12500} & /mralaysor/&$\\rightarrow$ & \\textit{15001} & /muralaysor/& \\\\
\\end{tabular}

\\vspace{20pt}\\hline

\\end{nopagebreak}
\\filbreak



\\vspace{15pt}
\\begin{nopagebreak}
\\noindent{\\fontsize{20pt}{10pt}\\textbf{nanalhen} } \\textit{79 of winter/February 19} (N)\\\\
\\noindent {\\tovian \\fontsize{20pt}{10pt} \\textbf{nanalhen} }\\\\
\\noindent /nan{\textprimstress}a{\textbeltl}en/\\\\
\\noindent Syllables: nan.aɬ.en\\\\
\\noindent lit. star+dream\\\\
\\noindent \\textit{Starble}\\\\


\\noindent History:

\\vspace{-0pt}
\\hspace{40pt}
\\begin{tabular}{ccc}
\\textit{6000} & /nnaana{\textbeltl}ena/&$\\rightarrow$ & \\textit{6500} & /naana{\textbeltl}ena/&$\\rightarrow$ & \\textit{9500} & /naana{\textbeltl}en/&$\\rightarrow$ & \\textit{12005} & /nana{\textbeltl}en/& \\\\
\\end{tabular}

\\vspace{20pt}\\hline

\\end{nopagebreak}
\\filbreak



\\vspace{15pt}
\\begin{nopagebreak}
\\noindent{\\fontsize{20pt}{10pt}\\textbf{munorhanapal} } \\textit{8 of autumn/September 8} (N)\\\\
\\noindent {\\tovian \\fontsize{20pt}{10pt} \\textbf{munorhanapal} }\\\\
\\noindent /munorħan{\textprimstress}apal/\\\\
\\noindent Syllables: mun.orħ.an.ap.al\\\\
\\noindent lit. dusk+stopping\\\\
\\noindent \\textit{Duskshalt}\\\\


\\noindent History:

\\vspace{-0pt}
\\hspace{40pt}
\\begin{tabular}{ccc}
\\textit{6000} & /pnorhanapalr/&$\\rightarrow$ & \\textit{6503} & /pnorhanaparl/&$\\rightarrow$ & \\textit{7600} & /pnorhanapal/&$\\rightarrow$ & \\textit{11000} & /mnorhanapal/&$\\rightarrow$ & \\textit{11996} & /munorhanapal/&$\\rightarrow$ & \\textit{14006} & /munorħanapal/& \\\\
\\end{tabular}

\\vspace{20pt}\\hline

\\end{nopagebreak}
\\filbreak



\\vspace{15pt}
\\begin{nopagebreak}
\\noindent{\\fontsize{20pt}{10pt}\\textbf{todakesh} } \\textit{8 of spring/March 8} (N)\\\\
\\noindent {\\tovian \\fontsize{20pt}{10pt} \\textbf{todakesh} }\\\\
\\noindent /{\textsubbridge{t}}od{\textprimstress}ake{\textesh}/\\\\
\\noindent Syllables: tod.ak.eʃ\\\\
\\noindent lit. eternal+fire\\\\
\\noindent \\textit{Everflame}\\\\


\\noindent History:

\\vspace{-0pt}
\\hspace{40pt}
\\begin{tabular}{ccc}
\\textit{6000} & /{\textsubbridge{t}}odake{\textesh}/& \\\\
\\end{tabular}

\\vspace{20pt}\\hline

\\end{nopagebreak}
\\filbreak



\\vspace{15pt}
\\begin{nopagebreak}
\\noindent{\\fontsize{20pt}{10pt}\\textbf{nalhom} } \\textit{8 of summer/June 8} (N)\\\\
\\noindent {\\tovian \\fontsize{20pt}{10pt} \\textbf{nalhom} }\\\\
\\noindent /n{\textprimstress}a{\textbeltl}om/\\\\
\\noindent Syllables: naɬ.om\\\\
\\noindent lit. light+container\\\\
\\noindent \\textit{Lightcup}\\\\


\\noindent History:

\\vspace{-0pt}
\\hspace{40pt}
\\begin{tabular}{ccc}
\\textit{6000} & /ana{\textbeltl}oma/&$\\rightarrow$ & \\textit{7500} & /na{\textbeltl}oma/&$\\rightarrow$ & \\textit{9500} & /na{\textbeltl}om/& \\\\
\\end{tabular}

\\vspace{20pt}\\hline

\\end{nopagebreak}
\\filbreak



\\vspace{15pt}
\\begin{nopagebreak}
\\noindent{\\fontsize{20pt}{10pt}\\textbf{dolrem} } \\textit{8 of winter/december 8} (N)\\\\
\\noindent {\\tovian \\fontsize{20pt}{10pt} \\textbf{dolrem} }\\\\
\\noindent /d{\textprimstress}olrem/\\\\
\\noindent Syllables: dolr.em\\\\
\\noindent lit. death+flower\\\\
\\noindent \\textit{Dappledviolet}\\\\


\\noindent History:

\\vspace{-0pt}
\\hspace{40pt}
\\begin{tabular}{ccc}
\\textit{6000} & /doldremp/&$\\rightarrow$ & \\textit{8750} & /dolrem/& \\\\
\\end{tabular}

\\vspace{20pt}\\hline

\\end{nopagebreak}
\\filbreak



\\vspace{15pt}
\\begin{nopagebreak}
\\noindent{\\fontsize{20pt}{10pt}\\textbf{nathan} } \\textit{80 of autumn/November 20} (N)\\\\
\\noindent {\\tovian \\fontsize{20pt}{10pt} \\textbf{nathan} }\\\\
\\noindent /n{\textprimstress}a{\texttheta}an/\\\\
\\noindent Syllables: naθ.an\\\\
\\noindent lit. light+tool\\\\
\\noindent \\textit{Glemfork}\\\\


\\noindent History:

\\vspace{-0pt}
\\hspace{40pt}
\\begin{tabular}{ccc}
\\textit{6000} & /ana{\dh}an/&$\\rightarrow$ & \\textit{7500} & /na{\dh}an/&$\\rightarrow$ & \\textit{12002} & /na{\texttheta}an/& \\\\
\\end{tabular}

\\vspace{20pt}\\hline

\\end{nopagebreak}
\\filbreak



\\vspace{15pt}
\\begin{nopagebreak}
\\noindent{\\fontsize{20pt}{10pt}\\textbf{shotlleth} } \\textit{80 of spring/May 20} (N)\\\\
\\noindent {\\tovian \\fontsize{20pt}{10pt} \\textbf{shotlleth} }\\\\
\\noindent /{\textesh}{\textprimstress}o{\texttoptiebar{t\textbeltl}}le{\texttheta}/\\\\
\\noindent Syllables: ʃoƛl.eθ\\\\
\\noindent lit. flowing+hoe\\\\
\\noindent \\textit{Mistsickle}\\\\


\\noindent History:

\\vspace{-0pt}
\\hspace{40pt}
\\begin{tabular}{ccc}
\\textit{6000} & /{\textesh}o{\texttoptiebar{t\textbeltl}}mllre{\dh}n/&$\\rightarrow$ & \\textit{6500} & /{\textesh}o{\texttoptiebar{t\textbeltl}}mlre{\dh}n/&$\\rightarrow$ & \\textit{6503} & /{\textesh}o{\texttoptiebar{t\textbeltl}}mrle{\dh}n/&$\\rightarrow$ & \\textit{7600} & /{\textesh}o{\texttoptiebar{t\textbeltl}}rle{\dh}n/&$\\rightarrow$ & \\textit{12002} & /{\textesh}o{\texttoptiebar{t\textbeltl}}rle{\texttheta}n/&$\\rightarrow$ & \\textit{12505} & /{\textesh}o{\texttoptiebar{t\textbeltl}}le{\texttheta}n/&$\\rightarrow$ & \\textit{13001} & /{\textesh}o{\texttoptiebar{t\textbeltl}}le{\texttheta}/& \\\\
\\end{tabular}

\\vspace{20pt}\\hline

\\end{nopagebreak}
\\filbreak



\\vspace{15pt}
\\begin{nopagebreak}
\\noindent{\\fontsize{20pt}{10pt}\\textbf{thangethan} } \\textit{80 of summer/August 20} (N)\\\\
\\noindent {\\tovian \\fontsize{20pt}{10pt} \\textbf{thangethan} }\\\\
\\noindent /{\texttheta}a{\ng}{\textprimstress}e{\texttheta}an/\\\\
\\noindent Syllables: θaŋ.eθ.an\\\\
\\noindent lit. leaf+tool\\\\
\\noindent \\textit{Leafcut}\\\\


\\noindent History:

\\vspace{-0pt}
\\hspace{40pt}
\\begin{tabular}{ccc}
\\textit{6000} & /a{\dh}a{\ng}ge{\dh}an/&$\\rightarrow$ & \\textit{7500} & /{\dh}a{\ng}ge{\dh}an/&$\\rightarrow$ & \\textit{8750} & /{\dh}a{\ng}e{\dh}an/&$\\rightarrow$ & \\textit{12002} & /{\texttheta}a{\ng}e{\texttheta}an/& \\\\
\\end{tabular}

\\vspace{20pt}\\hline

\\end{nopagebreak}
\\filbreak



\\vspace{15pt}
\\begin{nopagebreak}
\\noindent{\\fontsize{20pt}{10pt}\\textbf{murathan} } \\textit{80 of winter/February 20} (N)\\\\
\\noindent {\\tovian \\fontsize{20pt}{10pt} \\textbf{murathan} }\\\\
\\noindent /mur{\textprimstress}a{\texttheta}an/\\\\
\\noindent Syllables: mur.aθ.an\\\\
\\noindent lit. darkness+tool\\\\
\\noindent \\textit{Nighbasket}\\\\


\\noindent History:

\\vspace{-0pt}
\\hspace{40pt}
\\begin{tabular}{ccc}
\\textit{6000} & /bral{\dh}an/&$\\rightarrow$ & \\textit{11000} & /mral{\dh}an/&$\\rightarrow$ & \\textit{12002} & /mral{\texttheta}an/&$\\rightarrow$ & \\textit{12502} & /mra{\texttheta}an/&$\\rightarrow$ & \\textit{15001} & /mura{\texttheta}an/& \\\\
\\end{tabular}

\\vspace{20pt}\\hline

\\end{nopagebreak}
\\filbreak



\\vspace{15pt}
\\begin{nopagebreak}
\\noindent{\\fontsize{20pt}{10pt}\\textbf{???} } \\textit{81 of autumn/November 21} (N)\\\\
\\noindent {\\tovian \\fontsize{20pt}{10pt} \\textbf{???} }\\\\
\\noindent /???/\\\\
\\noindent Syllables: ???\\\\
\\noindent lit. ???\\\\
\\noindent \\textit{Skoble}\\\\


\\noindent History:

\\vspace{-0pt}
\\hspace{40pt}
\\begin{tabular}{ccc}
\\textit{6000} & /???/& \\\\
\\end{tabular}

\\vspace{20pt}\\hline

\\end{nopagebreak}
\\filbreak



\\vspace{15pt}
\\begin{nopagebreak}
\\noindent{\\fontsize{20pt}{10pt}\\textbf{sethok} } \\textit{81 of spring/May 21} (N)\\\\
\\noindent {\\tovian \\fontsize{20pt}{10pt} \\textbf{sethok} }\\\\
\\noindent /s{\textprimstress}e{\texttheta}ok/\\\\
\\noindent Syllables: seθ.ok\\\\
\\noindent lit. water+down\\\\
\\noindent \\textit{Raincress}\\\\


\\noindent History:

\\vspace{-0pt}
\\hspace{40pt}
\\begin{tabular}{ccc}
\\textit{6000} & /zel{\texttheta}ok/&$\\rightarrow$ & \\textit{12000} & /sel{\texttheta}ok/&$\\rightarrow$ & \\textit{12502} & /se{\texttheta}ok/& \\\\
\\end{tabular}

\\vspace{20pt}\\hline

\\end{nopagebreak}
\\filbreak



\\vspace{15pt}
\\begin{nopagebreak}
\\noindent{\\fontsize{20pt}{10pt}\\textbf{lhenashotlodrem} } \\textit{81 of summer/August 21} (N)\\\\
\\noindent {\\tovian \\fontsize{20pt}{10pt} \\textbf{lhenashotlodrem} }\\\\
\\noindent /{\textbeltl}ena{\textesh}o{\texttoptiebar{t\textbeltl}}{\textprimstress}odrem/\\\\
\\noindent Syllables: ɬen.aʃ.oƛ.odr.em\\\\
\\noindent lit. dream+flowing+flower\\\\
\\noindent \\textit{Tideflower}\\\\


\\noindent History:

\\vspace{-0pt}
\\hspace{40pt}
\\begin{tabular}{ccc}
\\textit{6000} & /{\textbeltl}ena{\textesh}o{\texttoptiebar{t\textbeltl}}dremp/&$\\rightarrow$ & \\textit{8750} & /{\textbeltl}ena{\textesh}o{\texttoptiebar{t\textbeltl}}drem/&$\\rightarrow$ & \\textit{9300} & /{\textbeltl}ena{\textesh}o{\texttoptiebar{t\textbeltl}}odrem/& \\\\
\\end{tabular}

\\vspace{20pt}\\hline

\\end{nopagebreak}
\\filbreak



\\vspace{15pt}
\\begin{nopagebreak}
\\noindent{\\fontsize{20pt}{10pt}\\textbf{deremen} } \\textit{81 of winter/February 1} (N)\\\\
\\noindent {\\tovian \\fontsize{20pt}{10pt} \\textbf{deremen} }\\\\
\\noindent /der{\textprimstress}emen/\\\\
\\noindent Syllables: der.em.en\\\\
\\noindent lit. flower+stone\\\\
\\noindent \\textit{Dremgypsum}\\\\


\\noindent History:

\\vspace{-0pt}
\\hspace{40pt}
\\begin{tabular}{ccc}
\\textit{6000} & /drempmern/&$\\rightarrow$ & \\textit{7600} & /drempmen/&$\\rightarrow$ & \\textit{8500} & /dremmen/&$\\rightarrow$ & \\textit{12007} & /dremen/&$\\rightarrow$ & \\textit{15001} & /deremen/& \\\\
\\end{tabular}

\\vspace{20pt}\\hline

\\end{nopagebreak}
\\filbreak



\\vspace{15pt}
\\begin{nopagebreak}
\\noindent{\\fontsize{20pt}{10pt}\\textbf{???} } \\textit{82 of autumn/November 22} (N)\\\\
\\noindent {\\tovian \\fontsize{20pt}{10pt} \\textbf{???} }\\\\
\\noindent /???/\\\\
\\noindent Syllables: ???\\\\
\\noindent lit. ???\\\\
\\noindent \\textit{Muncal}\\\\


\\noindent History:

\\vspace{-0pt}
\\hspace{40pt}
\\begin{tabular}{ccc}
\\textit{6000} & /???/& \\\\
\\end{tabular}

\\vspace{20pt}\\hline

\\end{nopagebreak}
\\filbreak



\\vspace{15pt}
\\begin{nopagebreak}
\\noindent{\\fontsize{20pt}{10pt}\\textbf{shotlom} } \\textit{82 of spring/May 22} (N)\\\\
\\noindent {\\tovian \\fontsize{20pt}{10pt} \\textbf{shotlom} }\\\\
\\noindent /{\textesh}{\textprimstress}o{\texttoptiebar{t\textbeltl}}om/\\\\
\\noindent Syllables: ʃoƛ.om\\\\
\\noindent lit. flowing+bloom\\\\
\\noindent \\textit{Mistbloom}\\\\


\\noindent History:

\\vspace{-0pt}
\\hspace{40pt}
\\begin{tabular}{ccc}
\\textit{6000} & /{\textesh}o{\texttoptiebar{t\textbeltl}}drmpfede/&$\\rightarrow$ & \\textit{7600} & /{\textesh}o{\texttoptiebar{t\textbeltl}}dmpfede/&$\\rightarrow$ & \\textit{8500} & /{\textesh}o{\texttoptiebar{t\textbeltl}}dmfede/&$\\rightarrow$ & \\textit{9300} & /{\textesh}o{\texttoptiebar{t\textbeltl}}odmfede/&$\\rightarrow$ & \\textit{9500} & /{\textesh}o{\texttoptiebar{t\textbeltl}}odmfed/&$\\rightarrow$ & \\textit{12003} & /{\textesh}o{\texttoptiebar{t\textbeltl}}odmd/&$\\rightarrow$ & \\textit{13001} & /{\textesh}o{\texttoptiebar{t\textbeltl}}odm/&$\\rightarrow$ & \\textit{13004} & /{\textesh}o{\texttoptiebar{t\textbeltl}}om/& \\\\
\\end{tabular}

\\vspace{20pt}\\hline

\\end{nopagebreak}
\\filbreak



\\vspace{15pt}
\\begin{nopagebreak}
\\noindent{\\fontsize{20pt}{10pt}\\textbf{munorhanatesh} } \\textit{82 of summer/August 22} (N)\\\\
\\noindent {\\tovian \\fontsize{20pt}{10pt} \\textbf{munorhanatesh} }\\\\
\\noindent /munorħan{\textprimstress}a{\textsubbridge{t}}e{\textesh}/\\\\
\\noindent Syllables: mun.orħ.an.at.eʃ\\\\
\\noindent lit. dusk+sound\\\\
\\noindent \\textit{Duskbell}\\\\


\\noindent History:

\\vspace{-0pt}
\\hspace{40pt}
\\begin{tabular}{ccc}
\\textit{6000} & /pnorhana{\textsubbridge{t}}e{\textyogh}/&$\\rightarrow$ & \\textit{11000} & /mnorhana{\textsubbridge{t}}e{\textyogh}/&$\\rightarrow$ & \\textit{11996} & /munorhana{\textsubbridge{t}}e{\textyogh}/&$\\rightarrow$ & \\textit{12001} & /munorhana{\textsubbridge{t}}e{\textesh}/&$\\rightarrow$ & \\textit{14006} & /munorħana{\textsubbridge{t}}e{\textesh}/& \\\\
\\end{tabular}

\\vspace{20pt}\\hline

\\end{nopagebreak}
\\filbreak



\\vspace{15pt}
\\begin{nopagebreak}
\\noindent{\\fontsize{20pt}{10pt}\\textbf{seyalethanan} } \\textit{82 of winter/February 22} (N)\\\\
\\noindent {\\tovian \\fontsize{20pt}{10pt} \\textbf{seyalethanan} }\\\\
\\noindent /seyale{\texttheta}{\textprimstress}anan/\\\\
\\noindent Syllables: sej.al.eθ.an.an\\\\
\\noindent lit. food+star\\\\
\\noindent \\textit{Starssalt}\\\\


\\noindent History:

\\vspace{-0pt}
\\hspace{40pt}
\\begin{tabular}{ccc}
\\textit{6000} & /syale{\dh}annaana/&$\\rightarrow$ & \\textit{6500} & /syale{\dh}anaana/&$\\rightarrow$ & \\textit{9500} & /syale{\dh}anaan/&$\\rightarrow$ & \\textit{12002} & /syale{\texttheta}anaan/&$\\rightarrow$ & \\textit{12005} & /syale{\texttheta}anan/&$\\rightarrow$ & \\textit{15001} & /seyale{\texttheta}anan/& \\\\
\\end{tabular}

\\vspace{20pt}\\hline

\\end{nopagebreak}
\\filbreak



\\vspace{15pt}
\\begin{nopagebreak}
\\noindent{\\fontsize{20pt}{10pt}\\textbf{holamel} } \\textit{83 of autumn/November 23} (N)\\\\
\\noindent {\\tovian \\fontsize{20pt}{10pt} \\textbf{holamel} }\\\\
\\noindent /ħol{\textprimstress}amel/\\\\
\\noindent Syllables: ħol.am.el\\\\
\\noindent lit. sweet+life\\\\
\\noindent \\textit{Sweevin}\\\\


\\noindent History:

\\vspace{-0pt}
\\hspace{40pt}
\\begin{tabular}{ccc}
\\textit{6000} & /holamel/&$\\rightarrow$ & \\textit{14006} & /ħolamel/& \\\\
\\end{tabular}

\\vspace{20pt}\\hline

\\end{nopagebreak}
\\filbreak



\\vspace{15pt}
\\begin{nopagebreak}
\\noindent{\\fontsize{20pt}{10pt}\\textbf{muloanalor} } \\textit{83 of spring/May 23} (N)\\\\
\\noindent {\\tovian \\fontsize{20pt}{10pt} \\textbf{muloanalor} }\\\\
\\noindent /muloan{\textprimstress}alor/\\\\
\\noindent Syllables: mul.o.an.al.or\\\\
\\noindent lit. morning+plant\\\\
\\noindent \\textit{Dawngrass}\\\\


\\noindent History:

\\vspace{-0pt}
\\hspace{40pt}
\\begin{tabular}{ccc}
\\textit{6000} & /ploanamllore/&$\\rightarrow$ & \\textit{6500} & /ploanamlore/&$\\rightarrow$ & \\textit{7600} & /ploanalore/&$\\rightarrow$ & \\textit{9500} & /ploanalor/&$\\rightarrow$ & \\textit{11000} & /mloanalor/&$\\rightarrow$ & \\textit{15001} & /muloanalor/& \\\\
\\end{tabular}

\\vspace{20pt}\\hline

\\end{nopagebreak}
\\filbreak



\\vspace{15pt}
\\begin{nopagebreak}
\\noindent{\\fontsize{20pt}{10pt}\\textbf{yelan} } \\textit{83 of summer/August 23} (N)\\\\
\\noindent {\\tovian \\fontsize{20pt}{10pt} \\textbf{yelan} }\\\\
\\noindent /y{\textprimstress}elan/\\\\
\\noindent Syllables: jel.an\\\\
\\noindent lit. person+light\\\\
\\noindent \\textit{Selitha}\\\\


\\noindent History:

\\vspace{-0pt}
\\hspace{40pt}
\\begin{tabular}{ccc}
\\textit{6000} & /dryelaana/&$\\rightarrow$ & \\textit{7501} & /ryelaana/&$\\rightarrow$ & \\textit{7600} & /yelaana/&$\\rightarrow$ & \\textit{9500} & /yelaan/&$\\rightarrow$ & \\textit{12005} & /yelan/& \\\\
\\end{tabular}

\\vspace{20pt}\\hline

\\end{nopagebreak}
\\filbreak



\\vspace{15pt}
\\begin{nopagebreak}
\\noindent{\\fontsize{20pt}{10pt}\\textbf{aymen} } \\textit{83 of winter/February 23} (N)\\\\
\\noindent {\\tovian \\fontsize{20pt}{10pt} \\textbf{aymen} }\\\\
\\noindent /{\textprimstress}aymen/\\\\
\\noindent Syllables: ajm.en\\\\
\\noindent lit. moon+stone\\\\
\\noindent \\textit{Mooniron}\\\\


\\noindent History:

\\vspace{-0pt}
\\hspace{40pt}
\\begin{tabular}{ccc}
\\textit{6000} & /aliysmern/&$\\rightarrow$ & \\textit{6300} & /alysmern/&$\\rightarrow$ & \\textit{7600} & /aysmen/&$\\rightarrow$ & \\textit{13002} & /aymen/& \\\\
\\end{tabular}

\\vspace{20pt}\\hline

\\end{nopagebreak}
\\filbreak



\\vspace{15pt}
\\begin{nopagebreak}
\\noindent{\\fontsize{20pt}{10pt}\\textbf{rithdrem} } \\textit{84 of autumn/November 24} (N)\\\\
\\noindent {\\tovian \\fontsize{20pt}{10pt} \\textbf{rithdrem} }\\\\
\\noindent /r{\textprimstress}i{\texttheta}drem/\\\\
\\noindent Syllables: riθdr.em\\\\
\\noindent lit. air+flower\\\\
\\noindent \\textit{Frothfloer}\\\\


\\noindent History:

\\vspace{-0pt}
\\hspace{40pt}
\\begin{tabular}{ccc}
\\textit{6000} & /ri{\texttheta}dremp/&$\\rightarrow$ & \\textit{8750} & /ri{\texttheta}drem/& \\\\
\\end{tabular}

\\vspace{20pt}\\hline

\\end{nopagebreak}
\\filbreak



\\vspace{15pt}
\\begin{nopagebreak}
\\noindent{\\fontsize{20pt}{10pt}\\textbf{selumlor} } \\textit{84 of spring/May 24} (N)\\\\
\\noindent {\\tovian \\fontsize{20pt}{10pt} \\textbf{selumlor} }\\\\
\\noindent /sel{\textprimstress}umlor/\\\\
\\noindent Syllables: sel.uml.or\\\\
\\noindent lit. water+vine\\\\
\\noindent \\textit{Rainvin}\\\\


\\noindent History:

\\vspace{-0pt}
\\hspace{40pt}
\\begin{tabular}{ccc}
\\textit{6000} & /zelkluwmllore/&$\\rightarrow$ & \\textit{6500} & /zelkluwmlore/&$\\rightarrow$ & \\textit{7600} & /zelklumlore/&$\\rightarrow$ & \\textit{8750} & /zellumlore/&$\\rightarrow$ & \\textit{9500} & /zellumlor/&$\\rightarrow$ & \\textit{12000} & /sellumlor/&$\\rightarrow$ & \\textit{12007} & /selumlor/& \\\\
\\end{tabular}

\\vspace{20pt}\\hline

\\end{nopagebreak}
\\filbreak



\\vspace{15pt}
\\begin{nopagebreak}
\\noindent{\\fontsize{20pt}{10pt}\\textbf{lhenaglar} } \\textit{84 of summer/August 24} (N)\\\\
\\noindent {\\tovian \\fontsize{20pt}{10pt} \\textbf{lhenaglar} }\\\\
\\noindent /{\textbeltl}en{\textprimstress}aglar/\\\\
\\noindent Syllables: ɬen.agl.ar\\\\
\\noindent lit. dream+tree\\\\
\\noindent \\textit{Dremroot}\\\\


\\noindent History:

\\vspace{-0pt}
\\hspace{40pt}
\\begin{tabular}{ccc}
\\textit{6000} & /{\textbeltl}enaaglare/&$\\rightarrow$ & \\textit{9500} & /{\textbeltl}enaaglar/&$\\rightarrow$ & \\textit{12005} & /{\textbeltl}enaglar/& \\\\
\\end{tabular}

\\vspace{20pt}\\hline

\\end{nopagebreak}
\\filbreak



\\vspace{15pt}
\\begin{nopagebreak}
\\noindent{\\fontsize{20pt}{10pt}\\textbf{muloanamen} } \\textit{84 of winter/February 24} (N)\\\\
\\noindent {\\tovian \\fontsize{20pt}{10pt} \\textbf{muloanamen} }\\\\
\\noindent /muloan{\textprimstress}amen/\\\\
\\noindent Syllables: mul.o.an.am.en\\\\
\\noindent lit. morning+metal\\\\
\\noindent \\textit{Dawnbrass}\\\\


\\noindent History:

\\vspace{-0pt}
\\hspace{40pt}
\\begin{tabular}{ccc}
\\textit{6000} & /ploana{\dh}nhmern/&$\\rightarrow$ & \\textit{6400} & /ploana{\dh}hmern/&$\\rightarrow$ & \\textit{7600} & /ploana{\dh}hmen/&$\\rightarrow$ & \\textit{11000} & /mloana{\dh}hmen/&$\\rightarrow$ & \\textit{12002} & /mloana{\texttheta}hmen/&$\\rightarrow$ & \\textit{13000} & /mloana{\texttheta}men/&$\\rightarrow$ & \\textit{13002} & /mloanamen/&$\\rightarrow$ & \\textit{15001} & /muloanamen/& \\\\
\\end{tabular}

\\vspace{20pt}\\hline

\\end{nopagebreak}
\\filbreak



\\vspace{15pt}
\\begin{nopagebreak}
\\noindent{\\fontsize{20pt}{10pt}\\textbf{lhenarith} } \\textit{85 of autumn/November 25} (N)\\\\
\\noindent {\\tovian \\fontsize{20pt}{10pt} \\textbf{lhenarith} }\\\\
\\noindent /{\textbeltl}en{\textprimstress}ari{\texttheta}/\\\\
\\noindent Syllables: ɬen.ar.iθ\\\\
\\noindent lit. dream+air\\\\
\\noindent \\textit{Depswe}\\\\


\\noindent History:

\\vspace{-0pt}
\\hspace{40pt}
\\begin{tabular}{ccc}
\\textit{6000} & /{\textbeltl}enari{\texttheta}/& \\\\
\\end{tabular}

\\vspace{20pt}\\hline

\\end{nopagebreak}
\\filbreak



\\vspace{15pt}
\\begin{nopagebreak}
\\noindent{\\fontsize{20pt}{10pt}\\textbf{nanatlan} } \\textit{85 of spring/May 25} (N)\\\\
\\noindent {\\tovian \\fontsize{20pt}{10pt} \\textbf{nanatlan} }\\\\
\\noindent /nan{\textprimstress}a{\texttoptiebar{t\textbeltl}}an/\\\\
\\noindent Syllables: nan.aƛ.an\\\\
\\noindent lit. star+animal\\\\
\\noindent \\textit{Starhare}\\\\


\\noindent History:

\\vspace{-0pt}
\\hspace{40pt}
\\begin{tabular}{ccc}
\\textit{6000} & /nnaana{\textsubbridge{t}}lamn/&$\\rightarrow$ & \\textit{6401} & /nnaana{\texttoptiebar{t\textbeltl}}amn/&$\\rightarrow$ & \\textit{6500} & /naana{\texttoptiebar{t\textbeltl}}amn/&$\\rightarrow$ & \\textit{7600} & /naana{\texttoptiebar{t\textbeltl}}an/&$\\rightarrow$ & \\textit{12005} & /nana{\texttoptiebar{t\textbeltl}}an/& \\\\
\\end{tabular}

\\vspace{20pt}\\hline

\\end{nopagebreak}
\\filbreak



\\vspace{15pt}
\\begin{nopagebreak}
\\noindent{\\fontsize{20pt}{10pt}\\textbf{yiselom} } \\textit{85 of summer/August 25} (N)\\\\
\\noindent {\\tovian \\fontsize{20pt}{10pt} \\textbf{yiselom} }\\\\
\\noindent /yis{\textprimstress}elom/\\\\
\\noindent Syllables: jis.el.om\\\\
\\noindent lit. moon+memory\\\\
\\noindent \\textit{Tlunalorum}\\\\


\\noindent History:

\\vspace{-0pt}
\\hspace{40pt}
\\begin{tabular}{ccc}
\\textit{6000} & /aliyselorm/&$\\rightarrow$ & \\textit{6300} & /alyselorm/&$\\rightarrow$ & \\textit{7500} & /lyselorm/&$\\rightarrow$ & \\textit{7501} & /yselorm/&$\\rightarrow$ & \\textit{7600} & /yselom/&$\\rightarrow$ & \\textit{15001} & /yiselom/& \\\\
\\end{tabular}

\\vspace{20pt}\\hline

\\end{nopagebreak}
\\filbreak



\\vspace{15pt}
\\begin{nopagebreak}
\\noindent{\\fontsize{20pt}{10pt}\\textbf{nanatan} } \\textit{85 of winter/February 25} (N)\\\\
\\noindent {\\tovian \\fontsize{20pt}{10pt} \\textbf{nanatan} }\\\\
\\noindent /nan{\textprimstress}a{\textsubbridge{t}}an/\\\\
\\noindent Syllables: nan.at.an\\\\
\\noindent lit. star+face\\\\
\\noindent \\textit{Starwhisker}\\\\


\\noindent History:

\\vspace{-0pt}
\\hspace{40pt}
\\begin{tabular}{ccc}
\\textit{6000} & /nnaana{\textsubbridge{t}}arn/&$\\rightarrow$ & \\textit{6500} & /naana{\textsubbridge{t}}arn/&$\\rightarrow$ & \\textit{7600} & /naana{\textsubbridge{t}}an/&$\\rightarrow$ & \\textit{12005} & /nana{\textsubbridge{t}}an/& \\\\
\\end{tabular}

\\vspace{20pt}\\hline

\\end{nopagebreak}
\\filbreak



\\vspace{15pt}
\\begin{nopagebreak}
\\noindent{\\fontsize{20pt}{10pt}\\textbf{nunasil} } \\textit{86 of autumn/November 26} (N)\\\\
\\noindent {\\tovian \\fontsize{20pt}{10pt} \\textbf{nunasil} }\\\\
\\noindent /nun{\textprimstress}asil/\\\\
\\noindent Syllables: nun.as.il\\\\
\\noindent lit. small+existence\\\\
\\noindent \\textit{Niwhis}\\\\


\\noindent History:

\\vspace{-0pt}
\\hspace{40pt}
\\begin{tabular}{ccc}
\\textit{6000} & /ninasil/&$\\rightarrow$ & \\textit{9200} & /n{\textschwa}nasil/&$\\rightarrow$ & \\textit{11990} & /nnasil/&$\\rightarrow$ & \\textit{11996} & /nunasil/& \\\\
\\end{tabular}

\\vspace{20pt}\\hline

\\end{nopagebreak}
\\filbreak



\\vspace{15pt}
\\begin{nopagebreak}
\\noindent{\\fontsize{20pt}{10pt}\\textbf{dolor} } \\textit{86 of spring/May 26} (N)\\\\
\\noindent {\\tovian \\fontsize{20pt}{10pt} \\textbf{dolor} }\\\\
\\noindent /d{\textprimstress}olor/\\\\
\\noindent Syllables: dol.or\\\\
\\noindent lit. death+blue\\\\
\\noindent \\textit{Twiliwoad}\\\\


\\noindent History:

\\vspace{-0pt}
\\hspace{40pt}
\\begin{tabular}{ccc}
\\textit{6000} & /doldora/&$\\rightarrow$ & \\textit{8750} & /dolora/&$\\rightarrow$ & \\textit{9500} & /dolor/& \\\\
\\end{tabular}

\\vspace{20pt}\\hline

\\end{nopagebreak}
\\filbreak



\\vspace{15pt}
\\begin{nopagebreak}
\\noindent{\\fontsize{20pt}{10pt}\\textbf{keshaglar} } \\textit{86 of summer/August 26} (N)\\\\
\\noindent {\\tovian \\fontsize{20pt}{10pt} \\textbf{keshaglar} }\\\\
\\noindent /ke{\textesh}{\textprimstress}aglar/\\\\
\\noindent Syllables: keʃ.agl.ar\\\\
\\noindent lit. fire+tree\\\\
\\noindent \\textit{Charwood}\\\\


\\noindent History:

\\vspace{-0pt}
\\hspace{40pt}
\\begin{tabular}{ccc}
\\textit{6000} & /ke{\textesh}aglare/&$\\rightarrow$ & \\textit{9500} & /ke{\textesh}aglar/& \\\\
\\end{tabular}

\\vspace{20pt}\\hline

\\end{nopagebreak}
\\filbreak



\\vspace{15pt}
\\begin{nopagebreak}
\\noindent{\\fontsize{20pt}{10pt}\\textbf{lanamen} } \\textit{86 of winter/February 26} (N)\\\\
\\noindent {\\tovian \\fontsize{20pt}{10pt} \\textbf{lanamen} }\\\\
\\noindent /lan{\textprimstress}amen/\\\\
\\noindent Syllables: lan.am.en\\\\
\\noindent lit. twilight+metal\\\\
\\noindent \\textit{Twilighttin}\\\\


\\noindent History:

\\vspace{-0pt}
\\hspace{40pt}
\\begin{tabular}{ccc}
\\textit{6000} & /brlana{\dh}nhmern/&$\\rightarrow$ & \\textit{6400} & /brlana{\dh}hmern/&$\\rightarrow$ & \\textit{7501} & /rlana{\dh}hmern/&$\\rightarrow$ & \\textit{7600} & /lana{\dh}hmen/&$\\rightarrow$ & \\textit{12002} & /lana{\texttheta}hmen/&$\\rightarrow$ & \\textit{13000} & /lana{\texttheta}men/&$\\rightarrow$ & \\textit{13002} & /lanamen/& \\\\
\\end{tabular}

\\vspace{20pt}\\hline

\\end{nopagebreak}
\\filbreak



\\vspace{15pt}
\\begin{nopagebreak}
\\noindent{\\fontsize{20pt}{10pt}\\textbf{yisen} } \\textit{87 of autumn/November 27} (N)\\\\
\\noindent {\\tovian \\fontsize{20pt}{10pt} \\textbf{yisen} }\\\\
\\noindent /y{\textprimstress}isen/\\\\
\\noindent Syllables: jis.en\\\\
\\noindent lit. moon+dream\\\\
\\noindent \\textit{Munverbena}\\\\


\\noindent History:

\\vspace{-0pt}
\\hspace{40pt}
\\begin{tabular}{ccc}
\\textit{6000} & /aliys{\textbeltl}ena/&$\\rightarrow$ & \\textit{6300} & /alys{\textbeltl}ena/&$\\rightarrow$ & \\textit{7500} & /lys{\textbeltl}ena/&$\\rightarrow$ & \\textit{7501} & /ys{\textbeltl}ena/&$\\rightarrow$ & \\textit{9500} & /ys{\textbeltl}en/&$\\rightarrow$ & \\textit{13000} & /ysen/&$\\rightarrow$ & \\textit{15001} & /yisen/& \\\\
\\end{tabular}

\\vspace{20pt}\\hline

\\end{nopagebreak}
\\filbreak



\\vspace{15pt}
\\begin{nopagebreak}
\\noindent{\\fontsize{20pt}{10pt}\\textbf{yiskel} } \\textit{87 of spring/May 27} (N)\\\\
\\noindent {\\tovian \\fontsize{20pt}{10pt} \\textbf{yiskel} }\\\\
\\noindent /y{\textprimstress}iskel/\\\\
\\noindent Syllables: jisk.el\\\\
\\noindent lit. moon+branch\\\\
\\noindent \\textit{Moonhazel}\\\\


\\noindent History:

\\vspace{-0pt}
\\hspace{40pt}
\\begin{tabular}{ccc}
\\textit{6000} & /aliyskelu/&$\\rightarrow$ & \\textit{6300} & /alyskelu/&$\\rightarrow$ & \\textit{7500} & /lyskelu/&$\\rightarrow$ & \\textit{7501} & /yskelu/&$\\rightarrow$ & \\textit{9500} & /yskel/&$\\rightarrow$ & \\textit{15001} & /yiskel/& \\\\
\\end{tabular}

\\vspace{20pt}\\hline

\\end{nopagebreak}
\\filbreak



\\vspace{15pt}
\\begin{nopagebreak}
\\noindent{\\fontsize{20pt}{10pt}\\textbf{rithenalor} } \\textit{87 of summer/August 27} (N)\\\\
\\noindent {\\tovian \\fontsize{20pt}{10pt} \\textbf{rithenalor} }\\\\
\\noindent /ri{\texttheta}en{\textprimstress}alor/\\\\
\\noindent Syllables: riθ.en.al.or\\\\
\\noindent lit. air+dream+plant\\\\
\\noindent \\textit{Zephyrherb}\\\\


\\noindent History:

\\vspace{-0pt}
\\hspace{40pt}
\\begin{tabular}{ccc}
\\textit{6000} & /ri{\texttheta}{\textbeltl}enamllore/&$\\rightarrow$ & \\textit{6500} & /ri{\texttheta}{\textbeltl}enamlore/&$\\rightarrow$ & \\textit{7600} & /ri{\texttheta}{\textbeltl}enalore/&$\\rightarrow$ & \\textit{9500} & /ri{\texttheta}{\textbeltl}enalor/&$\\rightarrow$ & \\textit{13000} & /ri{\texttheta}enalor/& \\\\
\\end{tabular}

\\vspace{20pt}\\hline

\\end{nopagebreak}
\\filbreak



\\vspace{15pt}
\\begin{nopagebreak}
\\noindent{\\fontsize{20pt}{10pt}\\textbf{muramen} } \\textit{87 of winter/February 27} (N)\\\\
\\noindent {\\tovian \\fontsize{20pt}{10pt} \\textbf{muramen} }\\\\
\\noindent /mur{\textprimstress}amen/\\\\
\\noindent Syllables: mur.am.en\\\\
\\noindent lit. darkness+metal\\\\
\\noindent \\textit{Shadowlead}\\\\


\\noindent History:

\\vspace{-0pt}
\\hspace{40pt}
\\begin{tabular}{ccc}
\\textit{6000} & /bral{\dh}nhmern/&$\\rightarrow$ & \\textit{6400} & /bral{\dh}hmern/&$\\rightarrow$ & \\textit{7600} & /bral{\dh}hmen/&$\\rightarrow$ & \\textit{11000} & /mral{\dh}hmen/&$\\rightarrow$ & \\textit{12002} & /mral{\texttheta}hmen/&$\\rightarrow$ & \\textit{12502} & /mra{\texttheta}hmen/&$\\rightarrow$ & \\textit{13000} & /mra{\texttheta}men/&$\\rightarrow$ & \\textit{13002} & /mramen/&$\\rightarrow$ & \\textit{15001} & /muramen/& \\\\
\\end{tabular}

\\vspace{20pt}\\hline

\\end{nopagebreak}
\\filbreak



\\vspace{15pt}
\\begin{nopagebreak}
\\noindent{\\fontsize{20pt}{10pt}\\textbf{engothr} } \\textit{88 of autumn/November 28} (N)\\\\
\\noindent {\\tovian \\fontsize{20pt}{10pt} \\textbf{engothr} }\\\\
\\noindent /{\textprimstress}e{\ng}o{\texttheta}r/\\\\
\\noindent Syllables: eŋ.oθr\\\\
\\noindent lit. sun+time\\\\
\\noindent \\textit{Sunethime}\\\\


\\noindent History:

\\vspace{-0pt}
\\hspace{40pt}
\\begin{tabular}{ccc}
\\textit{6000} & /efl{\ng}go{\texttheta}r/&$\\rightarrow$ & \\textit{7600} & /ef{\ng}go{\texttheta}r/&$\\rightarrow$ & \\textit{8750} & /ef{\ng}o{\texttheta}r/&$\\rightarrow$ & \\textit{13002} & /e{\ng}o{\texttheta}r/& \\\\
\\end{tabular}

\\vspace{20pt}\\hline

\\end{nopagebreak}
\\filbreak



\\vspace{15pt}
\\begin{nopagebreak}
\\noindent{\\fontsize{20pt}{10pt}\\textbf{muloanakel} } \\textit{88 of spring/May 28} (N)\\\\
\\noindent {\\tovian \\fontsize{20pt}{10pt} \\textbf{muloanakel} }\\\\
\\noindent /muloan{\textprimstress}akel/\\\\
\\noindent Syllables: mul.o.an.ak.el\\\\
\\noindent lit. dawn+branch\\\\
\\noindent \\textit{Duslower}\\\\


\\noindent History:

\\vspace{-0pt}
\\hspace{40pt}
\\begin{tabular}{ccc}
\\textit{6000} & /ploanakelu/&$\\rightarrow$ & \\textit{9500} & /ploanakel/&$\\rightarrow$ & \\textit{11000} & /mloanakel/&$\\rightarrow$ & \\textit{15001} & /muloanakel/& \\\\
\\end{tabular}

\\vspace{20pt}\\hline

\\end{nopagebreak}
\\filbreak



\\vspace{15pt}
\\begin{nopagebreak}
\\noindent{\\fontsize{20pt}{10pt}\\textbf{muralor} } \\textit{88 of summer/August 28} (N)\\\\
\\noindent {\\tovian \\fontsize{20pt}{10pt} \\textbf{muralor} }\\\\
\\noindent /mur{\textprimstress}alor/\\\\
\\noindent Syllables: mur.al.or\\\\
\\noindent lit. darkness+blue\\\\
\\noindent \\textit{Nightblue}\\\\


\\noindent History:

\\vspace{-0pt}
\\hspace{40pt}
\\begin{tabular}{ccc}
\\textit{6000} & /braldora/&$\\rightarrow$ & \\textit{8750} & /bralora/&$\\rightarrow$ & \\textit{9500} & /bralor/&$\\rightarrow$ & \\textit{11000} & /mralor/&$\\rightarrow$ & \\textit{15001} & /muralor/& \\\\
\\end{tabular}

\\vspace{20pt}\\hline

\\end{nopagebreak}
\\filbreak



\\vspace{15pt}
\\begin{nopagebreak}
\\noindent{\\fontsize{20pt}{10pt}\\textbf{lan} } \\textit{88 of winter/February 28} (N)\\\\
\\noindent {\\tovian \\fontsize{20pt}{10pt} \\textbf{lan} }\\\\
\\noindent /l{\textprimstress}an/\\\\
\\noindent Syllables: lan\\\\
\\noindent lit. twilight+???\\\\
\\noindent \\textit{Duskin}\\\\


\\noindent History:

\\vspace{-0pt}
\\hspace{40pt}
\\begin{tabular}{ccc}
\\textit{6000} & /brlana???/&$\\rightarrow$ & \\textit{7501} & /rlana???/&$\\rightarrow$ & \\textit{7600} & /lana???/&$\\rightarrow$ & \\textit{8500} & /lana/&$\\rightarrow$ & \\textit{9500} & /lan/& \\\\
\\end{tabular}

\\vspace{20pt}\\hline

\\end{nopagebreak}
\\filbreak



\\vspace{15pt}
\\begin{nopagebreak}
\\noindent{\\fontsize{20pt}{10pt}\\textbf{haradrem} } \\textit{89 of autumn/November 29} (N)\\\\
\\noindent {\\tovian \\fontsize{20pt}{10pt} \\textbf{haradrem} }\\\\
\\noindent /ħar{\textprimstress}adrem/\\\\
\\noindent Syllables: ħar.adr.em\\\\
\\noindent lit. pain+flower\\\\
\\noindent \\textit{Harpetal}\\\\


\\noindent History:

\\vspace{-0pt}
\\hspace{40pt}
\\begin{tabular}{ccc}
\\textit{6000} & /haradremp/&$\\rightarrow$ & \\textit{8750} & /haradrem/&$\\rightarrow$ & \\textit{14006} & /ħaradrem/& \\\\
\\end{tabular}

\\vspace{20pt}\\hline

\\end{nopagebreak}
\\filbreak



\\vspace{15pt}
\\begin{nopagebreak}
\\noindent{\\fontsize{20pt}{10pt}\\textbf{shotleon} } \\textit{89 of spring/May 29} (N)\\\\
\\noindent {\\tovian \\fontsize{20pt}{10pt} \\textbf{shotleon} }\\\\
\\noindent /{\textesh}o{\texttoptiebar{t\textbeltl}}{\textprimstress}eon/\\\\
\\noindent Syllables: ʃoƛ.e.on\\\\
\\noindent lit. flowing+mouth\\\\
\\noindent \\textit{Starwort}\\\\


\\noindent History:

\\vspace{-0pt}
\\hspace{40pt}
\\begin{tabular}{ccc}
\\textit{6000} & /{\textesh}o{\texttoptiebar{t\textbeltl}}eoln/&$\\rightarrow$ & \\textit{7600} & /{\textesh}o{\texttoptiebar{t\textbeltl}}eon/& \\\\
\\end{tabular}

\\vspace{20pt}\\hline

\\end{nopagebreak}
\\filbreak



\\vspace{15pt}
\\begin{nopagebreak}
\\noindent{\\fontsize{20pt}{10pt}\\textbf{fidol} } \\textit{89 of summer/August 29} (N)\\\\
\\noindent {\\tovian \\fontsize{20pt}{10pt} \\textbf{fidol} }\\\\
\\noindent /f{\textprimstress}idol/\\\\
\\noindent Syllables: fid.ol\\\\
\\noindent lit. sun+GEN+death\\\\
\\noindent \\textit{Sunseed}\\\\


\\noindent History:

\\vspace{-0pt}
\\hspace{40pt}
\\begin{tabular}{ccc}
\\textit{6000} & /eflhi-dol/&$\\rightarrow$ & \\textit{6402} & /ef{\textbeltl}i-dol/&$\\rightarrow$ & \\textit{7500} & /f{\textbeltl}i-dol/&$\\rightarrow$ & \\textit{8500} & /f{\textbeltl}idol/&$\\rightarrow$ & \\textit{13000} & /fidol/& \\\\
\\end{tabular}

\\vspace{20pt}\\hline

\\end{nopagebreak}
\\filbreak



\\vspace{15pt}
\\begin{nopagebreak}
\\noindent{\\fontsize{20pt}{10pt}\\textbf{tleshmen} } \\textit{89 of winter/February 29} (N)\\\\
\\noindent {\\tovian \\fontsize{20pt}{10pt} \\textbf{tleshmen} }\\\\
\\noindent /{\texttoptiebar{t\textbeltl}}{\textprimstress}e{\textesh}men/\\\\
\\noindent Syllables: ƛeʃm.en\\\\
\\noindent lit. speed+metal\\\\
\\noindent \\textit{Quicksilver}\\\\


\\noindent History:

\\vspace{-0pt}
\\hspace{40pt}
\\begin{tabular}{ccc}
\\textit{6000} & /{\texttoptiebar{t\textbeltl}}e{\textyogh}{\dh}nhmern/&$\\rightarrow$ & \\textit{6400} & /{\texttoptiebar{t\textbeltl}}e{\textyogh}{\dh}hmern/&$\\rightarrow$ & \\textit{7600} & /{\texttoptiebar{t\textbeltl}}e{\textyogh}{\dh}hmen/&$\\rightarrow$ & \\textit{12001} & /{\texttoptiebar{t\textbeltl}}e{\textesh}{\dh}hmen/&$\\rightarrow$ & \\textit{12002} & /{\texttoptiebar{t\textbeltl}}e{\textesh}{\texttheta}hmen/&$\\rightarrow$ & \\textit{13000} & /{\texttoptiebar{t\textbeltl}}e{\textesh}hmen/&$\\rightarrow$ & \\textit{13002} & /{\texttoptiebar{t\textbeltl}}e{\textesh}men/& \\\\
\\end{tabular}

\\vspace{20pt}\\hline

\\end{nopagebreak}
\\filbreak



\\vspace{15pt}
\\begin{nopagebreak}
\\noindent{\\fontsize{20pt}{10pt}\\textbf{lhenashith} } \\textit{9 of autumn/September 9} (N)\\\\
\\noindent {\\tovian \\fontsize{20pt}{10pt} \\textbf{lhenashith} }\\\\
\\noindent /{\textbeltl}en{\textprimstress}a{\textesh}i{\texttheta}/\\\\
\\noindent Syllables: ɬen.aʃ.iθ\\\\
\\noindent lit. dream+sharpness\\\\
\\noindent \\textit{Dremthorn}\\\\


\\noindent History:

\\vspace{-0pt}
\\hspace{40pt}
\\begin{tabular}{ccc}
\\textit{6000} & /{\textbeltl}ena{\textyogh}i{\texttheta}/&$\\rightarrow$ & \\textit{12001} & /{\textbeltl}ena{\textesh}i{\texttheta}/& \\\\
\\end{tabular}

\\vspace{20pt}\\hline

\\end{nopagebreak}
\\filbreak



\\vspace{15pt}
\\begin{nopagebreak}
\\noindent{\\fontsize{20pt}{10pt}\\textbf{todawen} } \\textit{9 of spring/March 9} (N)\\\\
\\noindent {\\tovian \\fontsize{20pt}{10pt} \\textbf{todawen} }\\\\
\\noindent /{\textsubbridge{t}}od{\textprimstress}awen/\\\\
\\noindent Syllables: tod.aw.en\\\\
\\noindent lit. eternal+beauty\\\\
\\noindent \\textit{Rashelwen}\\\\


\\noindent History:

\\vspace{-0pt}
\\hspace{40pt}
\\begin{tabular}{ccc}
\\textit{6000} & /{\textsubbridge{t}}odawen/& \\\\
\\end{tabular}

\\vspace{20pt}\\hline

\\end{nopagebreak}
\\filbreak



\\vspace{15pt}
\\begin{nopagebreak}
\\noindent{\\fontsize{20pt}{10pt}\\textbf{temen} } \\textit{9 of summer/June 9} (N)\\\\
\\noindent {\\tovian \\fontsize{20pt}{10pt} \\textbf{temen} }\\\\
\\noindent /{\textsubbridge{t}}{\textprimstress}emen/\\\\
\\noindent Syllables: tem.en\\\\
\\noindent lit. sound+stone\\\\
\\noindent \\textit{Whisperald}\\\\


\\noindent History:

\\vspace{-0pt}
\\hspace{40pt}
\\begin{tabular}{ccc}
\\textit{6000} & /{\textsubbridge{t}}e{\textyogh}mern/&$\\rightarrow$ & \\textit{7600} & /{\textsubbridge{t}}e{\textyogh}men/&$\\rightarrow$ & \\textit{12001} & /{\textsubbridge{t}}e{\textesh}men/&$\\rightarrow$ & \\textit{13002} & /{\textsubbridge{t}}emen/& \\\\
\\end{tabular}

\\vspace{20pt}\\hline

\\end{nopagebreak}
\\filbreak



\\vspace{15pt}
\\begin{nopagebreak}
\\noindent{\\fontsize{20pt}{10pt}\\textbf{lerithaglar} } \\textit{9 of winter/december 9} (N)\\\\
\\noindent {\\tovian \\fontsize{20pt}{10pt} \\textbf{lerithaglar} }\\\\
\\noindent /leri{\texttheta}{\textprimstress}aglar/\\\\
\\noindent Syllables: ler.iθ.agl.ar\\\\
\\noindent lit. breath+tree\\\\
\\noindent \\textit{Windtwig}\\\\


\\noindent History:

\\vspace{-0pt}
\\hspace{40pt}
\\begin{tabular}{ccc}
\\textit{6000} & /kwlyri{\texttheta}aglare/&$\\rightarrow$ & \\textit{7501} & /wlyri{\texttheta}aglare/&$\\rightarrow$ & \\textit{7600} & /lri{\texttheta}aglare/&$\\rightarrow$ & \\textit{9500} & /lri{\texttheta}aglar/&$\\rightarrow$ & \\textit{15001} & /leri{\texttheta}aglar/& \\\\
\\end{tabular}

\\vspace{20pt}\\hline

\\end{nopagebreak}
\\filbreak



\\vspace{15pt}
\\begin{nopagebreak}
\\noindent{\\fontsize{20pt}{10pt}\\textbf{meloresyalethalor} } \\textit{90 of autumn/November 30} (N)\\\\
\\noindent {\\tovian \\fontsize{20pt}{10pt} \\textbf{meloresyalethalor} }\\\\
\\noindent /meloresyale{\texttheta}{\textprimstress}alor/\\\\
\\noindent Syllables: mel.or.esj.al.eθ.al.or\\\\
\\noindent lit. grain+plant\\\\
\\noindent \\textit{Sorgho}\\\\


\\noindent History:

\\vspace{-0pt}
\\hspace{40pt}
\\begin{tabular}{ccc}
\\textit{6000} & /mlloresyale{\dh}amllore/&$\\rightarrow$ & \\textit{6500} & /mloresyale{\dh}amlore/&$\\rightarrow$ & \\textit{7502} & /meloresyale{\dh}amlore/&$\\rightarrow$ & \\textit{7600} & /meloresyale{\dh}alore/&$\\rightarrow$ & \\textit{9500} & /meloresyale{\dh}alor/&$\\rightarrow$ & \\textit{12002} & /meloresyale{\texttheta}alor/& \\\\
\\end{tabular}

\\vspace{20pt}\\hline

\\end{nopagebreak}
\\filbreak



\\vspace{15pt}
\\begin{nopagebreak}
\\noindent{\\fontsize{20pt}{10pt}\\textbf{shakwil} } \\textit{90 of spring/May 30} (N)\\\\
\\noindent {\\tovian \\fontsize{20pt}{10pt} \\textbf{shakwil} }\\\\
\\noindent /{\textesh}{\textprimstress}akwil/\\\\
\\noindent Syllables: ʃakw.il\\\\
\\noindent lit. ice+movement\\\\
\\noindent \\textit{Frostslide}\\\\


\\noindent History:

\\vspace{-0pt}
\\hspace{40pt}
\\begin{tabular}{ccc}
\\textit{6000} & /{\textesh}u{\textesh}akwil/&$\\rightarrow$ & \\textit{9200} & /{\textesh}{\textschwa}{\textesh}akwil/&$\\rightarrow$ & \\textit{11990} & /{\textesh}{\textesh}akwil/&$\\rightarrow$ & \\textit{12007} & /{\textesh}akwil/& \\\\
\\end{tabular}

\\vspace{20pt}\\hline

\\end{nopagebreak}
\\filbreak



\\vspace{15pt}
\\begin{nopagebreak}
\\noindent{\\fontsize{20pt}{10pt}\\textbf{falnan} } \\textit{90 of summer/August 30} (N)\\\\
\\noindent {\\tovian \\fontsize{20pt}{10pt} \\textbf{falnan} }\\\\
\\noindent /f{\textprimstress}alnan/\\\\
\\noindent Syllables: faln.an\\\\
\\noindent lit. leader+star\\\\
\\noindent \\textit{Shepherd's Star}\\\\


\\noindent History:

\\vspace{-0pt}
\\hspace{40pt}
\\begin{tabular}{ccc}
\\textit{6000} & /falrnnaana/&$\\rightarrow$ & \\textit{6500} & /falrnaana/&$\\rightarrow$ & \\textit{6503} & /farlnaana/&$\\rightarrow$ & \\textit{7600} & /falnaana/&$\\rightarrow$ & \\textit{9500} & /falnaan/&$\\rightarrow$ & \\textit{12005} & /falnan/& \\\\
\\end{tabular}

\\vspace{20pt}\\hline

\\end{nopagebreak}
\\filbreak



\\vspace{15pt}
\\begin{nopagebreak}
\\noindent{\\fontsize{20pt}{10pt}\\textbf{shathan} } \\textit{90 of winter/February 30} (N)\\\\
\\noindent {\\tovian \\fontsize{20pt}{10pt} \\textbf{shathan} }\\\\
\\noindent /{\textesh}{\textprimstress}a{\texttheta}an/\\\\
\\noindent Syllables: ʃaθ.an\\\\
\\noindent lit. ice+tool\\\\
\\noindent \\textit{Frostscreen}\\\\


\\noindent History:

\\vspace{-0pt}
\\hspace{40pt}
\\begin{tabular}{ccc}
\\textit{6000} & /{\textesh}u{\textesh}a{\dh}an/&$\\rightarrow$ & \\textit{9200} & /{\textesh}{\textschwa}{\textesh}a{\dh}an/&$\\rightarrow$ & \\textit{11990} & /{\textesh}{\textesh}a{\dh}an/&$\\rightarrow$ & \\textit{12002} & /{\textesh}{\textesh}a{\texttheta}an/&$\\rightarrow$ & \\textit{12007} & /{\textesh}a{\texttheta}an/& \\\\
\\end{tabular}

\\vspace{20pt}\\hline

\\end{nopagebreak}
\\filbreak



\\vspace{15pt}
\\begin{nopagebreak}
\\noindent{\\fontsize{20pt}{10pt}\\textbf{???} } \\textit{???} (PLACEHOLDER)\\\\
\\noindent {\\tovian \\fontsize{20pt}{10pt} \\textbf{???} }\\\\
\\noindent /???/\\\\
\\noindent Syllables: ???\\\\


\\noindent History:

\\vspace{-0pt}
\\hspace{40pt}
\\begin{tabular}{ccc}
\\textit{0} & /???/& \\\\
\\end{tabular}

\\vspace{20pt}\\hline

\\end{nopagebreak}
\\filbreak



\\vspace{15pt}
\\begin{nopagebreak}
\\noindent{\\fontsize{20pt}{10pt}\\textbf{remhil} } \\textit{abandoned} (N)\\\\
\\noindent {\\tovian \\fontsize{20pt}{10pt} \\textbf{remhil} }\\\\
\\noindent /r{\textprimstress}emħil/\\\\
\\noindent Syllables: remħ.il\\\\
\\noindent lit. empty+place\\\\


\\noindent History:

\\vspace{-0pt}
\\hspace{40pt}
\\begin{tabular}{ccc}
\\textit{0} & /oneremiha{\texttheta}ile/&$\\rightarrow$ & \\textit{1000} & /oneremih{\texttheta}ile/&$\\rightarrow$ & \\textit{3200} & /onremh{\texttheta}ile/&$\\rightarrow$ & \\textit{4502} & /onremhile/&$\\rightarrow$ & \\textit{7500} & /nremhile/&$\\rightarrow$ & \\textit{7600} & /remhile/&$\\rightarrow$ & \\textit{9500} & /remhil/&$\\rightarrow$ & \\textit{14006} & /remħil/& \\\\
\\end{tabular}

\\vspace{20pt}\\hline

\\end{nopagebreak}
\\filbreak



\\vspace{15pt}
\\begin{nopagebreak}
\\noindent{\\fontsize{20pt}{10pt}\\textbf{di-} } \\textit{ABL/ablative case marker/away from} (CASE)\\\\
\\noindent {\\tovian \\fontsize{20pt}{10pt} \\textbf{di-} }\\\\
\\noindent /d{\textprimstress}i/\\\\
\\noindent Syllables: di\\\\


\\noindent History:

\\vspace{-0pt}
\\hspace{40pt}
\\begin{tabular}{ccc}
\\textit{0} & /di-/& \\\\
\\end{tabular}

\\vspace{20pt}\\hline

\\end{nopagebreak}
\\filbreak



\\vspace{15pt}
\\begin{nopagebreak}
\\noindent{\\fontsize{20pt}{10pt}\\textbf{le-} } \\textit{ABS/abstract noun class} (CLASS)\\\\
\\noindent {\\tovian \\fontsize{20pt}{10pt} \\textbf{le-} }\\\\
\\noindent /l{\textprimstress}e/\\\\
\\noindent Syllables: le\\\\


\\noindent History:

\\vspace{-0pt}
\\hspace{40pt}
\\begin{tabular}{ccc}
\\textit{0} & /le-/& \\\\
\\end{tabular}

\\vspace{20pt}\\hline

\\end{nopagebreak}
\\filbreak



\\vspace{15pt}
\\begin{nopagebreak}
\\noindent{\\fontsize{20pt}{10pt}\\textbf{yu-} } \\textit{ACC/accusative case marker} (CASE)\\\\
\\noindent {\\tovian \\fontsize{20pt}{10pt} \\textbf{yu-} }\\\\
\\noindent /y{\textprimstress}u/\\\\
\\noindent Syllables: ju\\\\


\\noindent History:

\\vspace{-0pt}
\\hspace{40pt}
\\begin{tabular}{ccc}
\\textit{0} & /yu-/& \\\\
\\end{tabular}

\\vspace{20pt}\\hline

\\end{nopagebreak}
\\filbreak



\\vspace{15pt}
\\begin{nopagebreak}
\\noindent{\\fontsize{20pt}{10pt}\\textbf{lor} } \\textit{action/doing} (N)\\\\
\\noindent {\\tovian \\fontsize{20pt}{10pt} \\textbf{lor} }\\\\
\\noindent /l{\textprimstress}or/\\\\
\\noindent Syllables: lor\\\\


\\noindent History:

\\vspace{-0pt}
\\hspace{40pt}
\\begin{tabular}{ccc}
\\textit{0} & /lore/&$\\rightarrow$ & \\textit{9500} & /lor/& \\\\
\\end{tabular}

\\vspace{20pt}\\hline

\\end{nopagebreak}
\\filbreak



\\vspace{15pt}
\\begin{nopagebreak}
\\noindent{\\fontsize{20pt}{10pt}\\textbf{thim} } \\textit{adventure} (N)\\\\
\\noindent {\\tovian \\fontsize{20pt}{10pt} \\textbf{thim} }\\\\
\\noindent /{\texttheta}{\textprimstress}im/\\\\
\\noindent Syllables: θim\\\\
\\noindent lit. going+unknown\\\\


\\noindent History:

\\vspace{-0pt}
\\hspace{40pt}
\\begin{tabular}{ccc}
\\textit{1000} & /ki{\texttheta}sa{\textesh}imi/&$\\rightarrow$ & \\textit{1000} & /k{\texttheta}s{\textesh}imi/&$\\rightarrow$ & \\textit{4502} & /k{\texttheta}{\textesh}imi/&$\\rightarrow$ & \\textit{7501} & /{\texttheta}{\textesh}imi/&$\\rightarrow$ & \\textit{9500} & /{\texttheta}{\textesh}im/&$\\rightarrow$ & \\textit{13000} & /{\texttheta}im/& \\\\
\\end{tabular}

\\vspace{20pt}\\hline

\\end{nopagebreak}
\\filbreak



\\vspace{15pt}
\\begin{nopagebreak}
\\noindent{\\fontsize{20pt}{10pt}\\textbf{rith} } \\textit{air} (N)\\\\
\\noindent {\\tovian \\fontsize{20pt}{10pt} \\textbf{rith} }\\\\
\\noindent /r{\textprimstress}i{\texttheta}/\\\\
\\noindent Syllables: riθ\\\\


\\noindent History:

\\vspace{-0pt}
\\hspace{40pt}
\\begin{tabular}{ccc}
\\textit{0} & /rai{\texttheta}/&$\\rightarrow$ & \\textit{5500} & /ri{\texttheta}/& \\\\
\\end{tabular}

\\vspace{20pt}\\hline

\\end{nopagebreak}
\\filbreak



\\vspace{15pt}
\\begin{nopagebreak}
\\noindent{\\fontsize{20pt}{10pt}\\textbf{su-} } \\textit{ALL/allative case marker} (CASE)\\\\
\\noindent {\\tovian \\fontsize{20pt}{10pt} \\textbf{su-} }\\\\
\\noindent /s{\textprimstress}u/\\\\
\\noindent Syllables: su\\\\


\\noindent History:

\\vspace{-0pt}
\\hspace{40pt}
\\begin{tabular}{ccc}
\\textit{0} & /su-/& \\\\
\\end{tabular}

\\vspace{20pt}\\hline

\\end{nopagebreak}
\\filbreak



\\vspace{15pt}
\\begin{nopagebreak}
\\noindent{\\fontsize{20pt}{10pt}\\textbf{la-} } \\textit{ANA/animate noun class} (CLASS)\\\\
\\noindent {\\tovian \\fontsize{20pt}{10pt} \\textbf{la-} }\\\\
\\noindent /l{\textprimstress}a/\\\\
\\noindent Syllables: la\\\\


\\noindent History:

\\vspace{-0pt}
\\hspace{40pt}
\\begin{tabular}{ccc}
\\textit{0} & /la-/& \\\\
\\end{tabular}

\\vspace{20pt}\\hline

\\end{nopagebreak}
\\filbreak



\\vspace{15pt}
\\begin{nopagebreak}
\\noindent{\\fontsize{20pt}{10pt}\\textbf{nan} } \\textit{ancientness/antiquity/oldness} (N)\\\\
\\noindent {\\tovian \\fontsize{20pt}{10pt} \\textbf{nan} }\\\\
\\noindent /n{\textprimstress}an/\\\\
\\noindent Syllables: nan\\\\


\\noindent History:

\\vspace{-0pt}
\\hspace{40pt}
\\begin{tabular}{ccc}
\\textit{0} & /nan/& \\\\
\\end{tabular}

\\vspace{20pt}\\hline

\\end{nopagebreak}
\\filbreak



\\vspace{15pt}
\\begin{nopagebreak}
\\noindent{\\fontsize{20pt}{10pt}\\textbf{la-} } \\textit{ANI/animate noun class} (CLASS)\\\\
\\noindent {\\tovian \\fontsize{20pt}{10pt} \\textbf{la-} }\\\\
\\noindent /l{\textprimstress}a/\\\\
\\noindent Syllables: la\\\\


\\noindent History:

\\vspace{-0pt}
\\hspace{40pt}
\\begin{tabular}{ccc}
\\textit{0} & /la-/& \\\\
\\end{tabular}

\\vspace{20pt}\\hline

\\end{nopagebreak}
\\filbreak



\\vspace{15pt}
\\begin{nopagebreak}
\\noindent{\\fontsize{20pt}{10pt}\\textbf{tlan} } \\textit{animal/creature/beast} (N)\\\\
\\noindent {\\tovian \\fontsize{20pt}{10pt} \\textbf{tlan} }\\\\
\\noindent /{\texttoptiebar{t\textbeltl}}{\textprimstress}an/\\\\
\\noindent Syllables: ƛan\\\\


\\noindent History:

\\vspace{-0pt}
\\hspace{40pt}
\\begin{tabular}{ccc}
\\textit{0} & /{\textsubbridge{t}}elaman/&$\\rightarrow$ & \\textit{3200} & /{\textsubbridge{t}}lamn/&$\\rightarrow$ & \\textit{6401} & /{\texttoptiebar{t\textbeltl}}amn/&$\\rightarrow$ & \\textit{7600} & /{\texttoptiebar{t\textbeltl}}an/& \\\\
\\end{tabular}

\\vspace{20pt}\\hline

\\end{nopagebreak}
\\filbreak



\\vspace{15pt}
\\begin{nopagebreak}
\\noindent{\\fontsize{20pt}{10pt}\\textbf{gasatl} } \\textit{April} (N)\\\\
\\noindent {\\tovian \\fontsize{20pt}{10pt} \\textbf{gasatl} }\\\\
\\noindent /g{\textprimstress}asa{\texttoptiebar{t\textbeltl}}/\\\\
\\noindent Syllables: gas.aƛ\\\\
\\noindent lit. heat+rain\\\\
\\noindent \\textit{Foegvel}\\\\


\\noindent History:

\\vspace{-0pt}
\\hspace{40pt}
\\begin{tabular}{ccc}
\\textit{6000} & /agazl{\texttoptiebar{t\textbeltl}}a{\texttoptiebar{t\textbeltl}}a/&$\\rightarrow$ & \\textit{7500} & /gazl{\texttoptiebar{t\textbeltl}}a{\texttoptiebar{t\textbeltl}}a/&$\\rightarrow$ & \\textit{9500} & /gazl{\texttoptiebar{t\textbeltl}}a{\texttoptiebar{t\textbeltl}}/&$\\rightarrow$ & \\textit{12000} & /gasl{\texttoptiebar{t\textbeltl}}a{\texttoptiebar{t\textbeltl}}/&$\\rightarrow$ & \\textit{12500} & /gas{\texttoptiebar{t\textbeltl}}a{\texttoptiebar{t\textbeltl}}/&$\\rightarrow$ & \\textit{14001} & /gasa{\texttoptiebar{t\textbeltl}}/& \\\\
\\end{tabular}

\\vspace{20pt}\\hline

\\end{nopagebreak}
\\filbreak



\\vspace{15pt}
\\begin{nopagebreak}
\\noindent{\\fontsize{20pt}{10pt}\\textbf{kewinef} } \\textit{August} (N)\\\\
\\noindent {\\tovian \\fontsize{20pt}{10pt} \\textbf{kewinef} }\\\\
\\noindent /kew{\textprimstress}inef/\\\\
\\noindent Syllables: kew.in.ef\\\\
\\noindent lit. dancing+sun\\\\
\\noindent \\textit{Floreart}\\\\


\\noindent History:

\\vspace{-0pt}
\\hspace{40pt}
\\begin{tabular}{ccc}
\\textit{6000} & /kwilwnefl/&$\\rightarrow$ & \\textit{7600} & /kwiwnefl/&$\\rightarrow$ & \\textit{12500} & /kwiwnef/&$\\rightarrow$ & \\textit{13003} & /kwinnef/&$\\rightarrow$ & \\textit{15000} & /kwinef/&$\\rightarrow$ & \\textit{15001} & /kewinef/& \\\\
\\end{tabular}

\\vspace{20pt}\\hline

\\end{nopagebreak}
\\filbreak



\\vspace{15pt}
\\begin{nopagebreak}
\\noindent{\\fontsize{20pt}{10pt}\\textbf{mulhelesh} } \\textit{Autumn} (N)\\\\
\\noindent {\\tovian \\fontsize{20pt}{10pt} \\textbf{mulhelesh} }\\\\
\\noindent /mu{\textbeltl}{\textprimstress}ele{\textesh}/\\\\
\\noindent Syllables: muɬ.el.eʃ\\\\
\\noindent lit. present+harvest\\\\


\\noindent History:

\\vspace{-0pt}
\\hspace{40pt}
\\begin{tabular}{ccc}
\\textit{0} & /num{\textbeltl}ele{\textschwa}mellore{\textyogh}i{\texttheta}/&$\\rightarrow$ & \\textit{2300} & /num{\textbeltl}elemellore{\textyogh}i{\texttheta}/&$\\rightarrow$ & \\textit{3200} & /nm{\textbeltl}lemllre{\textyogh}{\texttheta}/&$\\rightarrow$ & \\textit{4502} & /nm{\textbeltl}lemllre{\textyogh}/&$\\rightarrow$ & \\textit{6500} & /nm{\textbeltl}lemlre{\textyogh}/&$\\rightarrow$ & \\textit{6503} & /nm{\textbeltl}lemrle{\textyogh}/&$\\rightarrow$ & \\textit{7501} & /m{\textbeltl}lemrle{\textyogh}/&$\\rightarrow$ & \\textit{7600} & /m{\textbeltl}lerle{\textyogh}/&$\\rightarrow$ & \\textit{12001} & /m{\textbeltl}lerle{\textesh}/&$\\rightarrow$ & \\textit{12500} & /m{\textbeltl}erle{\textesh}/&$\\rightarrow$ & \\textit{12505} & /m{\textbeltl}ele{\textesh}/&$\\rightarrow$ & \\textit{15001} & /mu{\textbeltl}ele{\textesh}/& \\\\
\\end{tabular}

\\vspace{20pt}\\hline

\\end{nopagebreak}
\\filbreak



\\vspace{15pt}
\\begin{nopagebreak}
\\noindent{\\fontsize{20pt}{10pt}\\textbf{kar} } \\textit{balance} (N)\\\\
\\noindent {\\tovian \\fontsize{20pt}{10pt} \\textbf{kar} }\\\\
\\noindent /k{\textprimstress}ar/\\\\
\\noindent Syllables: kar\\\\


\\noindent History:

\\vspace{-0pt}
\\hspace{40pt}
\\begin{tabular}{ccc}
\\textit{0} & /kara/&$\\rightarrow$ & \\textit{9500} & /kar/& \\\\
\\end{tabular}

\\vspace{20pt}\\hline

\\end{nopagebreak}
\\filbreak



\\vspace{15pt}
\\begin{nopagebreak}
\\noindent{\\fontsize{20pt}{10pt}\\textbf{wen} } \\textit{beauty} (N)\\\\
\\noindent {\\tovian \\fontsize{20pt}{10pt} \\textbf{wen} }\\\\
\\noindent /w{\textprimstress}en/\\\\
\\noindent Syllables: wen\\\\


\\noindent History:

\\vspace{-0pt}
\\hspace{40pt}
\\begin{tabular}{ccc}
\\textit{0} & /wen/& \\\\
\\end{tabular}

\\vspace{20pt}\\hline

\\end{nopagebreak}
\\filbreak



\\vspace{15pt}
\\begin{nopagebreak}
\\noindent{\\fontsize{20pt}{10pt}\\textbf{kan} } \\textit{becoming} (N)\\\\
\\noindent {\\tovian \\fontsize{20pt}{10pt} \\textbf{kan} }\\\\
\\noindent /k{\textprimstress}an/\\\\
\\noindent Syllables: kan\\\\
\\noindent lit. arriving+change\\\\


\\noindent History:

\\vspace{-0pt}
\\hspace{40pt}
\\begin{tabular}{ccc}
\\textit{0} & /ki{\texttheta}sakana/&$\\rightarrow$ & \\textit{1000} & /k{\texttheta}skana/&$\\rightarrow$ & \\textit{3000} & /k{\texttheta}sana/&$\\rightarrow$ & \\textit{4502} & /k{\texttheta}ana/&$\\rightarrow$ & \\textit{9500} & /k{\texttheta}an/&$\\rightarrow$ & \\textit{10700} & /k{\textsubbridge{t}}an/&$\\rightarrow$ & \\textit{11500} & /kan/& \\\\
\\end{tabular}

\\vspace{20pt}\\hline

\\end{nopagebreak}
\\filbreak



\\vspace{15pt}
\\begin{nopagebreak}
\\noindent{\\fontsize{20pt}{10pt}\\textbf{lon} } \\textit{beginning/starting} (N)\\\\
\\noindent {\\tovian \\fontsize{20pt}{10pt} \\textbf{lon} }\\\\
\\noindent /l{\textprimstress}on/\\\\
\\noindent Syllables: lon\\\\


\\noindent History:

\\vspace{-0pt}
\\hspace{40pt}
\\begin{tabular}{ccc}
\\textit{0} & /yaloran/&$\\rightarrow$ & \\textit{3200} & /ylorn/&$\\rightarrow$ & \\textit{7600} & /lon/& \\\\
\\end{tabular}

\\vspace{20pt}\\hline

\\end{nopagebreak}
\\filbreak



\\vspace{15pt}
\\begin{nopagebreak}
\\noindent{\\fontsize{20pt}{10pt}\\textbf{hul} } \\textit{behind} (N)\\\\
\\noindent {\\tovian \\fontsize{20pt}{10pt} \\textbf{hul} }\\\\
\\noindent /ħ{\textprimstress}ul/\\\\
\\noindent Syllables: ħul\\\\


\\noindent History:

\\vspace{-0pt}
\\hspace{40pt}
\\begin{tabular}{ccc}
\\textit{0} & /hul/&$\\rightarrow$ & \\textit{14006} & /ħul/& \\\\
\\end{tabular}

\\vspace{20pt}\\hline

\\end{nopagebreak}
\\filbreak



\\vspace{15pt}
\\begin{nopagebreak}
\\noindent{\\fontsize{20pt}{10pt}\\textbf{sil} } \\textit{being/existence/created thing} (N)\\\\
\\noindent {\\tovian \\fontsize{20pt}{10pt} \\textbf{sil} }\\\\
\\noindent /s{\textprimstress}il/\\\\
\\noindent Syllables: sil\\\\


\\noindent History:

\\vspace{-0pt}
\\hspace{40pt}
\\begin{tabular}{ccc}
\\textit{0} & /sil/& \\\\
\\end{tabular}

\\vspace{20pt}\\hline

\\end{nopagebreak}
\\filbreak



\\vspace{15pt}
\\begin{nopagebreak}
\\noindent{\\fontsize{20pt}{10pt}\\textbf{shanin} } \\textit{berry} (N)\\\\
\\noindent {\\tovian \\fontsize{20pt}{10pt} \\textbf{shanin} }\\\\
\\noindent /{\textesh}{\textprimstress}anin/\\\\
\\noindent Syllables: ʃan.in\\\\
\\noindent lit. fruit+small\\\\


\\noindent History:

\\vspace{-0pt}
\\hspace{40pt}
\\begin{tabular}{ccc}
\\textit{1000} & /{\textyogh}olanina/&$\\rightarrow$ & \\textit{3200} & /{\textyogh}lanina/&$\\rightarrow$ & \\textit{9500} & /{\textyogh}lanin/&$\\rightarrow$ & \\textit{12001} & /{\textesh}lanin/&$\\rightarrow$ & \\textit{12500} & /{\textesh}anin/& \\\\
\\end{tabular}

\\vspace{20pt}\\hline

\\end{nopagebreak}
\\filbreak



\\vspace{15pt}
\\begin{nopagebreak}
\\noindent{\\fontsize{20pt}{10pt}\\textbf{mot} } \\textit{big/large} (N)\\\\
\\noindent {\\tovian \\fontsize{20pt}{10pt} \\textbf{mot} }\\\\
\\noindent /m{\textprimstress}o{\textsubbridge{t}}/\\\\
\\noindent Syllables: mot\\\\


\\noindent History:

\\vspace{-0pt}
\\hspace{40pt}
\\begin{tabular}{ccc}
\\textit{0} & /mo{\textsubbridge{t}}/& \\\\
\\end{tabular}

\\vspace{20pt}\\hline

\\end{nopagebreak}
\\filbreak



\\vspace{15pt}
\\begin{nopagebreak}
\\noindent{\\fontsize{20pt}{10pt}\\textbf{rithan} } \\textit{bird} (N)\\\\
\\noindent {\\tovian \\fontsize{20pt}{10pt} \\textbf{rithan} }\\\\
\\noindent /r{\textprimstress}i{\texttheta}an/\\\\
\\noindent Syllables: riθ.an\\\\
\\noindent lit. air+animal\\\\


\\noindent History:

\\vspace{-0pt}
\\hspace{40pt}
\\begin{tabular}{ccc}
\\textit{100} & /rai{\texttheta}{\textsubbridge{t}}elaman/&$\\rightarrow$ & \\textit{3000} & /rai{\texttheta}elaman/&$\\rightarrow$ & \\textit{3200} & /rai{\texttheta}lamn/&$\\rightarrow$ & \\textit{5500} & /ri{\texttheta}lamn/&$\\rightarrow$ & \\textit{7600} & /ri{\texttheta}lan/&$\\rightarrow$ & \\textit{12500} & /ri{\texttheta}an/& \\\\
\\end{tabular}

\\vspace{20pt}\\hline

\\end{nopagebreak}
\\filbreak



\\vspace{15pt}
\\begin{nopagebreak}
\\noindent{\\fontsize{20pt}{10pt}\\textbf{lhaswan} } \\textit{birth} (N)\\\\
\\noindent {\\tovian \\fontsize{20pt}{10pt} \\textbf{lhaswan} }\\\\
\\noindent /{\textbeltl}{\textprimstress}aswan/\\\\
\\noindent Syllables: ɬasw.an\\\\
\\noindent lit. bringing+ALL+light\\\\


\\noindent History:

\\vspace{-0pt}
\\hspace{40pt}
\\begin{tabular}{ccc}
\\textit{50} & /{\textbeltl}ogasu-ahana/&$\\rightarrow$ & \\textit{1000} & /{\textbeltl}ogasuahana/&$\\rightarrow$ & \\textit{2202} & /{\textbeltl}ogasuwahana/&$\\rightarrow$ & \\textit{3200} & /{\textbeltl}gaswahana/&$\\rightarrow$ & \\textit{4500} & /{\textbeltl}aswahana/&$\\rightarrow$ & \\textit{5000} & /{\textbeltl}aswaana/&$\\rightarrow$ & \\textit{5500} & /{\textbeltl}aswana/&$\\rightarrow$ & \\textit{9500} & /{\textbeltl}aswan/& \\\\
\\end{tabular}

\\vspace{20pt}\\hline

\\end{nopagebreak}
\\filbreak



\\vspace{15pt}
\\begin{nopagebreak}
\\noindent{\\fontsize{20pt}{10pt}\\textbf{derik} } \\textit{blood} (N)\\\\
\\noindent {\\tovian \\fontsize{20pt}{10pt} \\textbf{derik} }\\\\
\\noindent /d{\textprimstress}erik/\\\\
\\noindent Syllables: der.ik\\\\


\\noindent History:

\\vspace{-0pt}
\\hspace{40pt}
\\begin{tabular}{ccc}
\\textit{0} & /driks/&$\\rightarrow$ & \\textit{13001} & /drik/&$\\rightarrow$ & \\textit{15001} & /derik/& \\\\
\\end{tabular}

\\vspace{20pt}\\hline

\\end{nopagebreak}
\\filbreak



\\vspace{15pt}
\\begin{nopagebreak}
\\noindent{\\fontsize{20pt}{10pt}\\textbf{med} } \\textit{bloom} (N)\\\\
\\noindent {\\tovian \\fontsize{20pt}{10pt} \\textbf{med} }\\\\
\\noindent /m{\textprimstress}ed/\\\\
\\noindent Syllables: med\\\\
\\noindent lit. flower+showing\\\\


\\noindent History:

\\vspace{-0pt}
\\hspace{40pt}
\\begin{tabular}{ccc}
\\textit{1000} & /drenpkfe{\textsubbridge{t}}e/&$\\rightarrow$ & \\textit{2000} & /drenpkfede/&$\\rightarrow$ & \\textit{3000} & /drenpfede/&$\\rightarrow$ & \\textit{3200} & /drnpfede/&$\\rightarrow$ & \\textit{6000} & /drmpfede/&$\\rightarrow$ & \\textit{7501} & /rmpfede/&$\\rightarrow$ & \\textit{7600} & /mpfede/&$\\rightarrow$ & \\textit{8750} & /mfede/&$\\rightarrow$ & \\textit{9500} & /mfed/&$\\rightarrow$ & \\textit{11995} & /med/& \\\\
\\end{tabular}

\\vspace{20pt}\\hline

\\end{nopagebreak}
\\filbreak



\\vspace{15pt}
\\begin{nopagebreak}
\\noindent{\\fontsize{20pt}{10pt}\\textbf{dor} } \\textit{blueness/blue} (N)\\\\
\\noindent {\\tovian \\fontsize{20pt}{10pt} \\textbf{dor} }\\\\
\\noindent /d{\textprimstress}or/\\\\
\\noindent Syllables: dor\\\\


\\noindent History:

\\vspace{-0pt}
\\hspace{40pt}
\\begin{tabular}{ccc}
\\textit{0} & /dora/&$\\rightarrow$ & \\textit{9500} & /dor/& \\\\
\\end{tabular}

\\vspace{20pt}\\hline

\\end{nopagebreak}
\\filbreak



\\vspace{15pt}
\\begin{nopagebreak}
\\noindent{\\fontsize{20pt}{10pt}\\textbf{sith} } \\textit{boat/ship/canoe} (N)\\\\
\\noindent {\\tovian \\fontsize{20pt}{10pt} \\textbf{sith} }\\\\
\\noindent /s{\textprimstress}i{\texttheta}/\\\\
\\noindent Syllables: siθ\\\\
\\noindent lit. water+traveling\\\\


\\noindent History:

\\vspace{-0pt}
\\hspace{40pt}
\\begin{tabular}{ccc}
\\textit{3000} & /zdelki{\texttheta}sa/&$\\rightarrow$ & \\textit{3000} & /zdeli{\texttheta}sa/&$\\rightarrow$ & \\textit{3200} & /zdli{\texttheta}sa/&$\\rightarrow$ & \\textit{4500} & /zli{\texttheta}sa/&$\\rightarrow$ & \\textit{4502} & /zli{\texttheta}a/&$\\rightarrow$ & \\textit{9500} & /zli{\texttheta}/&$\\rightarrow$ & \\textit{12000} & /sli{\texttheta}/&$\\rightarrow$ & \\textit{12500} & /si{\texttheta}/& \\\\
\\end{tabular}

\\vspace{20pt}\\hline

\\end{nopagebreak}
\\filbreak



\\vspace{15pt}
\\begin{nopagebreak}
\\noindent{\\fontsize{20pt}{10pt}\\textbf{adlel} } \\textit{body/trunk} (N)\\\\
\\noindent {\\tovian \\fontsize{20pt}{10pt} \\textbf{adlel} }\\\\
\\noindent /{\textprimstress}adlel/\\\\
\\noindent Syllables: adl.el\\\\


\\noindent History:

\\vspace{-0pt}
\\hspace{40pt}
\\begin{tabular}{ccc}
\\textit{0} & /a{\textsubbridge{t}}ellenal/&$\\rightarrow$ & \\textit{2000} & /adellenal/&$\\rightarrow$ & \\textit{3200} & /adllenl/&$\\rightarrow$ & \\textit{6500} & /adlenl/&$\\rightarrow$ & \\textit{7600} & /adlel/& \\\\
\\end{tabular}

\\vspace{20pt}\\hline

\\end{nopagebreak}
\\filbreak



\\vspace{15pt}
\\begin{nopagebreak}
\\noindent{\\fontsize{20pt}{10pt}\\textbf{men} } \\textit{brick} (N)\\\\
\\noindent {\\tovian \\fontsize{20pt}{10pt} \\textbf{men} }\\\\
\\noindent /m{\textprimstress}en/\\\\
\\noindent Syllables: men\\\\
\\noindent lit. fire+rock\\\\


\\noindent History:

\\vspace{-0pt}
\\hspace{40pt}
\\begin{tabular}{ccc}
\\textit{1300} & /ke{\textesh}meron/&$\\rightarrow$ & \\textit{3200} & /k{\textesh}mern/&$\\rightarrow$ & \\textit{7501} & /{\textesh}mern/&$\\rightarrow$ & \\textit{7600} & /{\textesh}men/&$\\rightarrow$ & \\textit{13002} & /men/& \\\\
\\end{tabular}

\\vspace{20pt}\\hline

\\end{nopagebreak}
\\filbreak



\\vspace{15pt}
\\begin{nopagebreak}
\\noindent{\\fontsize{20pt}{10pt}\\textbf{lhog} } \\textit{bringing/carrying} (N)\\\\
\\noindent {\\tovian \\fontsize{20pt}{10pt} \\textbf{lhog} }\\\\
\\noindent /{\textbeltl}{\textprimstress}og/\\\\
\\noindent Syllables: ɬog\\\\


\\noindent History:

\\vspace{-0pt}
\\hspace{40pt}
\\begin{tabular}{ccc}
\\textit{0} & /{\textbeltl}oga/&$\\rightarrow$ & \\textit{9500} & /{\textbeltl}og/& \\\\
\\end{tabular}

\\vspace{20pt}\\hline

\\end{nopagebreak}
\\filbreak



\\vspace{15pt}
\\begin{nopagebreak}
\\noindent{\\fontsize{20pt}{10pt}\\textbf{ke-} } \\textit{CAUS/causative voice marker} (VOICE)\\\\
\\noindent {\\tovian \\fontsize{20pt}{10pt} \\textbf{ke-} }\\\\
\\noindent /k{\textprimstress}e/\\\\
\\noindent Syllables: ke\\\\


\\noindent History:

\\vspace{-0pt}
\\hspace{40pt}
\\begin{tabular}{ccc}
\\textit{0} & /ke-/& \\\\
\\end{tabular}

\\vspace{20pt}\\hline

\\end{nopagebreak}
\\filbreak



\\vspace{15pt}
\\begin{nopagebreak}
\\noindent{\\fontsize{20pt}{10pt}\\textbf{lhel} } \\textit{center} (N)\\\\
\\noindent {\\tovian \\fontsize{20pt}{10pt} \\textbf{lhel} }\\\\
\\noindent /{\textbeltl}{\textprimstress}el/\\\\
\\noindent Syllables: ɬel\\\\


\\noindent History:

\\vspace{-0pt}
\\hspace{40pt}
\\begin{tabular}{ccc}
\\textit{0} & /{\textbeltl}el/& \\\\
\\end{tabular}

\\vspace{20pt}\\hline

\\end{nopagebreak}
\\filbreak



\\vspace{15pt}
\\begin{nopagebreak}
\\noindent{\\fontsize{20pt}{10pt}\\textbf{kan} } \\textit{change} (N)\\\\
\\noindent {\\tovian \\fontsize{20pt}{10pt} \\textbf{kan} }\\\\
\\noindent /k{\textprimstress}an/\\\\
\\noindent Syllables: kan\\\\


\\noindent History:

\\vspace{-0pt}
\\hspace{40pt}
\\begin{tabular}{ccc}
\\textit{0} & /kana/&$\\rightarrow$ & \\textit{9500} & /kan/& \\\\
\\end{tabular}

\\vspace{20pt}\\hline

\\end{nopagebreak}
\\filbreak



\\vspace{15pt}
\\begin{nopagebreak}
\\noindent{\\fontsize{20pt}{10pt}\\textbf{melhen} } \\textit{child} (N)\\\\
\\noindent {\\tovian \\fontsize{20pt}{10pt} \\textbf{melhen} }\\\\
\\noindent /m{\textprimstress}e{\textbeltl}en/\\\\
\\noindent Syllables: meɬ.en\\\\
\\noindent lit. newness+dream\\\\


\\noindent History:

\\vspace{-0pt}
\\hspace{40pt}
\\begin{tabular}{ccc}
\\textit{0} & /{\textyogh}emeno{\textbeltl}ena/&$\\rightarrow$ & \\textit{3200} & /{\textyogh}men{\textbeltl}ena/&$\\rightarrow$ & \\textit{6400} & /{\textyogh}me{\textbeltl}ena/&$\\rightarrow$ & \\textit{9500} & /{\textyogh}me{\textbeltl}en/&$\\rightarrow$ & \\textit{12001} & /{\textesh}me{\textbeltl}en/&$\\rightarrow$ & \\textit{13002} & /me{\textbeltl}en/& \\\\
\\end{tabular}

\\vspace{20pt}\\hline

\\end{nopagebreak}
\\filbreak



\\vspace{15pt}
\\begin{nopagebreak}
\\noindent{\\fontsize{20pt}{10pt}\\textbf{shithil} } \\textit{city} (N)\\\\
\\noindent {\\tovian \\fontsize{20pt}{10pt} \\textbf{shithil} }\\\\
\\noindent /{\textesh}{\textprimstress}i{\texttheta}il/\\\\
\\noindent Syllables: ʃiθ.il\\\\
\\noindent lit. wakefulness+place\\\\


\\noindent History:

\\vspace{-0pt}
\\hspace{40pt}
\\begin{tabular}{ccc}
\\textit{6000} & /{\textyogh}i{\texttheta}{\texttheta}ile/&$\\rightarrow$ & \\textit{6500} & /{\textyogh}i{\texttheta}ile/&$\\rightarrow$ & \\textit{9500} & /{\textyogh}i{\texttheta}il/&$\\rightarrow$ & \\textit{12001} & /{\textesh}i{\texttheta}il/& \\\\
\\end{tabular}

\\vspace{20pt}\\hline

\\end{nopagebreak}
\\filbreak



\\vspace{15pt}
\\begin{nopagebreak}
\\noindent{\\fontsize{20pt}{10pt}\\textbf{tam} } \\textit{closure/shutting} (N)\\\\
\\noindent {\\tovian \\fontsize{20pt}{10pt} \\textbf{tam} }\\\\
\\noindent /{\textsubbridge{t}}{\textprimstress}am/\\\\
\\noindent Syllables: tam\\\\


\\noindent History:

\\vspace{-0pt}
\\hspace{40pt}
\\begin{tabular}{ccc}
\\textit{0} & /{\textsubbridge{t}}anum/&$\\rightarrow$ & \\textit{3200} & /{\textsubbridge{t}}anm/&$\\rightarrow$ & \\textit{7600} & /{\textsubbridge{t}}am/& \\\\
\\end{tabular}

\\vspace{20pt}\\hline

\\end{nopagebreak}
\\filbreak



\\vspace{15pt}
\\begin{nopagebreak}
\\noindent{\\fontsize{20pt}{10pt}\\textbf{tush} } \\textit{coldness/cold} (N)\\\\
\\noindent {\\tovian \\fontsize{20pt}{10pt} \\textbf{tush} }\\\\
\\noindent /{\textsubbridge{t}}{\textprimstress}u{\textesh}/\\\\
\\noindent Syllables: tuʃ\\\\


\\noindent History:

\\vspace{-0pt}
\\hspace{40pt}
\\begin{tabular}{ccc}
\\textit{0} & /{\textsubbridge{t}}ku{\textesh}u/&$\\rightarrow$ & \\textit{3000} & /{\textsubbridge{t}}u{\textesh}u/&$\\rightarrow$ & \\textit{9500} & /{\textsubbridge{t}}u{\textesh}/& \\\\
\\end{tabular}

\\vspace{20pt}\\hline

\\end{nopagebreak}
\\filbreak



\\vspace{15pt}
\\begin{nopagebreak}
\\noindent{\\fontsize{20pt}{10pt}\\textbf{yi-} } \\textit{COM/comitative case marker/with/together with} (CASE)\\\\
\\noindent {\\tovian \\fontsize{20pt}{10pt} \\textbf{yi-} }\\\\
\\noindent /y{\textprimstress}i/\\\\
\\noindent Syllables: ji\\\\


\\noindent History:

\\vspace{-0pt}
\\hspace{40pt}
\\begin{tabular}{ccc}
\\textit{0} & /yi-/& \\\\
\\end{tabular}

\\vspace{20pt}\\hline

\\end{nopagebreak}
\\filbreak



\\vspace{15pt}
\\begin{nopagebreak}
\\noindent{\\fontsize{20pt}{10pt}\\textbf{yalrith} } \\textit{common bird} (N)\\\\
\\noindent {\\tovian \\fontsize{20pt}{10pt} \\textbf{yalrith} }\\\\
\\noindent /y{\textprimstress}alri{\texttheta}/\\\\
\\noindent Syllables: jalr.iθ\\\\
\\noindent lit. normality+air\\\\


\\noindent History:

\\vspace{-0pt}
\\hspace{40pt}
\\begin{tabular}{ccc}
\\textit{12000} & /yalri{\texttheta}/& \\\\
\\end{tabular}

\\vspace{20pt}\\hline

\\end{nopagebreak}
\\filbreak



\\vspace{15pt}
\\begin{nopagebreak}
\\noindent{\\fontsize{20pt}{10pt}\\textbf{yal} } \\textit{commonness/mundanity/normality} (N)\\\\
\\noindent {\\tovian \\fontsize{20pt}{10pt} \\textbf{yal} }\\\\
\\noindent /y{\textprimstress}al/\\\\
\\noindent Syllables: jal\\\\


\\noindent History:

\\vspace{-0pt}
\\hspace{40pt}
\\begin{tabular}{ccc}
\\textit{0} & /ileriala/&$\\rightarrow$ & \\textit{2201} & /ileriyala/&$\\rightarrow$ & \\textit{3200} & /ilriyala/&$\\rightarrow$ & \\textit{6300} & /ilryala/&$\\rightarrow$ & \\textit{6503} & /irlyala/&$\\rightarrow$ & \\textit{7500} & /rlyala/&$\\rightarrow$ & \\textit{7501} & /lyala/&$\\rightarrow$ & \\textit{7600} & /yala/&$\\rightarrow$ & \\textit{9500} & /yal/& \\\\
\\end{tabular}

\\vspace{20pt}\\hline

\\end{nopagebreak}
\\filbreak



\\vspace{15pt}
\\begin{nopagebreak}
\\noindent{\\fontsize{20pt}{10pt}\\textbf{lad} } \\textit{comparison} (N)\\\\
\\noindent {\\tovian \\fontsize{20pt}{10pt} \\textbf{lad} }\\\\
\\noindent /l{\textprimstress}ad/\\\\
\\noindent Syllables: lad\\\\


\\noindent History:

\\vspace{-0pt}
\\hspace{40pt}
\\begin{tabular}{ccc}
\\textit{0} & /roladil/&$\\rightarrow$ & \\textit{3200} & /rladl/&$\\rightarrow$ & \\textit{7600} & /ladl/&$\\rightarrow$ & \\textit{13001} & /lad/& \\\\
\\end{tabular}

\\vspace{20pt}\\hline

\\end{nopagebreak}
\\filbreak



\\vspace{15pt}
\\begin{nopagebreak}
\\noindent{\\fontsize{20pt}{10pt}\\textbf{om} } \\textit{completeness/wholeness} (N)\\\\
\\noindent {\\tovian \\fontsize{20pt}{10pt} \\textbf{om} }\\\\
\\noindent /{\textprimstress}om/\\\\
\\noindent Syllables: om\\\\


\\noindent History:

\\vspace{-0pt}
\\hspace{40pt}
\\begin{tabular}{ccc}
\\textit{0} & /oma/&$\\rightarrow$ & \\textit{9500} & /om/& \\\\
\\end{tabular}

\\vspace{20pt}\\hline

\\end{nopagebreak}
\\filbreak



\\vspace{15pt}
\\begin{nopagebreak}
\\noindent{\\fontsize{20pt}{10pt}\\textbf{ko-} } \\textit{COND/conditional mood marker} (MOOD)\\\\
\\noindent {\\tovian \\fontsize{20pt}{10pt} \\textbf{ko-} }\\\\
\\noindent /k{\textprimstress}o/\\\\
\\noindent Syllables: ko\\\\


\\noindent History:

\\vspace{-0pt}
\\hspace{40pt}
\\begin{tabular}{ccc}
\\textit{0} & /ko-/& \\\\
\\end{tabular}

\\vspace{20pt}\\hline

\\end{nopagebreak}
\\filbreak



\\vspace{15pt}
\\begin{nopagebreak}
\\noindent{\\fontsize{20pt}{10pt}\\textbf{gelor} } \\textit{conflict/fighting} (N)\\\\
\\noindent {\\tovian \\fontsize{20pt}{10pt} \\textbf{gelor} }\\\\
\\noindent /g{\textprimstress}elor/\\\\
\\noindent Syllables: gel.or\\\\


\\noindent History:

\\vspace{-0pt}
\\hspace{40pt}
\\begin{tabular}{ccc}
\\textit{0} & /gilore/&$\\rightarrow$ & \\textit{3200} & /glore/&$\\rightarrow$ & \\textit{9500} & /glor/&$\\rightarrow$ & \\textit{15001} & /gelor/& \\\\
\\end{tabular}

\\vspace{20pt}\\hline

\\end{nopagebreak}
\\filbreak



\\vspace{15pt}
\\begin{nopagebreak}
\\noindent{\\fontsize{20pt}{10pt}\\textbf{leth} } \\textit{consumption/tasting/eating} (N)\\\\
\\noindent {\\tovian \\fontsize{20pt}{10pt} \\textbf{leth} }\\\\
\\noindent /l{\textprimstress}e{\texttheta}/\\\\
\\noindent Syllables: leθ\\\\


\\noindent History:

\\vspace{-0pt}
\\hspace{40pt}
\\begin{tabular}{ccc}
\\textit{0} & /ale{\dh}a/&$\\rightarrow$ & \\textit{7500} & /le{\dh}a/&$\\rightarrow$ & \\textit{9500} & /le{\dh}/&$\\rightarrow$ & \\textit{12002} & /le{\texttheta}/& \\\\
\\end{tabular}

\\vspace{20pt}\\hline

\\end{nopagebreak}
\\filbreak



\\vspace{15pt}
\\begin{nopagebreak}
\\noindent{\\fontsize{20pt}{10pt}\\textbf{lhom} } \\textit{container/vessel} (N)\\\\
\\noindent {\\tovian \\fontsize{20pt}{10pt} \\textbf{lhom} }\\\\
\\noindent /{\textbeltl}{\textprimstress}om/\\\\
\\noindent Syllables: ɬom\\\\


\\noindent History:

\\vspace{-0pt}
\\hspace{40pt}
\\begin{tabular}{ccc}
\\textit{0} & /{\textbeltl}oma/&$\\rightarrow$ & \\textit{9500} & /{\textbeltl}om/& \\\\
\\end{tabular}

\\vspace{20pt}\\hline

\\end{nopagebreak}
\\filbreak



\\vspace{15pt}
\\begin{nopagebreak}
\\noindent{\\fontsize{20pt}{10pt}\\textbf{lhun} } \\textit{council} (N)\\\\
\\noindent {\\tovian \\fontsize{20pt}{10pt} \\textbf{lhun} }\\\\
\\noindent /{\textbeltl}{\textprimstress}un/\\\\
\\noindent Syllables: ɬun\\\\
\\noindent lit. thinking+joining\\\\


\\noindent History:

\\vspace{-0pt}
\\hspace{40pt}
\\begin{tabular}{ccc}
\\textit{0} & /na{\textbeltl}{\texttoptiebar{t\textbeltl}}una/&$\\rightarrow$ & \\textit{3200} & /n{\textbeltl}{\texttoptiebar{t\textbeltl}}una/&$\\rightarrow$ & \\textit{6400} & /{\textbeltl}{\texttoptiebar{t\textbeltl}}una/&$\\rightarrow$ & \\textit{9500} & /{\textbeltl}{\texttoptiebar{t\textbeltl}}un/&$\\rightarrow$ & \\textit{14001} & /{\textbeltl}un/& \\\\
\\end{tabular}

\\vspace{20pt}\\hline

\\end{nopagebreak}
\\filbreak



\\vspace{15pt}
\\begin{nopagebreak}
\\noindent{\\fontsize{20pt}{10pt}\\textbf{kso-} } \\textit{COUNTER/counterfactual mood marker} (MOOD)\\\\
\\noindent {\\tovian \\fontsize{20pt}{10pt} \\textbf{kso-} }\\\\
\\noindent /ks{\textprimstress}o/\\\\
\\noindent Syllables: kso\\\\


\\noindent History:

\\vspace{-0pt}
\\hspace{40pt}
\\begin{tabular}{ccc}
\\textit{0} & /kso-/& \\\\
\\end{tabular}

\\vspace{20pt}\\hline

\\end{nopagebreak}
\\filbreak



\\vspace{15pt}
\\begin{nopagebreak}
\\noindent{\\fontsize{20pt}{10pt}\\textbf{lor} } \\textit{dance} (N)\\\\
\\noindent {\\tovian \\fontsize{20pt}{10pt} \\textbf{lor} }\\\\
\\noindent /l{\textprimstress}or/\\\\
\\noindent Syllables: lor\\\\
\\noindent lit. movement+action\\\\


\\noindent History:

\\vspace{-0pt}
\\hspace{40pt}
\\begin{tabular}{ccc}
\\textit{0} & /kwillore/&$\\rightarrow$ & \\textit{3200} & /kwllore/&$\\rightarrow$ & \\textit{6500} & /kwlore/&$\\rightarrow$ & \\textit{7501} & /wlore/&$\\rightarrow$ & \\textit{7600} & /lore/&$\\rightarrow$ & \\textit{9500} & /lor/& \\\\
\\end{tabular}

\\vspace{20pt}\\hline

\\end{nopagebreak}
\\filbreak



\\vspace{15pt}
\\begin{nopagebreak}
\\noindent{\\fontsize{20pt}{10pt}\\textbf{kewin} } \\textit{dancing} (N)\\\\
\\noindent {\\tovian \\fontsize{20pt}{10pt} \\textbf{kewin} }\\\\
\\noindent /k{\textprimstress}ewin/\\\\
\\noindent Syllables: kew.in\\\\
\\noindent lit. movement+beauty\\\\


\\noindent History:

\\vspace{-0pt}
\\hspace{40pt}
\\begin{tabular}{ccc}
\\textit{0} & /kwilwen/&$\\rightarrow$ & \\textit{3200} & /kwilwn/&$\\rightarrow$ & \\textit{7600} & /kwiwn/&$\\rightarrow$ & \\textit{13001} & /kwiw/&$\\rightarrow$ & \\textit{13003} & /kwin/&$\\rightarrow$ & \\textit{15001} & /kewin/& \\\\
\\end{tabular}

\\vspace{20pt}\\hline

\\end{nopagebreak}
\\filbreak



\\vspace{15pt}
\\begin{nopagebreak}
\\noindent{\\fontsize{20pt}{10pt}\\textbf{mural} } \\textit{darkness/dark} (N)\\\\
\\noindent {\\tovian \\fontsize{20pt}{10pt} \\textbf{mural} }\\\\
\\noindent /m{\textprimstress}ural/\\\\
\\noindent Syllables: mur.al\\\\


\\noindent History:

\\vspace{-0pt}
\\hspace{40pt}
\\begin{tabular}{ccc}
\\textit{0} & /bral/&$\\rightarrow$ & \\textit{11000} & /mral/&$\\rightarrow$ & \\textit{15001} & /mural/& \\\\
\\end{tabular}

\\vspace{20pt}\\hline

\\end{nopagebreak}
\\filbreak



\\vspace{15pt}
\\begin{nopagebreak}
\\noindent{\\fontsize{20pt}{10pt}\\textbf{mi-} } \\textit{DAT/dative case marker} (CASE)\\\\
\\noindent {\\tovian \\fontsize{20pt}{10pt} \\textbf{mi-} }\\\\
\\noindent /m{\textprimstress}i/\\\\
\\noindent Syllables: mi\\\\


\\noindent History:

\\vspace{-0pt}
\\hspace{40pt}
\\begin{tabular}{ccc}
\\textit{0} & /mi-/& \\\\
\\end{tabular}

\\vspace{20pt}\\hline

\\end{nopagebreak}
\\filbreak



\\vspace{15pt}
\\begin{nopagebreak}
\\noindent{\\fontsize{20pt}{10pt}\\textbf{muloan} } \\textit{dawn/morning} (N)\\\\
\\noindent {\\tovian \\fontsize{20pt}{10pt} \\textbf{muloan} }\\\\
\\noindent /mul{\textprimstress}oan/\\\\
\\noindent Syllables: mul.o.an\\\\
\\noindent lit. first+light\\\\


\\noindent History:

\\vspace{-0pt}
\\hspace{40pt}
\\begin{tabular}{ccc}
\\textit{500} & /pu-loahana/&$\\rightarrow$ & \\textit{1000} & /puloahana/&$\\rightarrow$ & \\textit{3200} & /ploahana/&$\\rightarrow$ & \\textit{5000} & /ploaana/&$\\rightarrow$ & \\textit{5500} & /ploana/&$\\rightarrow$ & \\textit{9500} & /ploan/&$\\rightarrow$ & \\textit{11000} & /mloan/&$\\rightarrow$ & \\textit{15001} & /muloan/& \\\\
\\end{tabular}

\\vspace{20pt}\\hline

\\end{nopagebreak}
\\filbreak



\\vspace{15pt}
\\begin{nopagebreak}
\\noindent{\\fontsize{20pt}{10pt}\\textbf{lhom} } \\textit{day} (N)\\\\
\\noindent {\\tovian \\fontsize{20pt}{10pt} \\textbf{lhom} }\\\\
\\noindent /{\textbeltl}{\textprimstress}om/\\\\
\\noindent Syllables: ɬom\\\\
\\noindent \\textit{Period of time}\\\\


\\noindent History:

\\vspace{-0pt}
\\hspace{40pt}
\\begin{tabular}{ccc}
\\textit{0} & /{\textbeltl}onum/&$\\rightarrow$ & \\textit{3200} & /{\textbeltl}onm/&$\\rightarrow$ & \\textit{7600} & /{\textbeltl}om/& \\\\
\\end{tabular}

\\vspace{20pt}\\hline

\\end{nopagebreak}
\\filbreak



\\vspace{15pt}
\\begin{nopagebreak}
\\noindent{\\fontsize{20pt}{10pt}\\textbf{tewimen} } \\textit{daytime} (N)\\\\
\\noindent {\\tovian \\fontsize{20pt}{10pt} \\textbf{tewimen} }\\\\
\\noindent /{\textsubbridge{t}}ew{\textprimstress}imen/\\\\
\\noindent Syllables: tew.im.en\\\\
\\noindent \\textit{Time of day}\\\\


\\noindent History:

\\vspace{-0pt}
\\hspace{40pt}
\\begin{tabular}{ccc}
\\textit{0} & /{\textsubbridge{t}}uimena/&$\\rightarrow$ & \\textit{2202} & /{\textsubbridge{t}}uwimena/&$\\rightarrow$ & \\textit{3200} & /{\textsubbridge{t}}wimena/&$\\rightarrow$ & \\textit{9500} & /{\textsubbridge{t}}wimen/&$\\rightarrow$ & \\textit{15001} & /{\textsubbridge{t}}ewimen/& \\\\
\\end{tabular}

\\vspace{20pt}\\hline

\\end{nopagebreak}
\\filbreak



\\vspace{15pt}
\\begin{nopagebreak}
\\noindent{\\fontsize{20pt}{10pt}\\textbf{dol} } \\textit{death/dying} (N)\\\\
\\noindent {\\tovian \\fontsize{20pt}{10pt} \\textbf{dol} }\\\\
\\noindent /d{\textprimstress}ol/\\\\
\\noindent Syllables: dol\\\\


\\noindent History:

\\vspace{-0pt}
\\hspace{40pt}
\\begin{tabular}{ccc}
\\textit{0} & /{\textschwa}{\textsubbridge{t}}ol/&$\\rightarrow$ & \\textit{2000} & /{\textschwa}dol/&$\\rightarrow$ & \\textit{2300} & /dol/& \\\\
\\end{tabular}

\\vspace{20pt}\\hline

\\end{nopagebreak}
\\filbreak



\\vspace{15pt}
\\begin{nopagebreak}
\\noindent{\\fontsize{20pt}{10pt}\\textbf{ngothrimral} } \\textit{December} (N)\\\\
\\noindent {\\tovian \\fontsize{20pt}{10pt} \\textbf{ngothrimral} }\\\\
\\noindent /{\ng}o{\texttheta}r{\textprimstress}imral/\\\\
\\noindent Syllables: ŋoθr.imr.al\\\\
\\noindent lit. time+GEN+dark\\\\
\\noindent \\textit{Ventiosh}\\\\


\\noindent History:

\\vspace{-0pt}
\\hspace{40pt}
\\begin{tabular}{ccc}
\\textit{6000} & /{\ng}go{\texttheta}rhi-bral/&$\\rightarrow$ & \\textit{8500} & /{\ng}go{\texttheta}rhibral/&$\\rightarrow$ & \\textit{8750} & /{\ng}o{\texttheta}rhibral/&$\\rightarrow$ & \\textit{11000} & /{\ng}o{\texttheta}rhimral/&$\\rightarrow$ & \\textit{11002} & /{\ng}o{\texttheta}rimral/& \\\\
\\end{tabular}

\\vspace{20pt}\\hline

\\end{nopagebreak}
\\filbreak



\\vspace{15pt}
\\begin{nopagebreak}
\\noindent{\\fontsize{20pt}{10pt}\\textbf{i-} } \\textit{DEF/definite marker} (ARTICLE)\\\\
\\noindent {\\tovian \\fontsize{20pt}{10pt} \\textbf{i-} }\\\\
\\noindent /{\textprimstress}i/\\\\
\\noindent Syllables: i\\\\


\\noindent History:

\\vspace{-0pt}
\\hspace{40pt}
\\begin{tabular}{ccc}
\\textit{0} & /i-/& \\\\
\\end{tabular}

\\vspace{20pt}\\hline

\\end{nopagebreak}
\\filbreak



\\vspace{15pt}
\\begin{nopagebreak}
\\noindent{\\fontsize{20pt}{10pt}\\textbf{led} } \\textit{demonstration} (N)\\\\
\\noindent {\\tovian \\fontsize{20pt}{10pt} \\textbf{led} }\\\\
\\noindent /l{\textprimstress}ed/\\\\
\\noindent Syllables: led\\\\
\\noindent lit. action+showing\\\\


\\noindent History:

\\vspace{-0pt}
\\hspace{40pt}
\\begin{tabular}{ccc}
\\textit{0} & /lorefe{\textsubbridge{t}}e/&$\\rightarrow$ & \\textit{2000} & /lorefede/&$\\rightarrow$ & \\textit{3200} & /lrefede/&$\\rightarrow$ & \\textit{6503} & /rlefede/&$\\rightarrow$ & \\textit{7600} & /lefede/&$\\rightarrow$ & \\textit{9500} & /lefed/&$\\rightarrow$ & \\textit{12003} & /led/& \\\\
\\end{tabular}

\\vspace{20pt}\\hline

\\end{nopagebreak}
\\filbreak



\\vspace{15pt}
\\begin{nopagebreak}
\\noindent{\\fontsize{20pt}{10pt}\\textbf{shil} } \\textit{desert} (N)\\\\
\\noindent {\\tovian \\fontsize{20pt}{10pt} \\textbf{shil} }\\\\
\\noindent /{\textesh}{\textprimstress}il/\\\\
\\noindent Syllables: ʃil\\\\
\\noindent lit. sand+place\\\\


\\noindent History:

\\vspace{-0pt}
\\hspace{40pt}
\\begin{tabular}{ccc}
\\textit{1000} & /{\texttheta}ok{\textesh}{\textsubbridge{t}}{\textesh}{\texttoptiebar{t\textbeltl}}{\texttheta}ile/&$\\rightarrow$ & \\textit{1000} & /{\texttheta}k{\textesh}{\textsubbridge{t}}{\textesh}{\texttoptiebar{t\textbeltl}}{\texttheta}ile/&$\\rightarrow$ & \\textit{3000} & /{\texttheta}k{\textesh}{\textesh}{\texttoptiebar{t\textbeltl}}{\texttheta}ile/&$\\rightarrow$ & \\textit{4500} & /{\texttheta}{\textesh}{\textesh}{\texttoptiebar{t\textbeltl}}{\texttheta}ile/&$\\rightarrow$ & \\textit{4502} & /{\texttheta}{\textesh}{\texttoptiebar{t\textbeltl}}{\texttheta}ile/&$\\rightarrow$ & \\textit{7501} & /{\textesh}{\texttoptiebar{t\textbeltl}}{\texttheta}ile/&$\\rightarrow$ & \\textit{9500} & /{\textesh}{\texttoptiebar{t\textbeltl}}{\texttheta}il/&$\\rightarrow$ & \\textit{14000} & /{\textesh}{\texttoptiebar{t\textbeltl}}il/&$\\rightarrow$ & \\textit{14001} & /{\textesh}il/& \\\\
\\end{tabular}

\\vspace{20pt}\\hline

\\end{nopagebreak}
\\filbreak



\\vspace{15pt}
\\begin{nopagebreak}
\\noindent{\\fontsize{20pt}{10pt}\\textbf{tenod} } \\textit{desire} (N)\\\\
\\noindent {\\tovian \\fontsize{20pt}{10pt} \\textbf{tenod} }\\\\
\\noindent /{\textsubbridge{t}}{\textprimstress}enod/\\\\
\\noindent Syllables: ten.od\\\\


\\noindent History:

\\vspace{-0pt}
\\hspace{40pt}
\\begin{tabular}{ccc}
\\textit{0} & /{\textsubbridge{t}}ino{\textsubbridge{t}}am/&$\\rightarrow$ & \\textit{2000} & /{\textsubbridge{t}}inodam/&$\\rightarrow$ & \\textit{3200} & /{\textsubbridge{t}}nodm/&$\\rightarrow$ & \\textit{13001} & /{\textsubbridge{t}}nod/&$\\rightarrow$ & \\textit{15001} & /{\textsubbridge{t}}enod/& \\\\
\\end{tabular}

\\vspace{20pt}\\hline

\\end{nopagebreak}
\\filbreak



\\vspace{15pt}
\\begin{nopagebreak}
\\noindent{\\fontsize{20pt}{10pt}\\textbf{foretum} } \\textit{desk} (N)\\\\
\\noindent {\\tovian \\fontsize{20pt}{10pt} \\textbf{foretum} }\\\\
\\noindent /for{\textprimstress}e{\textsubbridge{t}}um/\\\\
\\noindent Syllables: for.et.um\\\\
\\noindent lit. writing+flatness\\\\


\\noindent History:

\\vspace{-0pt}
\\hspace{40pt}
\\begin{tabular}{ccc}
\\textit{7000} & /flore{\textsubbridge{t}}um/&$\\rightarrow$ & \\textit{12500} & /fore{\textsubbridge{t}}um/& \\\\
\\end{tabular}

\\vspace{20pt}\\hline

\\end{nopagebreak}
\\filbreak



\\vspace{15pt}
\\begin{nopagebreak}
\\noindent{\\fontsize{20pt}{10pt}\\textbf{sahyeol} } \\textit{discovery} (N)\\\\
\\noindent {\\tovian \\fontsize{20pt}{10pt} \\textbf{sahyeol} }\\\\
\\noindent /saħy{\textprimstress}eol/\\\\
\\noindent Syllables: saħj.e.ol\\\\
\\noindent lit. make+wisdom\\\\


\\noindent History:

\\vspace{-0pt}
\\hspace{40pt}
\\begin{tabular}{ccc}
\\textit{50} & /silahe{\textsubbridge{t}}arayeol/&$\\rightarrow$ & \\textit{1000} & /silah{\textsubbridge{t}}arayeol/&$\\rightarrow$ & \\textit{2100} & /silah{\textsubbridge{t}}aryeol/&$\\rightarrow$ & \\textit{3000} & /silaharyeol/&$\\rightarrow$ & \\textit{3200} & /slahryeol/&$\\rightarrow$ & \\textit{7600} & /slahyeol/&$\\rightarrow$ & \\textit{12500} & /sahyeol/&$\\rightarrow$ & \\textit{14006} & /saħyeol/& \\\\
\\end{tabular}

\\vspace{20pt}\\hline

\\end{nopagebreak}
\\filbreak



\\vspace{15pt}
\\begin{nopagebreak}
\\noindent{\\fontsize{20pt}{10pt}\\textbf{anon} } \\textit{dog} (N)\\\\
\\noindent {\\tovian \\fontsize{20pt}{10pt} \\textbf{anon} }\\\\
\\noindent /{\textprimstress}anon/\\\\
\\noindent Syllables: an.on\\\\


\\noindent History:

\\vspace{-0pt}
\\hspace{40pt}
\\begin{tabular}{ccc}
\\textit{0} & /am{\textsubbridge{t}}orin/&$\\rightarrow$ & \\textit{2000} & /amdorin/&$\\rightarrow$ & \\textit{3200} & /amdorn/&$\\rightarrow$ & \\textit{6000} & /andorn/&$\\rightarrow$ & \\textit{7600} & /andon/&$\\rightarrow$ & \\textit{8750} & /anon/& \\\\
\\end{tabular}

\\vspace{20pt}\\hline

\\end{nopagebreak}
\\filbreak



\\vspace{15pt}
\\begin{nopagebreak}
\\noindent{\\fontsize{20pt}{10pt}\\textbf{thok} } \\textit{down/underside/bottom} (N)\\\\
\\noindent {\\tovian \\fontsize{20pt}{10pt} \\textbf{thok} }\\\\
\\noindent /{\texttheta}{\textprimstress}ok/\\\\
\\noindent Syllables: θok\\\\


\\noindent History:

\\vspace{-0pt}
\\hspace{40pt}
\\begin{tabular}{ccc}
\\textit{0} & /{\texttheta}ok/& \\\\
\\end{tabular}

\\vspace{20pt}\\hline

\\end{nopagebreak}
\\filbreak



\\vspace{15pt}
\\begin{nopagebreak}
\\noindent{\\fontsize{20pt}{10pt}\\textbf{lhen} } \\textit{dream/spirit/soul} (N)\\\\
\\noindent {\\tovian \\fontsize{20pt}{10pt} \\textbf{lhen} }\\\\
\\noindent /{\textbeltl}{\textprimstress}en/\\\\
\\noindent Syllables: ɬen\\\\


\\noindent History:

\\vspace{-0pt}
\\hspace{40pt}
\\begin{tabular}{ccc}
\\textit{0} & /{\textbeltl}ena/&$\\rightarrow$ & \\textit{9500} & /{\textbeltl}en/& \\\\
\\end{tabular}

\\vspace{20pt}\\hline

\\end{nopagebreak}
\\filbreak



\\vspace{15pt}
\\begin{nopagebreak}
\\noindent{\\fontsize{20pt}{10pt}\\textbf{anolhyel} } \\textit{dreambringer} (N)\\\\
\\noindent {\\tovian \\fontsize{20pt}{10pt} \\textbf{anolhyel} }\\\\
\\noindent /an{\textprimstress}o{\textbeltl}yel/\\\\
\\noindent Syllables: an.oɬj.el\\\\
\\noindent lit. dream+COM+walking+person\\\\


\\noindent History:

\\vspace{-0pt}
\\hspace{40pt}
\\begin{tabular}{ccc}
\\textit{10} & /{\textbeltl}enayi-no{\textbeltl}driela/&$\\rightarrow$ & \\textit{1000} & /{\textbeltl}enayino{\textbeltl}driela/&$\\rightarrow$ & \\textit{2100} & /{\textbeltl}enyino{\textbeltl}driela/&$\\rightarrow$ & \\textit{2201} & /{\textbeltl}enyino{\textbeltl}driyela/&$\\rightarrow$ & \\textit{3200} & /{\textbeltl}nyno{\textbeltl}dryela/&$\\rightarrow$ & \\textit{4500} & /{\textbeltl}nyno{\textbeltl}ryela/&$\\rightarrow$ & \\textit{6000} & /{\textbeltl}nd{\textyogh}no{\textbeltl}ryela/&$\\rightarrow$ & \\textit{7501} & /nd{\textyogh}no{\textbeltl}ryela/&$\\rightarrow$ & \\textit{7600} & /nd{\textyogh}no{\textbeltl}yela/&$\\rightarrow$ & \\textit{8750} & /n{\textyogh}no{\textbeltl}yela/&$\\rightarrow$ & \\textit{8760} & /an{\textyogh}no{\textbeltl}yela/&$\\rightarrow$ & \\textit{9500} & /an{\textyogh}no{\textbeltl}yel/&$\\rightarrow$ & \\textit{12001} & /an{\textesh}no{\textbeltl}yel/&$\\rightarrow$ & \\textit{13002} & /anno{\textbeltl}yel/&$\\rightarrow$ & \\textit{15000} & /ano{\textbeltl}yel/& \\\\
\\end{tabular}

\\vspace{20pt}\\hline

\\end{nopagebreak}
\\filbreak



\\vspace{15pt}
\\begin{nopagebreak}
\\noindent{\\fontsize{20pt}{10pt}\\textbf{munorhan} } \\textit{dusk/evening/sunset} (N)\\\\
\\noindent {\\tovian \\fontsize{20pt}{10pt} \\textbf{munorhan} }\\\\
\\noindent /mun{\textprimstress}orħan/\\\\
\\noindent Syllables: mun.orħ.an\\\\
\\noindent lit. last+light\\\\


\\noindent History:

\\vspace{-0pt}
\\hspace{40pt}
\\begin{tabular}{ccc}
\\textit{500} & /panorahana/&$\\rightarrow$ & \\textit{3200} & /pnorhana/&$\\rightarrow$ & \\textit{9500} & /pnorhan/&$\\rightarrow$ & \\textit{11000} & /mnorhan/&$\\rightarrow$ & \\textit{11996} & /munorhan/&$\\rightarrow$ & \\textit{14006} & /munorħan/& \\\\
\\end{tabular}

\\vspace{20pt}\\hline

\\end{nopagebreak}
\\filbreak



\\vspace{15pt}
\\begin{nopagebreak}
\\noindent{\\fontsize{20pt}{10pt}\\textbf{thek} } \\textit{ear} (N)\\\\
\\noindent {\\tovian \\fontsize{20pt}{10pt} \\textbf{thek} }\\\\
\\noindent /{\texttheta}{\textprimstress}ek/\\\\
\\noindent Syllables: θek\\\\


\\noindent History:

\\vspace{-0pt}
\\hspace{40pt}
\\begin{tabular}{ccc}
\\textit{0} & /{\texttheta}ek/& \\\\
\\end{tabular}

\\vspace{20pt}\\hline

\\end{nopagebreak}
\\filbreak



\\vspace{15pt}
\\begin{nopagebreak}
\\noindent{\\fontsize{20pt}{10pt}\\textbf{tal} } \\textit{east} (N)\\\\
\\noindent {\\tovian \\fontsize{20pt}{10pt} \\textbf{tal} }\\\\
\\noindent /{\textsubbridge{t}}{\textprimstress}al/\\\\
\\noindent Syllables: tal\\\\


\\noindent History:

\\vspace{-0pt}
\\hspace{40pt}
\\begin{tabular}{ccc}
\\textit{0} & /{\textsubbridge{t}}al/& \\\\
\\end{tabular}

\\vspace{20pt}\\hline

\\end{nopagebreak}
\\filbreak



\\vspace{15pt}
\\begin{nopagebreak}
\\noindent{\\fontsize{20pt}{10pt}\\textbf{ath} } \\textit{eight/8} (N)\\\\
\\noindent {\\tovian \\fontsize{20pt}{10pt} \\textbf{ath} }\\\\
\\noindent /{\textprimstress}a{\texttheta}/\\\\
\\noindent Syllables: aθ\\\\


\\noindent History:

\\vspace{-0pt}
\\hspace{40pt}
\\begin{tabular}{ccc}
\\textit{0} & /a{\texttheta}/& \\\\
\\end{tabular}

\\vspace{20pt}\\hline

\\end{nopagebreak}
\\filbreak



\\vspace{15pt}
\\begin{nopagebreak}
\\noindent{\\fontsize{20pt}{10pt}\\textbf{muyar} } \\textit{eleven/11} (N)\\\\
\\noindent {\\tovian \\fontsize{20pt}{10pt} \\textbf{muyar} }\\\\
\\noindent /m{\textprimstress}uyar/\\\\
\\noindent Syllables: muj.ar\\\\


\\noindent History:

\\vspace{-0pt}
\\hspace{40pt}
\\begin{tabular}{ccc}
\\textit{0} & /beyar/&$\\rightarrow$ & \\textit{2100} & /byar/&$\\rightarrow$ & \\textit{11000} & /myar/&$\\rightarrow$ & \\textit{15001} & /muyar/& \\\\
\\end{tabular}

\\vspace{20pt}\\hline

\\end{nopagebreak}
\\filbreak



\\vspace{15pt}
\\begin{nopagebreak}
\\noindent{\\fontsize{20pt}{10pt}\\textbf{remi} } \\textit{empty/hollow} (N)\\\\
\\noindent {\\tovian \\fontsize{20pt}{10pt} \\textbf{remi} }\\\\
\\noindent /r{\textprimstress}emi/\\\\
\\noindent Syllables: rem.i\\\\


\\noindent History:

\\vspace{-0pt}
\\hspace{40pt}
\\begin{tabular}{ccc}
\\textit{0} & /oneremiha/&$\\rightarrow$ & \\textit{3200} & /onremiha/&$\\rightarrow$ & \\textit{5000} & /onremia/&$\\rightarrow$ & \\textit{7500} & /nremia/&$\\rightarrow$ & \\textit{7600} & /remia/&$\\rightarrow$ & \\textit{9500} & /remi/& \\\\
\\end{tabular}

\\vspace{20pt}\\hline

\\end{nopagebreak}
\\filbreak



\\vspace{15pt}
\\begin{nopagebreak}
\\noindent{\\fontsize{20pt}{10pt}\\textbf{par} } \\textit{end/final/last} (N)\\\\
\\noindent {\\tovian \\fontsize{20pt}{10pt} \\textbf{par} }\\\\
\\noindent /p{\textprimstress}ar/\\\\
\\noindent Syllables: par\\\\


\\noindent History:

\\vspace{-0pt}
\\hspace{40pt}
\\begin{tabular}{ccc}
\\textit{0} & /panor/&$\\rightarrow$ & \\textit{3200} & /panr/&$\\rightarrow$ & \\textit{7600} & /par/& \\\\
\\end{tabular}

\\vspace{20pt}\\hline

\\end{nopagebreak}
\\filbreak



\\vspace{15pt}
\\begin{nopagebreak}
\\noindent{\\fontsize{20pt}{10pt}\\textbf{yey} } \\textit{enemy} (N)\\\\
\\noindent {\\tovian \\fontsize{20pt}{10pt} \\textbf{yey} }\\\\
\\noindent /y{\textprimstress}ey/\\\\
\\noindent Syllables: jej\\\\
\\noindent lit. GEN+person+evil\\\\


\\noindent History:

\\vspace{-0pt}
\\hspace{40pt}
\\begin{tabular}{ccc}
\\textit{0} & /hi-drielaoyal/&$\\rightarrow$ & \\textit{1000} & /hidrielaoyal/&$\\rightarrow$ & \\textit{2100} & /hidrielayal/&$\\rightarrow$ & \\textit{2201} & /hidriyelayal/&$\\rightarrow$ & \\textit{3200} & /hdryelayl/&$\\rightarrow$ & \\textit{4500} & /hryelayl/&$\\rightarrow$ & \\textit{6300} & /hryelyl/&$\\rightarrow$ & \\textit{7501} & /ryelyl/&$\\rightarrow$ & \\textit{7600} & /yeyl/&$\\rightarrow$ & \\textit{13001} & /yey/& \\\\
\\end{tabular}

\\vspace{20pt}\\hline

\\end{nopagebreak}
\\filbreak



\\vspace{15pt}
\\begin{nopagebreak}
\\noindent{\\fontsize{20pt}{10pt}\\textbf{pu-} } \\textit{ESS/essive case marker} (CASE)\\\\
\\noindent {\\tovian \\fontsize{20pt}{10pt} \\textbf{pu-} }\\\\
\\noindent /p{\textprimstress}u/\\\\
\\noindent Syllables: pu\\\\


\\noindent History:

\\vspace{-0pt}
\\hspace{40pt}
\\begin{tabular}{ccc}
\\textit{0} & /pu-/& \\\\
\\end{tabular}

\\vspace{20pt}\\hline

\\end{nopagebreak}
\\filbreak



\\vspace{15pt}
\\begin{nopagebreak}
\\noindent{\\fontsize{20pt}{10pt}\\textbf{yal} } \\textit{evil/wickedness/badness} (N)\\\\
\\noindent {\\tovian \\fontsize{20pt}{10pt} \\textbf{yal} }\\\\
\\noindent /y{\textprimstress}al/\\\\
\\noindent Syllables: jal\\\\


\\noindent History:

\\vspace{-0pt}
\\hspace{40pt}
\\begin{tabular}{ccc}
\\textit{0} & /oyal/&$\\rightarrow$ & \\textit{2100} & /yal/& \\\\
\\end{tabular}

\\vspace{20pt}\\hline

\\end{nopagebreak}
\\filbreak



\\vspace{15pt}
\\begin{nopagebreak}
\\noindent{\\fontsize{20pt}{10pt}\\textbf{muril} } \\textit{eye/sight/seeing/watching} (N)\\\\
\\noindent {\\tovian \\fontsize{20pt}{10pt} \\textbf{muril} }\\\\
\\noindent /m{\textprimstress}uril/\\\\
\\noindent Syllables: mur.il\\\\


\\noindent History:

\\vspace{-0pt}
\\hspace{40pt}
\\begin{tabular}{ccc}
\\textit{0} & /em{\textschwa}rila/&$\\rightarrow$ & \\textit{2300} & /emrila/&$\\rightarrow$ & \\textit{7500} & /mrila/&$\\rightarrow$ & \\textit{9500} & /mril/&$\\rightarrow$ & \\textit{15001} & /muril/& \\\\
\\end{tabular}

\\vspace{20pt}\\hline

\\end{nopagebreak}
\\filbreak



\\vspace{15pt}
\\begin{nopagebreak}
\\noindent{\\fontsize{20pt}{10pt}\\textbf{tan} } \\textit{face} (N)\\\\
\\noindent {\\tovian \\fontsize{20pt}{10pt} \\textbf{tan} }\\\\
\\noindent /{\textsubbridge{t}}{\textprimstress}an/\\\\
\\noindent Syllables: tan\\\\


\\noindent History:

\\vspace{-0pt}
\\hspace{40pt}
\\begin{tabular}{ccc}
\\textit{0} & /{\textsubbridge{t}}aran/&$\\rightarrow$ & \\textit{3200} & /{\textsubbridge{t}}arn/&$\\rightarrow$ & \\textit{7600} & /{\textsubbridge{t}}an/& \\\\
\\end{tabular}

\\vspace{20pt}\\hline

\\end{nopagebreak}
\\filbreak



\\vspace{15pt}
\\begin{nopagebreak}
\\noindent{\\fontsize{20pt}{10pt}\\textbf{tlatl} } \\textit{falling} (N)\\\\
\\noindent {\\tovian \\fontsize{20pt}{10pt} \\textbf{tlatl} }\\\\
\\noindent /{\texttoptiebar{t\textbeltl}}{\textprimstress}a{\texttoptiebar{t\textbeltl}}/\\\\
\\noindent Syllables: ƛaƛ\\\\


\\noindent History:

\\vspace{-0pt}
\\hspace{40pt}
\\begin{tabular}{ccc}
\\textit{0} & /{\texttoptiebar{t\textbeltl}}a{\texttoptiebar{t\textbeltl}}a/&$\\rightarrow$ & \\textit{9500} & /{\texttoptiebar{t\textbeltl}}a{\texttoptiebar{t\textbeltl}}/& \\\\
\\end{tabular}

\\vspace{20pt}\\hline

\\end{nopagebreak}
\\filbreak



\\vspace{15pt}
\\begin{nopagebreak}
\\noindent{\\fontsize{20pt}{10pt}\\textbf{pan} } \\textit{father} (N)\\\\
\\noindent {\\tovian \\fontsize{20pt}{10pt} \\textbf{pan} }\\\\
\\noindent /p{\textprimstress}an/\\\\
\\noindent Syllables: pan\\\\


\\noindent History:

\\vspace{-0pt}
\\hspace{40pt}
\\begin{tabular}{ccc}
\\textit{0} & /pander/&$\\rightarrow$ & \\textit{3200} & /pandr/&$\\rightarrow$ & \\textit{8500} & /panr/&$\\rightarrow$ & \\textit{13001} & /pan/& \\\\
\\end{tabular}

\\vspace{20pt}\\hline

\\end{nopagebreak}
\\filbreak



\\vspace{15pt}
\\begin{nopagebreak}
\\noindent{\\fontsize{20pt}{10pt}\\textbf{ay} } \\textit{feather} (N)\\\\
\\noindent {\\tovian \\fontsize{20pt}{10pt} \\textbf{ay} }\\\\
\\noindent /{\textprimstress}ay/\\\\
\\noindent Syllables: aj\\\\


\\noindent History:

\\vspace{-0pt}
\\hspace{40pt}
\\begin{tabular}{ccc}
\\textit{0} & /aruya/&$\\rightarrow$ & \\textit{2100} & /arya/&$\\rightarrow$ & \\textit{7600} & /aya/&$\\rightarrow$ & \\textit{9500} & /ay/& \\\\
\\end{tabular}

\\vspace{20pt}\\hline

\\end{nopagebreak}
\\filbreak



\\vspace{15pt}
\\begin{nopagebreak}
\\noindent{\\fontsize{20pt}{10pt}\\textbf{nunaf} } \\textit{February} (N)\\\\
\\noindent {\\tovian \\fontsize{20pt}{10pt} \\textbf{nunaf} }\\\\
\\noindent /n{\textprimstress}unaf/\\\\
\\noindent Syllables: nun.af\\\\
\\noindent lit. small+sun\\\\
\\noindent \\textit{Nivost}\\\\


\\noindent History:

\\vspace{-0pt}
\\hspace{40pt}
\\begin{tabular}{ccc}
\\textit{6000} & /ninaefl/&$\\rightarrow$ & \\textit{9200} & /n{\textschwa}naefl/&$\\rightarrow$ & \\textit{10000} & /n{\textschwa}nafl/&$\\rightarrow$ & \\textit{11990} & /nnafl/&$\\rightarrow$ & \\textit{11996} & /nunafl/&$\\rightarrow$ & \\textit{12500} & /nunaf/& \\\\
\\end{tabular}

\\vspace{20pt}\\hline

\\end{nopagebreak}
\\filbreak



\\vspace{15pt}
\\begin{nopagebreak}
\\noindent{\\fontsize{20pt}{10pt}\\textbf{lheyam} } \\textit{fifteen/15} (N)\\\\
\\noindent {\\tovian \\fontsize{20pt}{10pt} \\textbf{lheyam} }\\\\
\\noindent /{\textbeltl}{\textprimstress}eyam/\\\\
\\noindent Syllables: ɬej.am\\\\
\\noindent lit. twelve+three\\\\


\\noindent History:

\\vspace{-0pt}
\\hspace{40pt}
\\begin{tabular}{ccc}
\\textit{4000} & /{\textbeltl}yafem/&$\\rightarrow$ & \\textit{12003} & /{\textbeltl}yam/&$\\rightarrow$ & \\textit{15001} & /{\textbeltl}eyam/& \\\\
\\end{tabular}

\\vspace{20pt}\\hline

\\end{nopagebreak}
\\filbreak



\\vspace{15pt}
\\begin{nopagebreak}
\\noindent{\\fontsize{20pt}{10pt}\\textbf{tom} } \\textit{finger/root} (N)\\\\
\\noindent {\\tovian \\fontsize{20pt}{10pt} \\textbf{tom} }\\\\
\\noindent /{\textsubbridge{t}}{\textprimstress}om/\\\\
\\noindent Syllables: tom\\\\


\\noindent History:

\\vspace{-0pt}
\\hspace{40pt}
\\begin{tabular}{ccc}
\\textit{0} & /{\textsubbridge{t}}omo/&$\\rightarrow$ & \\textit{9500} & /{\textsubbridge{t}}om/& \\\\
\\end{tabular}

\\vspace{20pt}\\hline

\\end{nopagebreak}
\\filbreak



\\vspace{15pt}
\\begin{nopagebreak}
\\noindent{\\fontsize{20pt}{10pt}\\textbf{kesh} } \\textit{fire/burning} (N)\\\\
\\noindent {\\tovian \\fontsize{20pt}{10pt} \\textbf{kesh} }\\\\
\\noindent /k{\textprimstress}e{\textesh}/\\\\
\\noindent Syllables: keʃ\\\\


\\noindent History:

\\vspace{-0pt}
\\hspace{40pt}
\\begin{tabular}{ccc}
\\textit{0} & /ke{\textesh}/& \\\\
\\end{tabular}

\\vspace{20pt}\\hline

\\end{nopagebreak}
\\filbreak



\\vspace{15pt}
\\begin{nopagebreak}
\\noindent{\\fontsize{20pt}{10pt}\\textbf{pul} } \\textit{first} (N)\\\\
\\noindent {\\tovian \\fontsize{20pt}{10pt} \\textbf{pul} }\\\\
\\noindent /p{\textprimstress}ul/\\\\
\\noindent Syllables: pul\\\\
\\noindent lit. ESS+one\\\\


\\noindent History:

\\vspace{-0pt}
\\hspace{40pt}
\\begin{tabular}{ccc}
\\textit{0} & /pu-lo/&$\\rightarrow$ & \\textit{1000} & /pulo/&$\\rightarrow$ & \\textit{9500} & /pul/& \\\\
\\end{tabular}

\\vspace{20pt}\\hline

\\end{nopagebreak}
\\filbreak



\\vspace{15pt}
\\begin{nopagebreak}
\\noindent{\\fontsize{20pt}{10pt}\\textbf{san} } \\textit{fish} (N)\\\\
\\noindent {\\tovian \\fontsize{20pt}{10pt} \\textbf{san} }\\\\
\\noindent /s{\textprimstress}an/\\\\
\\noindent Syllables: san\\\\
\\noindent lit. water+animal\\\\


\\noindent History:

\\vspace{-0pt}
\\hspace{40pt}
\\begin{tabular}{ccc}
\\textit{0} & /z{\textsubbridge{t}}el{\textsubbridge{t}}elaman/&$\\rightarrow$ & \\textit{2000} & /zdeldelaman/&$\\rightarrow$ & \\textit{3200} & /zdldlamn/&$\\rightarrow$ & \\textit{4500} & /zldlamn/&$\\rightarrow$ & \\textit{4501} & /zllamn/&$\\rightarrow$ & \\textit{6500} & /zlamn/&$\\rightarrow$ & \\textit{7600} & /zlan/&$\\rightarrow$ & \\textit{12000} & /slan/&$\\rightarrow$ & \\textit{12500} & /san/& \\\\
\\end{tabular}

\\vspace{20pt}\\hline

\\end{nopagebreak}
\\filbreak



\\vspace{15pt}
\\begin{nopagebreak}
\\noindent{\\fontsize{20pt}{10pt}\\textbf{rem} } \\textit{five/5} (N)\\\\
\\noindent {\\tovian \\fontsize{20pt}{10pt} \\textbf{rem} }\\\\
\\noindent /r{\textprimstress}em/\\\\
\\noindent Syllables: rem\\\\


\\noindent History:

\\vspace{-0pt}
\\hspace{40pt}
\\begin{tabular}{ccc}
\\textit{0} & /rep/&$\\rightarrow$ & \\textit{11000} & /rem/& \\\\
\\end{tabular}

\\vspace{20pt}\\hline

\\end{nopagebreak}
\\filbreak



\\vspace{15pt}
\\begin{nopagebreak}
\\noindent{\\fontsize{20pt}{10pt}\\textbf{tum} } \\textit{flatness} (_)\\\\
\\noindent {\\tovian \\fontsize{20pt}{10pt} \\textbf{tum} }\\\\
\\noindent /{\textsubbridge{t}}{\textprimstress}um/\\\\
\\noindent Syllables: tum\\\\
\\noindent lit. consistency/stability\\\\
\\noindent \\textit{N}\\\\


\\noindent History:

\\vspace{-0pt}
\\hspace{40pt}
\\begin{tabular}{ccc}
\\textit{0} & /{\textsubbridge{t}}um/& \\\\
\\end{tabular}

\\vspace{20pt}\\hline

\\end{nopagebreak}
\\filbreak



\\vspace{15pt}
\\begin{nopagebreak}
\\noindent{\\fontsize{20pt}{10pt}\\textbf{thok} } \\textit{floor/ground} (N)\\\\
\\noindent {\\tovian \\fontsize{20pt}{10pt} \\textbf{thok} }\\\\
\\noindent /{\texttheta}{\textprimstress}ok/\\\\
\\noindent Syllables: θok\\\\
\\noindent lit. down+foot\\\\


\\noindent History:

\\vspace{-0pt}
\\hspace{40pt}
\\begin{tabular}{ccc}
\\textit{10} & /{\texttheta}ok{\textesh}o{\textsubbridge{t}}/&$\\rightarrow$ & \\textit{1000} & /{\texttheta}ok{\textesh}{\textsubbridge{t}}/&$\\rightarrow$ & \\textit{3000} & /{\texttheta}ok{\textesh}/&$\\rightarrow$ & \\textit{13001} & /{\texttheta}ok/& \\\\
\\end{tabular}

\\vspace{20pt}\\hline

\\end{nopagebreak}
\\filbreak



\\vspace{15pt}
\\begin{nopagebreak}
\\noindent{\\fontsize{20pt}{10pt}\\textbf{shotl} } \\textit{flow} (N)\\\\
\\noindent {\\tovian \\fontsize{20pt}{10pt} \\textbf{shotl} }\\\\
\\noindent /{\textesh}{\textprimstress}o{\texttoptiebar{t\textbeltl}}/\\\\
\\noindent Syllables: ʃoƛ\\\\


\\noindent History:

\\vspace{-0pt}
\\hspace{40pt}
\\begin{tabular}{ccc}
\\textit{0} & /{\textesh}o{\texttoptiebar{t\textbeltl}}/& \\\\
\\end{tabular}

\\vspace{20pt}\\hline

\\end{nopagebreak}
\\filbreak



\\vspace{15pt}
\\begin{nopagebreak}
\\noindent{\\fontsize{20pt}{10pt}\\textbf{derem} } \\textit{flower} (N)\\\\
\\noindent {\\tovian \\fontsize{20pt}{10pt} \\textbf{derem} }\\\\
\\noindent /d{\textprimstress}erem/\\\\
\\noindent Syllables: der.em\\\\


\\noindent History:

\\vspace{-0pt}
\\hspace{40pt}
\\begin{tabular}{ccc}
\\textit{0} & /drenpok/&$\\rightarrow$ & \\textit{1000} & /drenpk/&$\\rightarrow$ & \\textit{3000} & /drenp/&$\\rightarrow$ & \\textit{6000} & /dremp/&$\\rightarrow$ & \\textit{8750} & /drem/&$\\rightarrow$ & \\textit{15001} & /derem/& \\\\
\\end{tabular}

\\vspace{20pt}\\hline

\\end{nopagebreak}
\\filbreak



\\vspace{15pt}
\\begin{nopagebreak}
\\noindent{\\fontsize{20pt}{10pt}\\textbf{shotl} } \\textit{flowing} (N)\\\\
\\noindent {\\tovian \\fontsize{20pt}{10pt} \\textbf{shotl} }\\\\
\\noindent /{\textesh}{\textprimstress}o{\texttoptiebar{t\textbeltl}}/\\\\
\\noindent Syllables: ʃoƛ\\\\


\\noindent History:

\\vspace{-0pt}
\\hspace{40pt}
\\begin{tabular}{ccc}
\\textit{0} & /{\textesh}o{\texttoptiebar{t\textbeltl}}/& \\\\
\\end{tabular}

\\vspace{20pt}\\hline

\\end{nopagebreak}
\\filbreak



\\vspace{15pt}
\\begin{nopagebreak}
\\noindent{\\fontsize{20pt}{10pt}\\textbf{lharis} } \\textit{flying} (N)\\\\
\\noindent {\\tovian \\fontsize{20pt}{10pt} \\textbf{lharis} }\\\\
\\noindent /{\textbeltl}{\textprimstress}aris/\\\\
\\noindent Syllables: ɬar.is\\\\


\\noindent History:

\\vspace{-0pt}
\\hspace{40pt}
\\begin{tabular}{ccc}
\\textit{998} & /{\textbeltl}e{\textyogh}ariza/&$\\rightarrow$ & \\textit{3200} & /{\textbeltl}{\textyogh}ariza/&$\\rightarrow$ & \\textit{4502} & /{\textbeltl}ariza/&$\\rightarrow$ & \\textit{9500} & /{\textbeltl}ariz/&$\\rightarrow$ & \\textit{12000} & /{\textbeltl}aris/& \\\\
\\end{tabular}

\\vspace{20pt}\\hline

\\end{nopagebreak}
\\filbreak



\\vspace{15pt}
\\begin{nopagebreak}
\\noindent{\\fontsize{20pt}{10pt}\\textbf{tlurith} } \\textit{fog} (N)\\\\
\\noindent {\\tovian \\fontsize{20pt}{10pt} \\textbf{tlurith} }\\\\
\\noindent /{\texttoptiebar{t\textbeltl}}{\textprimstress}uri{\texttheta}/\\\\
\\noindent Syllables: ƛur.iθ\\\\
\\noindent lit. tar+air\\\\


\\noindent History:

\\vspace{-0pt}
\\hspace{40pt}
\\begin{tabular}{ccc}
\\textit{1000} & /inga{\texttoptiebar{t\textbeltl}}unarai{\texttheta}/&$\\rightarrow$ & \\textit{3200} & /ing{\texttoptiebar{t\textbeltl}}unrai{\texttheta}/&$\\rightarrow$ & \\textit{5500} & /ing{\texttoptiebar{t\textbeltl}}unri{\texttheta}/&$\\rightarrow$ & \\textit{6000} & /i{\ng}g{\texttoptiebar{t\textbeltl}}unri{\texttheta}/&$\\rightarrow$ & \\textit{7500} & /{\ng}g{\texttoptiebar{t\textbeltl}}unri{\texttheta}/&$\\rightarrow$ & \\textit{7501} & /g{\texttoptiebar{t\textbeltl}}unri{\texttheta}/&$\\rightarrow$ & \\textit{7600} & /g{\texttoptiebar{t\textbeltl}}uri{\texttheta}/&$\\rightarrow$ & \\textit{11501} & /{\texttoptiebar{t\textbeltl}}uri{\texttheta}/& \\\\
\\end{tabular}

\\vspace{20pt}\\hline

\\end{nopagebreak}
\\filbreak



\\vspace{15pt}
\\begin{nopagebreak}
\\noindent{\\fontsize{20pt}{10pt}\\textbf{seyaleth} } \\textit{food} (N)\\\\
\\noindent {\\tovian \\fontsize{20pt}{10pt} \\textbf{seyaleth} }\\\\
\\noindent /sey{\textprimstress}ale{\texttheta}/\\\\
\\noindent Syllables: sej.al.eθ\\\\
\\noindent lit. INS+eating\\\\


\\noindent History:

\\vspace{-0pt}
\\hspace{40pt}
\\begin{tabular}{ccc}
\\textit{10} & /si-ale{\dh}a/&$\\rightarrow$ & \\textit{1000} & /siale{\dh}a/&$\\rightarrow$ & \\textit{2201} & /siyale{\dh}a/&$\\rightarrow$ & \\textit{3200} & /syale{\dh}a/&$\\rightarrow$ & \\textit{9500} & /syale{\dh}/&$\\rightarrow$ & \\textit{12002} & /syale{\texttheta}/&$\\rightarrow$ & \\textit{15001} & /seyale{\texttheta}/& \\\\
\\end{tabular}

\\vspace{20pt}\\hline

\\end{nopagebreak}
\\filbreak



\\vspace{15pt}
\\begin{nopagebreak}
\\noindent{\\fontsize{20pt}{10pt}\\textbf{shot} } \\textit{foot} (N)\\\\
\\noindent {\\tovian \\fontsize{20pt}{10pt} \\textbf{shot} }\\\\
\\noindent /{\textesh}{\textprimstress}o{\textsubbridge{t}}/\\\\
\\noindent Syllables: ʃot\\\\
\\noindent \\textit{_   }\\\\


\\noindent History:

\\vspace{-0pt}
\\hspace{40pt}
\\begin{tabular}{ccc}
\\textit{0} & /{\textesh}o{\textsubbridge{t}}/& \\\\
\\end{tabular}

\\vspace{20pt}\\hline

\\end{nopagebreak}
\\filbreak



\\vspace{15pt}
\\begin{nopagebreak}
\\noindent{\\fontsize{20pt}{10pt}\\textbf{lethil} } \\textit{forest} (N)\\\\
\\noindent {\\tovian \\fontsize{20pt}{10pt} \\textbf{lethil} }\\\\
\\noindent /l{\textprimstress}e{\texttheta}il/\\\\
\\noindent Syllables: leθ.il\\\\
\\noindent lit. tree+place\\\\


\\noindent History:

\\vspace{-0pt}
\\hspace{40pt}
\\begin{tabular}{ccc}
\\textit{1000} & /aklare{\texttheta}ile/&$\\rightarrow$ & \\textit{2000} & /aglare{\texttheta}ile/&$\\rightarrow$ & \\textit{3200} & /aglre{\texttheta}ile/&$\\rightarrow$ & \\textit{6503} & /agrle{\texttheta}ile/&$\\rightarrow$ & \\textit{7500} & /grle{\texttheta}ile/&$\\rightarrow$ & \\textit{7501} & /rle{\texttheta}ile/&$\\rightarrow$ & \\textit{7600} & /le{\texttheta}ile/&$\\rightarrow$ & \\textit{9500} & /le{\texttheta}il/& \\\\
\\end{tabular}

\\vspace{20pt}\\hline

\\end{nopagebreak}
\\filbreak



\\vspace{15pt}
\\begin{nopagebreak}
\\noindent{\\fontsize{20pt}{10pt}\\textbf{tod} } \\textit{forever/eternal/permanent} (N)\\\\
\\noindent {\\tovian \\fontsize{20pt}{10pt} \\textbf{tod} }\\\\
\\noindent /{\textsubbridge{t}}{\textprimstress}od/\\\\
\\noindent Syllables: tod\\\\


\\noindent History:

\\vspace{-0pt}
\\hspace{40pt}
\\begin{tabular}{ccc}
\\textit{0} & /{\textsubbridge{t}}oda/&$\\rightarrow$ & \\textit{9500} & /{\textsubbridge{t}}od/& \\\\
\\end{tabular}

\\vspace{20pt}\\hline

\\end{nopagebreak}
\\filbreak



\\vspace{15pt}
\\begin{nopagebreak}
\\noindent{\\fontsize{20pt}{10pt}\\textbf{tev} } \\textit{forty eight/48} (N)\\\\
\\noindent {\\tovian \\fontsize{20pt}{10pt} \\textbf{tev} }\\\\
\\noindent /{\textsubbridge{t}}{\textprimstress}ev/\\\\
\\noindent Syllables: tev\\\\
\\noindent lit. four+twelve\\\\


\\noindent History:

\\vspace{-0pt}
\\hspace{40pt}
\\begin{tabular}{ccc}
\\textit{4000} & /{\textsubbridge{t}}ev{\textbeltl}ya/&$\\rightarrow$ & \\textit{4502} & /{\textsubbridge{t}}evya/&$\\rightarrow$ & \\textit{9500} & /{\textsubbridge{t}}evy/&$\\rightarrow$ & \\textit{13001} & /{\textsubbridge{t}}ev/& \\\\
\\end{tabular}

\\vspace{20pt}\\hline

\\end{nopagebreak}
\\filbreak



\\vspace{15pt}
\\begin{nopagebreak}
\\noindent{\\fontsize{20pt}{10pt}\\textbf{tev} } \\textit{four/4} (N)\\\\
\\noindent {\\tovian \\fontsize{20pt}{10pt} \\textbf{tev} }\\\\
\\noindent /{\textsubbridge{t}}{\textprimstress}ev/\\\\
\\noindent Syllables: tev\\\\


\\noindent History:

\\vspace{-0pt}
\\hspace{40pt}
\\begin{tabular}{ccc}
\\textit{0} & /{\textsubbridge{t}}ev/& \\\\
\\end{tabular}

\\vspace{20pt}\\hline

\\end{nopagebreak}
\\filbreak



\\vspace{15pt}
\\begin{nopagebreak}
\\noindent{\\fontsize{20pt}{10pt}\\textbf{lheyalo} } \\textit{fourteen/14} (N)\\\\
\\noindent {\\tovian \\fontsize{20pt}{10pt} \\textbf{lheyalo} }\\\\
\\noindent /{\textbeltl}ey{\textprimstress}alo/\\\\
\\noindent Syllables: ɬej.al.o\\\\
\\noindent lit. twelve+two\\\\


\\noindent History:

\\vspace{-0pt}
\\hspace{40pt}
\\begin{tabular}{ccc}
\\textit{4000} & /{\textbeltl}yaloe/&$\\rightarrow$ & \\textit{9500} & /{\textbeltl}yalo/&$\\rightarrow$ & \\textit{15001} & /{\textbeltl}eyalo/& \\\\
\\end{tabular}

\\vspace{20pt}\\hline

\\end{nopagebreak}
\\filbreak



\\vspace{15pt}
\\begin{nopagebreak}
\\noindent{\\fontsize{20pt}{10pt}\\textbf{yelalh} } \\textit{friend/ally/friendship} (N)\\\\
\\noindent {\\tovian \\fontsize{20pt}{10pt} \\textbf{yelalh} }\\\\
\\noindent /y{\textprimstress}ela{\textbeltl}/\\\\
\\noindent Syllables: jel.aɬ\\\\
\\noindent lit. GEN+person+goodness\\\\


\\noindent History:

\\vspace{-0pt}
\\hspace{40pt}
\\begin{tabular}{ccc}
\\textit{0} & /hi-driela{\textbeltl}um/&$\\rightarrow$ & \\textit{1000} & /hidriela{\textbeltl}um/&$\\rightarrow$ & \\textit{2201} & /hidriyela{\textbeltl}um/&$\\rightarrow$ & \\textit{3200} & /hdryela{\textbeltl}m/&$\\rightarrow$ & \\textit{4500} & /hryela{\textbeltl}m/&$\\rightarrow$ & \\textit{7501} & /ryela{\textbeltl}m/&$\\rightarrow$ & \\textit{7600} & /yela{\textbeltl}m/&$\\rightarrow$ & \\textit{13001} & /yela{\textbeltl}/& \\\\
\\end{tabular}

\\vspace{20pt}\\hline

\\end{nopagebreak}
\\filbreak



\\vspace{15pt}
\\begin{nopagebreak}
\\noindent{\\fontsize{20pt}{10pt}\\textbf{heth} } \\textit{front} (N)\\\\
\\noindent {\\tovian \\fontsize{20pt}{10pt} \\textbf{heth} }\\\\
\\noindent /ħ{\textprimstress}e{\texttheta}/\\\\
\\noindent Syllables: ħeθ\\\\


\\noindent History:

\\vspace{-0pt}
\\hspace{40pt}
\\begin{tabular}{ccc}
\\textit{0} & /he{\texttheta}/&$\\rightarrow$ & \\textit{14006} & /ħe{\texttheta}/& \\\\
\\end{tabular}

\\vspace{20pt}\\hline

\\end{nopagebreak}
\\filbreak



\\vspace{15pt}
\\begin{nopagebreak}
\\noindent{\\fontsize{20pt}{10pt}\\textbf{laon} } \\textit{frown} (N)\\\\
\\noindent {\\tovian \\fontsize{20pt}{10pt} \\textbf{laon} }\\\\
\\noindent /l{\textprimstress}aon/\\\\
\\noindent Syllables: la.on\\\\
\\noindent lit. sad+mouth\\\\


\\noindent History:

\\vspace{-0pt}
\\hspace{40pt}
\\begin{tabular}{ccc}
\\textit{400} & /luraeolan/&$\\rightarrow$ & \\textit{3200} & /lraeoln/&$\\rightarrow$ & \\textit{6503} & /rlaeoln/&$\\rightarrow$ & \\textit{7600} & /laeon/&$\\rightarrow$ & \\textit{10000} & /laon/& \\\\
\\end{tabular}

\\vspace{20pt}\\hline

\\end{nopagebreak}
\\filbreak



\\vspace{15pt}
\\begin{nopagebreak}
\\noindent{\\fontsize{20pt}{10pt}\\textbf{shol} } \\textit{fruit} (N)\\\\
\\noindent {\\tovian \\fontsize{20pt}{10pt} \\textbf{shol} }\\\\
\\noindent /{\textesh}{\textprimstress}ol/\\\\
\\noindent Syllables: ʃol\\\\


\\noindent History:

\\vspace{-0pt}
\\hspace{40pt}
\\begin{tabular}{ccc}
\\textit{0} & /{\textyogh}ola/&$\\rightarrow$ & \\textit{9500} & /{\textyogh}ol/&$\\rightarrow$ & \\textit{12001} & /{\textesh}ol/& \\\\
\\end{tabular}

\\vspace{20pt}\\hline

\\end{nopagebreak}
\\filbreak



\\vspace{15pt}
\\begin{nopagebreak}
\\noindent{\\fontsize{20pt}{10pt}\\textbf{remhad} } \\textit{full} (N)\\\\
\\noindent {\\tovian \\fontsize{20pt}{10pt} \\textbf{remhad} }\\\\
\\noindent /r{\textprimstress}emħad/\\\\
\\noindent Syllables: remħ.ad\\\\
\\noindent lit. empty+not\\\\


\\noindent History:

\\vspace{-0pt}
\\hspace{40pt}
\\begin{tabular}{ccc}
\\textit{0} & /oneremihado/&$\\rightarrow$ & \\textit{3200} & /onremhado/&$\\rightarrow$ & \\textit{7500} & /nremhado/&$\\rightarrow$ & \\textit{7600} & /remhado/&$\\rightarrow$ & \\textit{9500} & /remhad/&$\\rightarrow$ & \\textit{14006} & /remħad/& \\\\
\\end{tabular}

\\vspace{20pt}\\hline

\\end{nopagebreak}
\\filbreak



\\vspace{15pt}
\\begin{nopagebreak}
\\noindent{\\fontsize{20pt}{10pt}\\textbf{heth} } \\textit{future} (Ns)\\\\
\\noindent {\\tovian \\fontsize{20pt}{10pt} \\textbf{heth} }\\\\
\\noindent /ħ{\textprimstress}e{\texttheta}/\\\\
\\noindent Syllables: ħeθ\\\\
\\noindent lit. front+path\\\\


\\noindent History:

\\vspace{-0pt}
\\hspace{40pt}
\\begin{tabular}{ccc}
\\textit{10} & /he{\texttheta}num/&$\\rightarrow$ & \\textit{3200} & /he{\texttheta}nm/&$\\rightarrow$ & \\textit{7600} & /he{\texttheta}m/&$\\rightarrow$ & \\textit{13001} & /he{\texttheta}/&$\\rightarrow$ & \\textit{14006} & /ħe{\texttheta}/& \\\\
\\end{tabular}

\\vspace{20pt}\\hline

\\end{nopagebreak}
\\filbreak



\\vspace{15pt}
\\begin{nopagebreak}
\\noindent{\\fontsize{20pt}{10pt}\\textbf{hi-} } \\textit{GEN/genitive case marker} (CASE)\\\\
\\noindent {\\tovian \\fontsize{20pt}{10pt} \\textbf{hi-} }\\\\
\\noindent /h{\textprimstress}i/\\\\
\\noindent Syllables: hi\\\\


\\noindent History:

\\vspace{-0pt}
\\hspace{40pt}
\\begin{tabular}{ccc}
\\textit{0} & /hi-/& \\\\
\\end{tabular}

\\vspace{20pt}\\hline

\\end{nopagebreak}
\\filbreak



\\vspace{15pt}
\\begin{nopagebreak}
\\noindent{\\fontsize{20pt}{10pt}\\textbf{misan} } \\textit{generosity} (N)\\\\
\\noindent {\\tovian \\fontsize{20pt}{10pt} \\textbf{misan} }\\\\
\\noindent /m{\textprimstress}isan/\\\\
\\noindent Syllables: mis.an\\\\
\\noindent lit. goodness+gift\\\\


\\noindent History:

\\vspace{-0pt}
\\hspace{40pt}
\\begin{tabular}{ccc}
\\textit{1300} & /{\textbeltl}umsalamin/&$\\rightarrow$ & \\textit{3200} & /{\textbeltl}mslamn/&$\\rightarrow$ & \\textit{6001} & /{\textbeltl}mislamn/&$\\rightarrow$ & \\textit{7600} & /{\textbeltl}mislan/&$\\rightarrow$ & \\textit{12500} & /{\textbeltl}misan/&$\\rightarrow$ & \\textit{13002} & /misan/& \\\\
\\end{tabular}

\\vspace{20pt}\\hline

\\end{nopagebreak}
\\filbreak



\\vspace{15pt}
\\begin{nopagebreak}
\\noindent{\\fontsize{20pt}{10pt}\\textbf{san} } \\textit{gift} (N)\\\\
\\noindent {\\tovian \\fontsize{20pt}{10pt} \\textbf{san} }\\\\
\\noindent /s{\textprimstress}an/\\\\
\\noindent Syllables: san\\\\


\\noindent History:

\\vspace{-0pt}
\\hspace{40pt}
\\begin{tabular}{ccc}
\\textit{0} & /salamin/&$\\rightarrow$ & \\textit{3200} & /slamn/&$\\rightarrow$ & \\textit{7600} & /slan/&$\\rightarrow$ & \\textit{12500} & /san/& \\\\
\\end{tabular}

\\vspace{20pt}\\hline

\\end{nopagebreak}
\\filbreak



\\vspace{15pt}
\\begin{nopagebreak}
\\noindent{\\fontsize{20pt}{10pt}\\textbf{elakorim} } \\textit{god of memory/cycle} (N)\\\\
\\noindent {\\tovian \\fontsize{20pt}{10pt} \\textbf{elakorim} }\\\\
\\noindent /elak{\textprimstress}orim/\\\\
\\noindent Syllables: el.ak.or.im\\\\


\\noindent History:

\\vspace{-0pt}
\\hspace{40pt}
\\begin{tabular}{ccc}
\\textit{0} & /elakorim/& \\\\
\\end{tabular}

\\vspace{20pt}\\hline

\\end{nopagebreak}
\\filbreak



\\vspace{15pt}
\\begin{nopagebreak}
\\noindent{\\fontsize{20pt}{10pt}\\textbf{lhum} } \\textit{goodness/virtue} (N)\\\\
\\noindent {\\tovian \\fontsize{20pt}{10pt} \\textbf{lhum} }\\\\
\\noindent /{\textbeltl}{\textprimstress}um/\\\\
\\noindent Syllables: ɬum\\\\


\\noindent History:

\\vspace{-0pt}
\\hspace{40pt}
\\begin{tabular}{ccc}
\\textit{0} & /{\textbeltl}um/& \\\\
\\end{tabular}

\\vspace{20pt}\\hline

\\end{nopagebreak}
\\filbreak



\\vspace{15pt}
\\begin{nopagebreak}
\\noindent{\\fontsize{20pt}{10pt}\\textbf{meloresyaleth} } \\textit{grain} (N)\\\\
\\noindent {\\tovian \\fontsize{20pt}{10pt} \\textbf{meloresyaleth} }\\\\
\\noindent /meloresy{\textprimstress}ale{\texttheta}/\\\\
\\noindent Syllables: mel.or.esj.al.eθ\\\\
\\noindent lit. plant+food\\\\


\\noindent History:

\\vspace{-0pt}
\\hspace{40pt}
\\begin{tabular}{ccc}
\\textit{6000} & /mlloresyale{\dh}a/&$\\rightarrow$ & \\textit{6500} & /mloresyale{\dh}a/&$\\rightarrow$ & \\textit{7502} & /meloresyale{\dh}a/&$\\rightarrow$ & \\textit{9500} & /meloresyale{\dh}/&$\\rightarrow$ & \\textit{12002} & /meloresyale{\texttheta}/& \\\\
\\end{tabular}

\\vspace{20pt}\\hline

\\end{nopagebreak}
\\filbreak



\\vspace{15pt}
\\begin{nopagebreak}
\\noindent{\\fontsize{20pt}{10pt}\\textbf{nanan} } \\textit{grandfather/grandmother/grandparent} (N)\\\\
\\noindent {\\tovian \\fontsize{20pt}{10pt} \\textbf{nanan} }\\\\
\\noindent /n{\textprimstress}anan/\\\\
\\noindent Syllables: nan.an\\\\
\\noindent lit. GEN+oldness+father\\\\


\\noindent History:

\\vspace{-0pt}
\\hspace{40pt}
\\begin{tabular}{ccc}
\\textit{642} & /hi-nanpander/&$\\rightarrow$ & \\textit{1000} & /hinanpander/&$\\rightarrow$ & \\textit{2000} & /hinanbander/&$\\rightarrow$ & \\textit{3200} & /hnanbandr/&$\\rightarrow$ & \\textit{6240} & /nanbandr/&$\\rightarrow$ & \\textit{8500} & /nanbanr/&$\\rightarrow$ & \\textit{8750} & /nananr/&$\\rightarrow$ & \\textit{13001} & /nanan/& \\\\
\\end{tabular}

\\vspace{20pt}\\hline

\\end{nopagebreak}
\\filbreak



\\vspace{15pt}
\\begin{nopagebreak}
\\noindent{\\fontsize{20pt}{10pt}\\textbf{lothr} } \\textit{grasp/grabbing} (N)\\\\
\\noindent {\\tovian \\fontsize{20pt}{10pt} \\textbf{lothr} }\\\\
\\noindent /l{\textprimstress}o{\texttheta}r/\\\\
\\noindent Syllables: loθr\\\\


\\noindent History:

\\vspace{-0pt}
\\hspace{40pt}
\\begin{tabular}{ccc}
\\textit{0} & /lo{\texttheta}ir/&$\\rightarrow$ & \\textit{3200} & /lo{\texttheta}r/& \\\\
\\end{tabular}

\\vspace{20pt}\\hline

\\end{nopagebreak}
\\filbreak



\\vspace{15pt}
\\begin{nopagebreak}
\\noindent{\\fontsize{20pt}{10pt}\\textbf{kelat} } \\textit{greenness/green} (N)\\\\
\\noindent {\\tovian \\fontsize{20pt}{10pt} \\textbf{kelat} }\\\\
\\noindent /k{\textprimstress}ela{\textsubbridge{t}}/\\\\
\\noindent Syllables: kel.at\\\\


\\noindent History:

\\vspace{-0pt}
\\hspace{40pt}
\\begin{tabular}{ccc}
\\textit{0} & /kla{\textsubbridge{t}}/&$\\rightarrow$ & \\textit{15001} & /kela{\textsubbridge{t}}/& \\\\
\\end{tabular}

\\vspace{20pt}\\hline

\\end{nopagebreak}
\\filbreak



\\vspace{15pt}
\\begin{nopagebreak}
\\noindent{\\fontsize{20pt}{10pt}\\textbf{amal} } \\textit{hair/person hair} (N)\\\\
\\noindent {\\tovian \\fontsize{20pt}{10pt} \\textbf{amal} }\\\\
\\noindent /{\textprimstress}amal/\\\\
\\noindent Syllables: am.al\\\\


\\noindent History:

\\vspace{-0pt}
\\hspace{40pt}
\\begin{tabular}{ccc}
\\textit{0} & /amalur/&$\\rightarrow$ & \\textit{3200} & /amalr/&$\\rightarrow$ & \\textit{6503} & /amarl/&$\\rightarrow$ & \\textit{7600} & /amal/& \\\\
\\end{tabular}

\\vspace{20pt}\\hline

\\end{nopagebreak}
\\filbreak



\\vspace{15pt}
\\begin{nopagebreak}
\\noindent{\\fontsize{20pt}{10pt}\\textbf{yeolath} } \\textit{half-truth/truth-shard/partly understood truth/selectively revealed truth} (N)\\\\
\\noindent {\\tovian \\fontsize{20pt}{10pt} \\textbf{yeolath} }\\\\
\\noindent /ye{\textprimstress}ola{\texttheta}/\\\\
\\noindent Syllables: je.ol.aθ\\\\
\\noindent lit. truth+shard\\\\
\\noindent \\textit{A truth partly understood or selectively revealed; the missing context may be unknown to the speaker or intentionally withheld.}\\\\


\\noindent History:

\\vspace{-0pt}
\\hspace{40pt}
\\begin{tabular}{ccc}
\\textit{0} & /arayeolkera{\texttheta}a/&$\\rightarrow$ & \\textit{2000} & /arayeolgera{\texttheta}a/&$\\rightarrow$ & \\textit{2100} & /aryeolgera{\texttheta}a/&$\\rightarrow$ & \\textit{3200} & /aryeolgra{\texttheta}a/&$\\rightarrow$ & \\textit{4501} & /aryeolra{\texttheta}a/&$\\rightarrow$ & \\textit{6503} & /aryeorla{\texttheta}a/&$\\rightarrow$ & \\textit{7500} & /ryeorla{\texttheta}a/&$\\rightarrow$ & \\textit{7600} & /yeola{\texttheta}a/&$\\rightarrow$ & \\textit{9500} & /yeola{\texttheta}/& \\\\
\\end{tabular}

\\vspace{20pt}\\hline

\\end{nopagebreak}
\\filbreak



\\vspace{15pt}
\\begin{nopagebreak}
\\noindent{\\fontsize{20pt}{10pt}\\textbf{thang} } \\textit{hand/leaf} (N)\\\\
\\noindent {\\tovian \\fontsize{20pt}{10pt} \\textbf{thang} }\\\\
\\noindent /{\texttheta}{\textprimstress}a{\ng}/\\\\
\\noindent Syllables: θaŋ\\\\


\\noindent History:

\\vspace{-0pt}
\\hspace{40pt}
\\begin{tabular}{ccc}
\\textit{0} & /a{\dh}ange/&$\\rightarrow$ & \\textit{6000} & /a{\dh}a{\ng}ge/&$\\rightarrow$ & \\textit{7500} & /{\dh}a{\ng}ge/&$\\rightarrow$ & \\textit{8750} & /{\dh}a{\ng}e/&$\\rightarrow$ & \\textit{9500} & /{\dh}a{\ng}/&$\\rightarrow$ & \\textit{12002} & /{\texttheta}a{\ng}/& \\\\
\\end{tabular}

\\vspace{20pt}\\hline

\\end{nopagebreak}
\\filbreak



\\vspace{15pt}
\\begin{nopagebreak}
\\noindent{\\fontsize{20pt}{10pt}\\textbf{nither} } \\textit{hanging} (N)\\\\
\\noindent {\\tovian \\fontsize{20pt}{10pt} \\textbf{nither} }\\\\
\\noindent /n{\textprimstress}i{\texttheta}er/\\\\
\\noindent Syllables: niθ.er\\\\
\\noindent lit. putting+up\\\\


\\noindent History:

\\vspace{-0pt}
\\hspace{40pt}
\\begin{tabular}{ccc}
\\textit{0} & /hemmani{\texttheta}era/&$\\rightarrow$ & \\textit{3200} & /hmmni{\texttheta}era/&$\\rightarrow$ & \\textit{6500} & /hmni{\texttheta}era/&$\\rightarrow$ & \\textit{7501} & /mni{\texttheta}era/&$\\rightarrow$ & \\textit{7600} & /ni{\texttheta}era/&$\\rightarrow$ & \\textit{9500} & /ni{\texttheta}er/& \\\\
\\end{tabular}

\\vspace{20pt}\\hline

\\end{nopagebreak}
\\filbreak



\\vspace{15pt}
\\begin{nopagebreak}
\\noindent{\\fontsize{20pt}{10pt}\\textbf{mat} } \\textit{happy} (N)\\\\
\\noindent {\\tovian \\fontsize{20pt}{10pt} \\textbf{mat} }\\\\
\\noindent /m{\textprimstress}a{\textsubbridge{t}}/\\\\
\\noindent Syllables: mat\\\\


\\noindent History:

\\vspace{-0pt}
\\hspace{40pt}
\\begin{tabular}{ccc}
\\textit{0} & /ma{\textsubbridge{t}}{\textsubbridge{t}}a/&$\\rightarrow$ & \\textit{3000} & /ma{\textsubbridge{t}}a/&$\\rightarrow$ & \\textit{9500} & /ma{\textsubbridge{t}}/& \\\\
\\end{tabular}

\\vspace{20pt}\\hline

\\end{nopagebreak}
\\filbreak



\\vspace{15pt}
\\begin{nopagebreak}
\\noindent{\\fontsize{20pt}{10pt}\\textbf{lha} } \\textit{hardness/strength} (N)\\\\
\\noindent {\\tovian \\fontsize{20pt}{10pt} \\textbf{lha} }\\\\
\\noindent /{\textbeltl}{\textprimstress}a/\\\\
\\noindent Syllables: ɬa\\\\


\\noindent History:

\\vspace{-0pt}
\\hspace{40pt}
\\begin{tabular}{ccc}
\\textit{0} & /{\textbeltl}aa/&$\\rightarrow$ & \\textit{5500} & /{\textbeltl}a/& \\\\
\\end{tabular}

\\vspace{20pt}\\hline

\\end{nopagebreak}
\\filbreak



\\vspace{15pt}
\\begin{nopagebreak}
\\noindent{\\fontsize{20pt}{10pt}\\textbf{lesh} } \\textit{harvest} (N)\\\\
\\noindent {\\tovian \\fontsize{20pt}{10pt} \\textbf{lesh} }\\\\
\\noindent /l{\textprimstress}e{\textesh}/\\\\
\\noindent Syllables: leʃ\\\\
\\noindent lit. plant+cutting\\\\


\\noindent History:

\\vspace{-0pt}
\\hspace{40pt}
\\begin{tabular}{ccc}
\\textit{1000} & /{\textschwa}mellore{\textyogh}i{\texttheta}/&$\\rightarrow$ & \\textit{2300} & /mellore{\textyogh}i{\texttheta}/&$\\rightarrow$ & \\textit{3200} & /mllre{\textyogh}{\texttheta}/&$\\rightarrow$ & \\textit{4502} & /mllre{\textyogh}/&$\\rightarrow$ & \\textit{6500} & /mlre{\textyogh}/&$\\rightarrow$ & \\textit{6503} & /mrle{\textyogh}/&$\\rightarrow$ & \\textit{7501} & /rle{\textyogh}/&$\\rightarrow$ & \\textit{7600} & /le{\textyogh}/&$\\rightarrow$ & \\textit{12001} & /le{\textesh}/& \\\\
\\end{tabular}

\\vspace{20pt}\\hline

\\end{nopagebreak}
\\filbreak



\\vspace{15pt}
\\begin{nopagebreak}
\\noindent{\\fontsize{20pt}{10pt}\\textbf{derunan} } \\textit{hat} (N)\\\\
\\noindent {\\tovian \\fontsize{20pt}{10pt} \\textbf{derunan} }\\\\
\\noindent /der{\textprimstress}unan/\\\\
\\noindent Syllables: der.un.an\\\\
\\noindent lit. head+covering\\\\


\\noindent History:

\\vspace{-0pt}
\\hspace{40pt}
\\begin{tabular}{ccc}
\\textit{3024} & /darunbenani/&$\\rightarrow$ & \\textit{3200} & /drunbnani/&$\\rightarrow$ & \\textit{8500} & /drunnani/&$\\rightarrow$ & \\textit{9500} & /drunnan/&$\\rightarrow$ & \\textit{12007} & /drunan/&$\\rightarrow$ & \\textit{15001} & /derunan/& \\\\
\\end{tabular}

\\vspace{20pt}\\hline

\\end{nopagebreak}
\\filbreak



\\vspace{15pt}
\\begin{nopagebreak}
\\noindent{\\fontsize{20pt}{10pt}\\textbf{ta} } \\textit{he/she/them} (P)\\\\
\\noindent {\\tovian \\fontsize{20pt}{10pt} \\textbf{ta} }\\\\
\\noindent /{\textsubbridge{t}}{\textprimstress}a/\\\\
\\noindent Syllables: ta\\\\


\\noindent History:

\\vspace{-0pt}
\\hspace{40pt}
\\begin{tabular}{ccc}
\\textit{0} & /{\textsubbridge{t}}a/& \\\\
\\end{tabular}

\\vspace{20pt}\\hline

\\end{nopagebreak}
\\filbreak



\\vspace{15pt}
\\begin{nopagebreak}
\\noindent{\\fontsize{20pt}{10pt}\\textbf{dan} } \\textit{head} (N)\\\\
\\noindent {\\tovian \\fontsize{20pt}{10pt} \\textbf{dan} }\\\\
\\noindent /d{\textprimstress}an/\\\\
\\noindent Syllables: dan\\\\


\\noindent History:

\\vspace{-0pt}
\\hspace{40pt}
\\begin{tabular}{ccc}
\\textit{0} & /darun/&$\\rightarrow$ & \\textit{3200} & /darn/&$\\rightarrow$ & \\textit{7600} & /dan/& \\\\
\\end{tabular}

\\vspace{20pt}\\hline

\\end{nopagebreak}
\\filbreak



\\vspace{15pt}
\\begin{nopagebreak}
\\noindent{\\fontsize{20pt}{10pt}\\textbf{yu-ta} } \\textit{her} (_)\\\\
\\noindent {\\tovian \\fontsize{20pt}{10pt} \\textbf{yu-ta} }\\\\
\\noindent /y{\textprimstress}u{\textsubbridge{t}}a/\\\\
\\noindent Syllables: ju.ta\\\\
\\noindent lit. ACC+she\\\\
\\noindent \\textit{P}\\\\


\\noindent History:

\\vspace{-0pt}
\\hspace{40pt}
\\begin{tabular}{ccc}
\\textit{0} & /yu-{\textsubbridge{t}}a/& \\\\
\\end{tabular}

\\vspace{20pt}\\hline

\\end{nopagebreak}
\\filbreak



\\vspace{15pt}
\\begin{nopagebreak}
\\noindent{\\fontsize{20pt}{10pt}\\textbf{dum} } \\textit{hiding} ()\\\\
\\noindent {\\tovian \\fontsize{20pt}{10pt} \\textbf{dum} }\\\\
\\noindent /d{\textprimstress}um/\\\\
\\noindent Syllables: dum\\\\
\\noindent lit. bbu\\\\


\\noindent History:

\\vspace{-0pt}
\\hspace{40pt}
\\begin{tabular}{ccc}
\\textit{220} & /bubu/&$\\rightarrow$ & \\textit{9500} & /bub/&$\\rightarrow$ & \\textit{11000} & /bum/&$\\rightarrow$ & \\textit{11001} & /dum/& \\\\
\\end{tabular}

\\vspace{20pt}\\hline

\\end{nopagebreak}
\\filbreak



\\vspace{15pt}
\\begin{nopagebreak}
\\noindent{\\fontsize{20pt}{10pt}\\textbf{yu-ta} } \\textit{him} (_)\\\\
\\noindent {\\tovian \\fontsize{20pt}{10pt} \\textbf{yu-ta} }\\\\
\\noindent /y{\textprimstress}u{\textsubbridge{t}}a/\\\\
\\noindent Syllables: ju.ta\\\\
\\noindent lit. ACC+he\\\\
\\noindent \\textit{P}\\\\


\\noindent History:

\\vspace{-0pt}
\\hspace{40pt}
\\begin{tabular}{ccc}
\\textit{0} & /yu-{\textsubbridge{t}}a/& \\\\
\\end{tabular}

\\vspace{20pt}\\hline

\\end{nopagebreak}
\\filbreak



\\vspace{15pt}
\\begin{nopagebreak}
\\noindent{\\fontsize{20pt}{10pt}\\textbf{hi-ta} } \\textit{his/her/their} (_)\\\\
\\noindent {\\tovian \\fontsize{20pt}{10pt} \\textbf{hi-ta} }\\\\
\\noindent /h{\textprimstress}i{\textsubbridge{t}}a/\\\\
\\noindent Syllables: hi.ta\\\\
\\noindent lit. GEN+she\\\\
\\noindent \\textit{P}\\\\


\\noindent History:

\\vspace{-0pt}
\\hspace{40pt}
\\begin{tabular}{ccc}
\\textit{0} & /hi-{\textsubbridge{t}}a/& \\\\
\\end{tabular}

\\vspace{20pt}\\hline

\\end{nopagebreak}
\\filbreak



\\vspace{15pt}
\\begin{nopagebreak}
\\noindent{\\fontsize{20pt}{10pt}\\textbf{leth} } \\textit{hoe} (N)\\\\
\\noindent {\\tovian \\fontsize{20pt}{10pt} \\textbf{leth} }\\\\
\\noindent /l{\textprimstress}e{\texttheta}/\\\\
\\noindent Syllables: leθ\\\\
\\noindent lit. plant+tool\\\\


\\noindent History:

\\vspace{-0pt}
\\hspace{40pt}
\\begin{tabular}{ccc}
\\textit{2000} & /{\textschwa}mellore{\dh}an/&$\\rightarrow$ & \\textit{2300} & /mellore{\dh}an/&$\\rightarrow$ & \\textit{3200} & /mllre{\dh}n/&$\\rightarrow$ & \\textit{6500} & /mlre{\dh}n/&$\\rightarrow$ & \\textit{6503} & /mrle{\dh}n/&$\\rightarrow$ & \\textit{7501} & /rle{\dh}n/&$\\rightarrow$ & \\textit{7600} & /le{\dh}n/&$\\rightarrow$ & \\textit{12002} & /le{\texttheta}n/&$\\rightarrow$ & \\textit{13001} & /le{\texttheta}/& \\\\
\\end{tabular}

\\vspace{20pt}\\hline

\\end{nopagebreak}
\\filbreak



\\vspace{15pt}
\\begin{nopagebreak}
\\noindent{\\fontsize{20pt}{10pt}\\textbf{ag} } \\textit{hotness/heat/warmth} (N)\\\\
\\noindent {\\tovian \\fontsize{20pt}{10pt} \\textbf{ag} }\\\\
\\noindent /{\textprimstress}ag/\\\\
\\noindent Syllables: ag\\\\


\\noindent History:

\\vspace{-0pt}
\\hspace{40pt}
\\begin{tabular}{ccc}
\\textit{0} & /asaka/&$\\rightarrow$ & \\textit{2000} & /asaga/&$\\rightarrow$ & \\textit{5100} & /aaga/&$\\rightarrow$ & \\textit{5500} & /aga/&$\\rightarrow$ & \\textit{9500} & /ag/& \\\\
\\end{tabular}

\\vspace{20pt}\\hline

\\end{nopagebreak}
\\filbreak



\\vspace{15pt}
\\begin{nopagebreak}
\\noindent{\\fontsize{20pt}{10pt}\\textbf{yethil} } \\textit{house/home} (N)\\\\
\\noindent {\\tovian \\fontsize{20pt}{10pt} \\textbf{yethil} }\\\\
\\noindent /y{\textprimstress}e{\texttheta}il/\\\\
\\noindent Syllables: jeθ.il\\\\
\\noindent lit. person+place\\\\


\\noindent History:

\\vspace{-0pt}
\\hspace{40pt}
\\begin{tabular}{ccc}
\\textit{10} & /driela{\texttheta}ile/&$\\rightarrow$ & \\textit{2201} & /driyela{\texttheta}ile/&$\\rightarrow$ & \\textit{3200} & /dryel{\texttheta}ile/&$\\rightarrow$ & \\textit{7501} & /ryel{\texttheta}ile/&$\\rightarrow$ & \\textit{7600} & /yel{\texttheta}ile/&$\\rightarrow$ & \\textit{9500} & /yel{\texttheta}il/&$\\rightarrow$ & \\textit{12502} & /ye{\texttheta}il/& \\\\
\\end{tabular}

\\vspace{20pt}\\hline

\\end{nopagebreak}
\\filbreak



\\vspace{15pt}
\\begin{nopagebreak}
\\noindent{\\fontsize{20pt}{10pt}\\textbf{si-lhan} } \\textit{how} (INT)\\\\
\\noindent {\\tovian \\fontsize{20pt}{10pt} \\textbf{si-lhan} }\\\\
\\noindent /s{\textprimstress}i{\textbeltl}an/\\\\
\\noindent Syllables: si.ɬan\\\\
\\noindent lit. INS+INT\\\\


\\noindent History:

\\vspace{-0pt}
\\hspace{40pt}
\\begin{tabular}{ccc}
\\textit{0} & /si-{\textbeltl}an/& \\\\
\\end{tabular}

\\vspace{20pt}\\hline

\\end{nopagebreak}
\\filbreak



\\vspace{15pt}
\\begin{nopagebreak}
\\noindent{\\fontsize{20pt}{10pt}\\textbf{shush} } \\textit{ice/frost} (N)\\\\
\\noindent {\\tovian \\fontsize{20pt}{10pt} \\textbf{shush} }\\\\
\\noindent /{\textesh}{\textprimstress}u{\textesh}/\\\\
\\noindent Syllables: ʃuʃ\\\\


\\noindent History:

\\vspace{-0pt}
\\hspace{40pt}
\\begin{tabular}{ccc}
\\textit{0} & /{\textesh}u{\textesh}a/&$\\rightarrow$ & \\textit{9500} & /{\textesh}u{\textesh}/& \\\\
\\end{tabular}

\\vspace{20pt}\\hline

\\end{nopagebreak}
\\filbreak



\\vspace{15pt}
\\begin{nopagebreak}
\\noindent{\\fontsize{20pt}{10pt}\\textbf{se-} } \\textit{IMP/imperative mood marker} (MOOD)\\\\
\\noindent {\\tovian \\fontsize{20pt}{10pt} \\textbf{se-} }\\\\
\\noindent /s{\textprimstress}e/\\\\
\\noindent Syllables: se\\\\


\\noindent History:

\\vspace{-0pt}
\\hspace{40pt}
\\begin{tabular}{ccc}
\\textit{0} & /se-/& \\\\
\\end{tabular}

\\vspace{20pt}\\hline

\\end{nopagebreak}
\\filbreak



\\vspace{15pt}
\\begin{nopagebreak}
\\noindent{\\fontsize{20pt}{10pt}\\textbf{gad} } \\textit{importance/grandness/greatness} (N)\\\\
\\noindent {\\tovian \\fontsize{20pt}{10pt} \\textbf{gad} }\\\\
\\noindent /g{\textprimstress}ad/\\\\
\\noindent Syllables: gad\\\\


\\noindent History:

\\vspace{-0pt}
\\hspace{40pt}
\\begin{tabular}{ccc}
\\textit{0} & /aga{\textsubbridge{t}}a/&$\\rightarrow$ & \\textit{2000} & /agada/&$\\rightarrow$ & \\textit{7500} & /gada/&$\\rightarrow$ & \\textit{9500} & /gad/& \\\\
\\end{tabular}

\\vspace{20pt}\\hline

\\end{nopagebreak}
\\filbreak



\\vspace{15pt}
\\begin{nopagebreak}
\\noindent{\\fontsize{20pt}{10pt}\\textbf{tid} } \\textit{in (obsolete)} (became locative "ti-")\\\\
\\noindent {\\tovian \\fontsize{20pt}{10pt} \\textbf{tid} }\\\\
\\noindent /{\textsubbridge{t}}{\textprimstress}id/\\\\
\\noindent Syllables: tid\\\\


\\noindent History:

\\vspace{-0pt}
\\hspace{40pt}
\\begin{tabular}{ccc}
\\textit{0} & /{\textsubbridge{t}}id/& \\\\
\\end{tabular}

\\vspace{20pt}\\hline

\\end{nopagebreak}
\\filbreak



\\vspace{15pt}
\\begin{nopagebreak}
\\noindent{\\fontsize{20pt}{10pt}\\textbf{lo-} } \\textit{INAN/inanimate noun class} (CLASS)\\\\
\\noindent {\\tovian \\fontsize{20pt}{10pt} \\textbf{lo-} }\\\\
\\noindent /l{\textprimstress}o/\\\\
\\noindent Syllables: lo\\\\


\\noindent History:

\\vspace{-0pt}
\\hspace{40pt}
\\begin{tabular}{ccc}
\\textit{0} & /lo-/& \\\\
\\end{tabular}

\\vspace{20pt}\\hline

\\end{nopagebreak}
\\filbreak



\\vspace{15pt}
\\begin{nopagebreak}
\\noindent{\\fontsize{20pt}{10pt}\\textbf{lo-} } \\textit{INANA/inanimate noun class} (CLASS)\\\\
\\noindent {\\tovian \\fontsize{20pt}{10pt} \\textbf{lo-} }\\\\
\\noindent /l{\textprimstress}o/\\\\
\\noindent Syllables: lo\\\\


\\noindent History:

\\vspace{-0pt}
\\hspace{40pt}
\\begin{tabular}{ccc}
\\textit{0} & /lo-/& \\\\
\\end{tabular}

\\vspace{20pt}\\hline

\\end{nopagebreak}
\\filbreak



\\vspace{15pt}
\\begin{nopagebreak}
\\noindent{\\fontsize{20pt}{10pt}\\textbf{e-} } \\textit{INDEF/indefinite marker} (ARTICLE)\\\\
\\noindent {\\tovian \\fontsize{20pt}{10pt} \\textbf{e-} }\\\\
\\noindent /{\textprimstress}e/\\\\
\\noindent Syllables: e\\\\


\\noindent History:

\\vspace{-0pt}
\\hspace{40pt}
\\begin{tabular}{ccc}
\\textit{0} & /e-/& \\\\
\\end{tabular}

\\vspace{20pt}\\hline

\\end{nopagebreak}
\\filbreak



\\vspace{15pt}
\\begin{nopagebreak}
\\noindent{\\fontsize{20pt}{10pt}\\textbf{lhanum} } \\textit{inheritance/inheriting/inherit/carrying forward a legacy} (N)\\\\
\\noindent {\\tovian \\fontsize{20pt}{10pt} \\textbf{lhanum} }\\\\
\\noindent /{\textbeltl}{\textprimstress}anum/\\\\
\\noindent Syllables: ɬan.um\\\\
\\noindent lit. carrying+past\\\\
\\noindent \\textit{Taking up and carrying forward what came before; applies to property traits duties traditions and ideas passed on from an earlier source.}\\\\


\\noindent History:

\\vspace{-0pt}
\\hspace{40pt}
\\begin{tabular}{ccc}
\\textit{500} & /{\textbeltl}ogahulnume/&$\\rightarrow$ & \\textit{3200} & /{\textbeltl}gahlnume/&$\\rightarrow$ & \\textit{4500} & /{\textbeltl}ahlnume/&$\\rightarrow$ & \\textit{7600} & /{\textbeltl}ahnume/&$\\rightarrow$ & \\textit{9500} & /{\textbeltl}ahnum/&$\\rightarrow$ & \\textit{13002} & /{\textbeltl}anum/& \\\\
\\end{tabular}

\\vspace{20pt}\\hline

\\end{nopagebreak}
\\filbreak



\\vspace{15pt}
\\begin{nopagebreak}
\\noindent{\\fontsize{20pt}{10pt}\\textbf{si-} } \\textit{INS/instrumental case marker} (CASE)\\\\
\\noindent {\\tovian \\fontsize{20pt}{10pt} \\textbf{si-} }\\\\
\\noindent /s{\textprimstress}i/\\\\
\\noindent Syllables: si\\\\


\\noindent History:

\\vspace{-0pt}
\\hspace{40pt}
\\begin{tabular}{ccc}
\\textit{0} & /si-/& \\\\
\\end{tabular}

\\vspace{20pt}\\hline

\\end{nopagebreak}
\\filbreak



\\vspace{15pt}
\\begin{nopagebreak}
\\noindent{\\fontsize{20pt}{10pt}\\textbf{teshan} } \\textit{instrument} (N)\\\\
\\noindent {\\tovian \\fontsize{20pt}{10pt} \\textbf{teshan} }\\\\
\\noindent /{\textsubbridge{t}}{\textprimstress}e{\textesh}an/\\\\
\\noindent Syllables: teʃ.an\\\\
\\noindent lit. sound+tool\\\\


\\noindent History:

\\vspace{-0pt}
\\hspace{40pt}
\\begin{tabular}{ccc}
\\textit{1500} & /{\textsubbridge{t}}{\textsubbridge{t}}e{\textyogh}ah{\dh}an/&$\\rightarrow$ & \\textit{3000} & /{\textsubbridge{t}}e{\textyogh}ah{\dh}an/&$\\rightarrow$ & \\textit{3200} & /{\textsubbridge{t}}{\textyogh}ah{\dh}n/&$\\rightarrow$ & \\textit{4502} & /{\textsubbridge{t}}{\textyogh}ahn/&$\\rightarrow$ & \\textit{6240} & /{\textsubbridge{t}}{\textyogh}an/&$\\rightarrow$ & \\textit{12001} & /{\textsubbridge{t}}{\textesh}an/&$\\rightarrow$ & \\textit{12400} & /{\textsubbridge{t}}e{\textesh}an/& \\\\
\\end{tabular}

\\vspace{20pt}\\hline

\\end{nopagebreak}
\\filbreak



\\vspace{15pt}
\\begin{nopagebreak}
\\noindent{\\fontsize{20pt}{10pt}\\textbf{lod} } \\textit{it/it (inanimate)} (P)\\\\
\\noindent {\\tovian \\fontsize{20pt}{10pt} \\textbf{lod} }\\\\
\\noindent /l{\textprimstress}od/\\\\
\\noindent Syllables: lod\\\\


\\noindent History:

\\vspace{-0pt}
\\hspace{40pt}
\\begin{tabular}{ccc}
\\textit{0} & /lo{\textsubbridge{t}}a/&$\\rightarrow$ & \\textit{2000} & /loda/&$\\rightarrow$ & \\textit{9500} & /lod/& \\\\
\\end{tabular}

\\vspace{20pt}\\hline

\\end{nopagebreak}
\\filbreak



\\vspace{15pt}
\\begin{nopagebreak}
\\noindent{\\fontsize{20pt}{10pt}\\textbf{hi-lota} } \\textit{its} (_)\\\\
\\noindent {\\tovian \\fontsize{20pt}{10pt} \\textbf{hi-lota} }\\\\
\\noindent /hil{\textprimstress}o{\textsubbridge{t}}a/\\\\
\\noindent Syllables: hi.lot.a\\\\
\\noindent lit. GEN+it\\\\
\\noindent \\textit{P}\\\\


\\noindent History:

\\vspace{-0pt}
\\hspace{40pt}
\\begin{tabular}{ccc}
\\textit{0} & /hi-lo{\textsubbridge{t}}a/& \\\\
\\end{tabular}

\\vspace{20pt}\\hline

\\end{nopagebreak}
\\filbreak



\\vspace{15pt}
\\begin{nopagebreak}
\\noindent{\\fontsize{20pt}{10pt}\\textbf{selushu} } \\textit{January} (N)\\\\
\\noindent {\\tovian \\fontsize{20pt}{10pt} \\textbf{selushu} }\\\\
\\noindent /sel{\textprimstress}u{\textesh}u/\\\\
\\noindent Syllables: sel.uʃ.u\\\\
\\noindent lit. water+cold+COM\\\\
\\noindent \\textit{Frohsten}\\\\


\\noindent History:

\\vspace{-0pt}
\\hspace{40pt}
\\begin{tabular}{ccc}
\\textit{6000} & /zel{\textsubbridge{t}}u{\textesh}uyi-/&$\\rightarrow$ & \\textit{6200} & /zel{\textsubbridge{t}}u{\textesh}ui-/&$\\rightarrow$ & \\textit{8500} & /zel{\textsubbridge{t}}u{\textesh}ui/&$\\rightarrow$ & \\textit{8750} & /zelu{\textesh}ui/&$\\rightarrow$ & \\textit{9500} & /zelu{\textesh}u/&$\\rightarrow$ & \\textit{12000} & /selu{\textesh}u/& \\\\
\\end{tabular}

\\vspace{20pt}\\hline

\\end{nopagebreak}
\\filbreak



\\vspace{15pt}
\\begin{nopagebreak}
\\noindent{\\fontsize{20pt}{10pt}\\textbf{lad} } \\textit{jealousy} (N)\\\\
\\noindent {\\tovian \\fontsize{20pt}{10pt} \\textbf{lad} }\\\\
\\noindent /l{\textprimstress}ad/\\\\
\\noindent Syllables: lad\\\\
\\noindent lit. wickedness+comparison\\\\


\\noindent History:

\\vspace{-0pt}
\\hspace{40pt}
\\begin{tabular}{ccc}
\\textit{402} & /oyalroladil/&$\\rightarrow$ & \\textit{2100} & /yalroladil/&$\\rightarrow$ & \\textit{3200} & /ylrladl/&$\\rightarrow$ & \\textit{6503} & /yrlladl/&$\\rightarrow$ & \\textit{7501} & /rlladl/&$\\rightarrow$ & \\textit{7600} & /lladl/&$\\rightarrow$ & \\textit{12007} & /ladl/&$\\rightarrow$ & \\textit{13001} & /lad/& \\\\
\\end{tabular}

\\vspace{20pt}\\hline

\\end{nopagebreak}
\\filbreak



\\vspace{15pt}
\\begin{nopagebreak}
\\noindent{\\fontsize{20pt}{10pt}\\textbf{tlun} } \\textit{joining} (N)\\\\
\\noindent {\\tovian \\fontsize{20pt}{10pt} \\textbf{tlun} }\\\\
\\noindent /{\texttoptiebar{t\textbeltl}}{\textprimstress}un/\\\\
\\noindent Syllables: ƛun\\\\


\\noindent History:

\\vspace{-0pt}
\\hspace{40pt}
\\begin{tabular}{ccc}
\\textit{0} & /{\texttoptiebar{t\textbeltl}}una/&$\\rightarrow$ & \\textit{9500} & /{\texttoptiebar{t\textbeltl}}un/& \\\\
\\end{tabular}

\\vspace{20pt}\\hline

\\end{nopagebreak}
\\filbreak



\\vspace{15pt}
\\begin{nopagebreak}
\\noindent{\\fontsize{20pt}{10pt}\\textbf{renef} } \\textit{July} (N)\\\\
\\noindent {\\tovian \\fontsize{20pt}{10pt} \\textbf{renef} }\\\\
\\noindent /r{\textprimstress}enef/\\\\
\\noindent Syllables: ren.ef\\\\
\\noindent lit. standing+sun\\\\
\\noindent \\textit{Thermiblase}\\\\


\\noindent History:

\\vspace{-0pt}
\\hspace{40pt}
\\begin{tabular}{ccc}
\\textit{6000} & /remnefl/&$\\rightarrow$ & \\textit{7600} & /renefl/&$\\rightarrow$ & \\textit{12500} & /renef/& \\\\
\\end{tabular}

\\vspace{20pt}\\hline

\\end{nopagebreak}
\\filbreak



\\vspace{15pt}
\\begin{nopagebreak}
\\noindent{\\fontsize{20pt}{10pt}\\textbf{romef} } \\textit{June} (N)\\\\
\\noindent {\\tovian \\fontsize{20pt}{10pt} \\textbf{romef} }\\\\
\\noindent /r{\textprimstress}omef/\\\\
\\noindent Syllables: rom.ef\\\\
\\noindent lit. sitting+sun\\\\
\\noindent \\textit{Thermiblase}\\\\


\\noindent History:

\\vspace{-0pt}
\\hspace{40pt}
\\begin{tabular}{ccc}
\\textit{6000} & /rom{\textsubbridge{t}}efl/&$\\rightarrow$ & \\textit{8750} & /romefl/&$\\rightarrow$ & \\textit{12500} & /romef/& \\\\
\\end{tabular}

\\vspace{20pt}\\hline

\\end{nopagebreak}
\\filbreak



\\vspace{15pt}
\\begin{nopagebreak}
\\noindent{\\fontsize{20pt}{10pt}\\textbf{vek} } \\textit{killing} (N)\\\\
\\noindent {\\tovian \\fontsize{20pt}{10pt} \\textbf{vek} }\\\\
\\noindent /v{\textprimstress}ek/\\\\
\\noindent Syllables: vek\\\\


\\noindent History:

\\vspace{-0pt}
\\hspace{40pt}
\\begin{tabular}{ccc}
\\textit{0} & /vekh/&$\\rightarrow$ & \\textit{5000} & /vek/& \\\\
\\end{tabular}

\\vspace{20pt}\\hline

\\end{nopagebreak}
\\filbreak



\\vspace{15pt}
\\begin{nopagebreak}
\\noindent{\\fontsize{20pt}{10pt}\\textbf{rish} } \\textit{kindness} (N)\\\\
\\noindent {\\tovian \\fontsize{20pt}{10pt} \\textbf{rish} }\\\\
\\noindent /r{\textprimstress}i{\textesh}/\\\\
\\noindent Syllables: riʃ\\\\


\\noindent History:

\\vspace{-0pt}
\\hspace{40pt}
\\begin{tabular}{ccc}
\\textit{0} & /rai{\textesh}el/&$\\rightarrow$ & \\textit{3200} & /rai{\textesh}l/&$\\rightarrow$ & \\textit{5500} & /ri{\textesh}l/&$\\rightarrow$ & \\textit{12500} & /ri{\textesh}/& \\\\
\\end{tabular}

\\vspace{20pt}\\hline

\\end{nopagebreak}
\\filbreak



\\vspace{15pt}
\\begin{nopagebreak}
\\noindent{\\fontsize{20pt}{10pt}\\textbf{nashith} } \\textit{knife} (N)\\\\
\\noindent {\\tovian \\fontsize{20pt}{10pt} \\textbf{nashith} }\\\\
\\noindent /n{\textprimstress}a{\textesh}i{\texttheta}/\\\\
\\noindent Syllables: naʃ.iθ\\\\
\\noindent lit. small+sharpness+tool\\\\


\\noindent History:

\\vspace{-0pt}
\\hspace{40pt}
\\begin{tabular}{ccc}
\\textit{50} & /nina{\textyogh}i{\texttheta}{\dh}an/&$\\rightarrow$ & \\textit{3200} & /nna{\textyogh}i{\texttheta}{\dh}n/&$\\rightarrow$ & \\textit{4502} & /nna{\textyogh}i{\texttheta}n/&$\\rightarrow$ & \\textit{6500} & /na{\textyogh}i{\texttheta}n/&$\\rightarrow$ & \\textit{12001} & /na{\textesh}i{\texttheta}n/&$\\rightarrow$ & \\textit{13001} & /na{\textesh}i{\texttheta}/& \\\\
\\end{tabular}

\\vspace{20pt}\\hline

\\end{nopagebreak}
\\filbreak



\\vspace{15pt}
\\begin{nopagebreak}
\\noindent{\\fontsize{20pt}{10pt}\\textbf{lhar} } \\textit{knowledge} (N)\\\\
\\noindent {\\tovian \\fontsize{20pt}{10pt} \\textbf{lhar} }\\\\
\\noindent /{\textbeltl}{\textprimstress}ar/\\\\
\\noindent Syllables: ɬar\\\\


\\noindent History:

\\vspace{-0pt}
\\hspace{40pt}
\\begin{tabular}{ccc}
\\textit{0} & /{\textbeltl}ara/&$\\rightarrow$ & \\textit{9500} & /{\textbeltl}ar/& \\\\
\\end{tabular}

\\vspace{20pt}\\hline

\\end{nopagebreak}
\\filbreak



\\vspace{15pt}
\\begin{nopagebreak}
\\noindent{\\fontsize{20pt}{10pt}\\textbf{tel} } \\textit{lake} (N)\\\\
\\noindent {\\tovian \\fontsize{20pt}{10pt} \\textbf{tel} }\\\\
\\noindent /{\textsubbridge{t}}{\textprimstress}el/\\\\
\\noindent Syllables: tel\\\\


\\noindent History:

\\vspace{-0pt}
\\hspace{40pt}
\\begin{tabular}{ccc}
\\textit{0} & /o{\textsubbridge{t}}kele/&$\\rightarrow$ & \\textit{3000} & /o{\textsubbridge{t}}ele/&$\\rightarrow$ & \\textit{7500} & /{\textsubbridge{t}}ele/&$\\rightarrow$ & \\textit{9500} & /{\textsubbridge{t}}el/& \\\\
\\end{tabular}

\\vspace{20pt}\\hline

\\end{nopagebreak}
\\filbreak



\\vspace{15pt}
\\begin{nopagebreak}
\\noindent{\\fontsize{20pt}{10pt}\\textbf{fal} } \\textit{leader} (N)\\\\
\\noindent {\\tovian \\fontsize{20pt}{10pt} \\textbf{fal} }\\\\
\\noindent /f{\textprimstress}al/\\\\
\\noindent Syllables: fal\\\\


\\noindent History:

\\vspace{-0pt}
\\hspace{40pt}
\\begin{tabular}{ccc}
\\textit{0} & /falor/&$\\rightarrow$ & \\textit{3200} & /falr/&$\\rightarrow$ & \\textit{6503} & /farl/&$\\rightarrow$ & \\textit{7600} & /fal/& \\\\
\\end{tabular}

\\vspace{20pt}\\hline

\\end{nopagebreak}
\\filbreak



\\vspace{15pt}
\\begin{nopagebreak}
\\noindent{\\fontsize{20pt}{10pt}\\textbf{hed} } \\textit{left} (N)\\\\
\\noindent {\\tovian \\fontsize{20pt}{10pt} \\textbf{hed} }\\\\
\\noindent /ħ{\textprimstress}ed/\\\\
\\noindent Syllables: ħed\\\\


\\noindent History:

\\vspace{-0pt}
\\hspace{40pt}
\\begin{tabular}{ccc}
\\textit{0} & /he{\textsubbridge{t}}a/&$\\rightarrow$ & \\textit{2000} & /heda/&$\\rightarrow$ & \\textit{9500} & /hed/&$\\rightarrow$ & \\textit{14006} & /ħed/& \\\\
\\end{tabular}

\\vspace{20pt}\\hline

\\end{nopagebreak}
\\filbreak



\\vspace{15pt}
\\begin{nopagebreak}
\\noindent{\\fontsize{20pt}{10pt}\\textbf{kel} } \\textit{leg/limb/branch/arm} (N)\\\\
\\noindent {\\tovian \\fontsize{20pt}{10pt} \\textbf{kel} }\\\\
\\noindent /k{\textprimstress}el/\\\\
\\noindent Syllables: kel\\\\


\\noindent History:

\\vspace{-0pt}
\\hspace{40pt}
\\begin{tabular}{ccc}
\\textit{0} & /kelu/&$\\rightarrow$ & \\textit{9500} & /kel/& \\\\
\\end{tabular}

\\vspace{20pt}\\hline

\\end{nopagebreak}
\\filbreak



\\vspace{15pt}
\\begin{nopagebreak}
\\noindent{\\fontsize{20pt}{10pt}\\textbf{mel} } \\textit{life} (N)\\\\
\\noindent {\\tovian \\fontsize{20pt}{10pt} \\textbf{mel} }\\\\
\\noindent /m{\textprimstress}el/\\\\
\\noindent Syllables: mel\\\\


\\noindent History:

\\vspace{-0pt}
\\hspace{40pt}
\\begin{tabular}{ccc}
\\textit{0} & /{\textschwa}mel/&$\\rightarrow$ & \\textit{2300} & /mel/& \\\\
\\end{tabular}

\\vspace{20pt}\\hline

\\end{nopagebreak}
\\filbreak



\\vspace{15pt}
\\begin{nopagebreak}
\\noindent{\\fontsize{20pt}{10pt}\\textbf{an} } \\textit{light} (N)\\\\
\\noindent {\\tovian \\fontsize{20pt}{10pt} \\textbf{an} }\\\\
\\noindent /{\textprimstress}an/\\\\
\\noindent Syllables: an\\\\


\\noindent History:

\\vspace{-0pt}
\\hspace{40pt}
\\begin{tabular}{ccc}
\\textit{0} & /ahana/&$\\rightarrow$ & \\textit{5000} & /aana/&$\\rightarrow$ & \\textit{5500} & /ana/&$\\rightarrow$ & \\textit{9500} & /an/& \\\\
\\end{tabular}

\\vspace{20pt}\\hline

\\end{nopagebreak}
\\filbreak



\\vspace{15pt}
\\begin{nopagebreak}
\\noindent{\\fontsize{20pt}{10pt}\\textbf{natlatl} } \\textit{lightning} (N)\\\\
\\noindent {\\tovian \\fontsize{20pt}{10pt} \\textbf{natlatl} }\\\\
\\noindent /n{\textprimstress}a{\texttoptiebar{t\textbeltl}}a{\texttoptiebar{t\textbeltl}}/\\\\
\\noindent Syllables: naƛ.aƛ\\\\
\\noindent lit. light+falling\\\\


\\noindent History:

\\vspace{-0pt}
\\hspace{40pt}
\\begin{tabular}{ccc}
\\textit{0} & /ahana{\texttoptiebar{t\textbeltl}}a{\texttoptiebar{t\textbeltl}}a/&$\\rightarrow$ & \\textit{3200} & /ahna{\texttoptiebar{t\textbeltl}}a{\texttoptiebar{t\textbeltl}}a/&$\\rightarrow$ & \\textit{6240} & /ana{\texttoptiebar{t\textbeltl}}a{\texttoptiebar{t\textbeltl}}a/&$\\rightarrow$ & \\textit{7500} & /na{\texttoptiebar{t\textbeltl}}a{\texttoptiebar{t\textbeltl}}a/&$\\rightarrow$ & \\textit{9500} & /na{\texttoptiebar{t\textbeltl}}a{\texttoptiebar{t\textbeltl}}/& \\\\
\\end{tabular}

\\vspace{20pt}\\hline

\\end{nopagebreak}
\\filbreak



\\vspace{15pt}
\\begin{nopagebreak}
\\noindent{\\fontsize{20pt}{10pt}\\textbf{ti-} } \\textit{LOC/locative case marker} (CASE)\\\\
\\noindent {\\tovian \\fontsize{20pt}{10pt} \\textbf{ti-} }\\\\
\\noindent /{\textsubbridge{t}}{\textprimstress}i/\\\\
\\noindent Syllables: ti\\\\


\\noindent History:

\\vspace{-0pt}
\\hspace{40pt}
\\begin{tabular}{ccc}
\\textit{0} & /{\textsubbridge{t}}i-/& \\\\
\\end{tabular}

\\vspace{20pt}\\hline

\\end{nopagebreak}
\\filbreak



\\vspace{15pt}
\\begin{nopagebreak}
\\noindent{\\fontsize{20pt}{10pt}\\textbf{talh} } \\textit{love} (N)\\\\
\\noindent {\\tovian \\fontsize{20pt}{10pt} \\textbf{talh} }\\\\
\\noindent /{\textsubbridge{t}}{\textprimstress}a{\textbeltl}/\\\\
\\noindent Syllables: taɬ\\\\


\\noindent History:

\\vspace{-0pt}
\\hspace{40pt}
\\begin{tabular}{ccc}
\\textit{0} & /{\textsubbridge{t}}a{\textbeltl}av/&$\\rightarrow$ & \\textit{3200} & /{\textsubbridge{t}}a{\textbeltl}v/&$\\rightarrow$ & \\textit{4502} & /{\textsubbridge{t}}a{\textbeltl}/& \\\\
\\end{tabular}

\\vspace{20pt}\\hline

\\end{nopagebreak}
\\filbreak



\\vspace{15pt}
\\begin{nopagebreak}
\\noindent{\\fontsize{20pt}{10pt}\\textbf{sil} } \\textit{make/create} (N)\\\\
\\noindent {\\tovian \\fontsize{20pt}{10pt} \\textbf{sil} }\\\\
\\noindent /s{\textprimstress}il/\\\\
\\noindent Syllables: sil\\\\


\\noindent History:

\\vspace{-0pt}
\\hspace{40pt}
\\begin{tabular}{ccc}
\\textit{0} & /silahe{\textsubbridge{t}}/&$\\rightarrow$ & \\textit{1000} & /silah{\textsubbridge{t}}/&$\\rightarrow$ & \\textit{3000} & /silah/&$\\rightarrow$ & \\textit{3200} & /silh/&$\\rightarrow$ & \\textit{5000} & /sil/& \\\\
\\end{tabular}

\\vspace{20pt}\\hline

\\end{nopagebreak}
\\filbreak



\\vspace{15pt}
\\begin{nopagebreak}
\\noindent{\\fontsize{20pt}{10pt}\\textbf{tlasel} } \\textit{March} (N)\\\\
\\noindent {\\tovian \\fontsize{20pt}{10pt} \\textbf{tlasel} }\\\\
\\noindent /{\texttoptiebar{t\textbeltl}}{\textprimstress}asel/\\\\
\\noindent Syllables: ƛas.el\\\\
\\noindent lit. falling+water\\\\
\\noindent \\textit{Vindrem}\\\\


\\noindent History:

\\vspace{-0pt}
\\hspace{40pt}
\\begin{tabular}{ccc}
\\textit{6000} & /{\texttoptiebar{t\textbeltl}}a{\texttoptiebar{t\textbeltl}}azel/&$\\rightarrow$ & \\textit{9200} & /{\texttoptiebar{t\textbeltl}}{\textschwa}{\texttoptiebar{t\textbeltl}}azel/&$\\rightarrow$ & \\textit{11990} & /{\texttoptiebar{t\textbeltl}}{\texttoptiebar{t\textbeltl}}azel/&$\\rightarrow$ & \\textit{12000} & /{\texttoptiebar{t\textbeltl}}{\texttoptiebar{t\textbeltl}}asel/&$\\rightarrow$ & \\textit{12007} & /{\texttoptiebar{t\textbeltl}}asel/& \\\\
\\end{tabular}

\\vspace{20pt}\\hline

\\end{nopagebreak}
\\filbreak



\\vspace{15pt}
\\begin{nopagebreak}
\\noindent{\\fontsize{20pt}{10pt}\\textbf{felor} } \\textit{May} (N)\\\\
\\noindent {\\tovian \\fontsize{20pt}{10pt} \\textbf{felor} }\\\\
\\noindent /f{\textprimstress}elor/\\\\
\\noindent Syllables: fel.or\\\\
\\noindent lit. three+plant\\\\
\\noindent \\textit{Luviosh}\\\\


\\noindent History:

\\vspace{-0pt}
\\hspace{40pt}
\\begin{tabular}{ccc}
\\textit{6000} & /femmllore/&$\\rightarrow$ & \\textit{6500} & /femlore/&$\\rightarrow$ & \\textit{7600} & /felore/&$\\rightarrow$ & \\textit{9500} & /felor/& \\\\
\\end{tabular}

\\vspace{20pt}\\hline

\\end{nopagebreak}
\\filbreak



\\vspace{15pt}
\\begin{nopagebreak}
\\noindent{\\fontsize{20pt}{10pt}\\textbf{na} } \\textit{me/I} (P)\\\\
\\noindent {\\tovian \\fontsize{20pt}{10pt} \\textbf{na} }\\\\
\\noindent /n{\textprimstress}a/\\\\
\\noindent Syllables: na\\\\


\\noindent History:

\\vspace{-0pt}
\\hspace{40pt}
\\begin{tabular}{ccc}
\\textit{0} & /na/& \\\\
\\end{tabular}

\\vspace{20pt}\\hline

\\end{nopagebreak}
\\filbreak



\\vspace{15pt}
\\begin{nopagebreak}
\\noindent{\\fontsize{20pt}{10pt}\\textbf{athir} } \\textit{meat} (N)\\\\
\\noindent {\\tovian \\fontsize{20pt}{10pt} \\textbf{athir} }\\\\
\\noindent /{\textprimstress}a{\texttheta}ir/\\\\
\\noindent Syllables: aθ.ir\\\\
\\noindent lit. tasting+GEN+blood\\\\


\\noindent History:

\\vspace{-0pt}
\\hspace{40pt}
\\begin{tabular}{ccc}
\\textit{1000} & /ale{\dh}ahi-driks/&$\\rightarrow$ & \\textit{1000} & /ale{\dh}ahidriks/&$\\rightarrow$ & \\textit{3200} & /al{\dh}ahidrks/&$\\rightarrow$ & \\textit{4501} & /al{\dh}ahidrs/&$\\rightarrow$ & \\textit{5000} & /al{\dh}aidrs/&$\\rightarrow$ & \\textit{5500} & /al{\dh}idrs/&$\\rightarrow$ & \\textit{12002} & /al{\texttheta}idrs/&$\\rightarrow$ & \\textit{12502} & /a{\texttheta}idrs/&$\\rightarrow$ & \\textit{13001} & /a{\texttheta}idr/&$\\rightarrow$ & \\textit{13004} & /a{\texttheta}ir/& \\\\
\\end{tabular}

\\vspace{20pt}\\hline

\\end{nopagebreak}
\\filbreak



\\vspace{15pt}
\\begin{nopagebreak}
\\noindent{\\fontsize{20pt}{10pt}\\textbf{elom} } \\textit{memory} (N)\\\\
\\noindent {\\tovian \\fontsize{20pt}{10pt} \\textbf{elom} }\\\\
\\noindent /{\textprimstress}elom/\\\\
\\noindent Syllables: el.om\\\\


\\noindent History:

\\vspace{-0pt}
\\hspace{40pt}
\\begin{tabular}{ccc}
\\textit{0} & /elakorim/&$\\rightarrow$ & \\textit{2000} & /elagorim/&$\\rightarrow$ & \\textit{3200} & /elgorm/&$\\rightarrow$ & \\textit{4501} & /elorm/&$\\rightarrow$ & \\textit{7600} & /elom/& \\\\
\\end{tabular}

\\vspace{20pt}\\hline

\\end{nopagebreak}
\\filbreak



\\vspace{15pt}
\\begin{nopagebreak}
\\noindent{\\fontsize{20pt}{10pt}\\textbf{korim} } \\textit{memory-soul} (N)\\\\
\\noindent {\\tovian \\fontsize{20pt}{10pt} \\textbf{korim} }\\\\
\\noindent /k{\textprimstress}orim/\\\\
\\noindent Syllables: kor.im\\\\


\\noindent History:

\\vspace{-0pt}
\\hspace{40pt}
\\begin{tabular}{ccc}
\\textit{0} & /korim/& \\\\
\\end{tabular}

\\vspace{20pt}\\hline

\\end{nopagebreak}
\\filbreak



\\vspace{15pt}
\\begin{nopagebreak}
\\noindent{\\fontsize{20pt}{10pt}\\textbf{men} } \\textit{metal} (N)\\\\
\\noindent {\\tovian \\fontsize{20pt}{10pt} \\textbf{men} }\\\\
\\noindent /m{\textprimstress}en/\\\\
\\noindent Syllables: men\\\\
\\noindent lit. tool+GEN+rock\\\\


\\noindent History:

\\vspace{-0pt}
\\hspace{40pt}
\\begin{tabular}{ccc}
\\textit{1000} & /{\dh}anhi-meron/&$\\rightarrow$ & \\textit{1000} & /{\dh}anhimeron/&$\\rightarrow$ & \\textit{3200} & /{\dh}nhmern/&$\\rightarrow$ & \\textit{6400} & /{\dh}hmern/&$\\rightarrow$ & \\textit{7501} & /hmern/&$\\rightarrow$ & \\textit{7600} & /hmen/&$\\rightarrow$ & \\textit{13002} & /men/& \\\\
\\end{tabular}

\\vspace{20pt}\\hline

\\end{nopagebreak}
\\filbreak



\\vspace{15pt}
\\begin{nopagebreak}
\\noindent{\\fontsize{20pt}{10pt}\\textbf{mo-} } \\textit{MID/middle voice marker} (VOICE)\\\\
\\noindent {\\tovian \\fontsize{20pt}{10pt} \\textbf{mo-} }\\\\
\\noindent /m{\textprimstress}o/\\\\
\\noindent Syllables: mo\\\\


\\noindent History:

\\vspace{-0pt}
\\hspace{40pt}
\\begin{tabular}{ccc}
\\textit{0} & /mo-/& \\\\
\\end{tabular}

\\vspace{20pt}\\hline

\\end{nopagebreak}
\\filbreak



\\vspace{15pt}
\\begin{nopagebreak}
\\noindent{\\fontsize{20pt}{10pt}\\textbf{hi-na} } \\textit{mine} (_)\\\\
\\noindent {\\tovian \\fontsize{20pt}{10pt} \\textbf{hi-na} }\\\\
\\noindent /h{\textprimstress}ina/\\\\
\\noindent Syllables: hi.na\\\\
\\noindent lit. GEN+I\\\\
\\noindent \\textit{P}\\\\


\\noindent History:

\\vspace{-0pt}
\\hspace{40pt}
\\begin{tabular}{ccc}
\\textit{0} & /hi-na/& \\\\
\\end{tabular}

\\vspace{20pt}\\hline

\\end{nopagebreak}
\\filbreak



\\vspace{15pt}
\\begin{nopagebreak}
\\noindent{\\fontsize{20pt}{10pt}\\textbf{tlurinin} } \\textit{mist} (N)\\\\
\\noindent {\\tovian \\fontsize{20pt}{10pt} \\textbf{tlurinin} }\\\\
\\noindent /{\texttoptiebar{t\textbeltl}}ur{\textprimstress}inin/\\\\
\\noindent Syllables: ƛur.in.in\\\\
\\noindent lit. fog+small\\\\


\\noindent History:

\\vspace{-0pt}
\\hspace{40pt}
\\begin{tabular}{ccc}
\\textit{100} & /inga{\texttoptiebar{t\textbeltl}}unarai{\texttheta}nina/&$\\rightarrow$ & \\textit{3200} & /ing{\texttoptiebar{t\textbeltl}}unrai{\texttheta}nina/&$\\rightarrow$ & \\textit{5500} & /ing{\texttoptiebar{t\textbeltl}}unri{\texttheta}nina/&$\\rightarrow$ & \\textit{6000} & /i{\ng}g{\texttoptiebar{t\textbeltl}}unri{\texttheta}nina/&$\\rightarrow$ & \\textit{7500} & /{\ng}g{\texttoptiebar{t\textbeltl}}unri{\texttheta}nina/&$\\rightarrow$ & \\textit{7501} & /g{\texttoptiebar{t\textbeltl}}unri{\texttheta}nina/&$\\rightarrow$ & \\textit{7600} & /g{\texttoptiebar{t\textbeltl}}uri{\texttheta}nina/&$\\rightarrow$ & \\textit{9500} & /g{\texttoptiebar{t\textbeltl}}uri{\texttheta}nin/&$\\rightarrow$ & \\textit{11501} & /{\texttoptiebar{t\textbeltl}}uri{\texttheta}nin/&$\\rightarrow$ & \\textit{13002} & /{\texttoptiebar{t\textbeltl}}urinin/& \\\\
\\end{tabular}

\\vspace{20pt}\\hline

\\end{nopagebreak}
\\filbreak



\\vspace{15pt}
\\begin{nopagebreak}
\\noindent{\\fontsize{20pt}{10pt}\\textbf{nunalom} } \\textit{month} (N)\\\\
\\noindent {\\tovian \\fontsize{20pt}{10pt} \\textbf{nunalom} }\\\\
\\noindent /nun{\textprimstress}alom/\\\\
\\noindent Syllables: nun.al.om\\\\
\\noindent lit. small+small+cycle\\\\


\\noindent History:

\\vspace{-0pt}
\\hspace{40pt}
\\begin{tabular}{ccc}
\\textit{0} & /ninaninaelakorim/&$\\rightarrow$ & \\textit{2000} & /ninaninaelagorim/&$\\rightarrow$ & \\textit{3200} & /nnannaelgorm/&$\\rightarrow$ & \\textit{4501} & /nnannaelorm/&$\\rightarrow$ & \\textit{6500} & /nanaelorm/&$\\rightarrow$ & \\textit{7600} & /nanaelom/&$\\rightarrow$ & \\textit{9200} & /n{\textschwa}naelom/&$\\rightarrow$ & \\textit{10000} & /n{\textschwa}nalom/&$\\rightarrow$ & \\textit{11990} & /nnalom/&$\\rightarrow$ & \\textit{11996} & /nunalom/& \\\\
\\end{tabular}

\\vspace{20pt}\\hline

\\end{nopagebreak}
\\filbreak



\\vspace{15pt}
\\begin{nopagebreak}
\\noindent{\\fontsize{20pt}{10pt}\\textbf{ay} } \\textit{moon} (N)\\\\
\\noindent {\\tovian \\fontsize{20pt}{10pt} \\textbf{ay} }\\\\
\\noindent /{\textprimstress}ay/\\\\
\\noindent Syllables: aj\\\\


\\noindent History:

\\vspace{-0pt}
\\hspace{40pt}
\\begin{tabular}{ccc}
\\textit{0} & /alius/&$\\rightarrow$ & \\textit{2201} & /aliyus/&$\\rightarrow$ & \\textit{3200} & /aliys/&$\\rightarrow$ & \\textit{6300} & /alys/&$\\rightarrow$ & \\textit{7600} & /ays/&$\\rightarrow$ & \\textit{13001} & /ay/& \\\\
\\end{tabular}

\\vspace{20pt}\\hline

\\end{nopagebreak}
\\filbreak



\\vspace{15pt}
\\begin{nopagebreak}
\\noindent{\\fontsize{20pt}{10pt}\\textbf{mar} } \\textit{mother} (N)\\\\
\\noindent {\\tovian \\fontsize{20pt}{10pt} \\textbf{mar} }\\\\
\\noindent /m{\textprimstress}ar/\\\\
\\noindent Syllables: mar\\\\


\\noindent History:

\\vspace{-0pt}
\\hspace{40pt}
\\begin{tabular}{ccc}
\\textit{0} & /mari/&$\\rightarrow$ & \\textit{9500} & /mar/& \\\\
\\end{tabular}

\\vspace{20pt}\\hline

\\end{nopagebreak}
\\filbreak



\\vspace{15pt}
\\begin{nopagebreak}
\\noindent{\\fontsize{20pt}{10pt}\\textbf{eon} } \\textit{mouth} (N)\\\\
\\noindent {\\tovian \\fontsize{20pt}{10pt} \\textbf{eon} }\\\\
\\noindent /{\textprimstress}eon/\\\\
\\noindent Syllables: e.on\\\\


\\noindent History:

\\vspace{-0pt}
\\hspace{40pt}
\\begin{tabular}{ccc}
\\textit{0} & /eolan/&$\\rightarrow$ & \\textit{3200} & /eoln/&$\\rightarrow$ & \\textit{7600} & /eon/& \\\\
\\end{tabular}

\\vspace{20pt}\\hline

\\end{nopagebreak}
\\filbreak



\\vspace{15pt}
\\begin{nopagebreak}
\\noindent{\\fontsize{20pt}{10pt}\\textbf{kewil} } \\textit{movement} (N)\\\\
\\noindent {\\tovian \\fontsize{20pt}{10pt} \\textbf{kewil} }\\\\
\\noindent /k{\textprimstress}ewil/\\\\
\\noindent Syllables: kew.il\\\\


\\noindent History:

\\vspace{-0pt}
\\hspace{40pt}
\\begin{tabular}{ccc}
\\textit{0} & /kwil/&$\\rightarrow$ & \\textit{15001} & /kewil/& \\\\
\\end{tabular}

\\vspace{20pt}\\hline

\\end{nopagebreak}
\\filbreak



\\vspace{15pt}
\\begin{nopagebreak}
\\noindent{\\fontsize{20pt}{10pt}\\textbf{lhemlor} } \\textit{narcotic} (N)\\\\
\\noindent {\\tovian \\fontsize{20pt}{10pt} \\textbf{lhemlor} }\\\\
\\noindent /{\textbeltl}{\textprimstress}emlor/\\\\
\\noindent Syllables: ɬeml.or\\\\
\\noindent lit. sleep+plant\\\\


\\noindent History:

\\vspace{-0pt}
\\hspace{40pt}
\\begin{tabular}{ccc}
\\textit{500} & /{\textbeltl}e{\dh}el{\textschwa}mellore/&$\\rightarrow$ & \\textit{2300} & /{\textbeltl}e{\dh}elmellore/&$\\rightarrow$ & \\textit{3200} & /{\textbeltl}{\dh}elmllore/&$\\rightarrow$ & \\textit{4502} & /{\textbeltl}elmllore/&$\\rightarrow$ & \\textit{6500} & /{\textbeltl}elmlore/&$\\rightarrow$ & \\textit{7600} & /{\textbeltl}emlore/&$\\rightarrow$ & \\textit{9500} & /{\textbeltl}emlor/& \\\\
\\end{tabular}

\\vspace{20pt}\\hline

\\end{nopagebreak}
\\filbreak



\\vspace{15pt}
\\begin{nopagebreak}
\\noindent{\\fontsize{20pt}{10pt}\\textbf{men} } \\textit{newness} (N)\\\\
\\noindent {\\tovian \\fontsize{20pt}{10pt} \\textbf{men} }\\\\
\\noindent /m{\textprimstress}en/\\\\
\\noindent Syllables: men\\\\


\\noindent History:

\\vspace{-0pt}
\\hspace{40pt}
\\begin{tabular}{ccc}
\\textit{0} & /{\textyogh}emeno/&$\\rightarrow$ & \\textit{3200} & /{\textyogh}meno/&$\\rightarrow$ & \\textit{9500} & /{\textyogh}men/&$\\rightarrow$ & \\textit{12001} & /{\textesh}men/&$\\rightarrow$ & \\textit{13002} & /men/& \\\\
\\end{tabular}

\\vspace{20pt}\\hline

\\end{nopagebreak}
\\filbreak



\\vspace{15pt}
\\begin{nopagebreak}
\\noindent{\\fontsize{20pt}{10pt}\\textbf{lhil} } \\textit{night/nighttime} (N)\\\\
\\noindent {\\tovian \\fontsize{20pt}{10pt} \\textbf{lhil} }\\\\
\\noindent /{\textbeltl}{\textprimstress}il/\\\\
\\noindent Syllables: ɬil\\\\
\\noindent \\textit{Time of day}\\\\


\\noindent History:

\\vspace{-0pt}
\\hspace{40pt}
\\begin{tabular}{ccc}
\\textit{0} & /a{\textbeltl}e{\textsubbridge{t}}inla/&$\\rightarrow$ & \\textit{1000} & /a{\textbeltl}{\textsubbridge{t}}inla/&$\\rightarrow$ & \\textit{3000} & /a{\textbeltl}inla/&$\\rightarrow$ & \\textit{7500} & /{\textbeltl}inla/&$\\rightarrow$ & \\textit{7600} & /{\textbeltl}ila/&$\\rightarrow$ & \\textit{9500} & /{\textbeltl}il/& \\\\
\\end{tabular}

\\vspace{20pt}\\hline

\\end{nopagebreak}
\\filbreak



\\vspace{15pt}
\\begin{nopagebreak}
\\noindent{\\fontsize{20pt}{10pt}\\textbf{ki} } \\textit{nine/9} (N)\\\\
\\noindent {\\tovian \\fontsize{20pt}{10pt} \\textbf{ki} }\\\\
\\noindent /k{\textprimstress}i/\\\\
\\noindent Syllables: ki\\\\


\\noindent History:

\\vspace{-0pt}
\\hspace{40pt}
\\begin{tabular}{ccc}
\\textit{0} & /ki/& \\\\
\\end{tabular}

\\vspace{20pt}\\hline

\\end{nopagebreak}
\\filbreak



\\vspace{15pt}
\\begin{nopagebreak}
\\noindent{\\fontsize{20pt}{10pt}\\textbf{i-} } \\textit{NOM/nomative case marker} (CASE)\\\\
\\noindent {\\tovian \\fontsize{20pt}{10pt} \\textbf{i-} }\\\\
\\noindent /{\textprimstress}i/\\\\
\\noindent Syllables: i\\\\


\\noindent History:

\\vspace{-0pt}
\\hspace{40pt}
\\begin{tabular}{ccc}
\\textit{0} & /i-/& \\\\
\\end{tabular}

\\vspace{20pt}\\hline

\\end{nopagebreak}
\\filbreak



\\vspace{15pt}
\\begin{nopagebreak}
\\noindent{\\fontsize{20pt}{10pt}\\textbf{anad} } \\textit{non-person hair/fur/hide/skin/scales} (N)\\\\
\\noindent {\\tovian \\fontsize{20pt}{10pt} \\textbf{anad} }\\\\
\\noindent /{\textprimstress}anad/\\\\
\\noindent Syllables: an.ad\\\\


\\noindent History:

\\vspace{-0pt}
\\hspace{40pt}
\\begin{tabular}{ccc}
\\textit{0} & /ana{\textsubbridge{t}}ur/&$\\rightarrow$ & \\textit{2000} & /anadur/&$\\rightarrow$ & \\textit{3200} & /anadr/&$\\rightarrow$ & \\textit{13001} & /anad/& \\\\
\\end{tabular}

\\vspace{20pt}\\hline

\\end{nopagebreak}
\\filbreak



\\vspace{15pt}
\\begin{nopagebreak}
\\noindent{\\fontsize{20pt}{10pt}\\textbf{lar} } \\textit{north} (N)\\\\
\\noindent {\\tovian \\fontsize{20pt}{10pt} \\textbf{lar} }\\\\
\\noindent /l{\textprimstress}ar/\\\\
\\noindent Syllables: lar\\\\


\\noindent History:

\\vspace{-0pt}
\\hspace{40pt}
\\begin{tabular}{ccc}
\\textit{0} & /lare/&$\\rightarrow$ & \\textit{9500} & /lar/& \\\\
\\end{tabular}

\\vspace{20pt}\\hline

\\end{nopagebreak}
\\filbreak



\\vspace{15pt}
\\begin{nopagebreak}
\\noindent{\\fontsize{20pt}{10pt}\\textbf{ven} } \\textit{nose} (N)\\\\
\\noindent {\\tovian \\fontsize{20pt}{10pt} \\textbf{ven} }\\\\
\\noindent /v{\textprimstress}en/\\\\
\\noindent Syllables: ven\\\\


\\noindent History:

\\vspace{-0pt}
\\hspace{40pt}
\\begin{tabular}{ccc}
\\textit{0} & /verun/&$\\rightarrow$ & \\textit{3200} & /vern/&$\\rightarrow$ & \\textit{7600} & /ven/& \\\\
\\end{tabular}

\\vspace{20pt}\\hline

\\end{nopagebreak}
\\filbreak



\\vspace{15pt}
\\begin{nopagebreak}
\\noindent{\\fontsize{20pt}{10pt}\\textbf{felesh} } \\textit{November} (N)\\\\
\\noindent {\\tovian \\fontsize{20pt}{10pt} \\textbf{felesh} }\\\\
\\noindent /f{\textprimstress}ele{\textesh}/\\\\
\\noindent Syllables: fel.eʃ\\\\
\\noindent lit. three+harvest\\\\
\\noindent \\textit{Prairilight}\\\\


\\noindent History:

\\vspace{-0pt}
\\hspace{40pt}
\\begin{tabular}{ccc}
\\textit{6000} & /femmllre{\textyogh}/&$\\rightarrow$ & \\textit{6500} & /femlre{\textyogh}/&$\\rightarrow$ & \\textit{6503} & /femrle{\textyogh}/&$\\rightarrow$ & \\textit{7600} & /ferle{\textyogh}/&$\\rightarrow$ & \\textit{12001} & /ferle{\textesh}/&$\\rightarrow$ & \\textit{12505} & /fele{\textesh}/& \\\\
\\end{tabular}

\\vspace{20pt}\\hline

\\end{nopagebreak}
\\filbreak



\\vspace{15pt}
\\begin{nopagebreak}
\\noindent{\\fontsize{20pt}{10pt}\\textbf{tlo-} } \\textit{OBL/obligative mood marker} (MOOD)\\\\
\\noindent {\\tovian \\fontsize{20pt}{10pt} \\textbf{tlo-} }\\\\
\\noindent /{\texttoptiebar{t\textbeltl}}{\textprimstress}o/\\\\
\\noindent Syllables: ƛo\\\\


\\noindent History:

\\vspace{-0pt}
\\hspace{40pt}
\\begin{tabular}{ccc}
\\textit{0} & /{\texttoptiebar{t\textbeltl}}o-/& \\\\
\\end{tabular}

\\vspace{20pt}\\hline

\\end{nopagebreak}
\\filbreak



\\vspace{15pt}
\\begin{nopagebreak}
\\noindent{\\fontsize{20pt}{10pt}\\textbf{renlesh} } \\textit{October} (N)\\\\
\\noindent {\\tovian \\fontsize{20pt}{10pt} \\textbf{renlesh} }\\\\
\\noindent /r{\textprimstress}enle{\textesh}/\\\\
\\noindent Syllables: renl.eʃ\\\\
\\noindent lit. standing+harvest\\\\
\\noindent \\textit{Fructiwane}\\\\


\\noindent History:

\\vspace{-0pt}
\\hspace{40pt}
\\begin{tabular}{ccc}
\\textit{6000} & /remnmllre{\textyogh}/&$\\rightarrow$ & \\textit{6500} & /remnmlre{\textyogh}/&$\\rightarrow$ & \\textit{6503} & /remnmrle{\textyogh}/&$\\rightarrow$ & \\textit{7600} & /renrle{\textyogh}/&$\\rightarrow$ & \\textit{12001} & /renrle{\textesh}/&$\\rightarrow$ & \\textit{12505} & /renle{\textesh}/& \\\\
\\end{tabular}

\\vspace{20pt}\\hline

\\end{nopagebreak}
\\filbreak



\\vspace{15pt}
\\begin{nopagebreak}
\\noindent{\\fontsize{20pt}{10pt}\\textbf{lheyalh} } \\textit{one hundred forty four/144} (N)\\\\
\\noindent {\\tovian \\fontsize{20pt}{10pt} \\textbf{lheyalh} }\\\\
\\noindent /{\textbeltl}{\textprimstress}eya{\textbeltl}/\\\\
\\noindent Syllables: ɬej.aɬ\\\\
\\noindent lit. twelve+twelve\\\\


\\noindent History:

\\vspace{-0pt}
\\hspace{40pt}
\\begin{tabular}{ccc}
\\textit{4000} & /{\textbeltl}ya{\textbeltl}ya/&$\\rightarrow$ & \\textit{9500} & /{\textbeltl}ya{\textbeltl}y/&$\\rightarrow$ & \\textit{13001} & /{\textbeltl}ya{\textbeltl}/&$\\rightarrow$ & \\textit{15001} & /{\textbeltl}eya{\textbeltl}/& \\\\
\\end{tabular}

\\vspace{20pt}\\hline

\\end{nopagebreak}
\\filbreak



\\vspace{15pt}
\\begin{nopagebreak}
\\noindent{\\fontsize{20pt}{10pt}\\textbf{lo} } \\textit{one/1} (N)\\\\
\\noindent {\\tovian \\fontsize{20pt}{10pt} \\textbf{lo} }\\\\
\\noindent /l{\textprimstress}o/\\\\
\\noindent Syllables: lo\\\\


\\noindent History:

\\vspace{-0pt}
\\hspace{40pt}
\\begin{tabular}{ccc}
\\textit{0} & /lo/& \\\\
\\end{tabular}

\\vspace{20pt}\\hline

\\end{nopagebreak}
\\filbreak



\\vspace{15pt}
\\begin{nopagebreak}
\\noindent{\\fontsize{20pt}{10pt}\\textbf{dam} } \\textit{openness} (N)\\\\
\\noindent {\\tovian \\fontsize{20pt}{10pt} \\textbf{dam} }\\\\
\\noindent /d{\textprimstress}am/\\\\
\\noindent Syllables: dam\\\\
\\noindent lit. not+closure\\\\


\\noindent History:

\\vspace{-0pt}
\\hspace{40pt}
\\begin{tabular}{ccc}
\\textit{1300} & /do{\textsubbridge{t}}anum/&$\\rightarrow$ & \\textit{2000} & /dodanum/&$\\rightarrow$ & \\textit{3200} & /ddanm/&$\\rightarrow$ & \\textit{6500} & /danm/&$\\rightarrow$ & \\textit{7600} & /dam/& \\\\
\\end{tabular}

\\vspace{20pt}\\hline

\\end{nopagebreak}
\\filbreak



\\vspace{15pt}
\\begin{nopagebreak}
\\noindent{\\fontsize{20pt}{10pt}\\textbf{do} } \\textit{opposite/not/un-} (N)\\\\
\\noindent {\\tovian \\fontsize{20pt}{10pt} \\textbf{do} }\\\\
\\noindent /d{\textprimstress}o/\\\\
\\noindent Syllables: do\\\\


\\noindent History:

\\vspace{-0pt}
\\hspace{40pt}
\\begin{tabular}{ccc}
\\textit{0} & /do/& \\\\
\\end{tabular}

\\vspace{20pt}\\hline

\\end{nopagebreak}
\\filbreak



\\vspace{15pt}
\\begin{nopagebreak}
\\noindent{\\fontsize{20pt}{10pt}\\textbf{lhe-} } \\textit{OPT/optative mood marker} (MOOD)\\\\
\\noindent {\\tovian \\fontsize{20pt}{10pt} \\textbf{lhe-} }\\\\
\\noindent /{\textbeltl}{\textprimstress}e/\\\\
\\noindent Syllables: ɬe\\\\


\\noindent History:

\\vspace{-0pt}
\\hspace{40pt}
\\begin{tabular}{ccc}
\\textit{0} & /{\textbeltl}e-/& \\\\
\\end{tabular}

\\vspace{20pt}\\hline

\\end{nopagebreak}
\\filbreak



\\vspace{15pt}
\\begin{nopagebreak}
\\noindent{\\fontsize{20pt}{10pt}\\textbf{remhedyel} } \\textit{orphan} (N)\\\\
\\noindent {\\tovian \\fontsize{20pt}{10pt} \\textbf{remhedyel} }\\\\
\\noindent /remħ{\textprimstress}edyel/\\\\
\\noindent Syllables: remħ.edj.el\\\\
\\noindent lit. abandoned+person\\\\


\\noindent History:

\\vspace{-0pt}
\\hspace{40pt}
\\begin{tabular}{ccc}
\\textit{1200} & /oneremih{\texttheta}iledriela/&$\\rightarrow$ & \\textit{2201} & /oneremih{\texttheta}iledriyela/&$\\rightarrow$ & \\textit{3200} & /onremh{\texttheta}ledryela/&$\\rightarrow$ & \\textit{4502} & /onremhledryela/&$\\rightarrow$ & \\textit{7500} & /nremhledryela/&$\\rightarrow$ & \\textit{7600} & /remhledyela/&$\\rightarrow$ & \\textit{9500} & /remhledyel/&$\\rightarrow$ & \\textit{12500} & /remhedyel/&$\\rightarrow$ & \\textit{14006} & /remħedyel/& \\\\
\\end{tabular}

\\vspace{20pt}\\hline

\\end{nopagebreak}
\\filbreak



\\vspace{15pt}
\\begin{nopagebreak}
\\noindent{\\fontsize{20pt}{10pt}\\textbf{hi-ne} } \\textit{our/ours} (_)\\\\
\\noindent {\\tovian \\fontsize{20pt}{10pt} \\textbf{hi-ne} }\\\\
\\noindent /h{\textprimstress}ine/\\\\
\\noindent Syllables: hi.ne\\\\
\\noindent lit. GEN+we\\\\
\\noindent \\textit{Ps}\\\\


\\noindent History:

\\vspace{-0pt}
\\hspace{40pt}
\\begin{tabular}{ccc}
\\textit{0} & /hi-ne/& \\\\
\\end{tabular}

\\vspace{20pt}\\hline

\\end{nopagebreak}
\\filbreak



\\vspace{15pt}
\\begin{nopagebreak}
\\noindent{\\fontsize{20pt}{10pt}\\textbf{har} } \\textit{pain} (N)\\\\
\\noindent {\\tovian \\fontsize{20pt}{10pt} \\textbf{har} }\\\\
\\noindent /ħ{\textprimstress}ar/\\\\
\\noindent Syllables: ħar\\\\


\\noindent History:

\\vspace{-0pt}
\\hspace{40pt}
\\begin{tabular}{ccc}
\\textit{0} & /hara/&$\\rightarrow$ & \\textit{9500} & /har/&$\\rightarrow$ & \\textit{14006} & /ħar/& \\\\
\\end{tabular}

\\vspace{20pt}\\hline

\\end{nopagebreak}
\\filbreak



\\vspace{15pt}
\\begin{nopagebreak}
\\noindent{\\fontsize{20pt}{10pt}\\textbf{pa-} } \\textit{PASS/passive voice marker} (VOICE)\\\\
\\noindent {\\tovian \\fontsize{20pt}{10pt} \\textbf{pa-} }\\\\
\\noindent /p{\textprimstress}a/\\\\
\\noindent Syllables: pa\\\\


\\noindent History:

\\vspace{-0pt}
\\hspace{40pt}
\\begin{tabular}{ccc}
\\textit{0} & /pa-/& \\\\
\\end{tabular}

\\vspace{20pt}\\hline

\\end{nopagebreak}
\\filbreak



\\vspace{15pt}
\\begin{nopagebreak}
\\noindent{\\fontsize{20pt}{10pt}\\textbf{hun} } \\textit{past} (Ns)\\\\
\\noindent {\\tovian \\fontsize{20pt}{10pt} \\textbf{hun} }\\\\
\\noindent /ħ{\textprimstress}un/\\\\
\\noindent Syllables: ħun\\\\
\\noindent lit. behind+path\\\\


\\noindent History:

\\vspace{-0pt}
\\hspace{40pt}
\\begin{tabular}{ccc}
\\textit{10} & /hulnum/&$\\rightarrow$ & \\textit{3200} & /hulnm/&$\\rightarrow$ & \\textit{7600} & /hunm/&$\\rightarrow$ & \\textit{13001} & /hun/&$\\rightarrow$ & \\textit{14006} & /ħun/& \\\\
\\end{tabular}

\\vspace{20pt}\\hline

\\end{nopagebreak}
\\filbreak



\\vspace{15pt}
\\begin{nopagebreak}
\\noindent{\\fontsize{20pt}{10pt}\\textbf{num} } \\textit{path} (N)\\\\
\\noindent {\\tovian \\fontsize{20pt}{10pt} \\textbf{num} }\\\\
\\noindent /n{\textprimstress}um/\\\\
\\noindent Syllables: num\\\\


\\noindent History:

\\vspace{-0pt}
\\hspace{40pt}
\\begin{tabular}{ccc}
\\textit{0} & /num/& \\\\
\\end{tabular}

\\vspace{20pt}\\hline

\\end{nopagebreak}
\\filbreak



\\vspace{15pt}
\\begin{nopagebreak}
\\noindent{\\fontsize{20pt}{10pt}\\textbf{tla} } \\textit{peace/peacefulness/calmness/calm} (N)\\\\
\\noindent {\\tovian \\fontsize{20pt}{10pt} \\textbf{tla} }\\\\
\\noindent /{\texttoptiebar{t\textbeltl}}{\textprimstress}a/\\\\
\\noindent Syllables: ƛa\\\\


\\noindent History:

\\vspace{-0pt}
\\hspace{40pt}
\\begin{tabular}{ccc}
\\textit{0} & /{\texttoptiebar{t\textbeltl}}af/&$\\rightarrow$ & \\textit{7501} & /{\texttoptiebar{t\textbeltl}}a{\textphi}/&$\\rightarrow$ & \\textit{8500} & /{\texttoptiebar{t\textbeltl}}a/& \\\\
\\end{tabular}

\\vspace{20pt}\\hline

\\end{nopagebreak}
\\filbreak



\\vspace{15pt}
\\begin{nopagebreak}
\\noindent{\\fontsize{20pt}{10pt}\\textbf{yel} } \\textit{person} (N)\\\\
\\noindent {\\tovian \\fontsize{20pt}{10pt} \\textbf{yel} }\\\\
\\noindent /y{\textprimstress}el/\\\\
\\noindent Syllables: jel\\\\


\\noindent History:

\\vspace{-0pt}
\\hspace{40pt}
\\begin{tabular}{ccc}
\\textit{0} & /driela/&$\\rightarrow$ & \\textit{2201} & /driyela/&$\\rightarrow$ & \\textit{3200} & /dryela/&$\\rightarrow$ & \\textit{7501} & /ryela/&$\\rightarrow$ & \\textit{7600} & /yela/&$\\rightarrow$ & \\textit{9500} & /yel/& \\\\
\\end{tabular}

\\vspace{20pt}\\hline

\\end{nopagebreak}
\\filbreak



\\vspace{15pt}
\\begin{nopagebreak}
\\noindent{\\fontsize{20pt}{10pt}\\textbf{terel} } \\textit{person skin} (N)\\\\
\\noindent {\\tovian \\fontsize{20pt}{10pt} \\textbf{terel} }\\\\
\\noindent /{\textsubbridge{t}}{\textprimstress}erel/\\\\
\\noindent Syllables: ter.el\\\\


\\noindent History:

\\vspace{-0pt}
\\hspace{40pt}
\\begin{tabular}{ccc}
\\textit{0} & /{\textsubbridge{t}}relahe{\textsubbridge{t}}/&$\\rightarrow$ & \\textit{1000} & /{\textsubbridge{t}}relah{\textsubbridge{t}}/&$\\rightarrow$ & \\textit{3000} & /{\textsubbridge{t}}relah/&$\\rightarrow$ & \\textit{3200} & /{\textsubbridge{t}}relh/&$\\rightarrow$ & \\textit{5000} & /{\textsubbridge{t}}rel/&$\\rightarrow$ & \\textit{15001} & /{\textsubbridge{t}}erel/& \\\\
\\end{tabular}

\\vspace{20pt}\\hline

\\end{nopagebreak}
\\filbreak



\\vspace{15pt}
\\begin{nopagebreak}
\\noindent{\\fontsize{20pt}{10pt}\\textbf{-e} } \\textit{PL/plural marker} (PLURAL)\\\\
\\noindent {\\tovian \\fontsize{20pt}{10pt} \\textbf{-e} }\\\\
\\noindent /{\textprimstress}e/\\\\
\\noindent Syllables: e\\\\


\\noindent History:

\\vspace{-0pt}
\\hspace{40pt}
\\begin{tabular}{ccc}
\\textit{0} & /-e/& \\\\
\\end{tabular}

\\vspace{20pt}\\hline

\\end{nopagebreak}
\\filbreak



\\vspace{15pt}
\\begin{nopagebreak}
\\noindent{\\fontsize{20pt}{10pt}\\textbf{thil} } \\textit{place} (N)\\\\
\\noindent {\\tovian \\fontsize{20pt}{10pt} \\textbf{thil} }\\\\
\\noindent /{\texttheta}{\textprimstress}il/\\\\
\\noindent Syllables: θil\\\\


\\noindent History:

\\vspace{-0pt}
\\hspace{40pt}
\\begin{tabular}{ccc}
\\textit{0} & /{\texttheta}ile/&$\\rightarrow$ & \\textit{9500} & /{\texttheta}il/& \\\\
\\end{tabular}

\\vspace{20pt}\\hline

\\end{nopagebreak}
\\filbreak



\\vspace{15pt}
\\begin{nopagebreak}
\\noindent{\\fontsize{20pt}{10pt}\\textbf{til} } \\textit{plain} (N)\\\\
\\noindent {\\tovian \\fontsize{20pt}{10pt} \\textbf{til} }\\\\
\\noindent /{\textsubbridge{t}}{\textprimstress}il/\\\\
\\noindent Syllables: til\\\\
\\noindent lit. flatness+place\\\\


\\noindent History:

\\vspace{-0pt}
\\hspace{40pt}
\\begin{tabular}{ccc}
\\textit{50} & /{\textsubbridge{t}}um{\texttheta}ile/&$\\rightarrow$ & \\textit{3200} & /{\textsubbridge{t}}m{\texttheta}ile/&$\\rightarrow$ & \\textit{3503} & /{\textsubbridge{t}}n{\texttheta}ile/&$\\rightarrow$ & \\textit{6400} & /{\textsubbridge{t}}{\texttheta}ile/&$\\rightarrow$ & \\textit{9500} & /{\textsubbridge{t}}{\texttheta}il/&$\\rightarrow$ & \\textit{10700} & /{\textsubbridge{t}}{\textsubbridge{t}}il/&$\\rightarrow$ & \\textit{12007} & /{\textsubbridge{t}}il/& \\\\
\\end{tabular}

\\vspace{20pt}\\hline

\\end{nopagebreak}
\\filbreak



\\vspace{15pt}
\\begin{nopagebreak}
\\noindent{\\fontsize{20pt}{10pt}\\textbf{melor} } \\textit{plant} (N)\\\\
\\noindent {\\tovian \\fontsize{20pt}{10pt} \\textbf{melor} }\\\\
\\noindent /m{\textprimstress}elor/\\\\
\\noindent Syllables: mel.or\\\\
\\noindent lit. life+action\\\\


\\noindent History:

\\vspace{-0pt}
\\hspace{40pt}
\\begin{tabular}{ccc}
\\textit{50} & /{\textschwa}mellore/&$\\rightarrow$ & \\textit{2300} & /mellore/&$\\rightarrow$ & \\textit{3200} & /mllore/&$\\rightarrow$ & \\textit{6500} & /mlore/&$\\rightarrow$ & \\textit{7502} & /melore/&$\\rightarrow$ & \\textit{9500} & /melor/& \\\\
\\end{tabular}

\\vspace{20pt}\\hline

\\end{nopagebreak}
\\filbreak



\\vspace{15pt}
\\begin{nopagebreak}
\\noindent{\\fontsize{20pt}{10pt}\\textbf{kog} } \\textit{playing} (N)\\\\
\\noindent {\\tovian \\fontsize{20pt}{10pt} \\textbf{kog} }\\\\
\\noindent /k{\textprimstress}og/\\\\
\\noindent Syllables: kog\\\\


\\noindent History:

\\vspace{-0pt}
\\hspace{40pt}
\\begin{tabular}{ccc}
\\textit{0} & /kokoki/&$\\rightarrow$ & \\textit{1000} & /kkoki/&$\\rightarrow$ & \\textit{2000} & /kkogi/&$\\rightarrow$ & \\textit{3000} & /kogi/&$\\rightarrow$ & \\textit{9500} & /kog/& \\\\
\\end{tabular}

\\vspace{20pt}\\hline

\\end{nopagebreak}
\\filbreak



\\vspace{15pt}
\\begin{nopagebreak}
\\noindent{\\fontsize{20pt}{10pt}\\textbf{el} } \\textit{possession} (N)\\\\
\\noindent {\\tovian \\fontsize{20pt}{10pt} \\textbf{el} }\\\\
\\noindent /{\textprimstress}el/\\\\
\\noindent Syllables: el\\\\


\\noindent History:

\\vspace{-0pt}
\\hspace{40pt}
\\begin{tabular}{ccc}
\\textit{0} & /elis/&$\\rightarrow$ & \\textit{3200} & /els/&$\\rightarrow$ & \\textit{13001} & /el/& \\\\
\\end{tabular}

\\vspace{20pt}\\hline

\\end{nopagebreak}
\\filbreak



\\vspace{15pt}
\\begin{nopagebreak}
\\noindent{\\fontsize{20pt}{10pt}\\textbf{sho} } \\textit{possibility} (N)\\\\
\\noindent {\\tovian \\fontsize{20pt}{10pt} \\textbf{sho} }\\\\
\\noindent /{\textesh}{\textprimstress}o/\\\\
\\noindent Syllables: ʃo\\\\


\\noindent History:

\\vspace{-0pt}
\\hspace{40pt}
\\begin{tabular}{ccc}
\\textit{0} & /{\textesh}ohe/&$\\rightarrow$ & \\textit{5000} & /{\textesh}oe/&$\\rightarrow$ & \\textit{9500} & /{\textesh}o/& \\\\
\\end{tabular}

\\vspace{20pt}\\hline

\\end{nopagebreak}
\\filbreak



\\vspace{15pt}
\\begin{nopagebreak}
\\noindent{\\fontsize{20pt}{10pt}\\textbf{sho-} } \\textit{POT/potential mood marker} (MOOD)\\\\
\\noindent {\\tovian \\fontsize{20pt}{10pt} \\textbf{sho-} }\\\\
\\noindent /{\textesh}{\textprimstress}o/\\\\
\\noindent Syllables: ʃo\\\\


\\noindent History:

\\vspace{-0pt}
\\hspace{40pt}
\\begin{tabular}{ccc}
\\textit{0} & /{\textesh}o-/& \\\\
\\end{tabular}

\\vspace{20pt}\\hline

\\end{nopagebreak}
\\filbreak



\\vspace{15pt}
\\begin{nopagebreak}
\\noindent{\\fontsize{20pt}{10pt}\\textbf{num} } \\textit{present} (Ns)\\\\
\\noindent {\\tovian \\fontsize{20pt}{10pt} \\textbf{num} }\\\\
\\noindent /n{\textprimstress}um/\\\\
\\noindent Syllables: num\\\\
\\noindent lit. path+center\\\\


\\noindent History:

\\vspace{-0pt}
\\hspace{40pt}
\\begin{tabular}{ccc}
\\textit{10} & /num{\textbeltl}el/&$\\rightarrow$ & \\textit{3200} & /num{\textbeltl}l/&$\\rightarrow$ & \\textit{12500} & /num{\textbeltl}/&$\\rightarrow$ & \\textit{13001} & /num/& \\\\
\\end{tabular}

\\vspace{20pt}\\hline

\\end{nopagebreak}
\\filbreak



\\vspace{15pt}
\\begin{nopagebreak}
\\noindent{\\fontsize{20pt}{10pt}\\textbf{sor} } \\textit{proximity/closeness} (N)\\\\
\\noindent {\\tovian \\fontsize{20pt}{10pt} \\textbf{sor} }\\\\
\\noindent /s{\textprimstress}or/\\\\
\\noindent Syllables: sor\\\\


\\noindent History:

\\vspace{-0pt}
\\hspace{40pt}
\\begin{tabular}{ccc}
\\textit{0} & /sor/& \\\\
\\end{tabular}

\\vspace{20pt}\\hline

\\end{nopagebreak}
\\filbreak



\\vspace{15pt}
\\begin{nopagebreak}
\\noindent{\\fontsize{20pt}{10pt}\\textbf{lhu-} } \\textit{PUR/purpose case marker/for} (CASE)\\\\
\\noindent {\\tovian \\fontsize{20pt}{10pt} \\textbf{lhu-} }\\\\
\\noindent /{\textbeltl}{\textprimstress}u/\\\\
\\noindent Syllables: ɬu\\\\


\\noindent History:

\\vspace{-0pt}
\\hspace{40pt}
\\begin{tabular}{ccc}
\\textit{0} & /{\textbeltl}u-/& \\\\
\\end{tabular}

\\vspace{20pt}\\hline

\\end{nopagebreak}
\\filbreak



\\vspace{15pt}
\\begin{nopagebreak}
\\noindent{\\fontsize{20pt}{10pt}\\textbf{man} } \\textit{putting/placing} (N)\\\\
\\noindent {\\tovian \\fontsize{20pt}{10pt} \\textbf{man} }\\\\
\\noindent /m{\textprimstress}an/\\\\
\\noindent Syllables: man\\\\


\\noindent History:

\\vspace{-0pt}
\\hspace{40pt}
\\begin{tabular}{ccc}
\\textit{0} & /hemmani/&$\\rightarrow$ & \\textit{3200} & /hmmani/&$\\rightarrow$ & \\textit{6500} & /hmani/&$\\rightarrow$ & \\textit{9500} & /hman/&$\\rightarrow$ & \\textit{13002} & /man/& \\\\
\\end{tabular}

\\vspace{20pt}\\hline

\\end{nopagebreak}
\\filbreak



\\vspace{15pt}
\\begin{nopagebreak}
\\noindent{\\fontsize{20pt}{10pt}\\textbf{lhan} } \\textit{question/INT} (INT)\\\\
\\noindent {\\tovian \\fontsize{20pt}{10pt} \\textbf{lhan} }\\\\
\\noindent /{\textbeltl}{\textprimstress}an/\\\\
\\noindent Syllables: ɬan\\\\


\\noindent History:

\\vspace{-0pt}
\\hspace{40pt}
\\begin{tabular}{ccc}
\\textit{0} & /{\textbeltl}an/& \\\\
\\end{tabular}

\\vspace{20pt}\\hline

\\end{nopagebreak}
\\filbreak



\\vspace{15pt}
\\begin{nopagebreak}
\\noindent{\\fontsize{20pt}{10pt}\\textbf{letlatl} } \\textit{rain} (N)\\\\
\\noindent {\\tovian \\fontsize{20pt}{10pt} \\textbf{letlatl} }\\\\
\\noindent /l{\textprimstress}e{\texttoptiebar{t\textbeltl}}a{\texttoptiebar{t\textbeltl}}/\\\\
\\noindent Syllables: leƛ.aƛ\\\\
\\noindent lit. water+falling\\\\


\\noindent History:

\\vspace{-0pt}
\\hspace{40pt}
\\begin{tabular}{ccc}
\\textit{0} & /z{\textsubbridge{t}}el{\texttoptiebar{t\textbeltl}}a{\texttoptiebar{t\textbeltl}}a/&$\\rightarrow$ & \\textit{2000} & /zdel{\texttoptiebar{t\textbeltl}}a{\texttoptiebar{t\textbeltl}}a/&$\\rightarrow$ & \\textit{3200} & /zdl{\texttoptiebar{t\textbeltl}}a{\texttoptiebar{t\textbeltl}}a/&$\\rightarrow$ & \\textit{4500} & /zl{\texttoptiebar{t\textbeltl}}a{\texttoptiebar{t\textbeltl}}a/&$\\rightarrow$ & \\textit{7501} & /l{\texttoptiebar{t\textbeltl}}a{\texttoptiebar{t\textbeltl}}a/&$\\rightarrow$ & \\textit{9500} & /l{\texttoptiebar{t\textbeltl}}a{\texttoptiebar{t\textbeltl}}/&$\\rightarrow$ & \\textit{15001} & /le{\texttoptiebar{t\textbeltl}}a{\texttoptiebar{t\textbeltl}}/& \\\\
\\end{tabular}

\\vspace{20pt}\\hline

\\end{nopagebreak}
\\filbreak



\\vspace{15pt}
\\begin{nopagebreak}
\\noindent{\\fontsize{20pt}{10pt}\\textbf{kith} } \\textit{reaching/arriving/traveling/going} (N)\\\\
\\noindent {\\tovian \\fontsize{20pt}{10pt} \\textbf{kith} }\\\\
\\noindent /k{\textprimstress}i{\texttheta}/\\\\
\\noindent Syllables: kiθ\\\\


\\noindent History:

\\vspace{-0pt}
\\hspace{40pt}
\\begin{tabular}{ccc}
\\textit{0} & /ki{\texttheta}sa/&$\\rightarrow$ & \\textit{4502} & /ki{\texttheta}a/&$\\rightarrow$ & \\textit{9500} & /ki{\texttheta}/& \\\\
\\end{tabular}

\\vspace{20pt}\\hline

\\end{nopagebreak}
\\filbreak



\\vspace{15pt}
\\begin{nopagebreak}
\\noindent{\\fontsize{20pt}{10pt}\\textbf{muloelhaswan} } \\textit{rebirth} (N)\\\\
\\noindent {\\tovian \\fontsize{20pt}{10pt} \\textbf{muloelhaswan} }\\\\
\\noindent /muloe{\textbeltl}{\textprimstress}aswan/\\\\
\\noindent Syllables: mul.o.eɬ.asw.an\\\\
\\noindent lit. second+birth\\\\


\\noindent History:

\\vspace{-0pt}
\\hspace{40pt}
\\begin{tabular}{ccc}
\\textit{1500} & /puloe{\textbeltl}ogasuahana/&$\\rightarrow$ & \\textit{2202} & /puloe{\textbeltl}ogasuwahana/&$\\rightarrow$ & \\textit{3200} & /ploe{\textbeltl}gaswahana/&$\\rightarrow$ & \\textit{4500} & /ploe{\textbeltl}aswahana/&$\\rightarrow$ & \\textit{5000} & /ploe{\textbeltl}aswaana/&$\\rightarrow$ & \\textit{5500} & /ploe{\textbeltl}aswana/&$\\rightarrow$ & \\textit{9500} & /ploe{\textbeltl}aswan/&$\\rightarrow$ & \\textit{11000} & /mloe{\textbeltl}aswan/&$\\rightarrow$ & \\textit{15001} & /muloe{\textbeltl}aswan/& \\\\
\\end{tabular}

\\vspace{20pt}\\hline

\\end{nopagebreak}
\\filbreak



\\vspace{15pt}
\\begin{nopagebreak}
\\noindent{\\fontsize{20pt}{10pt}\\textbf{ra-} } \\textit{REC/reciprocal voice marker} (VOICE)\\\\
\\noindent {\\tovian \\fontsize{20pt}{10pt} \\textbf{ra-} }\\\\
\\noindent /r{\textprimstress}a/\\\\
\\noindent Syllables: ra\\\\


\\noindent History:

\\vspace{-0pt}
\\hspace{40pt}
\\begin{tabular}{ccc}
\\textit{0} & /ra-/& \\\\
\\end{tabular}

\\vspace{20pt}\\hline

\\end{nopagebreak}
\\filbreak



\\vspace{15pt}
\\begin{nopagebreak}
\\noindent{\\fontsize{20pt}{10pt}\\textbf{lhod} } \\textit{redness/red} (N)\\\\
\\noindent {\\tovian \\fontsize{20pt}{10pt} \\textbf{lhod} }\\\\
\\noindent /{\textbeltl}{\textprimstress}od/\\\\
\\noindent Syllables: ɬod\\\\


\\noindent History:

\\vspace{-0pt}
\\hspace{40pt}
\\begin{tabular}{ccc}
\\textit{0} & /{\textbeltl}eho{\textsubbridge{t}}en/&$\\rightarrow$ & \\textit{1000} & /{\textbeltl}ho{\textsubbridge{t}}en/&$\\rightarrow$ & \\textit{2000} & /{\textbeltl}hoden/&$\\rightarrow$ & \\textit{3200} & /{\textbeltl}hodn/&$\\rightarrow$ & \\textit{4502} & /{\textbeltl}odn/&$\\rightarrow$ & \\textit{13001} & /{\textbeltl}od/& \\\\
\\end{tabular}

\\vspace{20pt}\\hline

\\end{nopagebreak}
\\filbreak



\\vspace{15pt}
\\begin{nopagebreak}
\\noindent{\\fontsize{20pt}{10pt}\\textbf{te-} } \\textit{REFL/reflexive voice marker} (VOICE)\\\\
\\noindent {\\tovian \\fontsize{20pt}{10pt} \\textbf{te-} }\\\\
\\noindent /{\textsubbridge{t}}{\textprimstress}e/\\\\
\\noindent Syllables: te\\\\


\\noindent History:

\\vspace{-0pt}
\\hspace{40pt}
\\begin{tabular}{ccc}
\\textit{0} & /{\textsubbridge{t}}e-/& \\\\
\\end{tabular}

\\vspace{20pt}\\hline

\\end{nopagebreak}
\\filbreak



\\vspace{15pt}
\\begin{nopagebreak}
\\noindent{\\fontsize{20pt}{10pt}\\textbf{lenum} } \\textit{repitition} (N)\\\\
\\noindent {\\tovian \\fontsize{20pt}{10pt} \\textbf{lenum} }\\\\
\\noindent /l{\textprimstress}enum/\\\\
\\noindent Syllables: len.um\\\\
\\noindent lit. action+future\\\\


\\noindent History:

\\vspace{-0pt}
\\hspace{40pt}
\\begin{tabular}{ccc}
\\textit{600} & /lorehe{\texttheta}nume/&$\\rightarrow$ & \\textit{1000} & /loreh{\texttheta}nume/&$\\rightarrow$ & \\textit{3200} & /lreh{\texttheta}nume/&$\\rightarrow$ & \\textit{4502} & /lrehnume/&$\\rightarrow$ & \\textit{6240} & /lrenume/&$\\rightarrow$ & \\textit{6503} & /rlenume/&$\\rightarrow$ & \\textit{7600} & /lenume/&$\\rightarrow$ & \\textit{9500} & /lenum/& \\\\
\\end{tabular}

\\vspace{20pt}\\hline

\\end{nopagebreak}
\\filbreak



\\vspace{15pt}
\\begin{nopagebreak}
\\noindent{\\fontsize{20pt}{10pt}\\textbf{yeolfethr} } \\textit{report/confession} (N)\\\\
\\noindent {\\tovian \\fontsize{20pt}{10pt} \\textbf{yeolfethr} }\\\\
\\noindent /ye{\textprimstress}olfe{\texttheta}r/\\\\
\\noindent Syllables: je.olf.eθr\\\\
\\noindent lit. truth+speech\\\\


\\noindent History:

\\vspace{-0pt}
\\hspace{40pt}
\\begin{tabular}{ccc}
\\textit{7200} & /aryeolfe{\dh}ra/&$\\rightarrow$ & \\textit{7500} & /ryeolfe{\dh}ra/&$\\rightarrow$ & \\textit{7600} & /yeolfe{\dh}ra/&$\\rightarrow$ & \\textit{9500} & /yeolfe{\dh}r/&$\\rightarrow$ & \\textit{12002} & /yeolfe{\texttheta}r/& \\\\
\\end{tabular}

\\vspace{20pt}\\hline

\\end{nopagebreak}
\\filbreak



\\vspace{15pt}
\\begin{nopagebreak}
\\noindent{\\fontsize{20pt}{10pt}\\textbf{gith} } \\textit{returning} (N)\\\\
\\noindent {\\tovian \\fontsize{20pt}{10pt} \\textbf{gith} }\\\\
\\noindent /g{\textprimstress}i{\texttheta}/\\\\
\\noindent Syllables: giθ\\\\
\\noindent lit. not+arriving\\\\


\\noindent History:

\\vspace{-0pt}
\\hspace{40pt}
\\begin{tabular}{ccc}
\\textit{0} & /doki{\texttheta}sa/&$\\rightarrow$ & \\textit{2000} & /dogi{\texttheta}sa/&$\\rightarrow$ & \\textit{3200} & /dgi{\texttheta}sa/&$\\rightarrow$ & \\textit{4502} & /dgi{\texttheta}a/&$\\rightarrow$ & \\textit{9500} & /dgi{\texttheta}/&$\\rightarrow$ & \\textit{12503} & /gi{\texttheta}/& \\\\
\\end{tabular}

\\vspace{20pt}\\hline

\\end{nopagebreak}
\\filbreak



\\vspace{15pt}
\\begin{nopagebreak}
\\noindent{\\fontsize{20pt}{10pt}\\textbf{lheya} } \\textit{right} (N)\\\\
\\noindent {\\tovian \\fontsize{20pt}{10pt} \\textbf{lheya} }\\\\
\\noindent /{\textbeltl}{\textprimstress}eya/\\\\
\\noindent Syllables: ɬej.a\\\\


\\noindent History:

\\vspace{-0pt}
\\hspace{40pt}
\\begin{tabular}{ccc}
\\textit{0} & /{\textbeltl}oya/&$\\rightarrow$ & \\textit{2100} & /{\textbeltl}ya/&$\\rightarrow$ & \\textit{15001} & /{\textbeltl}eya/& \\\\
\\end{tabular}

\\vspace{20pt}\\hline

\\end{nopagebreak}
\\filbreak



\\vspace{15pt}
\\begin{nopagebreak}
\\noindent{\\fontsize{20pt}{10pt}\\textbf{shotl} } \\textit{river} (N)\\\\
\\noindent {\\tovian \\fontsize{20pt}{10pt} \\textbf{shotl} }\\\\
\\noindent /{\textesh}{\textprimstress}o{\texttoptiebar{t\textbeltl}}/\\\\
\\noindent Syllables: ʃoƛ\\\\
\\noindent lit. flow+water\\\\


\\noindent History:

\\vspace{-0pt}
\\hspace{40pt}
\\begin{tabular}{ccc}
\\textit{500} & /{\textesh}o{\texttoptiebar{t\textbeltl}}z{\textsubbridge{t}}el/&$\\rightarrow$ & \\textit{2000} & /{\textesh}o{\texttoptiebar{t\textbeltl}}zdel/&$\\rightarrow$ & \\textit{3200} & /{\textesh}o{\texttoptiebar{t\textbeltl}}zdl/&$\\rightarrow$ & \\textit{4500} & /{\textesh}o{\texttoptiebar{t\textbeltl}}zl/&$\\rightarrow$ & \\textit{12000} & /{\textesh}o{\texttoptiebar{t\textbeltl}}sl/&$\\rightarrow$ & \\textit{12500} & /{\textesh}o{\texttoptiebar{t\textbeltl}}s/&$\\rightarrow$ & \\textit{13001} & /{\textesh}o{\texttoptiebar{t\textbeltl}}/& \\\\
\\end{tabular}

\\vspace{20pt}\\hline

\\end{nopagebreak}
\\filbreak



\\vspace{15pt}
\\begin{nopagebreak}
\\noindent{\\fontsize{20pt}{10pt}\\textbf{sethr} } \\textit{running/moving quickly/rushing} (N)\\\\
\\noindent {\\tovian \\fontsize{20pt}{10pt} \\textbf{sethr} }\\\\
\\noindent /s{\textprimstress}e{\texttheta}r/\\\\
\\noindent Syllables: seθr\\\\


\\noindent History:

\\vspace{-0pt}
\\hspace{40pt}
\\begin{tabular}{ccc}
\\textit{0} & /se{\texttheta}ir/&$\\rightarrow$ & \\textit{3200} & /se{\texttheta}r/& \\\\
\\end{tabular}

\\vspace{20pt}\\hline

\\end{nopagebreak}
\\filbreak



\\vspace{15pt}
\\begin{nopagebreak}
\\noindent{\\fontsize{20pt}{10pt}\\textbf{lur} } \\textit{sad} (N)\\\\
\\noindent {\\tovian \\fontsize{20pt}{10pt} \\textbf{lur} }\\\\
\\noindent /l{\textprimstress}ur/\\\\
\\noindent Syllables: lur\\\\


\\noindent History:

\\vspace{-0pt}
\\hspace{40pt}
\\begin{tabular}{ccc}
\\textit{0} & /lura/&$\\rightarrow$ & \\textit{9500} & /lur/& \\\\
\\end{tabular}

\\vspace{20pt}\\hline

\\end{nopagebreak}
\\filbreak



\\vspace{15pt}
\\begin{nopagebreak}
\\noindent{\\fontsize{20pt}{10pt}\\textbf{thoksh} } \\textit{sand} (N)\\\\
\\noindent {\\tovian \\fontsize{20pt}{10pt} \\textbf{thoksh} }\\\\
\\noindent /{\texttheta}{\textprimstress}ok{\textesh}/\\\\
\\noindent Syllables: θokʃ\\\\
\\noindent lit. ground+flow\\\\


\\noindent History:

\\vspace{-0pt}
\\hspace{40pt}
\\begin{tabular}{ccc}
\\textit{1000} & /{\texttheta}ok{\textesh}{\textsubbridge{t}}{\textesh}o{\texttoptiebar{t\textbeltl}}/&$\\rightarrow$ & \\textit{1000} & /{\texttheta}ok{\textesh}{\textsubbridge{t}}{\textesh}{\texttoptiebar{t\textbeltl}}/&$\\rightarrow$ & \\textit{3000} & /{\texttheta}ok{\textesh}{\textesh}{\texttoptiebar{t\textbeltl}}/&$\\rightarrow$ & \\textit{4502} & /{\texttheta}ok{\textesh}{\texttoptiebar{t\textbeltl}}/&$\\rightarrow$ & \\textit{13001} & /{\texttheta}ok{\textesh}/& \\\\
\\end{tabular}

\\vspace{20pt}\\hline

\\end{nopagebreak}
\\filbreak



\\vspace{15pt}
\\begin{nopagebreak}
\\noindent{\\fontsize{20pt}{10pt}\\textbf{nor} } \\textit{sea} (N)\\\\
\\noindent {\\tovian \\fontsize{20pt}{10pt} \\textbf{nor} }\\\\
\\noindent /n{\textprimstress}or/\\\\
\\noindent Syllables: nor\\\\


\\noindent History:

\\vspace{-0pt}
\\hspace{40pt}
\\begin{tabular}{ccc}
\\textit{0} & /nor/& \\\\
\\end{tabular}

\\vspace{20pt}\\hline

\\end{nopagebreak}
\\filbreak



\\vspace{15pt}
\\begin{nopagebreak}
\\noindent{\\fontsize{20pt}{10pt}\\textbf{rithagom} } \\textit{season} (N)\\\\
\\noindent {\\tovian \\fontsize{20pt}{10pt} \\textbf{rithagom} }\\\\
\\noindent /ri{\texttheta}{\textprimstress}agom/\\\\
\\noindent Syllables: riθ.ag.om\\\\
\\noindent lit. air+cycle\\\\


\\noindent History:

\\vspace{-0pt}
\\hspace{40pt}
\\begin{tabular}{ccc}
\\textit{1000} & /rai{\texttheta}elakorim/&$\\rightarrow$ & \\textit{2000} & /rai{\texttheta}elagorim/&$\\rightarrow$ & \\textit{3200} & /rai{\texttheta}lagorm/&$\\rightarrow$ & \\textit{5500} & /ri{\texttheta}lagorm/&$\\rightarrow$ & \\textit{7600} & /ri{\texttheta}lagom/&$\\rightarrow$ & \\textit{12500} & /ri{\texttheta}agom/& \\\\
\\end{tabular}

\\vspace{20pt}\\hline

\\end{nopagebreak}
\\filbreak



\\vspace{15pt}
\\begin{nopagebreak}
\\noindent{\\fontsize{20pt}{10pt}\\textbf{mulo} } \\textit{second} (Ns)\\\\
\\noindent {\\tovian \\fontsize{20pt}{10pt} \\textbf{mulo} }\\\\
\\noindent /m{\textprimstress}ulo/\\\\
\\noindent Syllables: mul.o\\\\
\\noindent lit. ESS+two\\\\


\\noindent History:

\\vspace{-0pt}
\\hspace{40pt}
\\begin{tabular}{ccc}
\\textit{0} & /pu-loe/&$\\rightarrow$ & \\textit{1000} & /puloe/&$\\rightarrow$ & \\textit{3200} & /ploe/&$\\rightarrow$ & \\textit{9500} & /plo/&$\\rightarrow$ & \\textit{11000} & /mlo/&$\\rightarrow$ & \\textit{15001} & /mulo/& \\\\
\\end{tabular}

\\vspace{20pt}\\hline

\\end{nopagebreak}
\\filbreak



\\vspace{15pt}
\\begin{nopagebreak}
\\noindent{\\fontsize{20pt}{10pt}\\textbf{lolesh} } \\textit{September} (N)\\\\
\\noindent {\\tovian \\fontsize{20pt}{10pt} \\textbf{lolesh} }\\\\
\\noindent /l{\textprimstress}ole{\textesh}/\\\\
\\noindent Syllables: lol.eʃ\\\\
\\noindent lit. one+harvest\\\\
\\noindent \\textit{Messihold}\\\\


\\noindent History:

\\vspace{-0pt}
\\hspace{40pt}
\\begin{tabular}{ccc}
\\textit{6000} & /lomllre{\textyogh}/&$\\rightarrow$ & \\textit{6500} & /lomlre{\textyogh}/&$\\rightarrow$ & \\textit{6503} & /lomrle{\textyogh}/&$\\rightarrow$ & \\textit{7600} & /lorle{\textyogh}/&$\\rightarrow$ & \\textit{12001} & /lorle{\textesh}/&$\\rightarrow$ & \\textit{12505} & /lole{\textesh}/& \\\\
\\end{tabular}

\\vspace{20pt}\\hline

\\end{nopagebreak}
\\filbreak



\\vspace{15pt}
\\begin{nopagebreak}
\\noindent{\\fontsize{20pt}{10pt}\\textbf{pol} } \\textit{seven/7} (N)\\\\
\\noindent {\\tovian \\fontsize{20pt}{10pt} \\textbf{pol} }\\\\
\\noindent /p{\textprimstress}ol/\\\\
\\noindent Syllables: pol\\\\


\\noindent History:

\\vspace{-0pt}
\\hspace{40pt}
\\begin{tabular}{ccc}
\\textit{0} & /pol/& \\\\
\\end{tabular}

\\vspace{20pt}\\hline

\\end{nopagebreak}
\\filbreak



\\vspace{15pt}
\\begin{nopagebreak}
\\noindent{\\fontsize{20pt}{10pt}\\textbf{lan} } \\textit{shadow/twilight} (N)\\\\
\\noindent {\\tovian \\fontsize{20pt}{10pt} \\textbf{lan} }\\\\
\\noindent /l{\textprimstress}an/\\\\
\\noindent Syllables: lan\\\\
\\noindent lit. darkness+light\\\\


\\noindent History:

\\vspace{-0pt}
\\hspace{40pt}
\\begin{tabular}{ccc}
\\textit{500} & /bralahana/&$\\rightarrow$ & \\textit{3200} & /brlahana/&$\\rightarrow$ & \\textit{5000} & /brlaana/&$\\rightarrow$ & \\textit{5500} & /brlana/&$\\rightarrow$ & \\textit{7501} & /rlana/&$\\rightarrow$ & \\textit{7600} & /lana/&$\\rightarrow$ & \\textit{9500} & /lan/& \\\\
\\end{tabular}

\\vspace{20pt}\\hline

\\end{nopagebreak}
\\filbreak



\\vspace{15pt}
\\begin{nopagebreak}
\\noindent{\\fontsize{20pt}{10pt}\\textbf{sul} } \\textit{shallows/shortness} (N)\\\\
\\noindent {\\tovian \\fontsize{20pt}{10pt} \\textbf{sul} }\\\\
\\noindent /s{\textprimstress}ul/\\\\
\\noindent Syllables: sul\\\\


\\noindent History:

\\vspace{-0pt}
\\hspace{40pt}
\\begin{tabular}{ccc}
\\textit{0} & /zul/&$\\rightarrow$ & \\textit{12000} & /sul/& \\\\
\\end{tabular}

\\vspace{20pt}\\hline

\\end{nopagebreak}
\\filbreak



\\vspace{15pt}
\\begin{nopagebreak}
\\noindent{\\fontsize{20pt}{10pt}\\textbf{kerath} } \\textit{shard/fragment} (N)\\\\
\\noindent {\\tovian \\fontsize{20pt}{10pt} \\textbf{kerath} }\\\\
\\noindent /k{\textprimstress}era{\texttheta}/\\\\
\\noindent Syllables: ker.aθ\\\\
\\noindent \\textit{A piece broken or separated from a whole; also a fragment of an idea or account.}\\\\


\\noindent History:

\\vspace{-0pt}
\\hspace{40pt}
\\begin{tabular}{ccc}
\\textit{0} & /kera{\texttheta}a/&$\\rightarrow$ & \\textit{3200} & /kra{\texttheta}a/&$\\rightarrow$ & \\textit{9500} & /kra{\texttheta}/&$\\rightarrow$ & \\textit{15001} & /kera{\texttheta}/& \\\\
\\end{tabular}

\\vspace{20pt}\\hline

\\end{nopagebreak}
\\filbreak



\\vspace{15pt}
\\begin{nopagebreak}
\\noindent{\\fontsize{20pt}{10pt}\\textbf{nal} } \\textit{sharing} (N)\\\\
\\noindent {\\tovian \\fontsize{20pt}{10pt} \\textbf{nal} }\\\\
\\noindent /n{\textprimstress}al/\\\\
\\noindent Syllables: nal\\\\
\\noindent lit. joining+possession\\\\


\\noindent History:

\\vspace{-0pt}
\\hspace{40pt}
\\begin{tabular}{ccc}
\\textit{500} & /{\texttoptiebar{t\textbeltl}}unaelis/&$\\rightarrow$ & \\textit{3200} & /{\texttoptiebar{t\textbeltl}}naels/&$\\rightarrow$ & \\textit{6501} & /naels/&$\\rightarrow$ & \\textit{10000} & /nals/&$\\rightarrow$ & \\textit{13001} & /nal/& \\\\
\\end{tabular}

\\vspace{20pt}\\hline

\\end{nopagebreak}
\\filbreak



\\vspace{15pt}
\\begin{nopagebreak}
\\noindent{\\fontsize{20pt}{10pt}\\textbf{shith} } \\textit{sharpness/alertness/wakefulness/cutting} (N)\\\\
\\noindent {\\tovian \\fontsize{20pt}{10pt} \\textbf{shith} }\\\\
\\noindent /{\textesh}{\textprimstress}i{\texttheta}/\\\\
\\noindent Syllables: ʃiθ\\\\


\\noindent History:

\\vspace{-0pt}
\\hspace{40pt}
\\begin{tabular}{ccc}
\\textit{0} & /{\textyogh}i{\texttheta}/&$\\rightarrow$ & \\textit{12001} & /{\textesh}i{\texttheta}/& \\\\
\\end{tabular}

\\vspace{20pt}\\hline

\\end{nopagebreak}
\\filbreak



\\vspace{15pt}
\\begin{nopagebreak}
\\noindent{\\fontsize{20pt}{10pt}\\textbf{munan} } \\textit{shelter/protection/covering} (N)\\\\
\\noindent {\\tovian \\fontsize{20pt}{10pt} \\textbf{munan} }\\\\
\\noindent /m{\textprimstress}unan/\\\\
\\noindent Syllables: mun.an\\\\


\\noindent History:

\\vspace{-0pt}
\\hspace{40pt}
\\begin{tabular}{ccc}
\\textit{0} & /benani/&$\\rightarrow$ & \\textit{3200} & /bnani/&$\\rightarrow$ & \\textit{9500} & /bnan/&$\\rightarrow$ & \\textit{11000} & /mnan/&$\\rightarrow$ & \\textit{11996} & /munan/& \\\\
\\end{tabular}

\\vspace{20pt}\\hline

\\end{nopagebreak}
\\filbreak



\\vspace{15pt}
\\begin{nopagebreak}
\\noindent{\\fontsize{20pt}{10pt}\\textbf{nalor} } \\textit{shine} (N)\\\\
\\noindent {\\tovian \\fontsize{20pt}{10pt} \\textbf{nalor} }\\\\
\\noindent /n{\textprimstress}alor/\\\\
\\noindent Syllables: nal.or\\\\
\\noindent lit. light+action\\\\


\\noindent History:

\\vspace{-0pt}
\\hspace{40pt}
\\begin{tabular}{ccc}
\\textit{5000} & /aanalore/&$\\rightarrow$ & \\textit{5500} & /analore/&$\\rightarrow$ & \\textit{7500} & /nalore/&$\\rightarrow$ & \\textit{9500} & /nalor/& \\\\
\\end{tabular}

\\vspace{20pt}\\hline

\\end{nopagebreak}
\\filbreak



\\vspace{15pt}
\\begin{nopagebreak}
\\noindent{\\fontsize{20pt}{10pt}\\textbf{fed} } \\textit{showing} (N)\\\\
\\noindent {\\tovian \\fontsize{20pt}{10pt} \\textbf{fed} }\\\\
\\noindent /f{\textprimstress}ed/\\\\
\\noindent Syllables: fed\\\\


\\noindent History:

\\vspace{-0pt}
\\hspace{40pt}
\\begin{tabular}{ccc}
\\textit{0} & /fe{\textsubbridge{t}}e/&$\\rightarrow$ & \\textit{2000} & /fede/&$\\rightarrow$ & \\textit{9500} & /fed/& \\\\
\\end{tabular}

\\vspace{20pt}\\hline

\\end{nopagebreak}
\\filbreak



\\vspace{15pt}
\\begin{nopagebreak}
\\noindent{\\fontsize{20pt}{10pt}\\textbf{menan} } \\textit{silver} (N)\\\\
\\noindent {\\tovian \\fontsize{20pt}{10pt} \\textbf{menan} }\\\\
\\noindent /m{\textprimstress}enan/\\\\
\\noindent Syllables: men.an\\\\
\\noindent lit. metal+light\\\\


\\noindent History:

\\vspace{-0pt}
\\hspace{40pt}
\\begin{tabular}{ccc}
\\textit{4000} & /{\dh}nhmernahana/&$\\rightarrow$ & \\textit{5000} & /{\dh}nhmernaana/&$\\rightarrow$ & \\textit{5500} & /{\dh}nhmernana/&$\\rightarrow$ & \\textit{6400} & /{\dh}hmernana/&$\\rightarrow$ & \\textit{7501} & /hmernana/&$\\rightarrow$ & \\textit{7600} & /hmenana/&$\\rightarrow$ & \\textit{9500} & /hmenan/&$\\rightarrow$ & \\textit{13002} & /menan/& \\\\
\\end{tabular}

\\vspace{20pt}\\hline

\\end{nopagebreak}
\\filbreak



\\vspace{15pt}
\\begin{nopagebreak}
\\noindent{\\fontsize{20pt}{10pt}\\textbf{fethiw} } \\textit{singing} (N)\\\\
\\noindent {\\tovian \\fontsize{20pt}{10pt} \\textbf{fethiw} }\\\\
\\noindent /f{\textprimstress}e{\texttheta}iw/\\\\
\\noindent Syllables: feθ.iw\\\\
\\noindent lit. speaking+beauty\\\\


\\noindent History:

\\vspace{-0pt}
\\hspace{40pt}
\\begin{tabular}{ccc}
\\textit{0} & /fe{\textsubbridge{t}}ei{\texttheta}irwen/&$\\rightarrow$ & \\textit{1000} & /f{\textsubbridge{t}}ei{\texttheta}irwen/&$\\rightarrow$ & \\textit{3000} & /fei{\texttheta}irwen/&$\\rightarrow$ & \\textit{3200} & /fei{\texttheta}irwn/&$\\rightarrow$ & \\textit{5500} & /fe{\texttheta}irwn/&$\\rightarrow$ & \\textit{7600} & /fe{\texttheta}iwn/&$\\rightarrow$ & \\textit{13001} & /fe{\texttheta}iw/& \\\\
\\end{tabular}

\\vspace{20pt}\\hline

\\end{nopagebreak}
\\filbreak



\\vspace{15pt}
\\begin{nopagebreak}
\\noindent{\\fontsize{20pt}{10pt}\\textbf{rom} } \\textit{sitting} (N)\\\\
\\noindent {\\tovian \\fontsize{20pt}{10pt} \\textbf{rom} }\\\\
\\noindent /r{\textprimstress}om/\\\\
\\noindent Syllables: rom\\\\


\\noindent History:

\\vspace{-0pt}
\\hspace{40pt}
\\begin{tabular}{ccc}
\\textit{0} & /roma{\textsubbridge{t}}/&$\\rightarrow$ & \\textit{3200} & /rom{\textsubbridge{t}}/&$\\rightarrow$ & \\textit{8750} & /rom/& \\\\
\\end{tabular}

\\vspace{20pt}\\hline

\\end{nopagebreak}
\\filbreak



\\vspace{15pt}
\\begin{nopagebreak}
\\noindent{\\fontsize{20pt}{10pt}\\textbf{tem} } \\textit{six/6} (N)\\\\
\\noindent {\\tovian \\fontsize{20pt}{10pt} \\textbf{tem} }\\\\
\\noindent /{\textsubbridge{t}}{\textprimstress}em/\\\\
\\noindent Syllables: tem\\\\


\\noindent History:

\\vspace{-0pt}
\\hspace{40pt}
\\begin{tabular}{ccc}
\\textit{0} & /{\textsubbridge{t}}ep/&$\\rightarrow$ & \\textit{11000} & /{\textsubbridge{t}}em/& \\\\
\\end{tabular}

\\vspace{20pt}\\hline

\\end{nopagebreak}
\\filbreak



\\vspace{15pt}
\\begin{nopagebreak}
\\noindent{\\fontsize{20pt}{10pt}\\textbf{lheyatev} } \\textit{sixteen/16} (N)\\\\
\\noindent {\\tovian \\fontsize{20pt}{10pt} \\textbf{lheyatev} }\\\\
\\noindent /{\textbeltl}ey{\textprimstress}a{\textsubbridge{t}}ev/\\\\
\\noindent Syllables: ɬej.at.ev\\\\
\\noindent lit. twelve+four\\\\


\\noindent History:

\\vspace{-0pt}
\\hspace{40pt}
\\begin{tabular}{ccc}
\\textit{4000} & /{\textbeltl}ya{\textsubbridge{t}}ev/&$\\rightarrow$ & \\textit{15001} & /{\textbeltl}eya{\textsubbridge{t}}ev/& \\\\
\\end{tabular}

\\vspace{20pt}\\hline

\\end{nopagebreak}
\\filbreak



\\vspace{15pt}
\\begin{nopagebreak}
\\noindent{\\fontsize{20pt}{10pt}\\textbf{rithil} } \\textit{sky} (N)\\\\
\\noindent {\\tovian \\fontsize{20pt}{10pt} \\textbf{rithil} }\\\\
\\noindent /r{\textprimstress}i{\texttheta}il/\\\\
\\noindent Syllables: riθ.il\\\\
\\noindent lit. air+place\\\\


\\noindent History:

\\vspace{-0pt}
\\hspace{40pt}
\\begin{tabular}{ccc}
\\textit{10} & /rai{\texttheta}{\texttheta}ile/&$\\rightarrow$ & \\textit{4502} & /rai{\texttheta}ile/&$\\rightarrow$ & \\textit{5500} & /ri{\texttheta}ile/&$\\rightarrow$ & \\textit{9500} & /ri{\texttheta}il/& \\\\
\\end{tabular}

\\vspace{20pt}\\hline

\\end{nopagebreak}
\\filbreak



\\vspace{15pt}
\\begin{nopagebreak}
\\noindent{\\fontsize{20pt}{10pt}\\textbf{nin} } \\textit{small} (N)\\\\
\\noindent {\\tovian \\fontsize{20pt}{10pt} \\textbf{nin} }\\\\
\\noindent /n{\textprimstress}in/\\\\
\\noindent Syllables: nin\\\\


\\noindent History:

\\vspace{-0pt}
\\hspace{40pt}
\\begin{tabular}{ccc}
\\textit{0} & /nina/&$\\rightarrow$ & \\textit{9500} & /nin/& \\\\
\\end{tabular}

\\vspace{20pt}\\hline

\\end{nopagebreak}
\\filbreak



\\vspace{15pt}
\\begin{nopagebreak}
\\noindent{\\fontsize{20pt}{10pt}\\textbf{rithor} } \\textit{smell} (N)\\\\
\\noindent {\\tovian \\fontsize{20pt}{10pt} \\textbf{rithor} }\\\\
\\noindent /r{\textprimstress}i{\texttheta}or/\\\\
\\noindent Syllables: riθ.or\\\\
\\noindent lit. air+action\\\\


\\noindent History:

\\vspace{-0pt}
\\hspace{40pt}
\\begin{tabular}{ccc}
\\textit{500} & /rai{\texttheta}lore/&$\\rightarrow$ & \\textit{5500} & /ri{\texttheta}lore/&$\\rightarrow$ & \\textit{9500} & /ri{\texttheta}lor/&$\\rightarrow$ & \\textit{12500} & /ri{\texttheta}or/& \\\\
\\end{tabular}

\\vspace{20pt}\\hline

\\end{nopagebreak}
\\filbreak



\\vspace{15pt}
\\begin{nopagebreak}
\\noindent{\\fontsize{20pt}{10pt}\\textbf{maon} } \\textit{smile} (N)\\\\
\\noindent {\\tovian \\fontsize{20pt}{10pt} \\textbf{maon} }\\\\
\\noindent /m{\textprimstress}aon/\\\\
\\noindent Syllables: ma.on\\\\
\\noindent lit. happy+mouth\\\\


\\noindent History:

\\vspace{-0pt}
\\hspace{40pt}
\\begin{tabular}{ccc}
\\textit{400} & /ma{\textsubbridge{t}}{\textsubbridge{t}}aeolan/&$\\rightarrow$ & \\textit{3000} & /ma{\textsubbridge{t}}aeolan/&$\\rightarrow$ & \\textit{3200} & /m{\textsubbridge{t}}aeoln/&$\\rightarrow$ & \\textit{7600} & /m{\textsubbridge{t}}aeon/&$\\rightarrow$ & \\textit{8750} & /maeon/&$\\rightarrow$ & \\textit{10000} & /maon/& \\\\
\\end{tabular}

\\vspace{20pt}\\hline

\\end{nopagebreak}
\\filbreak



\\vspace{15pt}
\\begin{nopagebreak}
\\noindent{\\fontsize{20pt}{10pt}\\textbf{lheth} } \\textit{softness/gentleness/smoothness} (N)\\\\
\\noindent {\\tovian \\fontsize{20pt}{10pt} \\textbf{lheth} }\\\\
\\noindent /{\textbeltl}{\textprimstress}e{\texttheta}/\\\\
\\noindent Syllables: ɬeθ\\\\


\\noindent History:

\\vspace{-0pt}
\\hspace{40pt}
\\begin{tabular}{ccc}
\\textit{0} & /{\textbeltl}e{\texttheta}/& \\\\
\\end{tabular}

\\vspace{20pt}\\hline

\\end{nopagebreak}
\\filbreak



\\vspace{15pt}
\\begin{nopagebreak}
\\noindent{\\fontsize{20pt}{10pt}\\textbf{tesh} } \\textit{sound} (N)\\\\
\\noindent {\\tovian \\fontsize{20pt}{10pt} \\textbf{tesh} }\\\\
\\noindent /{\textsubbridge{t}}{\textprimstress}e{\textesh}/\\\\
\\noindent Syllables: teʃ\\\\


\\noindent History:

\\vspace{-0pt}
\\hspace{40pt}
\\begin{tabular}{ccc}
\\textit{0} & /{\textsubbridge{t}}e{\textsubbridge{t}}e{\textyogh}ah/&$\\rightarrow$ & \\textit{1000} & /{\textsubbridge{t}}{\textsubbridge{t}}e{\textyogh}ah/&$\\rightarrow$ & \\textit{3000} & /{\textsubbridge{t}}e{\textyogh}ah/&$\\rightarrow$ & \\textit{3200} & /{\textsubbridge{t}}e{\textyogh}h/&$\\rightarrow$ & \\textit{4502} & /{\textsubbridge{t}}e{\textyogh}/&$\\rightarrow$ & \\textit{12001} & /{\textsubbridge{t}}e{\textesh}/& \\\\
\\end{tabular}

\\vspace{20pt}\\hline

\\end{nopagebreak}
\\filbreak



\\vspace{15pt}
\\begin{nopagebreak}
\\noindent{\\fontsize{20pt}{10pt}\\textbf{rar} } \\textit{south} (N)\\\\
\\noindent {\\tovian \\fontsize{20pt}{10pt} \\textbf{rar} }\\\\
\\noindent /r{\textprimstress}ar/\\\\
\\noindent Syllables: rar\\\\


\\noindent History:

\\vspace{-0pt}
\\hspace{40pt}
\\begin{tabular}{ccc}
\\textit{0} & /rar/& \\\\
\\end{tabular}

\\vspace{20pt}\\hline

\\end{nopagebreak}
\\filbreak



\\vspace{15pt}
\\begin{nopagebreak}
\\noindent{\\fontsize{20pt}{10pt}\\textbf{fethr} } \\textit{speaking} (N)\\\\
\\noindent {\\tovian \\fontsize{20pt}{10pt} \\textbf{fethr} }\\\\
\\noindent /f{\textprimstress}e{\texttheta}r/\\\\
\\noindent Syllables: feθr\\\\
\\noindent lit. showing+thought\\\\


\\noindent History:

\\vspace{-0pt}
\\hspace{40pt}
\\begin{tabular}{ccc}
\\textit{100} & /fe{\textsubbridge{t}}ei{\texttheta}ir/&$\\rightarrow$ & \\textit{1000} & /f{\textsubbridge{t}}ei{\texttheta}ir/&$\\rightarrow$ & \\textit{3000} & /fei{\texttheta}ir/&$\\rightarrow$ & \\textit{3200} & /fei{\texttheta}r/&$\\rightarrow$ & \\textit{5500} & /fe{\texttheta}r/& \\\\
\\end{tabular}

\\vspace{20pt}\\hline

\\end{nopagebreak}
\\filbreak



\\vspace{15pt}
\\begin{nopagebreak}
\\noindent{\\fontsize{20pt}{10pt}\\textbf{fethr} } \\textit{speech} (N)\\\\
\\noindent {\\tovian \\fontsize{20pt}{10pt} \\textbf{fethr} }\\\\
\\noindent /f{\textprimstress}e{\texttheta}r/\\\\
\\noindent Syllables: feθr\\\\


\\noindent History:

\\vspace{-0pt}
\\hspace{40pt}
\\begin{tabular}{ccc}
\\textit{0} & /fe{\dh}ra/&$\\rightarrow$ & \\textit{9500} & /fe{\dh}r/&$\\rightarrow$ & \\textit{12002} & /fe{\texttheta}r/& \\\\
\\end{tabular}

\\vspace{20pt}\\hline

\\end{nopagebreak}
\\filbreak



\\vspace{15pt}
\\begin{nopagebreak}
\\noindent{\\fontsize{20pt}{10pt}\\textbf{tlesh} } \\textit{speed} (N)\\\\
\\noindent {\\tovian \\fontsize{20pt}{10pt} \\textbf{tlesh} }\\\\
\\noindent /{\texttoptiebar{t\textbeltl}}{\textprimstress}e{\textesh}/\\\\
\\noindent Syllables: ƛeʃ\\\\


\\noindent History:

\\vspace{-0pt}
\\hspace{40pt}
\\begin{tabular}{ccc}
\\textit{0} & /{\texttoptiebar{t\textbeltl}}e{\textyogh}/&$\\rightarrow$ & \\textit{12001} & /{\texttoptiebar{t\textbeltl}}e{\textesh}/& \\\\
\\end{tabular}

\\vspace{20pt}\\hline

\\end{nopagebreak}
\\filbreak



\\vspace{15pt}
\\begin{nopagebreak}
\\noindent{\\fontsize{20pt}{10pt}\\textbf{mulhelor} } \\textit{Spring} (N)\\\\
\\noindent {\\tovian \\fontsize{20pt}{10pt} \\textbf{mulhelor} }\\\\
\\noindent /mu{\textbeltl}{\textprimstress}elor/\\\\
\\noindent Syllables: muɬ.el.or\\\\
\\noindent lit. present+plant\\\\


\\noindent History:

\\vspace{-0pt}
\\hspace{40pt}
\\begin{tabular}{ccc}
\\textit{0} & /num{\textbeltl}ele{\textschwa}mellore/&$\\rightarrow$ & \\textit{2300} & /num{\textbeltl}elemellore/&$\\rightarrow$ & \\textit{3200} & /nm{\textbeltl}lemllore/&$\\rightarrow$ & \\textit{6500} & /nm{\textbeltl}lemlore/&$\\rightarrow$ & \\textit{7501} & /m{\textbeltl}lemlore/&$\\rightarrow$ & \\textit{7600} & /m{\textbeltl}lelore/&$\\rightarrow$ & \\textit{9500} & /m{\textbeltl}lelor/&$\\rightarrow$ & \\textit{12500} & /m{\textbeltl}elor/&$\\rightarrow$ & \\textit{15001} & /mu{\textbeltl}elor/& \\\\
\\end{tabular}

\\vspace{20pt}\\hline

\\end{nopagebreak}
\\filbreak



\\vspace{15pt}
\\begin{nopagebreak}
\\noindent{\\fontsize{20pt}{10pt}\\textbf{ledul} } \\textit{spring vegetable} (N)\\\\
\\noindent {\\tovian \\fontsize{20pt}{10pt} \\textbf{ledul} }\\\\
\\noindent /l{\textprimstress}edul/\\\\
\\noindent Syllables: led.ul\\\\
\\noindent lit. plant+first\\\\


\\noindent History:

\\vspace{-0pt}
\\hspace{40pt}
\\begin{tabular}{ccc}
\\textit{100} & /{\textschwa}mellorepu-lo/&$\\rightarrow$ & \\textit{1000} & /{\textschwa}mellorepulo/&$\\rightarrow$ & \\textit{2000} & /{\textschwa}mellorebulo/&$\\rightarrow$ & \\textit{2300} & /mellorebulo/&$\\rightarrow$ & \\textit{3200} & /mllrebulo/&$\\rightarrow$ & \\textit{6500} & /mlrebulo/&$\\rightarrow$ & \\textit{6503} & /mrlebulo/&$\\rightarrow$ & \\textit{7501} & /rlebulo/&$\\rightarrow$ & \\textit{7600} & /lebulo/&$\\rightarrow$ & \\textit{9500} & /lebul/&$\\rightarrow$ & \\textit{11001} & /ledul/& \\\\
\\end{tabular}

\\vspace{20pt}\\hline

\\end{nopagebreak}
\\filbreak



\\vspace{15pt}
\\begin{nopagebreak}
\\noindent{\\fontsize{20pt}{10pt}\\textbf{ren} } \\textit{standing} (N)\\\\
\\noindent {\\tovian \\fontsize{20pt}{10pt} \\textbf{ren} }\\\\
\\noindent /r{\textprimstress}en/\\\\
\\noindent Syllables: ren\\\\


\\noindent History:

\\vspace{-0pt}
\\hspace{40pt}
\\begin{tabular}{ccc}
\\textit{0} & /reman/&$\\rightarrow$ & \\textit{3200} & /remn/&$\\rightarrow$ & \\textit{7600} & /ren/& \\\\
\\end{tabular}

\\vspace{20pt}\\hline

\\end{nopagebreak}
\\filbreak



\\vspace{15pt}
\\begin{nopagebreak}
\\noindent{\\fontsize{20pt}{10pt}\\textbf{nan} } \\textit{star} (N)\\\\
\\noindent {\\tovian \\fontsize{20pt}{10pt} \\textbf{nan} }\\\\
\\noindent /n{\textprimstress}an/\\\\
\\noindent Syllables: nan\\\\
\\noindent lit. small+light\\\\


\\noindent History:

\\vspace{-0pt}
\\hspace{40pt}
\\begin{tabular}{ccc}
\\textit{500} & /ninaahana/&$\\rightarrow$ & \\textit{3200} & /nnaahana/&$\\rightarrow$ & \\textit{5000} & /nnaaana/&$\\rightarrow$ & \\textit{5500} & /nnaana/&$\\rightarrow$ & \\textit{6500} & /naana/&$\\rightarrow$ & \\textit{9500} & /naan/&$\\rightarrow$ & \\textit{12005} & /nan/& \\\\
\\end{tabular}

\\vspace{20pt}\\hline

\\end{nopagebreak}
\\filbreak



\\vspace{15pt}
\\begin{nopagebreak}
\\noindent{\\fontsize{20pt}{10pt}\\textbf{ing} } \\textit{stickiness} (N)\\\\
\\noindent {\\tovian \\fontsize{20pt}{10pt} \\textbf{ing} }\\\\
\\noindent /{\textprimstress}i{\ng}/\\\\
\\noindent Syllables: iŋ\\\\


\\noindent History:

\\vspace{-0pt}
\\hspace{40pt}
\\begin{tabular}{ccc}
\\textit{0} & /inga/&$\\rightarrow$ & \\textit{6000} & /i{\ng}ga/&$\\rightarrow$ & \\textit{8750} & /i{\ng}a/&$\\rightarrow$ & \\textit{9500} & /i{\ng}/& \\\\
\\end{tabular}

\\vspace{20pt}\\hline

\\end{nopagebreak}
\\filbreak



\\vspace{15pt}
\\begin{nopagebreak}
\\noindent{\\fontsize{20pt}{10pt}\\textbf{lor} } \\textit{stomach} (N)\\\\
\\noindent {\\tovian \\fontsize{20pt}{10pt} \\textbf{lor} }\\\\
\\noindent /l{\textprimstress}or/\\\\
\\noindent Syllables: lor\\\\


\\noindent History:

\\vspace{-0pt}
\\hspace{40pt}
\\begin{tabular}{ccc}
\\textit{0} & /sildabore/&$\\rightarrow$ & \\textit{3200} & /sldbore/&$\\rightarrow$ & \\textit{4501} & /slbore/&$\\rightarrow$ & \\textit{7501} & /lbore/&$\\rightarrow$ & \\textit{8750} & /lore/&$\\rightarrow$ & \\textit{9500} & /lor/& \\\\
\\end{tabular}

\\vspace{20pt}\\hline

\\end{nopagebreak}
\\filbreak



\\vspace{15pt}
\\begin{nopagebreak}
\\noindent{\\fontsize{20pt}{10pt}\\textbf{men} } \\textit{stone/rock/earth} (N)\\\\
\\noindent {\\tovian \\fontsize{20pt}{10pt} \\textbf{men} }\\\\
\\noindent /m{\textprimstress}en/\\\\
\\noindent Syllables: men\\\\


\\noindent History:

\\vspace{-0pt}
\\hspace{40pt}
\\begin{tabular}{ccc}
\\textit{0} & /meron/&$\\rightarrow$ & \\textit{3200} & /mern/&$\\rightarrow$ & \\textit{7600} & /men/& \\\\
\\end{tabular}

\\vspace{20pt}\\hline

\\end{nopagebreak}
\\filbreak



\\vspace{15pt}
\\begin{nopagebreak}
\\noindent{\\fontsize{20pt}{10pt}\\textbf{pal} } \\textit{stopping/ceasing} (N)\\\\
\\noindent {\\tovian \\fontsize{20pt}{10pt} \\textbf{pal} }\\\\
\\noindent /p{\textprimstress}al/\\\\
\\noindent Syllables: pal\\\\


\\noindent History:

\\vspace{-0pt}
\\hspace{40pt}
\\begin{tabular}{ccc}
\\textit{0} & /palor/&$\\rightarrow$ & \\textit{3200} & /palr/&$\\rightarrow$ & \\textit{6503} & /parl/&$\\rightarrow$ & \\textit{7600} & /pal/& \\\\
\\end{tabular}

\\vspace{20pt}\\hline

\\end{nopagebreak}
\\filbreak



\\vspace{15pt}
\\begin{nopagebreak}
\\noindent{\\fontsize{20pt}{10pt}\\textbf{lerith} } \\textit{storm} (N)\\\\
\\noindent {\\tovian \\fontsize{20pt}{10pt} \\textbf{lerith} }\\\\
\\noindent /l{\textprimstress}eri{\texttheta}/\\\\
\\noindent Syllables: ler.iθ\\\\
\\noindent lit. wind+water\\\\


\\noindent History:

\\vspace{-0pt}
\\hspace{40pt}
\\begin{tabular}{ccc}
\\textit{50} & /kwilyi-rai{\texttheta}z{\textsubbridge{t}}el/&$\\rightarrow$ & \\textit{1000} & /kwilyirai{\texttheta}z{\textsubbridge{t}}el/&$\\rightarrow$ & \\textit{2000} & /kwilyirai{\texttheta}zdel/&$\\rightarrow$ & \\textit{3200} & /kwlyrai{\texttheta}zdl/&$\\rightarrow$ & \\textit{4500} & /kwlyrai{\texttheta}zl/&$\\rightarrow$ & \\textit{4502} & /kwlyrai{\texttheta}l/&$\\rightarrow$ & \\textit{5500} & /kwlyri{\texttheta}l/&$\\rightarrow$ & \\textit{7501} & /wlyri{\texttheta}l/&$\\rightarrow$ & \\textit{7600} & /lri{\texttheta}l/&$\\rightarrow$ & \\textit{12500} & /lri{\texttheta}/&$\\rightarrow$ & \\textit{15001} & /leri{\texttheta}/& \\\\
\\end{tabular}

\\vspace{20pt}\\hline

\\end{nopagebreak}
\\filbreak



\\vspace{15pt}
\\begin{nopagebreak}
\\noindent{\\fontsize{20pt}{10pt}\\textbf{fethim} } \\textit{story} (N)\\\\
\\noindent {\\tovian \\fontsize{20pt}{10pt} \\textbf{fethim} }\\\\
\\noindent /f{\textprimstress}e{\texttheta}im/\\\\
\\noindent Syllables: feθ.im\\\\
\\noindent lit. speaking+life\\\\


\\noindent History:

\\vspace{-0pt}
\\hspace{40pt}
\\begin{tabular}{ccc}
\\textit{200} & /fe{\textsubbridge{t}}ei{\texttheta}ir{\textschwa}mel/&$\\rightarrow$ & \\textit{1000} & /f{\textsubbridge{t}}ei{\texttheta}ir{\textschwa}mel/&$\\rightarrow$ & \\textit{2300} & /f{\textsubbridge{t}}ei{\texttheta}irmel/&$\\rightarrow$ & \\textit{3000} & /fei{\texttheta}irmel/&$\\rightarrow$ & \\textit{3200} & /fei{\texttheta}irml/&$\\rightarrow$ & \\textit{5500} & /fe{\texttheta}irml/&$\\rightarrow$ & \\textit{7600} & /fe{\texttheta}iml/&$\\rightarrow$ & \\textit{13001} & /fe{\texttheta}im/& \\\\
\\end{tabular}

\\vspace{20pt}\\hline

\\end{nopagebreak}
\\filbreak



\\vspace{15pt}
\\begin{nopagebreak}
\\noindent{\\fontsize{20pt}{10pt}\\textbf{wi-} } \\textit{SUB/subjunctive mood marker} (MOOD)\\\\
\\noindent {\\tovian \\fontsize{20pt}{10pt} \\textbf{wi-} }\\\\
\\noindent /w{\textprimstress}i/\\\\
\\noindent Syllables: wi\\\\


\\noindent History:

\\vspace{-0pt}
\\hspace{40pt}
\\begin{tabular}{ccc}
\\textit{0} & /wi-/& \\\\
\\end{tabular}

\\vspace{20pt}\\hline

\\end{nopagebreak}
\\filbreak



\\vspace{15pt}
\\begin{nopagebreak}
\\noindent{\\fontsize{20pt}{10pt}\\textbf{mulheag} } \\textit{Summer} (N)\\\\
\\noindent {\\tovian \\fontsize{20pt}{10pt} \\textbf{mulheag} }\\\\
\\noindent /mu{\textbeltl}{\textprimstress}eag/\\\\
\\noindent Syllables: muɬ.e.ag\\\\
\\noindent lit. present+heat\\\\


\\noindent History:

\\vspace{-0pt}
\\hspace{40pt}
\\begin{tabular}{ccc}
\\textit{0} & /num{\textbeltl}eleasaka/&$\\rightarrow$ & \\textit{2000} & /num{\textbeltl}eleasaga/&$\\rightarrow$ & \\textit{3200} & /nm{\textbeltl}leasaga/&$\\rightarrow$ & \\textit{5100} & /nm{\textbeltl}leaaga/&$\\rightarrow$ & \\textit{5500} & /nm{\textbeltl}leaga/&$\\rightarrow$ & \\textit{7501} & /m{\textbeltl}leaga/&$\\rightarrow$ & \\textit{9500} & /m{\textbeltl}leag/&$\\rightarrow$ & \\textit{12500} & /m{\textbeltl}eag/&$\\rightarrow$ & \\textit{15001} & /mu{\textbeltl}eag/& \\\\
\\end{tabular}

\\vspace{20pt}\\hline

\\end{nopagebreak}
\\filbreak



\\vspace{15pt}
\\begin{nopagebreak}
\\noindent{\\fontsize{20pt}{10pt}\\textbf{ef} } \\textit{sun} (N)\\\\
\\noindent {\\tovian \\fontsize{20pt}{10pt} \\textbf{ef} }\\\\
\\noindent /{\textprimstress}ef/\\\\
\\noindent Syllables: ef\\\\


\\noindent History:

\\vspace{-0pt}
\\hspace{40pt}
\\begin{tabular}{ccc}
\\textit{0} & /eifol/&$\\rightarrow$ & \\textit{3200} & /eifl/&$\\rightarrow$ & \\textit{5500} & /efl/&$\\rightarrow$ & \\textit{12500} & /ef/& \\\\
\\end{tabular}

\\vspace{20pt}\\hline

\\end{nopagebreak}
\\filbreak



\\vspace{15pt}
\\begin{nopagebreak}
\\noindent{\\fontsize{20pt}{10pt}\\textbf{efwyrith} } \\textit{sunshower} (N)\\\\
\\noindent {\\tovian \\fontsize{20pt}{10pt} \\textbf{efwyrith} }\\\\
\\noindent /{\textprimstress}efwyri{\texttheta}/\\\\
\\noindent Syllables: efwjr.iθ\\\\
\\noindent lit. sun+storm\\\\


\\noindent History:

\\vspace{-0pt}
\\hspace{40pt}
\\begin{tabular}{ccc}
\\textit{1200} & /eifolkwilyirai{\texttheta}z{\textsubbridge{t}}el/&$\\rightarrow$ & \\textit{2000} & /eifolkwilyirai{\texttheta}zdel/&$\\rightarrow$ & \\textit{3000} & /eifolwilyirai{\texttheta}zdel/&$\\rightarrow$ & \\textit{3200} & /eiflwlyrai{\texttheta}zdl/&$\\rightarrow$ & \\textit{4500} & /eiflwlyrai{\texttheta}zl/&$\\rightarrow$ & \\textit{4502} & /eiflwlyrai{\texttheta}l/&$\\rightarrow$ & \\textit{5500} & /eflwlyri{\texttheta}l/&$\\rightarrow$ & \\textit{7600} & /efwyri{\texttheta}l/&$\\rightarrow$ & \\textit{12500} & /efwyri{\texttheta}/& \\\\
\\end{tabular}

\\vspace{20pt}\\hline

\\end{nopagebreak}
\\filbreak



\\vspace{15pt}
\\begin{nopagebreak}
\\noindent{\\fontsize{20pt}{10pt}\\textbf{ngam} } \\textit{surpassing/exceeding/overwhelming} (N)\\\\
\\noindent {\\tovian \\fontsize{20pt}{10pt} \\textbf{ngam} }\\\\
\\noindent /{\ng}{\textprimstress}am/\\\\
\\noindent Syllables: ŋam\\\\


\\noindent History:

\\vspace{-0pt}
\\hspace{40pt}
\\begin{tabular}{ccc}
\\textit{0} & /{\ng}o{\ng}abi/&$\\rightarrow$ & \\textit{3200} & /{\ng}{\ng}abi/&$\\rightarrow$ & \\textit{6500} & /{\ng}abi/&$\\rightarrow$ & \\textit{9500} & /{\ng}ab/&$\\rightarrow$ & \\textit{11000} & /{\ng}am/& \\\\
\\end{tabular}

\\vspace{20pt}\\hline

\\end{nopagebreak}
\\filbreak



\\vspace{15pt}
\\begin{nopagebreak}
\\noindent{\\fontsize{20pt}{10pt}\\textbf{hol} } \\textit{sweet} (N)\\\\
\\noindent {\\tovian \\fontsize{20pt}{10pt} \\textbf{hol} }\\\\
\\noindent /ħ{\textprimstress}ol/\\\\
\\noindent Syllables: ħol\\\\
\\noindent lit. GEN+fruit\\\\


\\noindent History:

\\vspace{-0pt}
\\hspace{40pt}
\\begin{tabular}{ccc}
\\textit{400} & /hi-{\textyogh}ola/&$\\rightarrow$ & \\textit{1000} & /hi{\textyogh}ola/&$\\rightarrow$ & \\textit{3200} & /h{\textyogh}ola/&$\\rightarrow$ & \\textit{4502} & /hola/&$\\rightarrow$ & \\textit{9500} & /hol/&$\\rightarrow$ & \\textit{14006} & /ħol/& \\\\
\\end{tabular}

\\vspace{20pt}\\hline

\\end{nopagebreak}
\\filbreak



\\vspace{15pt}
\\begin{nopagebreak}
\\noindent{\\fontsize{20pt}{10pt}\\textbf{thuv} } \\textit{sword} (N)\\\\
\\noindent {\\tovian \\fontsize{20pt}{10pt} \\textbf{thuv} }\\\\
\\noindent /{\texttheta}{\textprimstress}uv/\\\\
\\noindent Syllables: θuv\\\\
\\noindent lit. tool+PUR+killing\\\\


\\noindent History:

\\vspace{-0pt}
\\hspace{40pt}
\\begin{tabular}{ccc}
\\textit{50} & /{\dh}an{\textbeltl}u-vekh/&$\\rightarrow$ & \\textit{1000} & /{\dh}an{\textbeltl}uvekh/&$\\rightarrow$ & \\textit{3200} & /{\dh}n{\textbeltl}uvkh/&$\\rightarrow$ & \\textit{4500} & /{\dh}n{\textbeltl}uvh/&$\\rightarrow$ & \\textit{4502} & /{\dh}n{\textbeltl}uv/&$\\rightarrow$ & \\textit{6400} & /{\dh}{\textbeltl}uv/&$\\rightarrow$ & \\textit{12002} & /{\texttheta}{\textbeltl}uv/&$\\rightarrow$ & \\textit{13000} & /{\texttheta}uv/& \\\\
\\end{tabular}

\\vspace{20pt}\\hline

\\end{nopagebreak}
\\filbreak



\\vspace{15pt}
\\begin{nopagebreak}
\\noindent{\\fontsize{20pt}{10pt}\\textbf{tlun} } \\textit{tar} (N)\\\\
\\noindent {\\tovian \\fontsize{20pt}{10pt} \\textbf{tlun} }\\\\
\\noindent /{\texttoptiebar{t\textbeltl}}{\textprimstress}un/\\\\
\\noindent Syllables: ƛun\\\\
\\noindent lit. stickiness+joining\\\\


\\noindent History:

\\vspace{-0pt}
\\hspace{40pt}
\\begin{tabular}{ccc}
\\textit{50} & /inga{\texttoptiebar{t\textbeltl}}una/&$\\rightarrow$ & \\textit{3200} & /ing{\texttoptiebar{t\textbeltl}}una/&$\\rightarrow$ & \\textit{6000} & /i{\ng}g{\texttoptiebar{t\textbeltl}}una/&$\\rightarrow$ & \\textit{7500} & /{\ng}g{\texttoptiebar{t\textbeltl}}una/&$\\rightarrow$ & \\textit{7501} & /g{\texttoptiebar{t\textbeltl}}una/&$\\rightarrow$ & \\textit{9500} & /g{\texttoptiebar{t\textbeltl}}un/&$\\rightarrow$ & \\textit{11501} & /{\texttoptiebar{t\textbeltl}}un/& \\\\
\\end{tabular}

\\vspace{20pt}\\hline

\\end{nopagebreak}
\\filbreak



\\vspace{15pt}
\\begin{nopagebreak}
\\noindent{\\fontsize{20pt}{10pt}\\textbf{lhad} } \\textit{teaching} (N)\\\\
\\noindent {\\tovian \\fontsize{20pt}{10pt} \\textbf{lhad} }\\\\
\\noindent /{\textbeltl}{\textprimstress}ad/\\\\
\\noindent Syllables: ɬad\\\\
\\noindent lit. knowledge+showing\\\\


\\noindent History:

\\vspace{-0pt}
\\hspace{40pt}
\\begin{tabular}{ccc}
\\textit{0} & /{\textbeltl}arafe{\textsubbridge{t}}e/&$\\rightarrow$ & \\textit{2000} & /{\textbeltl}arafede/&$\\rightarrow$ & \\textit{3200} & /{\textbeltl}rafede/&$\\rightarrow$ & \\textit{9500} & /{\textbeltl}rafed/&$\\rightarrow$ & \\textit{12003} & /{\textbeltl}rad/&$\\rightarrow$ & \\textit{12500} & /{\textbeltl}ad/& \\\\
\\end{tabular}

\\vspace{20pt}\\hline

\\end{nopagebreak}
\\filbreak



\\vspace{15pt}
\\begin{nopagebreak}
\\noindent{\\fontsize{20pt}{10pt}\\textbf{kogyal} } \\textit{tease} (N)\\\\
\\noindent {\\tovian \\fontsize{20pt}{10pt} \\textbf{kogyal} }\\\\
\\noindent /k{\textprimstress}ogyal/\\\\
\\noindent Syllables: kogj.al\\\\
\\noindent lit. playing+evil\\\\


\\noindent History:

\\vspace{-0pt}
\\hspace{40pt}
\\begin{tabular}{ccc}
\\textit{4700} & /kogiyal/&$\\rightarrow$ & \\textit{6300} & /kogyal/& \\\\
\\end{tabular}

\\vspace{20pt}\\hline

\\end{nopagebreak}
\\filbreak



\\vspace{15pt}
\\begin{nopagebreak}
\\noindent{\\fontsize{20pt}{10pt}\\textbf{tu-} } \\textit{TEMP/temporal case marker} (CASE)\\\\
\\noindent {\\tovian \\fontsize{20pt}{10pt} \\textbf{tu-} }\\\\
\\noindent /{\textsubbridge{t}}{\textprimstress}u/\\\\
\\noindent Syllables: tu\\\\


\\noindent History:

\\vspace{-0pt}
\\hspace{40pt}
\\begin{tabular}{ccc}
\\textit{0} & /{\textsubbridge{t}}u-/& \\\\
\\end{tabular}

\\vspace{20pt}\\hline

\\end{nopagebreak}
\\filbreak



\\vspace{15pt}
\\begin{nopagebreak}
\\noindent{\\fontsize{20pt}{10pt}\\textbf{tir} } \\textit{ten/10} (N)\\\\
\\noindent {\\tovian \\fontsize{20pt}{10pt} \\textbf{tir} }\\\\
\\noindent /{\textsubbridge{t}}{\textprimstress}ir/\\\\
\\noindent Syllables: tir\\\\


\\noindent History:

\\vspace{-0pt}
\\hspace{40pt}
\\begin{tabular}{ccc}
\\textit{0} & /{\textsubbridge{t}}ir/& \\\\
\\end{tabular}

\\vspace{20pt}\\hline

\\end{nopagebreak}
\\filbreak



\\vspace{15pt}
\\begin{nopagebreak}
\\noindent{\\fontsize{20pt}{10pt}\\textbf{hi-ta-e} } \\textit{their} (_)\\\\
\\noindent {\\tovian \\fontsize{20pt}{10pt} \\textbf{hi-ta-e} }\\\\
\\noindent /hi{\textsubbridge{t}}{\textprimstress}ae/\\\\
\\noindent Syllables: hi.ta.e\\\\
\\noindent lit. GEN+they\\\\
\\noindent \\textit{Ps}\\\\


\\noindent History:

\\vspace{-0pt}
\\hspace{40pt}
\\begin{tabular}{ccc}
\\textit{0} & /hi-{\textsubbridge{t}}a-e/& \\\\
\\end{tabular}

\\vspace{20pt}\\hline

\\end{nopagebreak}
\\filbreak



\\vspace{15pt}
\\begin{nopagebreak}
\\noindent{\\fontsize{20pt}{10pt}\\textbf{lod} } \\textit{them (inanimate)} (Ps)\\\\
\\noindent {\\tovian \\fontsize{20pt}{10pt} \\textbf{lod} }\\\\
\\noindent /l{\textprimstress}od/\\\\
\\noindent Syllables: lod\\\\


\\noindent History:

\\vspace{-0pt}
\\hspace{40pt}
\\begin{tabular}{ccc}
\\textit{0} & /lo{\textsubbridge{t}}e/&$\\rightarrow$ & \\textit{2000} & /lode/&$\\rightarrow$ & \\textit{9500} & /lod/& \\\\
\\end{tabular}

\\vspace{20pt}\\hline

\\end{nopagebreak}
\\filbreak



\\vspace{15pt}
\\begin{nopagebreak}
\\noindent{\\fontsize{20pt}{10pt}\\textbf{ta-e} } \\textit{they} (_)\\\\
\\noindent {\\tovian \\fontsize{20pt}{10pt} \\textbf{ta-e} }\\\\
\\noindent /{\textsubbridge{t}}{\textprimstress}ae/\\\\
\\noindent Syllables: ta.e\\\\
\\noindent lit. she+PL\\\\
\\noindent \\textit{Ps}\\\\


\\noindent History:

\\vspace{-0pt}
\\hspace{40pt}
\\begin{tabular}{ccc}
\\textit{0} & /{\textsubbridge{t}}a-e/& \\\\
\\end{tabular}

\\vspace{20pt}\\hline

\\end{nopagebreak}
\\filbreak



\\vspace{15pt}
\\begin{nopagebreak}
\\noindent{\\fontsize{20pt}{10pt}\\textbf{t} } \\textit{they (plural)} (Ps)\\\\
\\noindent {\\tovian \\fontsize{20pt}{10pt} \\textbf{t} }\\\\
\\noindent /{\textsubbridge{t}}/\\\\
\\noindent Syllables: t\\\\


\\noindent History:

\\vspace{-0pt}
\\hspace{40pt}
\\begin{tabular}{ccc}
\\textit{0} & /{\textsubbridge{t}}e/&$\\rightarrow$ & \\textit{12006} & /{\textsubbridge{t}}/& \\\\
\\end{tabular}

\\vspace{20pt}\\hline

\\end{nopagebreak}
\\filbreak



\\vspace{15pt}
\\begin{nopagebreak}
\\noindent{\\fontsize{20pt}{10pt}\\textbf{lhul} } \\textit{thickness/denseness} (N)\\\\
\\noindent {\\tovian \\fontsize{20pt}{10pt} \\textbf{lhul} }\\\\
\\noindent /{\textbeltl}{\textprimstress}ul/\\\\
\\noindent Syllables: ɬul\\\\


\\noindent History:

\\vspace{-0pt}
\\hspace{40pt}
\\begin{tabular}{ccc}
\\textit{0} & /{\textbeltl}unol/&$\\rightarrow$ & \\textit{3200} & /{\textbeltl}unl/&$\\rightarrow$ & \\textit{7600} & /{\textbeltl}ul/& \\\\
\\end{tabular}

\\vspace{20pt}\\hline

\\end{nopagebreak}
\\filbreak



\\vspace{15pt}
\\begin{nopagebreak}
\\noindent{\\fontsize{20pt}{10pt}\\textbf{l} } \\textit{thing} (N)\\\\
\\noindent {\\tovian \\fontsize{20pt}{10pt} \\textbf{l} }\\\\
\\noindent /l/\\\\
\\noindent Syllables: l\\\\


\\noindent History:

\\vspace{-0pt}
\\hspace{40pt}
\\begin{tabular}{ccc}
\\textit{0} & /{\textsubbridge{t}}iel/&$\\rightarrow$ & \\textit{2201} & /{\textsubbridge{t}}iyel/&$\\rightarrow$ & \\textit{3200} & /{\textsubbridge{t}}iyl/&$\\rightarrow$ & \\textit{6300} & /{\textsubbridge{t}}yl/&$\\rightarrow$ & \\textit{7501} & /yl/&$\\rightarrow$ & \\textit{7600} & /l/& \\\\
\\end{tabular}

\\vspace{20pt}\\hline

\\end{nopagebreak}
\\filbreak



\\vspace{15pt}
\\begin{nopagebreak}
\\noindent{\\fontsize{20pt}{10pt}\\textbf{nalh} } \\textit{thinking/believing/pondering} (N)\\\\
\\noindent {\\tovian \\fontsize{20pt}{10pt} \\textbf{nalh} }\\\\
\\noindent /n{\textprimstress}a{\textbeltl}/\\\\
\\noindent Syllables: naɬ\\\\


\\noindent History:

\\vspace{-0pt}
\\hspace{40pt}
\\begin{tabular}{ccc}
\\textit{0} & /na{\textbeltl}/& \\\\
\\end{tabular}

\\vspace{20pt}\\hline

\\end{nopagebreak}
\\filbreak



\\vspace{15pt}
\\begin{nopagebreak}
\\noindent{\\fontsize{20pt}{10pt}\\textbf{lheyal} } \\textit{thirteen/13} (N)\\\\
\\noindent {\\tovian \\fontsize{20pt}{10pt} \\textbf{lheyal} }\\\\
\\noindent /{\textbeltl}{\textprimstress}eyal/\\\\
\\noindent Syllables: ɬej.al\\\\
\\noindent lit. twelve+one\\\\


\\noindent History:

\\vspace{-0pt}
\\hspace{40pt}
\\begin{tabular}{ccc}
\\textit{4000} & /{\textbeltl}yalo/&$\\rightarrow$ & \\textit{9500} & /{\textbeltl}yal/&$\\rightarrow$ & \\textit{15001} & /{\textbeltl}eyal/& \\\\
\\end{tabular}

\\vspace{20pt}\\hline

\\end{nopagebreak}
\\filbreak



\\vspace{15pt}
\\begin{nopagebreak}
\\noindent{\\fontsize{20pt}{10pt}\\textbf{femlh} } \\textit{thirty six/36} (N)\\\\
\\noindent {\\tovian \\fontsize{20pt}{10pt} \\textbf{femlh} }\\\\
\\noindent /f{\textprimstress}em{\textbeltl}/\\\\
\\noindent Syllables: femɬ\\\\
\\noindent lit. three+twelve\\\\


\\noindent History:

\\vspace{-0pt}
\\hspace{40pt}
\\begin{tabular}{ccc}
\\textit{4000} & /fem{\textbeltl}ya/&$\\rightarrow$ & \\textit{9500} & /fem{\textbeltl}y/&$\\rightarrow$ & \\textit{13001} & /fem{\textbeltl}/& \\\\
\\end{tabular}

\\vspace{20pt}\\hline

\\end{nopagebreak}
\\filbreak



\\vspace{15pt}
\\begin{nopagebreak}
\\noindent{\\fontsize{20pt}{10pt}\\textbf{loelhyatem} } \\textit{thirty/30} (N)\\\\
\\noindent {\\tovian \\fontsize{20pt}{10pt} \\textbf{loelhyatem} }\\\\
\\noindent /loe{\textbeltl}y{\textprimstress}a{\textsubbridge{t}}em/\\\\
\\noindent Syllables: lo.eɬj.at.em\\\\
\\noindent lit. twenty four+six\\\\


\\noindent History:

\\vspace{-0pt}
\\hspace{40pt}
\\begin{tabular}{ccc}
\\textit{4000} & /loe{\textbeltl}ya{\textsubbridge{t}}ep/&$\\rightarrow$ & \\textit{11000} & /loe{\textbeltl}ya{\textsubbridge{t}}em/& \\\\
\\end{tabular}

\\vspace{20pt}\\hline

\\end{nopagebreak}
\\filbreak



\\vspace{15pt}
\\begin{nopagebreak}
\\noindent{\\fontsize{20pt}{10pt}\\textbf{ithr} } \\textit{thought} (N)\\\\
\\noindent {\\tovian \\fontsize{20pt}{10pt} \\textbf{ithr} }\\\\
\\noindent /{\textprimstress}i{\texttheta}r/\\\\
\\noindent Syllables: iθr\\\\


\\noindent History:

\\vspace{-0pt}
\\hspace{40pt}
\\begin{tabular}{ccc}
\\textit{0} & /i{\texttheta}ir/&$\\rightarrow$ & \\textit{3200} & /i{\texttheta}r/& \\\\
\\end{tabular}

\\vspace{20pt}\\hline

\\end{nopagebreak}
\\filbreak



\\vspace{15pt}
\\begin{nopagebreak}
\\noindent{\\fontsize{20pt}{10pt}\\textbf{fem} } \\textit{three/3} (N)\\\\
\\noindent {\\tovian \\fontsize{20pt}{10pt} \\textbf{fem} }\\\\
\\noindent /f{\textprimstress}em/\\\\
\\noindent Syllables: fem\\\\


\\noindent History:

\\vspace{-0pt}
\\hspace{40pt}
\\begin{tabular}{ccc}
\\textit{0} & /fem/& \\\\
\\end{tabular}

\\vspace{20pt}\\hline

\\end{nopagebreak}
\\filbreak



\\vspace{15pt}
\\begin{nopagebreak}
\\noindent{\\fontsize{20pt}{10pt}\\textbf{yisha} } \\textit{throwing} (N)\\\\
\\noindent {\\tovian \\fontsize{20pt}{10pt} \\textbf{yisha} }\\\\
\\noindent /y{\textprimstress}i{\textesh}a/\\\\
\\noindent Syllables: jiʃ.a\\\\


\\noindent History:

\\vspace{-0pt}
\\hspace{40pt}
\\begin{tabular}{ccc}
\\textit{0} & /huy{\textyogh}a/&$\\rightarrow$ & \\textit{2100} & /hy{\textyogh}a/&$\\rightarrow$ & \\textit{7501} & /y{\textyogh}a/&$\\rightarrow$ & \\textit{12001} & /y{\textesh}a/&$\\rightarrow$ & \\textit{15001} & /yi{\textesh}a/& \\\\
\\end{tabular}

\\vspace{20pt}\\hline

\\end{nopagebreak}
\\filbreak



\\vspace{15pt}
\\begin{nopagebreak}
\\noindent{\\fontsize{20pt}{10pt}\\textbf{lerithesh} } \\textit{thunder} (N)\\\\
\\noindent {\\tovian \\fontsize{20pt}{10pt} \\textbf{lerithesh} }\\\\
\\noindent /ler{\textprimstress}i{\texttheta}e{\textesh}/\\\\
\\noindent Syllables: ler.iθ.eʃ\\\\
\\noindent lit. storm+sound\\\\


\\noindent History:

\\vspace{-0pt}
\\hspace{40pt}
\\begin{tabular}{ccc}
\\textit{100} & /kwilyi-rai{\texttheta}z{\textsubbridge{t}}el{\textsubbridge{t}}e{\textsubbridge{t}}e{\textyogh}ah/&$\\rightarrow$ & \\textit{1000} & /kwilyirai{\texttheta}z{\textsubbridge{t}}el{\textsubbridge{t}}{\textsubbridge{t}}e{\textyogh}ah/&$\\rightarrow$ & \\textit{2000} & /kwilyirai{\texttheta}zdel{\textsubbridge{t}}{\textsubbridge{t}}e{\textyogh}ah/&$\\rightarrow$ & \\textit{3000} & /kwilyirai{\texttheta}zdel{\textsubbridge{t}}e{\textyogh}ah/&$\\rightarrow$ & \\textit{3200} & /kwlyrai{\texttheta}zdl{\textsubbridge{t}}e{\textyogh}h/&$\\rightarrow$ & \\textit{4500} & /kwlyrai{\texttheta}zl{\textsubbridge{t}}e{\textyogh}h/&$\\rightarrow$ & \\textit{4501} & /kwlyrai{\texttheta}zle{\textyogh}h/&$\\rightarrow$ & \\textit{4502} & /kwlyrai{\texttheta}le{\textyogh}/&$\\rightarrow$ & \\textit{5500} & /kwlyri{\texttheta}le{\textyogh}/&$\\rightarrow$ & \\textit{7501} & /wlyri{\texttheta}le{\textyogh}/&$\\rightarrow$ & \\textit{7600} & /lri{\texttheta}le{\textyogh}/&$\\rightarrow$ & \\textit{12001} & /lri{\texttheta}le{\textesh}/&$\\rightarrow$ & \\textit{12500} & /lri{\texttheta}e{\textesh}/&$\\rightarrow$ & \\textit{15001} & /leri{\texttheta}e{\textesh}/& \\\\
\\end{tabular}

\\vspace{20pt}\\hline

\\end{nopagebreak}
\\filbreak



\\vspace{15pt}
\\begin{nopagebreak}
\\noindent{\\fontsize{20pt}{10pt}\\textbf{nalor} } \\textit{tie/bind/knot} (N)\\\\
\\noindent {\\tovian \\fontsize{20pt}{10pt} \\textbf{nalor} }\\\\
\\noindent /n{\textprimstress}alor/\\\\
\\noindent Syllables: nal.or\\\\
\\noindent lit. joining+action\\\\


\\noindent History:

\\vspace{-0pt}
\\hspace{40pt}
\\begin{tabular}{ccc}
\\textit{10} & /{\texttoptiebar{t\textbeltl}}unalore/&$\\rightarrow$ & \\textit{3200} & /{\texttoptiebar{t\textbeltl}}nalore/&$\\rightarrow$ & \\textit{6501} & /nalore/&$\\rightarrow$ & \\textit{9500} & /nalor/& \\\\
\\end{tabular}

\\vspace{20pt}\\hline

\\end{nopagebreak}
\\filbreak



\\vspace{15pt}
\\begin{nopagebreak}
\\noindent{\\fontsize{20pt}{10pt}\\textbf{ngothr} } \\textit{time} (N)\\\\
\\noindent {\\tovian \\fontsize{20pt}{10pt} \\textbf{ngothr} }\\\\
\\noindent /{\ng}{\textprimstress}o{\texttheta}r/\\\\
\\noindent Syllables: ŋoθr\\\\


\\noindent History:

\\vspace{-0pt}
\\hspace{40pt}
\\begin{tabular}{ccc}
\\textit{0} & /nko{\texttheta}ir/&$\\rightarrow$ & \\textit{2000} & /ngo{\texttheta}ir/&$\\rightarrow$ & \\textit{3200} & /ngo{\texttheta}r/&$\\rightarrow$ & \\textit{6000} & /{\ng}go{\texttheta}r/&$\\rightarrow$ & \\textit{8750} & /{\ng}o{\texttheta}r/& \\\\
\\end{tabular}

\\vspace{20pt}\\hline

\\end{nopagebreak}
\\filbreak



\\vspace{15pt}
\\begin{nopagebreak}
\\noindent{\\fontsize{20pt}{10pt}\\textbf{than} } \\textit{tool} (N)\\\\
\\noindent {\\tovian \\fontsize{20pt}{10pt} \\textbf{than} }\\\\
\\noindent /{\texttheta}{\textprimstress}an/\\\\
\\noindent Syllables: θan\\\\


\\noindent History:

\\vspace{-0pt}
\\hspace{40pt}
\\begin{tabular}{ccc}
\\textit{0} & /{\dh}an/&$\\rightarrow$ & \\textit{12002} & /{\texttheta}an/& \\\\
\\end{tabular}

\\vspace{20pt}\\hline

\\end{nopagebreak}
\\filbreak



\\vspace{15pt}
\\begin{nopagebreak}
\\noindent{\\fontsize{20pt}{10pt}\\textbf{gang} } \\textit{tooth} (N)\\\\
\\noindent {\\tovian \\fontsize{20pt}{10pt} \\textbf{gang} }\\\\
\\noindent /g{\textprimstress}a{\ng}/\\\\
\\noindent Syllables: gaŋ\\\\


\\noindent History:

\\vspace{-0pt}
\\hspace{40pt}
\\begin{tabular}{ccc}
\\textit{0} & /ganga/&$\\rightarrow$ & \\textit{6000} & /ga{\ng}ga/&$\\rightarrow$ & \\textit{8750} & /ga{\ng}a/&$\\rightarrow$ & \\textit{9500} & /ga{\ng}/& \\\\
\\end{tabular}

\\vspace{20pt}\\hline

\\end{nopagebreak}
\\filbreak



\\vspace{15pt}
\\begin{nopagebreak}
\\noindent{\\fontsize{20pt}{10pt}\\textbf{thul} } \\textit{top/height} (N)\\\\
\\noindent {\\tovian \\fontsize{20pt}{10pt} \\textbf{thul} }\\\\
\\noindent /{\texttheta}{\textprimstress}ul/\\\\
\\noindent Syllables: θul\\\\


\\noindent History:

\\vspace{-0pt}
\\hspace{40pt}
\\begin{tabular}{ccc}
\\textit{0} & /{\dh}ul/&$\\rightarrow$ & \\textit{12002} & /{\texttheta}ul/& \\\\
\\end{tabular}

\\vspace{20pt}\\hline

\\end{nopagebreak}
\\filbreak



\\vspace{15pt}
\\begin{nopagebreak}
\\noindent{\\fontsize{20pt}{10pt}\\textbf{le-fethr} } \\textit{Tovian/language} (N)\\\\
\\noindent {\\tovian \\fontsize{20pt}{10pt} \\textbf{le-fethr} }\\\\
\\noindent /l{\textprimstress}efe{\texttheta}r/\\\\
\\noindent Syllables: le.feθr\\\\
\\noindent lit. ABS+speech\\\\


\\noindent History:

\\vspace{-0pt}
\\hspace{40pt}
\\begin{tabular}{ccc}
\\textit{9999999} & /le-fe{\texttheta}r/& \\\\
\\end{tabular}

\\vspace{20pt}\\hline

\\end{nopagebreak}
\\filbreak



\\vspace{15pt}
\\begin{nopagebreak}
\\noindent{\\fontsize{20pt}{10pt}\\textbf{pel} } \\textit{town/village} (N)\\\\
\\noindent {\\tovian \\fontsize{20pt}{10pt} \\textbf{pel} }\\\\
\\noindent /p{\textprimstress}el/\\\\
\\noindent Syllables: pel\\\\


\\noindent History:

\\vspace{-0pt}
\\hspace{40pt}
\\begin{tabular}{ccc}
\\textit{0} & /pelor/&$\\rightarrow$ & \\textit{3200} & /pelr/&$\\rightarrow$ & \\textit{6503} & /perl/&$\\rightarrow$ & \\textit{7600} & /pel/& \\\\
\\end{tabular}

\\vspace{20pt}\\hline

\\end{nopagebreak}
\\filbreak



\\vspace{15pt}
\\begin{nopagebreak}
\\noindent{\\fontsize{20pt}{10pt}\\textbf{lheth} } \\textit{trance/sleep} (N)\\\\
\\noindent {\\tovian \\fontsize{20pt}{10pt} \\textbf{lheth} }\\\\
\\noindent /{\textbeltl}{\textprimstress}e{\texttheta}/\\\\
\\noindent Syllables: ɬeθ\\\\


\\noindent History:

\\vspace{-0pt}
\\hspace{40pt}
\\begin{tabular}{ccc}
\\textit{0} & /{\textbeltl}e{\dh}el/&$\\rightarrow$ & \\textit{3200} & /{\textbeltl}e{\dh}l/&$\\rightarrow$ & \\textit{12002} & /{\textbeltl}e{\texttheta}l/&$\\rightarrow$ & \\textit{12500} & /{\textbeltl}e{\texttheta}/& \\\\
\\end{tabular}

\\vspace{20pt}\\hline

\\end{nopagebreak}
\\filbreak



\\vspace{15pt}
\\begin{nopagebreak}
\\noindent{\\fontsize{20pt}{10pt}\\textbf{gelar} } \\textit{tree} (N)\\\\
\\noindent {\\tovian \\fontsize{20pt}{10pt} \\textbf{gelar} }\\\\
\\noindent /g{\textprimstress}elar/\\\\
\\noindent Syllables: gel.ar\\\\


\\noindent History:

\\vspace{-0pt}
\\hspace{40pt}
\\begin{tabular}{ccc}
\\textit{0} & /aklare/&$\\rightarrow$ & \\textit{2000} & /aglare/&$\\rightarrow$ & \\textit{7500} & /glare/&$\\rightarrow$ & \\textit{9500} & /glar/&$\\rightarrow$ & \\textit{15001} & /gelar/& \\\\
\\end{tabular}

\\vspace{20pt}\\hline

\\end{nopagebreak}
\\filbreak



\\vspace{15pt}
\\begin{nopagebreak}
\\noindent{\\fontsize{20pt}{10pt}\\textbf{lelom} } \\textit{turn/rotate} (N)\\\\
\\noindent {\\tovian \\fontsize{20pt}{10pt} \\textbf{lelom} }\\\\
\\noindent /l{\textprimstress}elom/\\\\
\\noindent Syllables: lel.om\\\\
\\noindent lit. movement+cycle\\\\


\\noindent History:

\\vspace{-0pt}
\\hspace{40pt}
\\begin{tabular}{ccc}
\\textit{1000} & /kwilelakorim/&$\\rightarrow$ & \\textit{2000} & /kwilelagorim/&$\\rightarrow$ & \\textit{3200} & /kwlelgorm/&$\\rightarrow$ & \\textit{4501} & /kwlelorm/&$\\rightarrow$ & \\textit{7501} & /wlelorm/&$\\rightarrow$ & \\textit{7600} & /lelom/& \\\\
\\end{tabular}

\\vspace{20pt}\\hline

\\end{nopagebreak}
\\filbreak



\\vspace{15pt}
\\begin{nopagebreak}
\\noindent{\\fontsize{20pt}{10pt}\\textbf{lheya} } \\textit{twelve/12} (N)\\\\
\\noindent {\\tovian \\fontsize{20pt}{10pt} \\textbf{lheya} }\\\\
\\noindent /{\textbeltl}{\textprimstress}eya/\\\\
\\noindent Syllables: ɬej.a\\\\


\\noindent History:

\\vspace{-0pt}
\\hspace{40pt}
\\begin{tabular}{ccc}
\\textit{0} & /{\textbeltl}iya/&$\\rightarrow$ & \\textit{2100} & /{\textbeltl}ya/&$\\rightarrow$ & \\textit{15001} & /{\textbeltl}eya/& \\\\
\\end{tabular}

\\vspace{20pt}\\hline

\\end{nopagebreak}
\\filbreak



\\vspace{15pt}
\\begin{nopagebreak}
\\noindent{\\fontsize{20pt}{10pt}\\textbf{loelhyal} } \\textit{twenty five/25} (N)\\\\
\\noindent {\\tovian \\fontsize{20pt}{10pt} \\textbf{loelhyal} }\\\\
\\noindent /lo{\textprimstress}e{\textbeltl}yal/\\\\
\\noindent Syllables: lo.eɬj.al\\\\
\\noindent lit. twenty four+one\\\\


\\noindent History:

\\vspace{-0pt}
\\hspace{40pt}
\\begin{tabular}{ccc}
\\textit{4000} & /loe{\textbeltl}yalo/&$\\rightarrow$ & \\textit{9500} & /loe{\textbeltl}yal/& \\\\
\\end{tabular}

\\vspace{20pt}\\hline

\\end{nopagebreak}
\\filbreak



\\vspace{15pt}
\\begin{nopagebreak}
\\noindent{\\fontsize{20pt}{10pt}\\textbf{loelh} } \\textit{twenty four/24} (N)\\\\
\\noindent {\\tovian \\fontsize{20pt}{10pt} \\textbf{loelh} }\\\\
\\noindent /l{\textprimstress}oe{\textbeltl}/\\\\
\\noindent Syllables: lo.eɬ\\\\
\\noindent lit. two+twelve\\\\


\\noindent History:

\\vspace{-0pt}
\\hspace{40pt}
\\begin{tabular}{ccc}
\\textit{4000} & /loe{\textbeltl}ya/&$\\rightarrow$ & \\textit{9500} & /loe{\textbeltl}y/&$\\rightarrow$ & \\textit{13001} & /loe{\textbeltl}/& \\\\
\\end{tabular}

\\vspace{20pt}\\hline

\\end{nopagebreak}
\\filbreak



\\vspace{15pt}
\\begin{nopagebreak}
\\noindent{\\fontsize{20pt}{10pt}\\textbf{loelhyalo} } \\textit{twenty six/26} (N)\\\\
\\noindent {\\tovian \\fontsize{20pt}{10pt} \\textbf{loelhyalo} }\\\\
\\noindent /loe{\textbeltl}y{\textprimstress}alo/\\\\
\\noindent Syllables: lo.eɬj.al.o\\\\
\\noindent lit. twenty four+two\\\\


\\noindent History:

\\vspace{-0pt}
\\hspace{40pt}
\\begin{tabular}{ccc}
\\textit{4000} & /loe{\textbeltl}yaloe/&$\\rightarrow$ & \\textit{9500} & /loe{\textbeltl}yalo/& \\\\
\\end{tabular}

\\vspace{20pt}\\hline

\\end{nopagebreak}
\\filbreak



\\vspace{15pt}
\\begin{nopagebreak}
\\noindent{\\fontsize{20pt}{10pt}\\textbf{lheyath} } \\textit{twenty/20} (N)\\\\
\\noindent {\\tovian \\fontsize{20pt}{10pt} \\textbf{lheyath} }\\\\
\\noindent /{\textbeltl}{\textprimstress}eya{\texttheta}/\\\\
\\noindent Syllables: ɬej.aθ\\\\
\\noindent lit. twelve+eight\\\\


\\noindent History:

\\vspace{-0pt}
\\hspace{40pt}
\\begin{tabular}{ccc}
\\textit{4000} & /{\textbeltl}yaa{\texttheta}/&$\\rightarrow$ & \\textit{5500} & /{\textbeltl}ya{\texttheta}/&$\\rightarrow$ & \\textit{15001} & /{\textbeltl}eya{\texttheta}/& \\\\
\\end{tabular}

\\vspace{20pt}\\hline

\\end{nopagebreak}
\\filbreak



\\vspace{15pt}
\\begin{nopagebreak}
\\noindent{\\fontsize{20pt}{10pt}\\textbf{lo} } \\textit{two/2} (Ns)\\\\
\\noindent {\\tovian \\fontsize{20pt}{10pt} \\textbf{lo} }\\\\
\\noindent /l{\textprimstress}o/\\\\
\\noindent Syllables: lo\\\\


\\noindent History:

\\vspace{-0pt}
\\hspace{40pt}
\\begin{tabular}{ccc}
\\textit{0} & /lo/& \\\\
\\end{tabular}

\\vspace{20pt}\\hline

\\end{nopagebreak}
\\filbreak



\\vspace{15pt}
\\begin{nopagebreak}
\\noindent{\\fontsize{20pt}{10pt}\\textbf{hi-thokan} } \\textit{undershine} (N)\\\\
\\noindent {\\tovian \\fontsize{20pt}{10pt} \\textbf{hi-thokan} }\\\\
\\noindent /ħi{\texttheta}{\textprimstress}okan/\\\\
\\noindent Syllables: ħi.θok.an\\\\
\\noindent lit. GEN+bottom+light\\\\


\\noindent History:

\\vspace{-0pt}
\\hspace{40pt}
\\begin{tabular}{ccc}
\\textit{14000} & /hi-{\texttheta}okan/&$\\rightarrow$ & \\textit{14006} & /ħi-{\texttheta}okan/& \\\\
\\end{tabular}

\\vspace{20pt}\\hline

\\end{nopagebreak}
\\filbreak



\\vspace{15pt}
\\begin{nopagebreak}
\\noindent{\\fontsize{20pt}{10pt}\\textbf{shim} } \\textit{unknown/mystery} (N)\\\\
\\noindent {\\tovian \\fontsize{20pt}{10pt} \\textbf{shim} }\\\\
\\noindent /{\textesh}{\textprimstress}im/\\\\
\\noindent Syllables: ʃim\\\\


\\noindent History:

\\vspace{-0pt}
\\hspace{40pt}
\\begin{tabular}{ccc}
\\textit{0} & /{\textesh}imi/&$\\rightarrow$ & \\textit{9500} & /{\textesh}im/& \\\\
\\end{tabular}

\\vspace{20pt}\\hline

\\end{nopagebreak}
\\filbreak



\\vspace{15pt}
\\begin{nopagebreak}
\\noindent{\\fontsize{20pt}{10pt}\\textbf{ther} } \\textit{up} (N)\\\\
\\noindent {\\tovian \\fontsize{20pt}{10pt} \\textbf{ther} }\\\\
\\noindent /{\texttheta}{\textprimstress}er/\\\\
\\noindent Syllables: θer\\\\


\\noindent History:

\\vspace{-0pt}
\\hspace{40pt}
\\begin{tabular}{ccc}
\\textit{0} & /{\texttheta}era/&$\\rightarrow$ & \\textit{9500} & /{\texttheta}er/& \\\\
\\end{tabular}

\\vspace{20pt}\\hline

\\end{nopagebreak}
\\filbreak



\\vspace{15pt}
\\begin{nopagebreak}
\\noindent{\\fontsize{20pt}{10pt}\\textbf{n} } \\textit{us/we} (Ps)\\\\
\\noindent {\\tovian \\fontsize{20pt}{10pt} \\textbf{n} }\\\\
\\noindent /n/\\\\
\\noindent Syllables: n\\\\


\\noindent History:

\\vspace{-0pt}
\\hspace{40pt}
\\begin{tabular}{ccc}
\\textit{0} & /ne/&$\\rightarrow$ & \\textit{12006} & /n/& \\\\
\\end{tabular}

\\vspace{20pt}\\hline

\\end{nopagebreak}
\\filbreak



\\vspace{15pt}
\\begin{nopagebreak}
\\noindent{\\fontsize{20pt}{10pt}\\textbf{var} } \\textit{value} (N)\\\\
\\noindent {\\tovian \\fontsize{20pt}{10pt} \\textbf{var} }\\\\
\\noindent /v{\textprimstress}ar/\\\\
\\noindent Syllables: var\\\\


\\noindent History:

\\vspace{-0pt}
\\hspace{40pt}
\\begin{tabular}{ccc}
\\textit{0} & /velare/&$\\rightarrow$ & \\textit{3200} & /vlare/&$\\rightarrow$ & \\textit{9500} & /vlar/&$\\rightarrow$ & \\textit{12500} & /var/& \\\\
\\end{tabular}

\\vspace{20pt}\\hline

\\end{nopagebreak}
\\filbreak



\\vspace{15pt}
\\begin{nopagebreak}
\\noindent{\\fontsize{20pt}{10pt}\\textbf{kelumlor} } \\textit{vine} (N)\\\\
\\noindent {\\tovian \\fontsize{20pt}{10pt} \\textbf{kelumlor} }\\\\
\\noindent /kel{\textprimstress}umlor/\\\\
\\noindent Syllables: kel.uml.or\\\\
\\noindent lit. arm+plant\\\\


\\noindent History:

\\vspace{-0pt}
\\hspace{40pt}
\\begin{tabular}{ccc}
\\textit{200} & /kelu{\textschwa}mellore/&$\\rightarrow$ & \\textit{2202} & /keluw{\textschwa}mellore/&$\\rightarrow$ & \\textit{2300} & /keluwmellore/&$\\rightarrow$ & \\textit{3200} & /kluwmllore/&$\\rightarrow$ & \\textit{6500} & /kluwmlore/&$\\rightarrow$ & \\textit{7600} & /klumlore/&$\\rightarrow$ & \\textit{9500} & /klumlor/&$\\rightarrow$ & \\textit{15001} & /kelumlor/& \\\\
\\end{tabular}

\\vspace{20pt}\\hline

\\end{nopagebreak}
\\filbreak



\\vspace{15pt}
\\begin{nopagebreak}
\\noindent{\\fontsize{20pt}{10pt}\\textbf{anshish} } \\textit{waking dreambringer} (N)\\\\
\\noindent {\\tovian \\fontsize{20pt}{10pt} \\textbf{anshish} }\\\\
\\noindent /{\textprimstress}an{\textesh}i{\textesh}/\\\\
\\noindent Syllables: anʃ.iʃ\\\\
\\noindent lit. dream+COM+alertness\\\\


\\noindent History:

\\vspace{-0pt}
\\hspace{40pt}
\\begin{tabular}{ccc}
\\textit{10} & /{\textbeltl}enayi-{\textyogh}i{\texttheta}/&$\\rightarrow$ & \\textit{1000} & /{\textbeltl}enayi{\textyogh}i{\texttheta}/&$\\rightarrow$ & \\textit{2100} & /{\textbeltl}enyi{\textyogh}i{\texttheta}/&$\\rightarrow$ & \\textit{3200} & /{\textbeltl}nyi{\textyogh}{\texttheta}/&$\\rightarrow$ & \\textit{4502} & /{\textbeltl}nyi{\textyogh}/&$\\rightarrow$ & \\textit{6000} & /{\textbeltl}nd{\textyogh}i{\textyogh}/&$\\rightarrow$ & \\textit{7501} & /nd{\textyogh}i{\textyogh}/&$\\rightarrow$ & \\textit{8750} & /n{\textyogh}i{\textyogh}/&$\\rightarrow$ & \\textit{8760} & /an{\textyogh}i{\textyogh}/&$\\rightarrow$ & \\textit{12001} & /an{\textesh}i{\textesh}/& \\\\
\\end{tabular}

\\vspace{20pt}\\hline

\\end{nopagebreak}
\\filbreak



\\vspace{15pt}
\\begin{nopagebreak}
\\noindent{\\fontsize{20pt}{10pt}\\textbf{nolh} } \\textit{walking} (N)\\\\
\\noindent {\\tovian \\fontsize{20pt}{10pt} \\textbf{nolh} }\\\\
\\noindent /n{\textprimstress}o{\textbeltl}/\\\\
\\noindent Syllables: noɬ\\\\


\\noindent History:

\\vspace{-0pt}
\\hspace{40pt}
\\begin{tabular}{ccc}
\\textit{0} & /no{\textbeltl}/& \\\\
\\end{tabular}

\\vspace{20pt}\\hline

\\end{nopagebreak}
\\filbreak



\\vspace{15pt}
\\begin{nopagebreak}
\\noindent{\\fontsize{20pt}{10pt}\\textbf{sor} } \\textit{washing} (N)\\\\
\\noindent {\\tovian \\fontsize{20pt}{10pt} \\textbf{sor} }\\\\
\\noindent /s{\textprimstress}or/\\\\
\\noindent Syllables: sor\\\\
\\noindent lit. water+action\\\\


\\noindent History:

\\vspace{-0pt}
\\hspace{40pt}
\\begin{tabular}{ccc}
\\textit{0} & /z{\textsubbridge{t}}ellore/&$\\rightarrow$ & \\textit{2000} & /zdellore/&$\\rightarrow$ & \\textit{3200} & /zdllore/&$\\rightarrow$ & \\textit{4500} & /zllore/&$\\rightarrow$ & \\textit{6500} & /zlore/&$\\rightarrow$ & \\textit{9500} & /zlor/&$\\rightarrow$ & \\textit{12000} & /slor/&$\\rightarrow$ & \\textit{12500} & /sor/& \\\\
\\end{tabular}

\\vspace{20pt}\\hline

\\end{nopagebreak}
\\filbreak



\\vspace{15pt}
\\begin{nopagebreak}
\\noindent{\\fontsize{20pt}{10pt}\\textbf{sel} } \\textit{water} (N)\\\\
\\noindent {\\tovian \\fontsize{20pt}{10pt} \\textbf{sel} }\\\\
\\noindent /s{\textprimstress}el/\\\\
\\noindent Syllables: sel\\\\


\\noindent History:

\\vspace{-0pt}
\\hspace{40pt}
\\begin{tabular}{ccc}
\\textit{0} & /z{\textsubbridge{t}}el/&$\\rightarrow$ & \\textit{2000} & /zdel/&$\\rightarrow$ & \\textit{4500} & /zel/&$\\rightarrow$ & \\textit{12000} & /sel/& \\\\
\\end{tabular}

\\vspace{20pt}\\hline

\\end{nopagebreak}
\\filbreak



\\vspace{15pt}
\\begin{nopagebreak}
\\noindent{\\fontsize{20pt}{10pt}\\textbf{tirlhom} } \\textit{week} (N)\\\\
\\noindent {\\tovian \\fontsize{20pt}{10pt} \\textbf{tirlhom} }\\\\
\\noindent /{\textsubbridge{t}}{\textprimstress}ir{\textbeltl}om/\\\\
\\noindent Syllables: tirɬ.om\\\\
\\noindent lit. ten+day\\\\


\\noindent History:

\\vspace{-0pt}
\\hspace{40pt}
\\begin{tabular}{ccc}
\\textit{4000} & /{\textsubbridge{t}}ir{\textbeltl}onm/&$\\rightarrow$ & \\textit{7600} & /{\textsubbridge{t}}ir{\textbeltl}om/& \\\\
\\end{tabular}

\\vspace{20pt}\\hline

\\end{nopagebreak}
\\filbreak



\\vspace{15pt}
\\begin{nopagebreak}
\\noindent{\\fontsize{20pt}{10pt}\\textbf{tlem} } \\textit{weight} (N)\\\\
\\noindent {\\tovian \\fontsize{20pt}{10pt} \\textbf{tlem} }\\\\
\\noindent /{\texttoptiebar{t\textbeltl}}{\textprimstress}em/\\\\
\\noindent Syllables: ƛem\\\\


\\noindent History:

\\vspace{-0pt}
\\hspace{40pt}
\\begin{tabular}{ccc}
\\textit{0} & /a{\texttoptiebar{t\textbeltl}}emo/&$\\rightarrow$ & \\textit{7500} & /{\texttoptiebar{t\textbeltl}}emo/&$\\rightarrow$ & \\textit{9500} & /{\texttoptiebar{t\textbeltl}}em/& \\\\
\\end{tabular}

\\vspace{20pt}\\hline

\\end{nopagebreak}
\\filbreak



\\vspace{15pt}
\\begin{nopagebreak}
\\noindent{\\fontsize{20pt}{10pt}\\textbf{fil} } \\textit{west} (N)\\\\
\\noindent {\\tovian \\fontsize{20pt}{10pt} \\textbf{fil} }\\\\
\\noindent /f{\textprimstress}il/\\\\
\\noindent Syllables: fil\\\\


\\noindent History:

\\vspace{-0pt}
\\hspace{40pt}
\\begin{tabular}{ccc}
\\textit{0} & /fil/& \\\\
\\end{tabular}

\\vspace{20pt}\\hline

\\end{nopagebreak}
\\filbreak



\\vspace{15pt}
\\begin{nopagebreak}
\\noindent{\\fontsize{20pt}{10pt}\\textbf{lo-lhan} } \\textit{what} (INT)\\\\
\\noindent {\\tovian \\fontsize{20pt}{10pt} \\textbf{lo-lhan} }\\\\
\\noindent /l{\textprimstress}o{\textbeltl}an/\\\\
\\noindent Syllables: lo.ɬan\\\\
\\noindent lit. INAN+INT\\\\


\\noindent History:

\\vspace{-0pt}
\\hspace{40pt}
\\begin{tabular}{ccc}
\\textit{0} & /lo-{\textbeltl}an/& \\\\
\\end{tabular}

\\vspace{20pt}\\hline

\\end{nopagebreak}
\\filbreak



\\vspace{15pt}
\\begin{nopagebreak}
\\noindent{\\fontsize{20pt}{10pt}\\textbf{tu-lhan} } \\textit{when} (INT)\\\\
\\noindent {\\tovian \\fontsize{20pt}{10pt} \\textbf{tu-lhan} }\\\\
\\noindent /{\textsubbridge{t}}{\textprimstress}u{\textbeltl}an/\\\\
\\noindent Syllables: tu.ɬan\\\\
\\noindent lit. TEMP+INT\\\\


\\noindent History:

\\vspace{-0pt}
\\hspace{40pt}
\\begin{tabular}{ccc}
\\textit{0} & /{\textsubbridge{t}}u-{\textbeltl}an/& \\\\
\\end{tabular}

\\vspace{20pt}\\hline

\\end{nopagebreak}
\\filbreak



\\vspace{15pt}
\\begin{nopagebreak}
\\noindent{\\fontsize{20pt}{10pt}\\textbf{ti-lhan} } \\textit{where} (INT)\\\\
\\noindent {\\tovian \\fontsize{20pt}{10pt} \\textbf{ti-lhan} }\\\\
\\noindent /{\textsubbridge{t}}{\textprimstress}i{\textbeltl}an/\\\\
\\noindent Syllables: ti.ɬan\\\\
\\noindent lit. LOC+INT\\\\


\\noindent History:

\\vspace{-0pt}
\\hspace{40pt}
\\begin{tabular}{ccc}
\\textit{0} & /{\textsubbridge{t}}i-{\textbeltl}an/& \\\\
\\end{tabular}

\\vspace{20pt}\\hline

\\end{nopagebreak}
\\filbreak



\\vspace{15pt}
\\begin{nopagebreak}
\\noindent{\\fontsize{20pt}{10pt}\\textbf{hi-lhanshimi} } \\textit{whichever} (_)\\\\
\\noindent {\\tovian \\fontsize{20pt}{10pt} \\textbf{hi-lhanshimi} }\\\\
\\noindent /hi{\textbeltl}an{\textesh}{\textprimstress}imi/\\\\
\\noindent Syllables: hi.ɬanʃ.im.i\\\\
\\noindent lit. GEN+INT+unknown\\\\
\\noindent \\textit{INT}\\\\


\\noindent History:

\\vspace{-0pt}
\\hspace{40pt}
\\begin{tabular}{ccc}
\\textit{0} & /hi-{\textbeltl}an{\textesh}imi/& \\\\
\\end{tabular}

\\vspace{20pt}\\hline

\\end{nopagebreak}
\\filbreak



\\vspace{15pt}
\\begin{nopagebreak}
\\noindent{\\fontsize{20pt}{10pt}\\textbf{i-lhan} } \\textit{who} (INT)\\\\
\\noindent {\\tovian \\fontsize{20pt}{10pt} \\textbf{i-lhan} }\\\\
\\noindent /{\textprimstress}i{\textbeltl}an/\\\\
\\noindent Syllables: i.ɬan\\\\
\\noindent lit. NOM+INT\\\\


\\noindent History:

\\vspace{-0pt}
\\hspace{40pt}
\\begin{tabular}{ccc}
\\textit{0} & /i-{\textbeltl}an/& \\\\
\\end{tabular}

\\vspace{20pt}\\hline

\\end{nopagebreak}
\\filbreak



\\vspace{15pt}
\\begin{nopagebreak}
\\noindent{\\fontsize{20pt}{10pt}\\textbf{hi-lhanshimi} } \\textit{whoever} (_)\\\\
\\noindent {\\tovian \\fontsize{20pt}{10pt} \\textbf{hi-lhanshimi} }\\\\
\\noindent /hi{\textbeltl}an{\textesh}{\textprimstress}imi/\\\\
\\noindent Syllables: hi.ɬanʃ.im.i\\\\
\\noindent lit. GEN+INT+unknown\\\\
\\noindent \\textit{INT}\\\\


\\noindent History:

\\vspace{-0pt}
\\hspace{40pt}
\\begin{tabular}{ccc}
\\textit{0} & /hi-{\textbeltl}an{\textesh}imi/& \\\\
\\end{tabular}

\\vspace{20pt}\\hline

\\end{nopagebreak}
\\filbreak



\\vspace{15pt}
\\begin{nopagebreak}
\\noindent{\\fontsize{20pt}{10pt}\\textbf{yu-lhan} } \\textit{whom} (INT)\\\\
\\noindent {\\tovian \\fontsize{20pt}{10pt} \\textbf{yu-lhan} }\\\\
\\noindent /y{\textprimstress}u{\textbeltl}an/\\\\
\\noindent Syllables: ju.ɬan\\\\
\\noindent lit. ACC+INT\\\\


\\noindent History:

\\vspace{-0pt}
\\hspace{40pt}
\\begin{tabular}{ccc}
\\textit{0} & /yu-{\textbeltl}an/& \\\\
\\end{tabular}

\\vspace{20pt}\\hline

\\end{nopagebreak}
\\filbreak



\\vspace{15pt}
\\begin{nopagebreak}
\\noindent{\\fontsize{20pt}{10pt}\\textbf{hi-lhan} } \\textit{whose} (INT)\\\\
\\noindent {\\tovian \\fontsize{20pt}{10pt} \\textbf{hi-lhan} }\\\\
\\noindent /h{\textprimstress}i{\textbeltl}an/\\\\
\\noindent Syllables: hi.ɬan\\\\
\\noindent lit. GEN+INT\\\\


\\noindent History:

\\vspace{-0pt}
\\hspace{40pt}
\\begin{tabular}{ccc}
\\textit{0} & /hi-{\textbeltl}an/& \\\\
\\end{tabular}

\\vspace{20pt}\\hline

\\end{nopagebreak}
\\filbreak



\\vspace{15pt}
\\begin{nopagebreak}
\\noindent{\\fontsize{20pt}{10pt}\\textbf{lhu-lhan} } \\textit{why} (INT)\\\\
\\noindent {\\tovian \\fontsize{20pt}{10pt} \\textbf{lhu-lhan} }\\\\
\\noindent /{\textbeltl}{\textprimstress}u{\textbeltl}an/\\\\
\\noindent Syllables: ɬu.ɬan\\\\
\\noindent lit. PUR+INT\\\\


\\noindent History:

\\vspace{-0pt}
\\hspace{40pt}
\\begin{tabular}{ccc}
\\textit{0} & /{\textbeltl}u-{\textbeltl}an/& \\\\
\\end{tabular}

\\vspace{20pt}\\hline

\\end{nopagebreak}
\\filbreak



\\vspace{15pt}
\\begin{nopagebreak}
\\noindent{\\fontsize{20pt}{10pt}\\textbf{shad} } \\textit{width} (N)\\\\
\\noindent {\\tovian \\fontsize{20pt}{10pt} \\textbf{shad} }\\\\
\\noindent /{\textesh}{\textprimstress}ad/\\\\
\\noindent Syllables: ʃad\\\\


\\noindent History:

\\vspace{-0pt}
\\hspace{40pt}
\\begin{tabular}{ccc}
\\textit{0} & /{\textyogh}ila{\textsubbridge{t}}e/&$\\rightarrow$ & \\textit{2000} & /{\textyogh}ilade/&$\\rightarrow$ & \\textit{3200} & /{\textyogh}lade/&$\\rightarrow$ & \\textit{9500} & /{\textyogh}lad/&$\\rightarrow$ & \\textit{12001} & /{\textesh}lad/&$\\rightarrow$ & \\textit{12500} & /{\textesh}ad/& \\\\
\\end{tabular}

\\vspace{20pt}\\hline

\\end{nopagebreak}
\\filbreak



\\vspace{15pt}
\\begin{nopagebreak}
\\noindent{\\fontsize{20pt}{10pt}\\textbf{lerith} } \\textit{wind/breath} (N)\\\\
\\noindent {\\tovian \\fontsize{20pt}{10pt} \\textbf{lerith} }\\\\
\\noindent /l{\textprimstress}eri{\texttheta}/\\\\
\\noindent Syllables: ler.iθ\\\\
\\noindent lit. movement+COM+air\\\\


\\noindent History:

\\vspace{-0pt}
\\hspace{40pt}
\\begin{tabular}{ccc}
\\textit{10} & /kwilyi-rai{\texttheta}/&$\\rightarrow$ & \\textit{1000} & /kwilyirai{\texttheta}/&$\\rightarrow$ & \\textit{3200} & /kwlyrai{\texttheta}/&$\\rightarrow$ & \\textit{5500} & /kwlyri{\texttheta}/&$\\rightarrow$ & \\textit{7501} & /wlyri{\texttheta}/&$\\rightarrow$ & \\textit{7600} & /lri{\texttheta}/&$\\rightarrow$ & \\textit{15001} & /leri{\texttheta}/& \\\\
\\end{tabular}

\\vspace{20pt}\\hline

\\end{nopagebreak}
\\filbreak



\\vspace{15pt}
\\begin{nopagebreak}
\\noindent{\\fontsize{20pt}{10pt}\\textbf{mulhetush} } \\textit{Winter} (N)\\\\
\\noindent {\\tovian \\fontsize{20pt}{10pt} \\textbf{mulhetush} }\\\\
\\noindent /mu{\textbeltl}{\textprimstress}e{\textsubbridge{t}}u{\textesh}/\\\\
\\noindent Syllables: muɬ.et.uʃ\\\\
\\noindent lit. present+cold\\\\


\\noindent History:

\\vspace{-0pt}
\\hspace{40pt}
\\begin{tabular}{ccc}
\\textit{0} & /num{\textbeltl}ele{\textsubbridge{t}}ku{\textesh}u/&$\\rightarrow$ & \\textit{3000} & /num{\textbeltl}ele{\textsubbridge{t}}u{\textesh}u/&$\\rightarrow$ & \\textit{3200} & /nm{\textbeltl}le{\textsubbridge{t}}u{\textesh}u/&$\\rightarrow$ & \\textit{7501} & /m{\textbeltl}le{\textsubbridge{t}}u{\textesh}u/&$\\rightarrow$ & \\textit{9500} & /m{\textbeltl}le{\textsubbridge{t}}u{\textesh}/&$\\rightarrow$ & \\textit{12500} & /m{\textbeltl}e{\textsubbridge{t}}u{\textesh}/&$\\rightarrow$ & \\textit{15001} & /mu{\textbeltl}e{\textsubbridge{t}}u{\textesh}/& \\\\
\\end{tabular}

\\vspace{20pt}\\hline

\\end{nopagebreak}
\\filbreak



\\vspace{15pt}
\\begin{nopagebreak}
\\noindent{\\fontsize{20pt}{10pt}\\textbf{yeol} } \\textit{wisdom/truth} (N)\\\\
\\noindent {\\tovian \\fontsize{20pt}{10pt} \\textbf{yeol} }\\\\
\\noindent /y{\textprimstress}eol/\\\\
\\noindent Syllables: je.ol\\\\


\\noindent History:

\\vspace{-0pt}
\\hspace{40pt}
\\begin{tabular}{ccc}
\\textit{0} & /arayeol/&$\\rightarrow$ & \\textit{2100} & /aryeol/&$\\rightarrow$ & \\textit{7500} & /ryeol/&$\\rightarrow$ & \\textit{7600} & /yeol/& \\\\
\\end{tabular}

\\vspace{20pt}\\hline

\\end{nopagebreak}
\\filbreak



\\vspace{15pt}
\\begin{nopagebreak}
\\noindent{\\fontsize{20pt}{10pt}\\textbf{theraf} } \\textit{wondering/curiosity} (N)\\\\
\\noindent {\\tovian \\fontsize{20pt}{10pt} \\textbf{theraf} }\\\\
\\noindent /{\texttheta}{\textprimstress}eraf/\\\\
\\noindent Syllables: θer.af\\\\


\\noindent History:

\\vspace{-0pt}
\\hspace{40pt}
\\begin{tabular}{ccc}
\\textit{0} & /{\texttheta}irafa/&$\\rightarrow$ & \\textit{3200} & /{\texttheta}rafa/&$\\rightarrow$ & \\textit{9500} & /{\texttheta}raf/&$\\rightarrow$ & \\textit{15001} & /{\texttheta}eraf/& \\\\
\\end{tabular}

\\vspace{20pt}\\hline

\\end{nopagebreak}
\\filbreak



\\vspace{15pt}
\\begin{nopagebreak}
\\noindent{\\fontsize{20pt}{10pt}\\textbf{gelar} } \\textit{woods} (N)\\\\
\\noindent {\\tovian \\fontsize{20pt}{10pt} \\textbf{gelar} }\\\\
\\noindent /g{\textprimstress}elar/\\\\
\\noindent Syllables: gel.ar\\\\


\\noindent History:

\\vspace{-0pt}
\\hspace{40pt}
\\begin{tabular}{ccc}
\\textit{0} & /glar/&$\\rightarrow$ & \\textit{15001} & /gelar/& \\\\
\\end{tabular}

\\vspace{20pt}\\hline

\\end{nopagebreak}
\\filbreak



\\vspace{15pt}
\\begin{nopagebreak}
\\noindent{\\fontsize{20pt}{10pt}\\textbf{han} } \\textit{world} (N)\\\\
\\noindent {\\tovian \\fontsize{20pt}{10pt} \\textbf{han} }\\\\
\\noindent /ħ{\textprimstress}an/\\\\
\\noindent Syllables: ħan\\\\


\\noindent History:

\\vspace{-0pt}
\\hspace{40pt}
\\begin{tabular}{ccc}
\\textit{0} & /hu{\textsubbridge{t}}anna/&$\\rightarrow$ & \\textit{1000} & /h{\textsubbridge{t}}anna/&$\\rightarrow$ & \\textit{3000} & /hanna/&$\\rightarrow$ & \\textit{6500} & /hana/&$\\rightarrow$ & \\textit{9500} & /han/&$\\rightarrow$ & \\textit{14006} & /ħan/& \\\\
\\end{tabular}

\\vspace{20pt}\\hline

\\end{nopagebreak}
\\filbreak



\\vspace{15pt}
\\begin{nopagebreak}
\\noindent{\\fontsize{20pt}{10pt}\\textbf{for} } \\textit{writing/painting} (N)\\\\
\\noindent {\\tovian \\fontsize{20pt}{10pt} \\textbf{for} }\\\\
\\noindent /f{\textprimstress}or/\\\\
\\noindent Syllables: for\\\\
\\noindent lit. showing+action\\\\


\\noindent History:

\\vspace{-0pt}
\\hspace{40pt}
\\begin{tabular}{ccc}
\\textit{10} & /fe{\textsubbridge{t}}elore/&$\\rightarrow$ & \\textit{1000} & /f{\textsubbridge{t}}elore/&$\\rightarrow$ & \\textit{3000} & /felore/&$\\rightarrow$ & \\textit{3200} & /flore/&$\\rightarrow$ & \\textit{9500} & /flor/&$\\rightarrow$ & \\textit{12500} & /for/& \\\\
\\end{tabular}

\\vspace{20pt}\\hline

\\end{nopagebreak}
\\filbreak



\\vspace{15pt}
\\begin{nopagebreak}
\\noindent{\\fontsize{20pt}{10pt}\\textbf{nalom} } \\textit{year} (N)\\\\
\\noindent {\\tovian \\fontsize{20pt}{10pt} \\textbf{nalom} }\\\\
\\noindent /n{\textprimstress}alom/\\\\
\\noindent Syllables: nal.om\\\\
\\noindent lit. small+cycle\\\\


\\noindent History:

\\vspace{-0pt}
\\hspace{40pt}
\\begin{tabular}{ccc}
\\textit{10} & /ninaelakorim/&$\\rightarrow$ & \\textit{2000} & /ninaelagorim/&$\\rightarrow$ & \\textit{3200} & /nnaelgorm/&$\\rightarrow$ & \\textit{4501} & /nnaelorm/&$\\rightarrow$ & \\textit{6500} & /naelorm/&$\\rightarrow$ & \\textit{7600} & /naelom/&$\\rightarrow$ & \\textit{10000} & /nalom/& \\\\
\\end{tabular}

\\vspace{20pt}\\hline

\\end{nopagebreak}
\\filbreak



\\vspace{15pt}
\\begin{nopagebreak}
\\noindent{\\fontsize{20pt}{10pt}\\textbf{wa} } \\textit{you} (P)\\\\
\\noindent {\\tovian \\fontsize{20pt}{10pt} \\textbf{wa} }\\\\
\\noindent /w{\textprimstress}a/\\\\
\\noindent Syllables: wa\\\\


\\noindent History:

\\vspace{-0pt}
\\hspace{40pt}
\\begin{tabular}{ccc}
\\textit{0} & /wa/& \\\\
\\end{tabular}

\\vspace{20pt}\\hline

\\end{nopagebreak}
\\filbreak



\\vspace{15pt}
\\begin{nopagebreak}
\\noindent{\\fontsize{20pt}{10pt}\\textbf{w} } \\textit{you (plural)} (Ps)\\\\
\\noindent {\\tovian \\fontsize{20pt}{10pt} \\textbf{w} }\\\\
\\noindent /w/\\\\
\\noindent Syllables: w\\\\


\\noindent History:

\\vspace{-0pt}
\\hspace{40pt}
\\begin{tabular}{ccc}
\\textit{0} & /we/&$\\rightarrow$ & \\textit{12006} & /w/& \\\\
\\end{tabular}

\\vspace{20pt}\\hline

\\end{nopagebreak}
\\filbreak



\\vspace{15pt}
\\begin{nopagebreak}
\\noindent{\\fontsize{20pt}{10pt}\\textbf{don} } \\textit{young} (N)\\\\
\\noindent {\\tovian \\fontsize{20pt}{10pt} \\textbf{don} }\\\\
\\noindent /d{\textprimstress}on/\\\\
\\noindent Syllables: don\\\\
\\noindent lit. not+oldness\\\\


\\noindent History:

\\vspace{-0pt}
\\hspace{40pt}
\\begin{tabular}{ccc}
\\textit{0} & /donan/&$\\rightarrow$ & \\textit{3200} & /donn/&$\\rightarrow$ & \\textit{6500} & /don/& \\\\
\\end{tabular}

\\vspace{20pt}\\hline

\\end{nopagebreak}
\\filbreak



\\vspace{15pt}
\\begin{nopagebreak}
\\noindent{\\fontsize{20pt}{10pt}\\textbf{hi-wa} } \\textit{your} (_)\\\\
\\noindent {\\tovian \\fontsize{20pt}{10pt} \\textbf{hi-wa} }\\\\
\\noindent /h{\textprimstress}iwa/\\\\
\\noindent Syllables: hi.wa\\\\
\\noindent lit. GEN+you\\\\
\\noindent \\textit{P}\\\\


\\noindent History:

\\vspace{-0pt}
\\hspace{40pt}
\\begin{tabular}{ccc}
\\textit{0} & /hi-wa/& \\\\
\\end{tabular}

\\vspace{20pt}\\hline

\\end{nopagebreak}
\\filbreak



\\onecolumn


\end{document}
